\documentclass[11pt,letterpaper]{article}
\usepackage[T1]{fontenc}
\usepackage[dvipsnames]{xcolor}
\usepackage{amsmath,amsthm}
\usepackage{newtxtext,newtxmath}
\usepackage[margin=1in,headheight=15pt]{geometry}
\usepackage{microtype}
\usepackage{mathtools,booktabs,array,longtable}
\usepackage{graphicx,enumitem,fancyhdr}
\usepackage[most]{tcolorbox}
\usepackage{hyperref}
\hypersetup{colorlinks=true,linkcolor=MidnightBlue,citecolor=MidnightBlue,
 urlcolor=MidnightBlue,pdftitle={Fixed-height shifted rectangles: asymptotics and algebraicity},
 pdfsubject={Part I: asymptotics and algebraicity; Part II: rational diagonals and rare boundary phases; Parts III-IV: the A181198 and A181199 recurrences proved; Part V: two earlier independent routes},
 pdfauthor={Research report prepared for Vladimir Reshetnikov}}
\urlstyle{same}
\definecolor{inkblue}{HTML}{173752}
\definecolor{pale}{HTML}{F2F5F8}
\pagestyle{fancy}
\fancyhf{}
\fancyhead[L]{\small\textsc{Fixed-height shifted rectangles}}
\fancyhead[R]{\small OEIS research report}
\fancyfoot[C]{\thepage}
\renewcommand{\headrulewidth}{0.3pt}
\setlength{\parskip}{4pt}
\setlength{\parindent}{14pt}
\setlist{nosep,leftmargin=*}
\numberwithin{equation}{section}
\newtheorem{theorem}{Theorem}[section]
\newtheorem{proposition}[theorem]{Proposition}
\newtheorem{lemma}[theorem]{Lemma}
\newtheorem{corollary}[theorem]{Corollary}
\theoremstyle{definition}
\newtheorem{definition}[theorem]{Definition}
\newtheorem{remark}[theorem]{Remark}
\newtheorem{question}{Research question}
\newcommand{\E}{\mathbb E}
\newcommand{\Pp}{\mathbb P}
\newcommand{\R}{\mathbb R}
\newcommand{\Q}{\mathbb Q}
\newcommand{\N}{\mathbb N}
\newcommand{\1}{\mathbf 1}
\newcommand{\fall}[2]{#1^{\underline{#2}}}
\newcommand{\eps}{\varepsilon}
\DeclareMathOperator{\Tr}{Tr}
\DeclareMathOperator{\Var}{Var}
\DeclareMathOperator{\sgn}{sgn}
% Part II (batch 98, manuscript 02) macros, as delivered; Part I's \fall is used
\newcommand{\Z}{\mathbb Z}
\newcommand{\PP}{\mathbb P}
\newcommand{\ind}{\mathbf 1}
\newcommand{\rise}[2]{(#1)^{\overline{#2}}}
\newcommand{\Diag}{\operatorname{Diag}}
\newcommand{\Sym}{\mathfrak S}
\newcommand{\Stirling}[2]{\genfrac{\{}{\}}{0pt}{}{#1}{#2}}
\newcommand{\M}{\mathfrak M}
% Parts III-V (batch 101, Research Reports 231, 233, 85 and 91) macros, as
% delivered; Report 85's \tr (upright Tr) is typeset with \Tr
\newcommand{\one}{\mathbf1}
\newcommand{\ip}[2]{\langle #1,#2\rangle}
\DeclareMathOperator{\tr}{tr}
\DeclareMathOperator{\Pf}{Pf}
\DeclareMathOperator{\diver}{div}
\newcommand{\dd}{\,d}
\newcommand{\coef}[2]{[#1^{#2}]}
\newcommand{\C}{\mathbb C}
\newcommand{\D}{\Delta}
\newcommand{\diag}{\operatorname{diag}}
\newcommand{\Cat}{\operatorname{Cat}}
\newcommand{\del}{\delta}
\newcommand{\tht}{\vartheta}
\newtcolorbox{resultbox}{colback=pale,colframe=inkblue,boxrule=0.6pt,
 arc=1mm,left=9pt,right=9pt,top=7pt,bottom=7pt,breakable}
% [write] Dated editorial notes added when the report was written into
% ProveIt (batch 85); the delivered text is otherwise unchanged apart from
% the label prefix shr: and the page anchors around the title page.
% Part II (batch 98, manuscript 02) was added on 5 October 2026 with the
% label prefix shr:bp:; its editorial changes are listed at its head.
% Parts III-V (batch 101) were added on 5 October 2026 with the label
% prefixes shr:h4:, shr:h5:, shr:sa: and shr:sd:; their editorial changes
% are listed at their heads.
\newenvironment{writenote}[1][]{\par\smallskip\begingroup\small
 \noindent\textbf{[write]}\if\relax\detokenize{#1}\relax\else~\textbf{#1.}\fi\ }%
 {\par\endgroup\smallskip}

\begin{document}
\hypersetup{pageanchor=false}
\begin{titlepage}
\vspace*{0.45in}
{\color{inkblue}\large\textsc{Enumerative combinatorics \& asymptotic analysis}\par}
\vspace{0.3in}
{\color{inkblue}\Huge\bfseries Fixed-height\par shifted rectangles\par}
\vspace{0.25in}
{\Large An unconditional proof of the A181199 asymptotic,\par
all-order expansions, and an algebraicity dichotomy\par}
\vspace{0.35in}
{\large Prepared for Vladimir Reshetnikov\par}
{\large October 3, 2026\par}
\vspace{0.35in}
\begin{resultbox}
For every fixed number of rows $m$, let $T_m(n)$ count the $m\times n$
arrays increasing in rows, columns, diagonals, and downward antidiagonals.
The principal result is
\[
T_m(n)\sim
\frac{\sqrt m\,\prod_{j=0}^{m-1}j!}
 {2^{m(m-1)}(2\pi)^{(m-1)/2}}
\frac{m^{mn}}{n^{(m^2-1)/2}}.
\]
The proof supplies every inverse-power correction, proves the asymptotic
conjecture in OEIS A181199, and implies that
$\sum_{n\ge0}T_m(n)z^n$ is algebraic if and only if $m\le2$.
\end{resultbox}
\vspace{0.18in}
\noindent\textbf{Abstract.}
We split the order-polytope model of a shifted rectangle at a fixed numerical
threshold. Except for an explicitly bounded exponentially small boundary
contribution, its enumeration is an expectation under independent binomial
random variables, weighted by a squared Vandermonde and a rational pair
interaction. The shifted hook formula identifies the weight exactly.
Standardized binomial moments are polynomials in $1/n$, so the expectation
has an effective all-order expansion with rational coefficients. Its leading
Gaussian integral gives the constant above; the first two corrections are
proved as explicit polynomials in $m$. We obtain the conjectured constant
$9\sqrt5/(2^{17}\pi^2)$ for A181199 and five correction terms, a universal
limiting ratio to ordinary rectangular tableaux, nonalgebraicity at every
height at least three, and Gaussian-ensemble laws for random intermediate
row lengths. Complete proofs, an exact-arithmetic coefficient engine,
independent enumerators, and a source/status audit are included.

\vfill
\noindent\small\textbf{Status.} AI-assisted, unrefereed research manuscript.
The theorems have proofs below, but have not been independently peer reviewed
or proof-assistant verified. Historical priority is not certified. No
conjectured recurrence is used as a premise, and none is claimed proved here.
\end{titlepage}
\hypersetup{pageanchor=true}

\begin{writenote}[Added 5 October 2026, batch~98: two Parts]
Since 5 October 2026 this report has two Parts. Part~\ref{shr:part:one}
(Sections~1--13 and Appendices~A--B) is the batch-85 manuscript described
in the first note of Section~1.3, printed as before: its sections,
equations, theorems, research questions and labels keep their numbers, and
that note's ``a single manuscript, printed in full'' now describes
Part~\ref{shr:part:one}. Part~\ref{shr:bp:part} (Sections~14--25 and
Appendices~C--D) is manuscript~02 of batch~98, \emph{Rational diagonals and
rare boundary phases in shifted rectangles} (4~October 2026, pinned at
\texttt{db20eb379}), printed in full with the label prefix
\texttt{shr:bp:}. It proves that $F_m$ is the diagonal of a rational power
series, hence D-finite, for every fixed $m$
(Theorem~\ref{shr:bp:thm:diagonal-main}), which answers the first half of
Research question~7; the second half, the growth of the minimal order,
stays open. It also proves a sharp all-order expansion for the number
$R_m(n)$ of arrays whose first row is completed before the last row starts
(Theorem~\ref{shr:bp:thm:rare-main}). That count is a small part of the
boundary error of Theorem~\ref{shr:thm:cut}, not its asymptotic, so
Research question~4 is not answered by it (Remark~\ref{shr:bp:rem:q4});
the leading term of the boundary error is the write's
Proposition~\ref{shr:bp:prop:boundary}. Dated notes at the end of
Section~4.1, after Remark~8.1 and at Research questions~4 and~7 say so.
\end{writenote}

\begin{writenote}[Added 5 October 2026, batch~101: five Parts]
Since 5 October 2026 this report has five Parts; Parts~\ref{shr:part:one}
and~\ref{shr:bp:part} keep every number and label, and change only by
dated notes. Part~\ref{shr:h4:part} (Sections~26--33) is Research
Report~231 of the session bundle placed in batch~101: it proves the
order-two A181198 recurrence and Conjecture~18 of Kauers and Koutschan
\cite{KK} for every $n\ge1$ (Theorems~\ref{shr:h4:thm:main},
\ref{shr:h4:thm:oeis} and~\ref{shr:h4:thm:sum}). Part~\ref{shr:h5:part}
(Sections~34--42 and Appendix~\ref{shr:h5:app:literal}) is Research
Report~233: it proves Conjecture~19 and the order-three, degree-24
A181199 recurrence (Theorems~\ref{shr:h5:thm:sum}
and~\ref{shr:h5:thm:operator}). These are the recurrences that
Parts~\ref{shr:part:one} and~\ref{shr:bp:part}, and the OEIS entries, call
conjectural. Part~\ref{shr:sa:part} (Sections~43--45) prints Research
Reports~85 and~91, dated 1 and 2~October 2026, before the manuscripts of
Parts~\ref{shr:part:one} and~\ref{shr:bp:part}: second routes, earlier by
date line, to Theorems~\ref{shr:thm:main} and~\ref{shr:thm:algebraic} and
to Theorem~\ref{shr:bp:thm:diagonal-main}. The write adds
Proposition~\ref{shr:h4:prop:minimal} (the A181198 operator has minimal
order) and Proposition~\ref{shr:h5:prop:nothyper} (A181199 obeys no
first-order recurrence). Dated notes at the end of Section~1.3, after
Remark~8.1, at Research questions~1, 6, 7, 8 and~11, after the notes of
Sections~11.3 and~23.3, after Remark~\ref{shr:bp:rem:constructive} and in
Section~\ref{shr:bp:sec:further} say so. The intake's independent check
of the write (5~October 2026, end of Section~\ref{shr:sa:sec:further})
added Propositions~\ref{shr:h4:prop:higher} and~\ref{shr:h5:prop:higher}:
recurrences of lower degree and higher order for both sequences.
\end{writenote}

{\fontsize{9.5}{10.5}\selectfont\setlength{\parskip}{0pt}\tableofcontents}
\newpage

\part{An unconditional proof of the A181199 asymptotic, all-order expansions, and an algebraicity dichotomy}
\label{shr:part:one}

\section{The problem, the open target, and the source audit}

For positive integers $m,n$, consider arrays $(a_{i,j})$ containing each
integer in $\{1,\ldots,mn\}$ exactly once, subject to
\begin{align}
 a_{i,j}&<a_{i,j+1},&a_{i,j}&<a_{i+1,j},\\
 a_{i,j}&<a_{i+1,j+1},&a_{i,j+1}&<a_{i+1,j},
 \label{shr:eq:arrayconditions}
\end{align}
whenever the indicated cells exist. Write $T_m(n)$ for their number, and
set $T_m(0)=1$. The diagonal inequality is redundant once row and column
increase are imposed. The downward-antidiagonal inequality is not.

These numbers form the array A181196. Its fixed-height rows include
A181197 for $m=3$, A181198 for $m=4$, and A181199 for $m=5$
\cite{oeis196,oeis197,oeis198,oeis199}. Throughout this article, \emph{$m$ is
fixed and the row length $n$ tends to infinity}. Interchanging these two
parameters changes the problem: already $T_2(n)=C_{n-1}$ for $n\ge1$,
whereas $T_m(2)=1$ for every $m$.

\subsection{The conjecture resolved here}
The A181199 entry inspected on October 3, 2026, labels the following formula
as a conjecture of Vaclav Kotesovec, dated February 27, 2023:
\begin{equation}
 T_5(n)\sim \frac{9\,5^{5n+1/2}}{2^{17}\pi^2n^{12}}.
 \label{shr:eq:oeisconj}
\end{equation}
The entry says it was based on Christoph Koutschan's conjectured recurrence.
Kauers and Koutschan discuss the relevant guessed recurrences and finite-sum
conjectures for heights four and five in their 2023 article
\cite{KK,oeis199}. We do not assume those recurrences. We prove
\eqref{shr:eq:oeisconj} directly from the combinatorial definition and extend it
to every fixed height.

\subsection{What is established background, and what is contributed}
The shifted hook product and the order-statistics interpretation are
classical tools, not new inventions of this report. Sun's treatment of
shifted strips gives both tools and proves enumeration results in several
small cases \cite{Sun}. The known height-three leading asymptotic in
A181197 is recovered, not newly conjectured. Chan's transfer-matrix theory
concerns repeating strips with a fixed number of cells in each row; in the
notation here that is the regime with $n$ fixed and $m$ growing
\cite{Chan}. Its constant-coefficient recurrences do not supply a
fixed-height asymptotic in $n$.

The contributions developed in this manuscript are the exact
random-threshold reduction with a boundary bound, the resulting all-height
all-order theorem, universal first and second corrections, an algebraicity
classification, and intermediate-shape limit laws. These statements are
proved below. The source review did not locate these statements in the
works inspected, but that observation is not an exhaustive priority search.
In particular, a formula's continued ``conjecture'' label in OEIS is not by
itself proof that no published proof exists.

\subsection{Relationship to ProveIt}
The repository was inspected at revision
\texttt{6bf7f30d0352f7596e70928b3d4f304914075907}. Identifier searches for
A181198 and A181199 returned no matching repository report. This is a
limited search result, not a claim to have audited every repository file.
The existing path-forest reports distinguish unconditional expansions,
conjectured recurrences, and inverse-index approximations
\cite{ProveIt}. We retain that separation. No theorem in that report is a
mathematical dependency of the present proof, and the standard inverse-index
method in Section~\ref{shr:sec:inverse} is not presented as new.

\begin{resultbox}
\textbf{Scope of the resolution.} Equation~\eqref{shr:eq:oeisconj} is proved
here. The order-two recurrence for A181198 and the order-three recurrence
for A181199 are \emph{not} proved here. The asymptotic series is an all-order
Poincar\'e expansion, with a fixed truncation order; no Borel summability,
Stokes data, or complete exponentially improved transseries is asserted.
\end{resultbox}

\begin{writenote}[Added 3 October 2026, batch~85]
This report is a single manuscript, printed in full: manuscript~09 of
batch~85 of ProveIt's incoming-reports intake, delivered as the archive
\texttt{Shifted\_\allowbreak Rectangle\_\allowbreak Asymptotics.zip}
(343{,}338 bytes, wrapper directory
\texttt{Shifted\_\allowbreak Rectangle\_\allowbreak Asymptotics/}) in commit
\texttt{9d6968c8a}, and placed in commit \texttt{ddf8df5d5} in the
research-report collection, among the OEIS-sequence asymptotics reports, as
the report \texttt{a181199-\allowbreak shifted-\allowbreak rectangles}. The
title page reads ``Prepared for Vladimir Reshetnikov''; the PDF metadata
names ``Research report prepared for Vladimir Reshetnikov''. No human author
and no tool are named. The delivered text is printed unchanged except that
every label carries the prefix \texttt{shr:}, the title page no longer
creates a duplicate page anchor, and dated notes such as this one are added.

\emph{Pin and repository claims.} The revision quoted above,
\texttt{6bf7f30d0} (3 October 2026), is the ProveIt commit the manuscript
was written against. At placement, and again at this writing, the intake
searched the whole tracked tree for A181196--A181199, shifted rectangles,
shifted tableaux, shifted hook products and Sun's article; apart from the
arrival archive, the only hits are the batch-85 sibling named below, which
points to this report. This report continues no repository material, and
the sentence ``Identifier searches for A181198 and A181199 returned no matching
repository report'' remains true. ``The existing path-forest reports'' is
one report,
\texttt{oeis-sequence-asymptotics/\allowbreak a189281-path-forest-expansions},
in three Parts; as the manuscript says, none of its theorems is used here.
No statement of this report is formalized in Lean or Rocq, no formal
development in ProveIt treats shifted tableaux, and the report's place in
the collection confers no formal status.

\emph{Neighbours.} The companion reports of the same batch are
\texttt{a047874-\allowbreak long-\allowbreak increasing-\allowbreak subsequences}
(manuscript~08), whose fixed-sector problem also uses the hook count of the
rectangle with $m$ rows of length $n$ as a normalizer, and
\texttt{a182220-\allowbreak source-\allowbreak boundary} (manuscript~07,
unrelated mathematics). Other conjectures of the 2023 article of Kauers
and Koutschan \cite{KK} are proved in the collection reports
\texttt{enumerative-\allowbreak combinatorics/\allowbreak a181280-\allowbreak binary-\allowbreak matrix-\allowbreak formula}
(its Conjecture~20, Section~6.5, a C-finite formula),
\texttt{enumerative-\allowbreak combinatorics/\allowbreak a195806-\allowbreak hexagonal-\allowbreak lattice}
(Conjecture~11) and Part~IV of
\texttt{a215561-\allowbreak fixed-\allowbreak composition-\allowbreak excursions}
(Conjecture~15). The present report concerns that article's Section~6.4,
and leaves the guessed recurrences of that section unproved. Two notes below relate Corollary~2.4 and Section~\ref{shr:sec:inverse}
to the transseries volumes in
\texttt{Analysis/\allowbreak Transseries/\allowbreak docs/\allowbreak series-\allowbreak and-\allowbreak transseries/}.

\emph{Notation.} The manuscript reuses several letters with local meanings.
They are kept as delivered; the table fixes each reading. No symbol was
renamed.

\smallskip
\begin{center}\footnotesize
\begin{tabular}{@{}p{0.2\textwidth}p{0.74\textwidth}@{}}
\multicolumn{2}{@{}l}{\textbf{[write]} Notation table of the note in
Section~1.3 (added 3 October 2026)}\\
\toprule
Symbol & Meanings in this report\\
\midrule
$T_m(n)$ & the shifted-rectangle count (A181196 family), \emph{not} the
number $f^{(n^m)}$ of standard Young tableaux of the $m\times n$ rectangle;
the transseries volume writes $T_d(n)$ for the latter, with $d$ rows.\\
$d$ & $\binom m2$, the number of row pairs (\eqref{shr:eq:normalizations});
in the transseries volume $d$ is the number of rows, $m$ here.\\
$K_m$; $K$; $K_j(X)$ & the leading constant \eqref{shr:eq:Km}; $K=K_m$ in
Section~\ref{shr:sec:inverse}, but the number of samples below $p$ in the proof
of Theorem~\ref{shr:thm:tracezero}; Krawtchouk polynomials in Section~6.1.\\
$L$; $L_5(n)$ & a truncation order (Theorem~\ref{shr:thm:main}); $\log(y/K)$ in
Section~\ref{shr:sec:inverse}; the leading term in Section~7.3.\\
$h_r(q)$; $h_j$ & complete homogeneous polynomials \eqref{shr:eq:hdef};
squared norms $j!\,\fall n j/n^j$ in Section~6.1.\\
$\lambda$ & a state or strict partition (Section~3, Appendix~A); $m\log m$ in
Section~\ref{shr:sec:inverse}.\\
$R_n$; $R_m$; $R(z)$; $\mathcal R_{L,n}$ & the pair weight
\eqref{shr:eq:Rweight}; the radius $m^{-m}$ in Section~\ref{shr:sec:algebraicity};
the rational function of Appendix~A; the Taylor remainder
\eqref{shr:eq:taylorbound}.\\
$p$; $p_r$ & the threshold in $(0,1)$; the power sums $\sum_iY_i^r$ (Section~6).\\
$q_{ij}$; $q$ & $(y_i+y_j)^2$; a derivative order in
Section~\ref{shr:sec:algebraicity}.\\
$Z_m(n)$; $Z_i$ & the binomial Vandermonde normalizer \eqref{shr:eq:binomZ};
standardized counts in Section~\ref{shr:sec:GUE}.\\
$A$, $B$; $\mathcal A$ & $\binom m3$ and $m(m-1)^2(m-2)/12$ (Section~6.1);
the event that all comparisons hold.\\
$C(z)$; $C_{n-1}$; $C_{m,L}$; $\mathcal C_n$ & the Catalan generating
function; Catalan numbers; a constant in \eqref{shr:eq:taylorbound}; the
central event.\\
$E_{m,n}(p)$; $\E$; $\mathcal E_z$ & the boundary error \eqref{shr:eq:exactcut};
expectation; the moment functional \eqref{shr:eq:Ez}.\\
$H(\lambda)$; $H_j$; $H_{20}$ & the hook expression of Appendix~A; Hermite
polynomials; the constant \eqref{shr:eq:H20}.\\
$\beta_j$ & recurrence coefficients (Section~6.2); not the exponent
$\beta=-\alpha_m$ of the transseries volume.\\
\bottomrule
\end{tabular}
\end{center}
\end{writenote}

\begin{writenote}[Added 5 October 2026, batch~101: the recurrences]
The note above says that this report ``leaves the guessed recurrences of
that section unproved'', and the box ``Scope of the resolution'' says the
same of the manuscript. Both remain true of Part~\ref{shr:part:one}, but no
longer of the report: Part~\ref{shr:h4:part} proves the A181198 recurrence
and Conjecture~18 of \cite{KK} (Theorems~\ref{shr:h4:thm:oeis}
and~\ref{shr:h4:thm:sum}), and Part~\ref{shr:h5:part} the A181199
recurrence and Conjecture~19 (Theorems~\ref{shr:h5:thm:sum}
and~\ref{shr:h5:thm:operator}), each from the array definition. Neither
proof uses the asymptotic theorems of Part~\ref{shr:part:one}; both credit
them.
\end{writenote}

\section{Main results}
Set
\begin{equation}
 d=\binom m2,\qquad J_m=\prod_{j=0}^{m-1}j!,\qquad
 M_m(n)=\frac{(mn)!}{(n!)^m},\qquad \alpha_m=\frac{m^2-1}{2},
 \label{shr:eq:normalizations}
\end{equation}
and define
\begin{equation}
 K_m=\frac{\sqrt m\,J_m}{4^d(2\pi)^{(m-1)/2}}.
 \label{shr:eq:Km}
\end{equation}
The multinomial coefficient $M_m(n)$ counts all row-increasing arrays,
before imposing any comparisons between different rows.

\begin{theorem}[All orders at fixed height]\label{shr:thm:main}
For every fixed integer $m\ge1$ there is a uniquely determined sequence
$c_{m,j}\in\Q$, with $c_{m,0}=1$, such that, for every fixed $L\ge0$,
\begin{equation}
 T_m(n)=M_m(n)\frac{J_m}{(4n)^d}
 \left(\sum_{j=0}^{L}\frac{c_{m,j}}{n^j}
       +O_{m,L}(n^{-L-1})\right).
 \label{shr:eq:mainM}
\end{equation}
Every coefficient is given by the finite moment formula
\eqref{shr:eq:coefficientformula}. In particular,
\begin{equation}
 c_{m,1}=\frac{m(m^2-1)}{12},\qquad
 c_{m,2}=\frac{m^2(m^2-1)(m^2+2)}{288}.
 \label{shr:eq:c12}
\end{equation}
Equivalently, there are rational $b_{m,j}$, $b_{m,0}=1$, such that
\begin{equation}
 T_m(n)=K_m m^{mn}n^{-\alpha_m}
 \left(\sum_{j=0}^{L}\frac{b_{m,j}}{n^j}
       +O_{m,L}(n^{-L-1})\right),
 \label{shr:eq:mainK}
\end{equation}
where
\begin{align}
 b_{m,1}&=\frac{(m^2-1)^2}{12m},\\
 b_{m,2}&=\frac{(m^2-1)(m^6+3m^2-1)}{288m^2}.
 \label{shr:eq:b12}
\end{align}
\end{theorem}

\begin{corollary}[A181199, with five corrections]\label{shr:cor:A181199}
With $a(n)=T_5(n)$,
\begin{align}
 a(n)=\frac{9\,5^{5n+1/2}}{2^{17}\pi^2n^{12}}
 \bigg(&1+\frac{48}{5n}+\frac{5233}{100n^2}
       +\frac{1028391}{5000n^3}\notag\\
       &+\frac{60248377}{100000n^4}
       +\frac{273939939}{250000n^5}+O(n^{-6})\bigg).
 \label{shr:eq:A181199exp}
\end{align}
\end{corollary}

\begin{theorem}[Algebraicity dichotomy]\label{shr:thm:algebraic}
The ordinary generating function
\[
 F_m(z)=\sum_{n\ge0}T_m(n)z^n
\]
is algebraic over $\Q(z)$ if and only if $m=1$ or $m=2$.
The same classification holds for algebraicity over $\mathbb C(z)$.
\end{theorem}

\begin{corollary}[A universal comparison with ordinary rectangles]
Let $f^{(n^m)}$ count ordinary standard Young tableaux of rectangular shape
with $m$ rows of length $n$. For fixed $m$,
\begin{align}
\frac{T_m(n)}{f^{(n^m)}}=2^{-m(m-1)}\bigg(&1+
 \frac{m(m^2-1)}{4n}\notag\\
 &+\frac{m^2(m^2-1)(m^2-2)}{32n^2}+O_m(n^{-3})\bigg).
 \label{shr:eq:rectangle-ratio}
\end{align}
Thus the extra antidiagonal comparisons retain a positive limiting fraction
$2^{-m(m-1)}$ of ordinary rectangular tableaux.
\end{corollary}

The constant is not obtained by treating the comparisons as independent.
It comes from two shifted-hook denominators, one for each side of a
threshold cut. Their product is the essential feature of the proof.
The distributional counterpart is given in Section~\ref{shr:sec:GUE}.

\begin{writenote}[Added 3 October 2026, batch~85: the comparison constant]
The normalizer of Corollary~2.4 is the classical hook count of the
rectangle, recorded in ProveIt as equation~(3.19), label
\texttt{t2:eq:tableaux}, of the volume \emph{Combinatorial Transseries and
Their Inverses}
(\texttt{Analysis/\allowbreak Transseries/\allowbreak docs/\allowbreak series-\allowbreak and-\allowbreak transseries/\allowbreak Combinatorial\_\allowbreak Transseries\_\allowbreak Inverses/},
Section~3.5), which writes it $T_d(n)$ with $d$ rows; the hook product
displayed at the end of Section~6.4 here is the same formula with $d=m$.
That volume's constants (3.20) (\texttt{t2:eq:tableaux-constants}) give
$f^{(n^m)}\sim C\,m^{mn}n^{\beta}$ with $a=m\log m$,
$\beta=-(m^2-1)/2=-\alpha_m$ and
$C=\sqrt m\,(2\pi)^{(1-m)/2}\prod_{j=0}^{m-1}j!$. Comparing with
\eqref{shr:eq:Km}, $K_m=2^{-m(m-1)}C$ exactly, since $4^d=2^{m(m-1)}$: the
two families share the rate $m^m$ and the exponent $-\alpha_m$, and the
leading term of \eqref{shr:eq:rectangle-ratio} is the quotient of the two
constants. The intake checked this identity symbolically for
$m=2,\ldots,7$, and checked that \eqref{shr:eq:c12} and the rising-factorial
expansion of $\prod_{j<m}(n+j)!/n!$ give the two printed corrections of
\eqref{shr:eq:rectangle-ratio} for the same $m$. Neither the formula nor the
constant of the ordinary rectangle is new here; the new statement is the
shifted count and its ratio.
\end{writenote}

\section{Prefix states and the shifted hook product}

\subsection{The exact state space}
Fill an array by inserting $1,2,\ldots,mn$ in order. At any intermediate
stage let $x_i$ be the number of filled cells in row $i$. Row increase means
that the filled cells in each row are a prefix. The legal states satisfy
\begin{equation}
 n\ge x_1\ge x_2\ge\cdots\ge x_m\ge0,
 \qquad x_i=x_{i+1}\Longrightarrow x_i\in\{0,n\}.
 \label{shr:eq:states}
\end{equation}
Indeed, placing a filled cell in row $i+1$, column $j<n$, forces the cell
in row $i$, column $j+1$, to have already been filled. When the lower row
has length $n$, column increase forces the upper row to be full as well.
Conversely, these conditions make the filled set an order ideal of the
array poset. They therefore characterize the prefix states exactly.

A transition increases $x_i$ by one. For $i>1$ it is allowed precisely when
\begin{equation}
 x_{i-1}\ge x_i+2
 \quad\hbox{or}\quad (x_i=n-1\hbox{ and }x_{i-1}=n),
 \label{shr:eq:transition}
\end{equation}
with $x_i<n$; the first row has only the latter upper bound. Empty rows
can tie, and completed rows can tie. Forgetting either exception would
change the sequence.

For a legal state $\lambda=(\lambda_1,\ldots,\lambda_m)$ define its rotated
complement by
\begin{equation}
 \lambda^c=(n-\lambda_m,\ldots,n-\lambda_1).
 \label{shr:eq:complement}
\end{equation}
Rotation through $180$ degrees, together with reversal of the numerical
labels, preserves the comparisons. Thus suffix enumeration at $\lambda$
is prefix enumeration at $\lambda^c$.

\subsection{Interior states}
Call a state \emph{interior} when
\begin{equation}
 n>\lambda_1>\cdots>\lambda_m>0.
 \label{shr:eq:interior}
\end{equation}
Every path ending there has avoided the completed-row boundary. After
removing trailing zero parts, it is a chain of strict partitions obtained
by adding one box at a time. Such chains are standard shifted tableaux.
If $g^\lambda$ denotes their number and $k=|\lambda|$, the shifted hook
product is
\begin{equation}
 g^\lambda=\frac{k!}{\prod_i\lambda_i!}
 \prod_{i<j}\frac{\lambda_i-\lambda_j}{\lambda_i+\lambda_j}.
 \label{shr:eq:shiftedhook}
\end{equation}
This standard formula appears, for example, in Sun's account \cite{Sun}.
Appendix~\ref{shr:app:hook} proves it directly from the strict-partition
recursion; no shifted-tableau enumeration result is needed without proof.

For an interior state both $\lambda$ and $\lambda^c$ are strict positive
partitions. Consequently, the prefix and suffix counts are simultaneously
available in product form. This is why cutting \emph{in the middle}, rather
than tracking the first completed row, is effective for asymptotics.

\section{An exact random-threshold reduction}

\subsection{Order statistics and the cut}
Take $m$ independent rows, each consisting of $n$ independent uniform
random numbers in $(0,1)$, sorted increasingly. Let $\mathcal A$ be the event
that the comparisons between rows also hold. Every interleaving of these
sorted rows has the same probability, so
\begin{equation}
 \Pp(\mathcal A)=\frac{T_m(n)}{M_m(n)}.
 \label{shr:eq:probability}
\end{equation}
This is the order-statistics form of the order-polytope interpretation
\cite{Sun}; it also follows immediately by sorting all $mn$ samples at once.

Fix $p\in(0,1)$ and let $X_i$ count the entries in row $i$ below $p$.
Before conditioning on $\mathcal A$, the variables $X_i$ are independent
$\operatorname{Bin}(n,p)$. Conditional on these counts, the samples below
$p$ and above $p$ are independent sorted uniforms in their respective
intervals. When the count vector is an interior state $\lambda$, the
conditional probability of satisfying all comparisons below the cut is
\[
 \frac{g^\lambda}{k!/\prod_i\lambda_i!}.
\]
The analogous suffix probability uses $\lambda^c$. All comparisons across
the cut are automatic because the prefix is an order ideal. Multiplying
and using \eqref{shr:eq:shiftedhook} gives
\begin{equation}
 \Pp(\mathcal A\mid X=\lambda)=R_n(\lambda),\qquad
 R_n(x)=\prod_{i<j}
 \frac{(x_i-x_j)^2}{(x_i+x_j)(2n-x_i-x_j)}.
 \label{shr:eq:Rweight}
\end{equation}
The weight is symmetric in the coordinates. Equal interior coordinates
give zero, and all distinct coordinates have exactly one decreasing order.

\begin{theorem}[Interior sum and explicit boundary error]\label{shr:thm:cut}
Let $X_1,\ldots,X_m$ be independent $\operatorname{Bin}(n,p)$ variables and
$\mathcal I=\{1\le X_i\le n-1\text{ for all }i\}$. Then
\begin{equation}
 \frac{T_m(n)}{M_m(n)}=
 \frac1{m!}\E\left[\1_{\mathcal I}R_n(X)\right]+E_{m,n}(p),
 \label{shr:eq:exactcut}
\end{equation}
where the error is nonnegative and satisfies
\begin{equation}
 0\le E_{m,n}(p)\le m\bigl(p^n+(1-p)^n\bigr).
 \label{shr:eq:boundary}
\end{equation}
For $p=1/2$ the expectation is the finite rational sum
\begin{equation}
 I_m(n)=2^{-mn}
 \sum_{n>x_1>\cdots>x_m>0}
 \left(\prod_{i=1}^m\binom n{x_i}\right)
 \prod_{i<j}\frac{(x_i-x_j)^2}{(x_i+x_j)(2n-x_i-x_j)},
 \label{shr:eq:interiorsum}
\end{equation}
and
\begin{equation}
 0\le T_m(n)-M_m(n)I_m(n)\le 2m\,2^{-n}M_m(n).
 \label{shr:eq:exactbound}
\end{equation}
\end{theorem}

\begin{proof}
The contribution to $\Pp(\mathcal A)$ from interior count vectors is the
sum of their binomial probabilities times \eqref{shr:eq:Rweight}. Symmetry
converts this decreasing-vector sum into $1/m!$ times the unrestricted
expectation on $\mathcal I$. The remaining event is
$\mathcal A\cap\mathcal I^c$. It has nonnegative probability, at most the
probability that some row has all its samples on one side of the cut.
A union bound gives \eqref{shr:eq:boundary}. The specialization is immediate.
\end{proof}

\begin{remark}[A useful pointwise bound]\label{shr:rem:Rbound}
For $0\le x,y\le n$,
$|x-y|\le x+y$ and $|x-y|\le2n-x-y$. Thus every factor of
$R_n$ lies in $[0,1]$ whenever its denominator is nonzero. In particular,
$0\le R_n\le1$ on $\mathcal I$. No uncontrolled singularity is hidden
near the endpoints in \eqref{shr:eq:exactcut}.
\end{remark}

The boundary estimate is an exact, nonasymptotic statement. It is deliberately
simple and is not claimed sharp. Since the main probability is of order
$n^{-d}$, this exponentially small additive error is negligible relative
to every term of the inverse-power expansion.

\begin{writenote}[Added 5 October 2026, batch~98]
The rate $2^{-n}$ of the bound at $p=1/2$ is the true one, but the bound
is not sharp in its constant, nor, for $m\ge3$, in its power of $n$: the
write's Proposition~\ref{shr:bp:prop:boundary} in Part~\ref{shr:bp:part}
shows $E_{m,n}(1/2)\sim2^{1-n}3^{1-m}J_{m-1}(4n)^{-\binom{m-1}2}$, against
the bound $2m\,2^{-n}$.
\end{writenote}

\section{Gaussian reduction and an effective all-order theorem}

\subsection{The scaled integrand}
Use $p=1/2$ and put
\begin{equation}
 Y_i=\frac{2X_i-n}{\sqrt n},\qquad
 \Delta(y)=\prod_{i<j}(y_i-y_j),\qquad
 q_{ij}(y)=(y_i+y_j)^2.
\end{equation}
The pair weight becomes exactly
\begin{equation}
 R_n(X)=(4n)^{-d}\Delta(Y)^2
 \prod_{i<j}\left(1-\frac{q_{ij}(Y)}{4n}\right)^{-1}
 \quad\hbox{on }\mathcal I.
 \label{shr:eq:scaledweight}
\end{equation}
Let $h_r(q)$ be the complete homogeneous polynomial of degree $r$ in the
$d$ variables $q_{ij}$, defined by
\begin{equation}
 \prod_{i<j}(1-uq_{ij})^{-1}=\sum_{r\ge0}h_r(q)u^r.
 \label{shr:eq:hdef}
\end{equation}
In particular $h_0=1$. For $m=1$, the empty product is one and $h_r=0$
for $r>0$.

\subsection{Polynomial standardized moments}
A variable $Y_i$ has the same law as
$n^{-1/2}(\eps_1+\cdots+\eps_n)$, where the $\eps_j$ are independent
symmetric signs. Its moment generating function is
\begin{equation}
 \E e^{tY_i}=
 \left(\cosh(t/\sqrt n)\right)^n.
 \label{shr:eq:mgf}
\end{equation}
For every fixed nonnegative integer $k$, its $k$th moment is exactly a
polynomial $\mu_k(z)$ at $z=1/n$. One convenient formal definition is
\begin{equation}
 \sum_{k\ge0}\mu_k(z)\frac{t^k}{k!}
 =\exp\left(\frac1z\log\cosh(t\sqrt z)\right).
 \label{shr:eq:moments}
\end{equation}
The apparent singularity at $z=0$ is removable in the formal expansion.
Odd moments vanish, and
\begin{align}
 \mu_0&=1,&\mu_2&=1,&\mu_4&=3-2z,\\
 \mu_6&=15-30z+16z^2,&
 \mu_8&=105-420z+588z^2-272z^3.
 \label{shr:eq:momentexamples}
\end{align}
For $k=2s>0$, the degree is at most $s-1$. This follows either by grouping
terms according to equal sign indices, or from the moment--cumulant formula:
a partition into even blocks contributes
$z^{s-\#\text{blocks}}$ times a product of sign cumulants.

Define a linear functional $\mathcal E_z:\Q[y_1,\ldots,y_m]\to\Q[z]$ by
\begin{equation}
 \mathcal E_z(y_1^{k_1}\cdots y_m^{k_m})=
 \prod_{i=1}^m\mu_{k_i}(z).
 \label{shr:eq:Ez}
\end{equation}
Then $\mathcal E_{1/n}$ is the exact expectation for the independent
variables $Y_i$. At $z=0$, it is expectation under independent standard
normal variables.

\subsection{The coefficient formula}
\begin{proposition}[Finite formula for each coefficient]\label{shr:prop:formula}
For $j\ge0$, the coefficient in \eqref{shr:eq:mainM} is
\begin{equation}
 \boxed{\displaystyle
 c_{m,j}=\frac1{m!J_m}
 \sum_{r=0}^{j}\frac1{4^r}
 [z^{j-r}]\,\mathcal E_z\!\left(\Delta(y)^2h_r(q(y))\right).}
 \label{shr:eq:coefficientformula}
\end{equation}
The formula uses finite polynomials and rational arithmetic. It does not
require previously computed sequence terms or a guessed recurrence.
\end{proposition}

\subsection{Uniform error control}
We supply the estimates that justify the formal calculation.
Let
\[
 \mathcal C_n=\{|Y_i|\le\sqrt n/2\text{ for all }i\}.
\]
Since $\cosh t\le e^{t^2/2}$, Chernoff's inequality gives
\begin{equation}
 \Pp(\mathcal C_n^c)\le2m e^{-n/8}.
 \label{shr:eq:central-tail}
\end{equation}
The same moment generating function bound implies that the absolute moments
of every fixed degree of the $Y_i$ are bounded uniformly in $n$.

On $\mathcal C_n$, each $q_{ij}(Y)/(4n)$ lies in $[0,1/4]$. Taylor's theorem
applied to the finite product in \eqref{shr:eq:scaledweight} therefore gives,
for every fixed $L$,
\begin{align}
 \prod_{i<j}\left(1-\frac{q_{ij}(Y)}{4n}\right)^{-1}
 &=\sum_{r=0}^{L}\frac{h_r(q(Y))}{(4n)^r}+\mathcal R_{L,n}(Y),\\
 |\mathcal R_{L,n}(Y)|
 &\le C_{m,L}n^{-L-1}(1+\|Y\|)^{2L+2}.
 \label{shr:eq:taylorbound}
\end{align}
To see the second bound, differentiate the product $L+1$ times with respect
to $1/n$. Every differentiated term has total numerator degree $2L+2$;
all denominators are bounded below by fixed powers of $3/4$.
After multiplication by $\Delta(Y)^2$, the expected remainder is still
$O_{m,L}(n^{-L-1})$ by the uniform moment bounds.

We may replace polynomial expectations on $\mathcal C_n$ by unrestricted
ones with exponentially small error. For example, Cauchy--Schwarz,
\eqref{shr:eq:central-tail}, and a uniform moment bound control the difference
for every fixed polynomial. For the original rational weight, use
$R_n\le1$ from Remark~\ref{shr:rem:Rbound}. After restoring the factor
$(4n)^d$, its discarded tail is only a fixed polynomial in $n$ times an
exponential. The endpoint error \eqref{shr:eq:boundary} behaves similarly.
We have consequently proved
\begin{equation}
 \frac{T_m(n)}{M_m(n)}=
 \frac{(4n)^{-d}}{m!}
 \left(\sum_{r=0}^{L}\frac{
 \mathcal E_{1/n}(\Delta^2h_r)}{(4n)^r}
 +O_{m,L}(n^{-L-1})\right).
 \label{shr:eq:beforecollect}
\end{equation}
This establishes \eqref{shr:eq:coefficientformula} after collecting powers.
It also proves that parity effects cannot occur in the algebraic sector:
the moments are exact polynomials in $1/n$, on both parity subsequences.

\subsection{The leading Gaussian integral}
For independent standard normals $G_1,\ldots,G_m$,
\begin{equation}
 \E\Delta(G)^2=m!\prod_{j=0}^{m-1}j!=m!J_m.
 \label{shr:eq:gaussZ}
\end{equation}
Here is an elementary verification. Replace the monomials in the
Vandermonde determinant by the monic Hermite polynomials $H_j$, whose
squared norms for standard normal measure are $j!$. The determinant is
unchanged because this basis change is unit triangular. Expanding the
square and using orthogonality leaves the $m!$ matching permutations,
each contributing $\prod j!$.

Thus $c_{m,0}=1$. Finally, Stirling's expansion gives
\begin{equation}
 M_m(n)\sim \sqrt m\,(2\pi n)^{(1-m)/2}m^{mn},
\end{equation}
which proves the leading constant and exponent in \eqref{shr:eq:mainK}.
More explicitly, the complete conversion between the two normalizations is
\begin{equation}
 \sum_{j\ge0}b_{m,j}z^j=
 \left(\sum_{j\ge0}c_{m,j}z^j\right)
 \exp\left(\sum_{r\ge1}
 \frac{B_{2r}}{2r(2r-1)}
 (m^{1-2r}-m)z^{2r-1}\right),
 \label{shr:eq:stirling-convert}
\end{equation}
where $B_{2r}$ are Bernoulli numbers. All uses of this identity are formal,
or truncated to a fixed finite order with Stirling's remainder estimate
\cite{DLMFgamma}. This completes the proof of Theorem~\ref{shr:thm:main}, apart
from the explicit evaluations in the next section.

\section{Universal first and second corrections}

The finite formula is already an all-order algorithm. To obtain short
formulas valid at every height, it is better to keep the Vandermonde
weight intact rather than expand it into individual monomials.

\subsection{An exact binomial Vandermonde normalization}
Define the normalized binomial ensemble by weighting the independent $Y_i$
with $\Delta(Y)^2$. Its normalizing factor is
\begin{equation}
 Z_m(n)=\E\Delta(Y)^2
   =m!J_m\prod_{j=0}^{m-1}\frac{\fall n j}{n^j}.
 \label{shr:eq:binomZ}
\end{equation}
This identity is valid for $n\ge m-1$ (and interpreted as zero when the
support cannot accommodate $m$ distinct points).

For completeness, let $X\sim\operatorname{Bin}(n,1/2)$ and define $K_j(X)$
by
\[
 (1+t)^X(1-t)^{n-X}=\sum_{j=0}^n K_j(X)t^j.
\]
Multiplying two such generating functions and taking expectation gives
$(1+st)^n$. Hence the $K_j$ are orthogonal and
$\E K_j^2=\binom nj$. Their leading coefficient as polynomials in
$Y=(2X-n)/\sqrt n$ is $n^{j/2}/j!$. The corresponding monic polynomials
therefore have squared norms
\[
 h_j=j!\,\fall n j/n^j.
\]
The same determinant argument as in \eqref{shr:eq:gaussZ} proves
\eqref{shr:eq:binomZ}. These are the symmetric Krawtchouk polynomials; the
standard orthogonal-polynomial framework is recorded in \cite{DLMFpoly}.

Put
\begin{equation}
 A=\binom m3,\qquad B=\frac{m(m-1)^2(m-2)}{12}.
\end{equation}
Expanding the finite product in \eqref{shr:eq:binomZ} gives
\begin{equation}
 \frac{Z_m(n)}{m!J_m}=1-\frac A n+\frac{A^2-B}{2n^2}+O_m(n^{-3}).
 \label{shr:eq:Zexpand}
\end{equation}
Indeed, $A=\sum_{j<m}\sum_{k<j}k$ and
$B=\sum_{j<m}\sum_{k<j}k^2$.

\subsection{The first pair statistic}
Let $p_r=\sum_iY_i^r$. The first complete homogeneous polynomial is
\begin{equation}
 S_1=h_1(q)=\sum_{i<j}(Y_i+Y_j)^2=(m-2)p_2+p_1^2.
 \label{shr:eq:S1}
\end{equation}
Write $\langle\cdot\rangle_n$ for expectation in the normalized binomial
Vandermonde ensemble. The monic polynomial recurrence is
\[
 yP_j(y)=P_{j+1}(y)+\beta_jP_{j-1}(y),\qquad
 \beta_j=j\left(1-\frac{j-1}{n}\right).
\]
The diagonal coefficient is zero by symmetry. The norm quotient proves
the displayed value of $\beta_j$.

We need the standard finite-projection identities
\begin{align}
 \langle p_2\rangle_n&=\beta_m+2\sum_{j=1}^{m-1}\beta_j
 =m^2-\frac{m(m-1)(2m-1)}{3n},\\
 \langle p_1^2\rangle_n&=\beta_m
 =m-\frac{m(m-1)}n.
 \label{shr:eq:traceidentities}
\end{align}
They can be seen directly from the determinant ensemble: multiplication
by $y$ is tridiagonal in the orthonormal basis, with off-diagonal entries
$\sqrt{\beta_j}$. The mean of a linear statistic is its multiplication
operator's trace on degrees $0,\ldots,m-1$. Its variance is the squared
norm of the part crossing that projection. For $p_1$ only the edge from
degree $m-1$ to degree $m$ crosses, giving $\beta_m$.
These trace rules also follow by expanding
$\E\prod_i(1+t f(Y_i))$ as a determinant of the moment matrix and taking
the first two derivatives at $t=0$.

Consequently,
\begin{equation}
 \langle S_1\rangle_n=s_{10}+\frac{s_{11}}n,
 \quad
 s_{10}=m(m-1)^2,\quad
 s_{11}=-\frac{m(m-1)(2m^2-5m+5)}3.
 \label{shr:eq:s10s11}
\end{equation}

\subsection{The second pair statistic in the Gaussian limit}
Let $\langle\cdot\rangle_0$ denote the Gaussian Vandermonde ensemble.
Since
\[
 h_2(q)=\frac12\left(S_1^2+\sum_{i<j}(Y_i+Y_j)^4\right),
\]
and
\[
 \sum_{i<j}(Y_i+Y_j)^4=(m-8)p_4+4p_1p_3+3p_2^2,
\]
we need only a few Gaussian moments:
\begin{align}
 \langle p_2^2\rangle_0&=m^2(m^2+2),&
 \langle p_2p_1^2\rangle_0&=m(m^2+2),\\
 \langle p_1^4\rangle_0&=3m^2,&
 \langle p_4\rangle_0&=2m^3+m,&
 \langle p_1p_3\rangle_0&=3m^2.
 \label{shr:eq:gaussmoments}
\end{align}
There is no need to invoke random-matrix integration to establish these.
Scaling and translation of the Gaussian integral give
\begin{equation}
 \left\langle e^{t p_2+s p_1}\right\rangle_0
 =(1-2t)^{-m^2/2}
 \exp\left(\frac{ms^2}{2(1-2t)}\right).
\end{equation}
This proves the first three identities by differentiation. Translation
integration by parts, using $\Delta(y+c\mathbf1)=\Delta(y)$, gives
$\langle p_1p_3\rangle_0=3\langle p_2\rangle_0=3m^2$.
Finally, in the Hermite multiplication matrix, the degree-$j$ diagonal
entry of the fourth power is $6j^2+6j+3$; summing over $0\le j<m$ yields
$2m^3+m$. Thus
\begin{equation}
 H_{20}:=\langle h_2(q)\rangle_0
 =\frac{m(m-1)(m^4-3m^3+10m^2-18m+16)}2.
 \label{shr:eq:H20}
\end{equation}

\subsection{Collecting the terms}
Combining \eqref{shr:eq:Zexpand}, \eqref{shr:eq:s10s11}, and \eqref{shr:eq:H20},
\begin{align}
 c_{m,1}&=-A+\frac{s_{10}}4,\\
 c_{m,2}&=\frac{A^2-B}{2}+
           \frac{-As_{10}+s_{11}}4+\frac{H_{20}}{16}.
\end{align}
Straight polynomial simplification gives \eqref{shr:eq:c12}.
The first Stirling logarithm coefficient is
$u_1=(m^{-1}-m)/12$, so
$b_{m,1}=c_{m,1}+u_1$ and
$b_{m,2}=c_{m,2}+u_1c_{m,1}+u_1^2/2$. This proves \eqref{shr:eq:b12}.

The ordinary rectangular hook product is
\[
 f^{(n^m)}=M_m(n)J_m\prod_{j=0}^{m-1}\frac{n!}{(n+j)!}.
\]
Its finite rising-factorial expansion, combined with \eqref{shr:eq:c12},
proves \eqref{shr:eq:rectangle-ratio}.

\section{OEIS specializations and numerical checks}

\subsection{Low-height controls}
For $m=1$, $T_1(n)=1$. For $m=2$, the usual strict-ballot argument gives
$T_2(n)=C_{n-1}$ for $n\ge1$; equivalently,
\begin{equation}
 T_2(n)=\binom{2n}{n}\frac1{2(2n-1)}
 =\binom{2n}{n}\frac1{4n}
   \sum_{j\ge0}\frac1{(2n)^j}.
\end{equation}
Thus $c_{2,j}=2^{-j}$ at every order. This is a useful exact control of
both the normalization and the coefficient generator.

For height three, the theorem recovers the existing A181197 leading term
\cite{oeis197} and gives
\begin{equation}
 T_3(n)=\frac{\sqrt3\,27^n}{64\pi n^4}
 \left(1+\frac{16}{9n}+\frac{755}{324n^2}+O(n^{-3})\right).
\end{equation}
For height four, corresponding to A181198,
\begin{equation}
 T_4(n)=\frac{3\,256^n}{1024\sqrt2\,\pi^{3/2}n^{15/2}}
 \left(1+\frac{75}{16n}+\frac{6905}{512n^2}+O(n^{-3})\right).
 \label{shr:eq:A181198exp}
\end{equation}
The leading constants are consequences of the general theorem, not inputs
fitted from these sequences.

\subsection{Five corrections in multinomial normalization}
Table~\ref{shr:tab:coeffs} reports exact output of \eqref{shr:eq:coefficientformula}.
Only the compact multinomial normalization is used in this table.
The archive also supplies the corresponding $b_{m,j}$.

\begin{table}[ht]
\centering
\caption{Coefficients $c_{m,j}$ in \eqref{shr:eq:mainM}; $c_{m,0}=1$.}
\label{shr:tab:coeffs}
\renewcommand{\arraystretch}{1.3}
\begin{tabular}{c r r r r r}
\toprule
$m$ & $c_{m,1}$ & $c_{m,2}$ & $c_{m,3}$ & $c_{m,4}$ & $c_{m,5}$\\
\midrule
2 & $1/2$ & $1/4$ & $1/8$ & $1/16$ & $1/32$\\
3 & $2$ & $11/4$ & $23/8$ & $45/32$ & $-57/16$\\
4 & $5$ & $15$ & $33$ & $789/16$ & $93/16$\\
5 & $10$ & $225/4$ & $1819/8$ & $22045/32$ & $1353$\\
6 & $35/2$ & $665/4$ & $8939/8$ & $92351/16$ & $731621/32$\\
\bottomrule
\end{tabular}
\end{table}

\subsection{How much the corrections improve A181199}
For $m=5$, let $L_5(n)=K_5\,5^{5n}n^{-12}$ and
$A_{5,L}(n)=L_5(n)\sum_{j=0}^L b_{5,j}n^{-j}$.
Table~\ref{shr:tab:numerics} shows the signed relative error
$A_{5,L}(n)/T_5(n)-1$. Its values are numerical diagnostics, not rigorous
interval enclosures. The counts at $n=10,20,40$ were independently
regenerated by the row-state dynamic program and compared with OEIS;
$n=80$ is an OEIS reference value, not regenerated in this run
\cite{b199}.

\begin{table}[ht]
\centering
\caption{A181199: signed relative errors of successive truncations.}
\label{shr:tab:numerics}
\renewcommand{\arraystretch}{1.2}
\begin{tabular}{r r r r r}
\toprule
$n$ & $L=0$ & $L=1$ & $L=2$ & $L=5$\\
\midrule
10 & $-6.37190\cdot10^{-1}$ & $-2.88893\cdot10^{-1}$ & $-9.90345\cdot10^{-2}$ & $1.42184\cdot10^{-3}$\\
20 & $-3.90471\cdot10^{-1}$ & $-9.78967\cdot10^{-2}$ & $-1.81550\cdot10^{-2}$ & $1.97547\cdot10^{-5}$\\
40 & $-2.16403\cdot10^{-1}$ & $-2.83393\cdot10^{-2}$ & $-2.71080\cdot10^{-3}$ & $2.67266\cdot10^{-7}$\\
80 & $-1.13941\cdot10^{-1}$ & $-7.61418\cdot10^{-3}$ & $-3.69269\cdot10^{-4}$ & $3.76974\cdot10^{-9}$\\
\bottomrule
\end{tabular}
\end{table}

The leading term alone converges rather slowly at height five because
$b_{5,1}=9.6$. This explains why the constant is difficult to recognize
reliably from modestly sized data without correction terms. The proof,
however, does not use the numerical agreement.

\section{A complete algebraicity classification}\label{shr:sec:algebraicity}

We now prove Theorem~\ref{shr:thm:algebraic}. The argument uses the classical
Newton--Puiseux property: an algebraic function has a convergent fractional
power expansion at a finite point, with algebraic coefficients when both
the defining polynomial and the expansion point are algebraic. It has no
logarithmic leading singularity. This is standard algebraic-function
background; see the algebraic-function treatment in \cite{FS}.

Put $R_m=m^{-m}$ and rescale
\[
 G_m(z)=F_m(R_mz)=\sum_{n\ge0}v_nz^n,\qquad
 v_n=R_m^nT_m(n).
\]
For $m\ge2$, Theorem~\ref{shr:thm:main} gives
\begin{equation}
 v_n=K_m n^{-\alpha_m}(1+O(n^{-1})).
 \label{shr:eq:vn}
\end{equation}
In particular the radius of convergence of $G_m$ is exactly one.

\subsection{Odd heights: a logarithmic obstruction}
Suppose $m\ge3$ is odd. Then $\alpha_m$ is a positive integer. Set
$q=\alpha_m-1$. The coefficient of $z^k$ in $G_m^{(q)}(z)$ is
\[
 \frac{(k+q)!}{k!}v_{k+q}=\frac{K_m}{k}+O(k^{-2})
 \quad(k\longrightarrow\infty).
\]
After a harmless adjustment of finitely many terms, subtraction of
$K_m\sum_{k\ge1}z^k/k$ leaves an absolutely summable coefficient sequence.
Therefore, along the real interval $z\uparrow1$,
\begin{equation}
 G_m^{(q)}(z)=K_m\log\frac1{1-z}+O(1).
 \label{shr:eq:logsing}
\end{equation}
A nonzero logarithmic divergence is incompatible with a Puiseux expansion.
If $F_m$ were algebraic, so would be its rescaling and every derivative;
\eqref{shr:eq:logsing} is a contradiction.

\subsection{Even heights: a transcendental-amplitude obstruction}
Suppose $m\ge4$ is even. Now $q=\alpha_m-1/2=(m^2-2)/2$ is an integer.
The derivative coefficients satisfy
\[
 [z^k]G_m^{(q)}(z)=K_m k^{-1/2}+O(k^{-3/2}).
\]
Since
$[z^k](1-z)^{-1/2}=1/\sqrt{\pi k}+O(k^{-3/2})$, coefficient subtraction
again leaves an absolutely summable series. Hence
\begin{equation}
 G_m^{(q)}(z)=K_m\sqrt\pi\,(1-z)^{-1/2}+O(1),
 \qquad z\uparrow1.
 \label{shr:eq:halfsing}
\end{equation}
If the function were algebraic over $\Q(z)$, the leading coefficient in
this Puiseux expansion at the algebraic point $z=1$ would be algebraic.
But
\begin{equation}
 K_m\sqrt\pi=
 \frac{\sqrt m\,J_m}{4^d2^{(m-1)/2}}
 \pi^{-(m-2)/2}
 \label{shr:eq:amplitude}
\end{equation}
is a nonzero algebraic multiple of a nonzero integer power of $\pi^{-1}$.
It is transcendental: otherwise that power of $\pi$, and hence $\pi$
itself, would be algebraic. This contradiction uses only the classical
transcendence of $\pi$.

\subsection{The two algebraic cases and descent of constants}
The exceptional cases are explicit:
\[
 F_1(z)=\frac1{1-z},\qquad
 F_2(z)=1+zC(z),\qquad
 C(z)=\frac{1-\sqrt{1-4z}}{2z}.
\]
Thus $m=1,2$ are algebraic.

Finally, a rational-coefficient power series algebraic over
$\mathbb C(z)$ is algebraic over $\Q(z)$. Indeed, fix finite degree bounds
for a nonzero polynomial relation. Equating all coefficients gives a
homogeneous linear system with rational entries in finitely many unknown
polynomial coefficients. The descending chain of solution spaces
stabilizes after finitely many independent equations. A nonzero complex
solution therefore implies a nonzero rational solution. This proves the
asserted classification over either field.

\begin{remark}
Nonalgebraicity does not imply non-D-finiteness. Logarithmic singularities,
and half-integer singularities with transcendental connection constants,
are compatible with D-finite functions. The guessed recurrences for
heights four and five remain separate questions.
\end{remark}

\begin{writenote}[Added 5 October 2026, batch~98]
$F_m$ is in fact D-finite for every fixed $m$: it is the diagonal of a
rational power series over $\Q$ (Part~\ref{shr:bp:part},
Theorem~\ref{shr:bp:thm:diagonal-main}, through the multiple-binomial-sum
theorem of Bostan, Lairez and Salvy). With Theorem~\ref{shr:thm:algebraic},
$F_m$ is algebraic for $m\le2$ and transcendental but D-finite for $m\ge3$.
That proof is an existence proof; it produces no operator, and the guessed
recurrences for heights four and five remain unproved.
\end{writenote}

\begin{writenote}[Added 5 October 2026, batch~101]
The last sentence of the note above is out of date for heights four and
five: Parts~\ref{shr:h4:part} and~\ref{shr:h5:part} prove the guessed
recurrences, which are explicit recurrences of orders two and three for the
coefficients of $F_4$ and $F_5$. The order two at height four is minimal
(Proposition~\ref{shr:h4:prop:minimal}); at height five the minimal order
is two or three (Proposition~\ref{shr:h5:prop:nothyper}). For $m\ge6$ no
operator is known.
\end{writenote}

\section{Gaussian laws for intermediate shapes}\label{shr:sec:GUE}

The same reduction describes typical partially filled arrays, not only
their total number. We give both a randomized numerical cut and a fixed
number of inserted labels; these are different probability models.

\subsection{A cut at a fixed numerical level}
Choose a tableau uniformly among the $T_m(n)$ arrays. Independently choose
$mn$ uniform random numbers in $(0,1)$, sort them, and place them in the
array in increasing order of its integer labels. This produces the uniform
order-statistics model conditioned on $\mathcal A$.
For a fixed $p\in(0,1)$ let $X_i$ be the number of row-$i$ entries below
$p$ and set
\begin{equation}
 Z_i=\frac{X_i-np}{\sqrt{np(1-p)}}.
\end{equation}

\begin{theorem}[Gaussian Vandermonde limit]\label{shr:thm:gue}
For fixed $m$ and fixed $p\in(0,1)$, the decreasing vector $Z$ converges in
distribution to the probability density
\begin{equation}
 \frac{e^{-\|z\|^2/2}\Delta(z)^2}{(2\pi)^{m/2}J_m},
 \qquad z_1>z_2>\cdots>z_m,
 \label{shr:eq:GUEdensity}
\end{equation}
with respect to $m$-dimensional Lebesgue measure. Polynomial moments of
fixed degree converge as well.
\end{theorem}

\begin{proof}
For interior decreasing $x$, the exact conditional mass is
\[
 \frac{\prod_i\Pp(\operatorname{Bin}(n,p)=x_i)R_n(x)}
      {T_m(n)/M_m(n)}.
\]
In bounded sets of standardized coordinates, the local central limit
estimate for the binomial mass follows directly from Stirling's formula.
Moreover,
\[
 x_i-x_j=\sqrt{np(1-p)}(z_i-z_j),
\]
whereas the two denominator factors in \eqref{shr:eq:Rweight} are asymptotic
to $2np$ and $2n(1-p)$. Thus
\[
 R_n(x)=(4n)^{-d}\Delta(z)^2(1+o(1))
\]
uniformly on every bounded set. Divide by
$T_m(n)/M_m(n)\sim J_m(4n)^{-d}$ and use the lattice mesh
$1/\sqrt{np(1-p)}$ in every coordinate. The Riemann sums converge to
\eqref{shr:eq:GUEdensity}.

To justify passage beyond bounded sets, first restrict every $x_i/n$ to a
fixed compact interval inside $(0,1)$ containing $p$. There the pair
denominators are bounded below by a constant times $n^2$, and binomial
local bounds give a Gaussian upper bound times a fixed polynomial in the
standardized coordinates. This is integrable, with every polynomial
moment. Outside that macroscopic interval, Chernoff bounds and $R_n\le1$
give an exponentially small contribution even after division by the
order-$n^{-d}$ conditioning probability. The boundary term satisfies the
same conclusion by \eqref{shr:eq:boundary}. This proves tightness and moment
convergence. Normalization of the limit is \eqref{shr:eq:gaussZ}, with the
factor $m!$ removed on the decreasing chamber.
\end{proof}

The density \eqref{shr:eq:GUEdensity} is the ordered eigenvalue density of the
Gaussian unitary ensemble in the normalization with Gaussian factor
$e^{-\sum z_i^2/2}$. This terminology is optional: the theorem and its
normalization have already been specified by an explicit density.

\subsection{A deterministic fraction of the labels}
Let $\Lambda_i(k)$ be the number of cells in row $i$ occupied by labels
$1,\ldots,k$ of a uniformly random tableau. Fix $p\in(0,1)$ and take
$k=\lfloor mnp\rfloor$.

\begin{theorem}[Trace-zero limit at deterministic time]\label{shr:thm:tracezero}
The vector
\[
 \left(\frac{\Lambda_i(k)-np}{\sqrt{np(1-p)}}\right)_{i=1}^m
\]
converges to the density
\begin{equation}
 \frac{e^{-\|z\|^2/2}\Delta(z)^2}{(2\pi)^{(m-1)/2}J_m}
 \label{shr:eq:tracezerodensity}
\end{equation}
on $\{z_1>\cdots>z_m,\ \sum_i z_i=0\}$, with respect to Euclidean
$(m-1)$-dimensional measure on the trace-zero hyperplane.
\end{theorem}

\begin{proof}
In the preceding randomized-cut construction, the total number $K$ of
samples below $p$ is $\operatorname{Bin}(mn,p)$ and is independent of the
integer tableau. Conditioning on $K=k$ therefore gives precisely
$\Lambda(k)$. For an interior vector $x$ of total size $k$, divide the
mass used in the preceding proof by $\Pp(K=k)$.
The latter is asymptotic to $(2\pi mnp(1-p))^{-1/2}$.
The lattice $\{x\in\mathbb Z^m:\sum x_i=k\}$ has Euclidean covolume
$\sqrt m$. In standardized coordinates its cell volume is
$\sqrt m\,(np(1-p))^{-(m-1)/2}$. These factors give exactly the density
\eqref{shr:eq:tracezerodensity} under the local binomial estimate.

For domination, use the same compact macroscopic region as before.
The product of local binomial upper bounds on the hyperplane, after the
additional conditioning, is bounded by a fixed polynomial times a Gaussian
in the standardized coordinates, times the hyperplane lattice-cell
volume. Its tails are uniformly summable. Outside the macroscopic region,
the extra factor $1/\Pp(K=k)=O(\sqrt n)$ cannot overcome an exponential
Chernoff bound. This proves convergence of the hyperplane Riemann sums.

Finally, $\Delta$ is invariant under common translation. Orthogonally
splitting $\R^m$ into the trace-zero hyperplane and the unit trace direction
shows that the full Gaussian Vandermonde integral is $\sqrt{2\pi}$ times
the trace-zero integral. Equation~\eqref{shr:eq:gaussZ} then verifies the
normalization in \eqref{shr:eq:tracezerodensity}.
\end{proof}

This theorem concerns one fixed macroscopic time. It is not a claim of
process convergence to a noncolliding Brownian bridge; proving tightness
in a path space and identifying the joint time laws require additional
work.

\section{Ratios, eventual log-convexity, and inversion}\label{shr:sec:inverse}

\subsection{Consecutive ratios}
Write $\alpha=\alpha_m$ and $b_1=b_{m,1}$. For $m\ge2$, the all-order
expansion gives
\begin{equation}
 \frac{T_m(n+1)}{T_m(n)}=m^m
 \left(1-\frac\alpha n+
 \frac{\alpha(\alpha+1)/2-b_1}{n^2}+O_m(n^{-3})\right).
 \label{shr:eq:ratio}
\end{equation}
Hence $T_m(n)$ is eventually strictly increasing. Expanding
$\log T_m(n)$ at $n-1,n,n+1$ with enough terms also gives
\begin{equation}
 \log T_m(n+1)-2\log T_m(n)+\log T_m(n-1)
 =\frac\alpha{n^2}+O_m(n^{-3})>0
\end{equation}
for sufficiently large $n$. Thus it is eventually strictly log-convex.
These statements do not identify the smallest valid $n$.

\subsection{Inverting the asymptotic profile}
Put $\lambda=m\log m$, $K=K_m$, and $L=\log(y/K)$. The large solution of
\[
 \lambda x-\alpha\log x=L
\]
is
\begin{equation}
 x_0=-\frac\alpha\lambda
 W_{-1}\!\left(-\frac\lambda\alpha e^{-L/\alpha}\right),
 \label{shr:eq:inverseW}
\end{equation}
where $W_{-1}$ is the real branch taking values at most $-1$
\cite{DLMFW}. The branch choice matters: the principal branch describes
the small, irrelevant solution. In elementary logarithmic form,
\begin{equation}
 x_0=\frac{L+\alpha\log L-\alpha\log\lambda}{\lambda}
 +O_m\left(\frac{\log L}{L}\right).
\end{equation}

Let $\ell_2=b_{m,2}-b_{m,1}^2/2$. Formal inversion of the continuous
asymptotic profile gives
\begin{equation}
 x=x_0-\frac{b_{m,1}}{\lambda x_0}
 -\frac{\alpha b_{m,1}/\lambda^2+\ell_2/\lambda}{x_0^2}
 +O_m(x_0^{-3}).
 \label{shr:eq:inversecorrection}
\end{equation}
To verify it, insert $x=x_0+\delta$ into
$\lambda x-\alpha\log x+b_{m,1}/x+\ell_2/x^2=L$ and compare the
coefficients of $x_0^{-1}$ and $x_0^{-2}$.
Along the exact values $y=T_m(n)$, the approximation in
\eqref{shr:eq:inversecorrection} returns $n+O_m(n^{-3})$.
The higher inverse terms are obtained by the same recursive substitution.
This inversion technique is standard and also used elsewhere in ProveIt;
its application here is a consequence of the new expansion, not a new
inversion method.

For an arbitrary target $y$, the least integer index attaining it is a
staircase function. An asymptotic estimate alone must not be rounded and
advertised as an exact threshold certificate. Near an integer transition,
explicit upper and lower bounds for the neighboring sequence values are
required. The finite interior sum and its boundary error
\eqref{shr:eq:exactbound} supply one possible, though not always computationally
efficient, route to such bounds.

\begin{writenote}[Added 3 October 2026, batch~85: the inversion is an instance]
The inversion \eqref{shr:eq:inverseW} is a case of Theorem~J.23
(\texttt{p0:thm:lambert-core}, ``Exact solution of the dominant block'') of
the canonical transseries volume
\texttt{Analysis/\allowbreak Transseries/\allowbreak docs/\allowbreak series-\allowbreak and-\allowbreak transseries/\allowbreak Transseries\_\allowbreak And\_\allowbreak Inversion/},
which solves $aX+b\log X=L$ by $X=(b/a)\,W_k\bigl((a/b)e^{L/b}\bigr)$. Here
$a=\lambda=m\log m>0$ and $b=-\alpha_m<0$ for $m\ge2$, and that theorem's
branch rule for $b<0$ gives exactly \eqref{shr:eq:inverseW}: the large root is
the $W_{-1}$ root, and it exists precisely when
$L>L_c=\alpha_m\bigl(1-\log(\alpha_m/\lambda)\bigr)$ (for $m=5$,
$K_5e^{L_c}\approx0.021$, so every $y\ge1$ qualifies). The text's own
sentence that the technique ``is standard and also used elsewhere in
ProveIt'' is accurate; nothing in this section is a new inversion result.
The real substitution, branch rules, threshold and slope of Theorem~J.23
are formalized in
\texttt{Analysis/\allowbreak FabiusFunction/\allowbreak Lean/\allowbreak FabiusFunction/\allowbreak LinLogCoreInversion.lean}
(the volume's register rates the theorem ``Partial'': the asymptotic and
general complex-branch clauses are not formalized); that formalization concerns the general equation, and
nothing specific to shifted rectangles is formalized. The caution about
integer thresholds in the preceding paragraph is the separation condition
of the volume's staircase theorem, Theorem~J.43 (\texttt{p0:thm:staircase}).
The path-forest report cited in Section~1.3 inverts its sequences by the
same apparatus (its \texttt{spf:thm:inverse-first} and
\texttt{spf:thm:inverse-all}).
\end{writenote}

\section{Reproducibility and computational scope}

\subsection{Exact coefficient generation}
The program \texttt{code/coefficients.py} implements
\eqref{shr:eq:coefficientformula} using Python's exact rational arithmetic.
It avoids expanding the squared Vandermonde in full. For any symmetric
polynomial $P$ and any product measure with identical one-variable moments,
\begin{equation}
 \mathcal E_z(\Delta^2P)=m!
 \sum_{\sigma\in S_m}\sgn(\sigma)\,
 \mathcal E_z\left(P\prod_{i=0}^{m-1}y_{i+1}^{\,i+\sigma(i)}\right).
 \label{shr:eq:collapseperms}
\end{equation}
Expand the two Vandermonde determinants, permute the integration variables
in one expansion, and use symmetry of $P$ to prove this identity.
Only $m!$ permutations remain. Equal sorted exponent profiles are grouped
before evaluating products of moment polynomials.

The $h_r$ are generated from \eqref{shr:eq:hdef}; the standardized moments are
generated from sign cumulants. The Stirling conversion
\eqref{shr:eq:stirling-convert} is also done with rational numbers. The method
is practical for modest fixed $m$ and moderate order. Its permutation
factor makes it unsuitable as presently implemented for large $m$;
no claim of polynomial complexity in $m$ is made.

\subsection{Independent enumeration and exact finite identities}
The program \texttt{code/tableaux.py} supplies two enumerators: a row-state
dynamic program using \eqref{shr:eq:transition}, and an independent bitmask
linear-extension enumerator using the cell-level predecessor relation.
The latter is intentionally restricted to small arrays. Neither uses a
conjectured recurrence.

The verification run recorded in the archive passed 471 checks in the
following suites:
\begin{center}
\begin{tabular}{p{0.76\textwidth}r}
\toprule
Suite & Checks\\
\midrule
Bitmask-poset versus row-state enumeration & 27\\
Interior prefix counts versus the shifted hook product & 238\\
Threshold split, rational interior sum, and boundary inequality & 120\\
Exact Catalan control, $n=1,\ldots,30$ & 30\\
Selected A181198 values, independently regenerated & 12\\
Selected A181199 values, independently regenerated & 12\\
Universal corrections versus the coefficient engine, $m=1,\ldots,6$ & 30\\
Standardized binomial-moment controls & 2\\
\bottomrule
\end{tabular}
\end{center}
The independently regenerated OEIS comparisons use $n=1,\ldots,10,20,40$
for each of heights four and five \cite{b198,b199}. The threshold identity
is checked for $1\le m\le5$ and $1\le n\le8$.
These are finite corroborations of the proofs, not substitutes for them.

\subsection{Commands and data provenance}
From the archive's top-level directory:
\begin{verbatim}
python3 code/verify.py
python3 code/coefficients.py --height 5 --order 5
python3 code/numerics.py
bash build.sh
\end{verbatim}
The first two commands use the Python standard library only.
The numerical diagnostics additionally use \texttt{mpmath}; they were run
with 90 decimal digits. The PDF build requires a conventional LaTeX
installation with the packages listed in the preamble.

The reference integers in \texttt{data/oeis\_selected.json} were transcribed
from the entry pages and b-files inspected during the source audit. They
are not outputs of the coefficient engine. Counts regenerated independently
are saved separately in \texttt{data/exact\_regenerated.json}. Values at
$n=80$ appear only as reference inputs to numerical diagnostics. Full
third-party articles are not redistributed in this archive.

\begin{writenote}[Added 3 October 2026, batch~85: the shipped layout]
``The archive'' in this section is the delivered package. In ProveIt its
files keep the subdirectories \texttt{code/} and \texttt{data/}, except
four, which moved as follows:
\begin{center}\footnotesize
\begin{tabular}{@{}ll@{}}
\toprule
Delivered as & Shipped as\\
\midrule
\texttt{build.sh} & \texttt{code/build.sh}\\
\texttt{requirements-numerics.txt} & \texttt{data/requirements-numerics.txt}\\
\texttt{notes/SOURCES\_AND\_STATUS.md} & \texttt{notes-SOURCES\_AND\_STATUS.md}\\
\texttt{notes/PROPOSED\_OEIS\_ADDITIONS.txt} & \texttt{data/notes-PROPOSED\_OEIS\_ADDITIONS.txt}\\
\bottomrule
\end{tabular}
\end{center}
\texttt{code/build.sh} changes to its own directory, so in its shipped
place it no longer finds \texttt{article.tex}. The delivered PDF and
checksum list are not shipped; they survive in the arrival commit's
archive. \texttt{code/verify.py} and
\texttt{code/numerics.py} write their outputs into the report's own
\texttt{data/} directory, overwriting the shipped records, so they must be
rerun on a copy; the README of the report gives the commands. At intake
(3 October 2026, Python 3.14.4 on Windows) \texttt{code/verify.py} passed
all 471 checks on a copy, and every regenerated data file agreed with the
shipped one after removal of carriage returns, except the timing field
of \texttt{data/\allowbreak verification.json}; \texttt{code/coefficients.py}
and \texttt{code/numerics.py} (mpmath~1.3.0) reproduced
\texttt{data/\allowbreak coefficients\_m5.json} and \texttt{data/\allowbreak numerics.csv}.

\texttt{data/\allowbreak oeis\_selected.json} holds terms of A181198 and
A181199 copied from the OEIS entries and b-files (Koutschan, with earlier
terms by Heinz); OEIS content is published by The OEIS Foundation Inc.\
under the Creative Commons Attribution-ShareAlike 4.0 licence, so this file
is third-party data under that licence, not MIT-0. The draft OEIS text in
\texttt{data/\allowbreak notes-\allowbreak PROPOSED\_\allowbreak OEIS\_\allowbreak ADDITIONS.txt} is
marked ``DRAFT ONLY --- NOT SUBMITTED'' and stays so: nothing was submitted
to OEIS by the package's author or by the intake, and its own instruction
not to mark either conjectured recurrence as proved stands.
\end{writenote}

\begin{writenote}[Added 5 October 2026, batch~101: the OEIS draft]
The draft's instruction ``Do not mark either conjectured sequence
recurrence as proved on the strength of this article'' was right for this
manuscript, which proves no recurrence, and the shipped file stays
byte-identical. Since 5~October 2026 Parts~\ref{shr:h4:part}
and~\ref{shr:h5:part} prove both recurrences. The draft has not been
updated, nothing has been submitted to OEIS, and any change to the OEIS
entries is for a human editor after review.
\end{writenote}

\section{Further research questions and proposed approaches}

The following questions are not presented as solved. They are chosen to
continue the present line of work rather than to restate the resolved
leading-term conjecture.

\begin{question}[Prove the specific height-four and height-five recurrences]
Can the exact enumeration be transformed into a finite hypergeometric sum
or a Pfaffian identity that admits a compact creative-telescoping
certificate for the recurrences in A181198 and A181199?
\end{question}
The random-threshold expression is exact only after adding its boundary
part. For a recurrence proof that part cannot be discarded. A useful first
step is an exact decomposition according to the number of full and empty
rows at the cut, followed by shifted-skew-tableau formulas for those
boundary states. A proof should compare the resulting operator with the
published guessed operator, not merely generate more confirming terms.

\begin{writenote}[Added 5 October 2026, batch~101: Research question~1 answered]
Answered at both heights, through a Pfaffian identity rather than the
boundary decomposition suggested above. Part~\ref{shr:h4:part} (Report~231)
writes the order-polytope volume as an integral of a determinant of the
Volterra transfer $V^{n-1}$, evaluates it by a four-variable de~Bruijn
Pfaffian as two kernel traces, reduces both traces to one weighted triangle
period $Z_n$, and proves a first-order recurrence for $Z_n$ with a
six-coefficient polynomial divergence certificate
(Theorem~\ref{shr:h4:thm:Z}); elimination gives the published A181198
operator (Theorem~\ref{shr:h4:thm:oeis}) and variation of constants
Conjecture~18 of \cite{KK} (Theorem~\ref{shr:h4:thm:sum}).
Part~\ref{shr:h5:part} (Report~233) does the same at height five with an
odd augmented Pfaffian, an open-chain integral and a second period $X_n$,
closed by polynomial contiguity and a finite moment closure; it proves the
nested sum of Conjecture~19 and the published A181199 operator
(Theorems~\ref{shr:h5:thm:sum} and~\ref{shr:h5:thm:operator}). Both keep
every boundary contribution, compare their operators with the published
ones coefficient by coefficient, and settle the initial indices: the
operators hold for $n\ge1$ and fail at $n=0$ under $a_0=1$. Neither uses a
creative-telescoping certificate; a compact one at height five is
Research question~17. Nothing is known for $m\ge6$ (item~(F7) and
Research question~16).
\end{writenote}

\begin{question}[Polynomial dependence of every correction on height]
For each fixed $j$, is $c_{m,j}$ a polynomial in $m$ of degree at most $3j$?
Can its factorization and denominator structure be described uniformly?
\end{question}
The formulas for $j=1,2$ support this degree scale. The question beyond
those orders is not proved here. A promising approach is to reorganize
\eqref{shr:eq:coefficientformula} as Gaussian matrix moments or as connected
cumulant diagrams, in which the powers of $m$ count index loops. This could
replace factorial-in-$m$ computation with symbolic fixed-order formulas.

\begin{question}[A growing-height transition]
For what growth rates $m=m(n)$ does the leading fixed-height formula remain
relatively accurate? Is the natural first transition governed by $m^3/n$?
\end{question}
The coefficient $c_{m,1}\sim m^3/12$ suggests this parameter in multinomial
normalization. It does not establish a uniform regime. A proof would need
uniform Vandermonde moment bounds and a growing-dimension tail analysis.
For $m^3/n$ of order one, connected logarithmic corrections may require
resummation rather than ordinary fixed-order truncation.

\begin{question}[Exponentially small sectors and sharp boundary asymptotics]
What is the true leading asymptotic of the discarded boundary term
$E_{m,n}(1/2)$, and how does it combine with large deviations of the
interior expectation?
\end{question}
The bound $2m2^{-n}$ is sufficient for every inverse power but is not an
asymptotic formula. Configurations with a row empty or full give a natural
finite hierarchy of lower-dimensional boundary strata. Computing their
constants could produce a genuine transseries beyond the algebraic sector,
and could clarify the signs and optimal truncation of the coefficients.

\begin{writenote}[Added 5 October 2026, batch~98: Question~4 re-scoped]
Part~\ref{shr:bp:part} does not answer this question, although its
introduction counts the non-sharp boundary bound among ``the two gaps
addressed here''. Its rare count $R_m(n)$ is the part of the boundary event
in which the first row is completed before the last row starts: it
satisfies $R_m(n)\le M_m(n)E_{m,n}(1/2)$, and the ratio grows like $2^n$
times a power of $n$ (Remark~\ref{shr:bp:rem:q4}). The leading asymptotic
of $E_{m,n}(1/2)$ itself is the write's
Proposition~\ref{shr:bp:prop:boundary}: exactly, $E_{m,n}(1/2)$ is twice
the stratum with one boundary row (the last row empty at the cut, or, by
symmetry, the first row full) plus $O(4^{-n})$, and
\[
 E_{m,n}(1/2)\sim2^{1-n}3^{1-m}J_{m-1}(4n)^{-\binom{m-1}2}.
\]
That proposition is unrefereed and checked numerically only on finite
ranges. Still open: an all-order expansion of the boundary term (the
one-row stratum is an interior sum with $m-1$ rows and an extra factor
$x_i/(2n-x_i)$ per row, so the method of Section~5 should apply, but this
has not been carried out); the constants of the deeper strata, with two or
more rows empty or full at the cut, and of the rare sectors inside them
(compare Research question~10 of Part~\ref{shr:bp:part}); their
combination with large deviations of the interior expectation; and any
transseries reading.
\end{writenote}

\begin{question}[A full path-process limit]
Do the centered row-length processes of a uniform tableau converge to a
trace-zero noncolliding Gaussian bridge on compact time intervals inside
$(0,1)$? What are the entrance and exit scalings near the endpoints?
\end{question}
Theorem~\ref{shr:thm:tracezero} establishes one-time marginals only. A proof of
process convergence would require multi-cut enumeration, control of
increments, and tightness. The complement symmetry provides a useful
constraint on any proposed entrance/exit description.

\begin{question}[Certified asymptotic enclosures and integer thresholds]
Can the moment expansion be packaged into efficient explicit upper and
lower bounds with practical constants, avoiding the $\binom{n-1}{m}$-term
interior sum?
\end{question}
On a central region the pair interaction has a positive geometric
expansion. Combining its finite partial sums with concentration bounds and
exact moment determinants could give certified intervals. These would
turn the inverse approximations into exact index decisions away from a
quantified transition zone.

\begin{writenote}[Added 5 October 2026, batch~101: Research question~6 advanced at heights four and five]
At $m=4$ and $m=5$, Parts~\ref{shr:h4:part} and~\ref{shr:h5:part} turn the
proved recurrences into effective statements: rational error constants for
every fixed truncation (Theorems~\ref{shr:h4:thm:bexp}
and~\ref{shr:h5:asy:fixedorder}), explicit logarithmic remainders, finite
Lambert-$W_{-1}$ models with effective inverse errors, and two-ceiling
envelopes for the integer threshold (Theorems~\ref{shr:h4:thm:threshold}
and~\ref{shr:h5:inv:envelope}), which leave at most two adjacent candidates
near an integer. At height five the receipt evaluates constants at $K=0$
and $K=6$, for example $|b_n-q_6(n)|\le18n^{-17}$ for $n\ge64$
\eqref{shr:h5:asy:explicit6}. These bounds need no interior sum, but they
rest on the recurrences and do not reach $m\ge6$. Still open: practical
constants at general $m$, and the constants that Parts~III and~IV construct
without evaluating (items~(F6) and~(F10)).
\end{writenote}

\begin{question}[D-finiteness versus nonalgebraicity at arbitrary height]
Is $F_m(z)$ D-finite for every fixed $m$, and how does the minimal
annihilating order grow with $m$?
\end{question}
Theorem~\ref{shr:thm:algebraic} settles the algebraic/nonalgebraic distinction,
not the D-finite/non-D-finite one. An exact finite-dimensional differential
system in the fixed-height regime would be a stronger structural result.
The opposite-orientation constant-width transfer matrix does not answer
this question.

\begin{writenote}[Added 5 October 2026, batch~98: Question~7 answered in part]
The first half is answered: for every fixed $m$, $F_m$ is the diagonal of
a rational power series over $\Q$, hence D-finite, and $(T_m(n))$ is
P-recursive (Part~\ref{shr:bp:part}, Theorem~\ref{shr:bp:thm:diagonal-main}:
the exact sum of Theorem~\ref{shr:bp:thm:skeleton} over the $2m$ row births
and completions, a finite sum of products of reflection determinants, is a
multiple binomial sum in the sense of Bostan, Lairez and Salvy). The second
half, how the minimal annihilating order (and degree) grows with $m$,
remains open; it is Research question~11 of Part~\ref{shr:bp:part}. The
existence proof supplies no operator, so the guessed recurrences of heights
four and five (Research question~1) remain unproved.
\end{writenote}

\begin{writenote}[Added 5 October 2026, batch~101]
The last sentence of the note above is out of date:
Parts~\ref{shr:h4:part} and~\ref{shr:h5:part} prove the guessed recurrences
of heights four and five (note at Research question~1). On the second half
of this question, the minimal order of a recurrence is two at $m=4$
(Proposition~\ref{shr:h4:prop:minimal}) and two or three at $m=5$
(Proposition~\ref{shr:h5:prop:nothyper}).
\end{writenote}

\section{Conclusion}
A shifted rectangle admits a particularly effective asymptotic cut: both
sides of an interior threshold state are ordinary strict shifted shapes.
The two shifted hook products turn the count into a binomial Vandermonde
expectation, with an explicitly controlled boundary error. This gives a
proof of the A181199 leading constant independent of all guessed
recurrences, an effective all-order expansion for every fixed height,
and universal correction polynomials through second order.

The consequences go beyond adding coefficients to an OEIS entry. The
same constant yields a complete algebraicity dichotomy for the fixed-height
generating functions, while the same exact weight identifies Gaussian
laws for typical intermediate shapes. The remaining recurrence, uniformity,
and exponentially small questions are separate, concrete targets with
clear starting points.

\clearpage
\part{Rational diagonals and rare boundary phases in shifted rectangles}
\label{shr:bp:part}

\noindent{\Large D-finiteness at every fixed height and a sharp
exponential refinement for OEIS A181199\par}
\smallskip
\noindent{\large Prepared for Vladimir Reshetnikov, October 4, 2026\par}
\medskip

\begin{resultbox}
For every fixed height $m$, the generating function of the shifted-rectangle
counts $T_m(n)$ is a diagonal of a rational function, hence D-finite.
For $m\ge2$, let $R_m(n)$ count arrays in which row $1$ finishes before row $m$ starts. Then
\[
 R_m(n)\sim
 \sqrt{\frac m2}\,(2\pi)^{-(m-2)/2}
 \frac{\prod_{j=0}^{m-3}j!}{4^{\binom{m-2}{2}}9^{m-2}}
 \left(\frac{m^m}{4}\right)^n n^{-(m-2)^2/2}.
\]
Both statements hold for all fixed heights, not only the tabulated cases.
\end{resultbox}

\paragraph{Abstract.}
The fixed-height shifted-rectangle family combines strict inequalities
between active rows with exceptional ties at empty and completed rows.
We separate these boundary events from the intervening walks. The resulting
finite sum of reflection determinants is a multiple binomial sum, proving
rational diagonality and D-finiteness at every fixed height and resolving
an explicit question in the preceding ProveIt report. We also count the
rare failure of a common active phase. A first-completion cut gives a
shifted-hook sum with an exponentially smaller remainder. A beta--binomial
mixture then yields its sharp exponential rate, leading constant, effective
all-order expansion, and a universal first correction. We compute three
corrections at height five, describe the conditional Gaussian shape, and
invert the rare-event profile. Exact verification programs and a source
and status audit accompany the manuscript.

\paragraph{Status.} AI-assisted, unrefereed research manuscript.
The stated results have mathematical proofs below, but have not been
independently peer reviewed or proof-assistant verified. Historical priority
is not certified. The particular conjectured recurrences and finite-sum
identities for A181198 and A181199 are \emph{not} proved here.

\begin{writenote}[Added 5 October 2026, batch~98: provenance of Part~II]
Part~II is a single manuscript, printed in full: manuscript~02 of batch~98
of ProveIt's incoming-reports intake, \emph{Rational diagonals and rare
boundary phases in shifted rectangles}, subtitled ``D-finiteness at every
fixed height and a sharp exponential refinement for OEIS A181199'' and
dated October~4, 2026 (a 1{,}656-line source with 97 labels and a 24-page
PDF). Its title page reads ``Prepared for Vladimir Reshetnikov''; its PDF
metadata name ``Research report prepared for Vladimir Reshetnikov''. No
human author and no tool are named. It arrived as the archive
\texttt{Shifted\_\allowbreak Rectangle\_\allowbreak Boundary\_\allowbreak Research.zip}
(390{,}943 bytes, 22 files under a doubled wrapper directory of the same
name) in commit \texttt{2172df76a} and was placed in commit
\texttt{0fba5167f} (batch~98B) as an addition to this report. It pins the
repository at commit \texttt{db20eb37982abc0f163c6d308d392f5ac4c9bc62}
(4~October 2026; Section~\ref{shr:bp:sec:scope}). This report was not
changed between the pin and the placement, so the manuscript read exactly
the text printed as Part~\ref{shr:part:one}, and its citations of that text
(Theorems~2.1 and~2.3, Research question~7) carry the right numbers. Its
code, data and source notes are shipped with the prefix
\texttt{02-boundary-phases-}; its \LaTeX\ source, README, PDF and checksum
list are not shipped and survive in the archive of the arrival commit. The
manuscript is unrefereed and AI-assisted; its proofs are ordinary
mathematical proofs, nothing in it is formalized in Lean or Rocq, and its
place in this report confers no formal status.

\emph{What it answers.} Theorem~\ref{shr:bp:thm:diagonal-main} answers the
first half of Part~\ref{shr:part:one}'s Research question~7 (D-finiteness
at every fixed height), not its second half (Remark~\ref{shr:bp:rem:q7}).
Theorem~\ref{shr:bp:thm:rare-main} is a sharp result about a sub-sector of
the boundary event of Theorem~\ref{shr:thm:cut}; it does not answer
Research question~4, which asks for the asymptotic of the whole boundary
error (Remark~\ref{shr:bp:rem:q4}; the write's
Proposition~\ref{shr:bp:prop:boundary} gives its leading term). Dated notes
at both questions, at Remark~8.1 and at the end of Section~4.1 of
Part~\ref{shr:part:one} say so.

\emph{Overlap with Part~\ref{shr:part:one}, credited.} The two texts share
1.98\,\% of their word 8-grams (214 of 10{,}822). The state space
\eqref{shr:bp:eq:states} is \eqref{shr:eq:states}, re-derived; the shifted
hook product \eqref{shr:bp:eq:hook} is imported (Appendix~\ref{shr:app:hook}
proves it); \eqref{shr:bp:eq:gaussianZ} is \eqref{shr:eq:gaussZ} with $r$
variables; \eqref{shr:bp:eq:binomialVdm} extends \eqref{shr:eq:binomZ} to a
general success probability and is then mixed over the beta law;
\eqref{shr:bp:eq:gue-moments} are the Gaussian limits of
\eqref{shr:eq:traceidentities}; the Stirling factor \eqref{shr:bp:eq:outer}
is the analogue of \eqref{shr:eq:stirling-convert};
Lemma~\ref{shr:bp:lem:mixture} is a random-cut variant of
Theorem~\ref{shr:thm:cut}, and Theorem~\ref{shr:bp:thm:shape} the analogue,
for the rare event, of Theorem~\ref{shr:thm:gue}. The inversion of
Section~\ref{shr:bp:sec:inverse} uses the apparatus of
Section~\ref{shr:sec:inverse} with a different rate and power, and
Corollary~\ref{shr:bp:cor:probability} imports Theorem~\ref{shr:thm:main}
and Corollary~\ref{shr:cor:A181199} ($b_{5,1}=48/5$). Everything else is
new to the repository. ``The preceding (ProveIt, repository, asymptotic)
report'' in Part~II means Part~\ref{shr:part:one}.

\emph{Editorial changes.} Section~$k$ of the manuscript is
Section~$k+13$ here, its Appendices~A and~B are Appendices~C and~D, and its
Research questions~1--8 are Research questions~8--15 (one counter runs
through both Parts); theorem and equation numbers follow their sections.
Every label carries the prefix \texttt{shr:bp:}, and one label was added
(Remark~\ref{shr:bp:rem:constructive}). The manuscript's $d=\binom{m-2}2$
is renamed $d_{\rm R}$ throughout, because Part~\ref{shr:part:one}'s $d$ is
$\binom m2$; no normalization changes. Its three citations of its
predecessor (the bibliography entry \texttt{prior}, whose pin and path this
note records) now point to Part~\ref{shr:part:one}. Its title, subtitle,
result box, abstract and status paragraph, from its title page, open this
Part. Falling factorials are typeset as in Part~\ref{shr:part:one}
($n^{\underline k}$, the manuscript's $(n)_{\underline k}$); rising
factorials keep the manuscript's $(n)^{\overline k}$. The bibliographies
are merged: the OEIS, Kauers--Koutschan and Sun entries are
Part~\ref{shr:part:one}'s (the manuscript inspected the OEIS entries on
4~October 2026 and arXiv~v2 of Kauers--Koutschan, its Section~6.4 and
Conjectures~18--19), and Bostan--Lairez--Salvy \cite{BLS} is added. The
\textbf{[write]} notes, Remarks~\ref{shr:bp:rem:q7}
and~\ref{shr:bp:rem:q4}, Section~\ref{shr:bp:sec:boundary} with
Proposition~\ref{shr:bp:prop:boundary}, and
Section~\ref{shr:bp:sec:further} are new. No other statement, proof or
number of the manuscript was changed.
\end{writenote}

\begin{writenote}[Added 5 October 2026, batch~98: notation of Part~II]
Part~II keeps the manuscript's symbols except $d_{\rm R}$, and several of
them clash with Part~\ref{shr:part:one}'s. Within Part~II every symbol has
the meaning below; the right column gives Part~\ref{shr:part:one}'s
meaning, the tempting false reading. Shared without change: $T_m(n)$,
$F_m(z)$, $M_m(n)$, $\alpha_m$, $\Delta$, the hook count $g^\lambda$,
$J_k=\prod_{j<k}j!$ (Part~II uses it at $k=r$), $W_{-1}$, and $\E$.
{\footnotesize
\begin{longtable}{@{}>{\raggedright\arraybackslash}p{0.17\linewidth}>{\raggedright\arraybackslash}p{0.43\linewidth}>{\raggedright\arraybackslash}p{0.34\linewidth}@{}}
\toprule
symbol & meaning in Part~II & \emph{not}: meaning in Part~\ref{shr:part:one}\\
\midrule\endhead
$d_{\rm R}$ & $\binom r2=\binom{m-2}2$, the manuscript's $d$ & $d=\binom m2$\\
$d_s$, $b_s$, $r_s$ & completions, births and active rows among the first $s$ event letters (Section~\ref{shr:bp:sec:binomial}) & ---\\
$r$ & the size of an active block in Sections~\ref{shr:bp:sec:events}--\ref{shr:bp:sec:words}; $m-2$ from Section~\ref{shr:bp:sec:cut} on & the degree of $h_r$, the index of $p_r$\\
$R_m(n)$ & the rare count $\#\{D_1<B_m\}$ \eqref{shr:bp:eq:raredef} & $R_n(x)$ the pair weight; $R_m=m^{-m}$; $R(z)$\\
$H_m(n)$; $H(a)$ & the core count \eqref{shr:bp:eq:coredef}; the indicator $\binom aa$ of $a\ge0$ \eqref{shr:bp:eq:Hdelta} & $H(\lambda)$, Hermite $H_j$, $H_{20}$\\
$\mathcal K_m$; $K_r(u,v)$; $K$ & the rare constant \eqref{shr:bp:eq:rareK}; the chamber kernel; $\sum_ik_i$ in \eqref{shr:bp:eq:factorialmoment} & $K_m$ \eqref{shr:eq:Km}; Krawtchouk $K_j$\\
$b^{\rm R}_{m,j}$; $b^{\rm total}_{5,1}$ & rare coefficients; Part~I's $b_{5,1}=48/5$ & $b_{m,j}$ of \eqref{shr:eq:mainK}\\
$B_i$, $D_i$; $B$, $D$; $\mathsf B$, $\mathsf D$ & birth and completion ranks $a_{i,1}$, $a_{i,n}$; event letters; the birth and completion maps & $B=m(m-1)^2(m-2)/12$; Bernoulli $B_{2r}$ (also in \eqref{shr:bp:eq:outer})\\
$\mathcal C_r(n)$ & the strict chamber \eqref{shr:bp:eq:chamber} & $\mathcal C_n$ the central event\\
$\mathcal I$; $X_i$ & $\{1\le X_i\le n-2\}$; $\operatorname{Bin}(n,U)$ counts mixed over $U\sim\operatorname{Beta}(n,n+1)$ (Section~\ref{shr:bp:sec:beta}) & $\{1\le X_i\le n-1\}$; independent $\operatorname{Bin}(n,p)$\\
$W_n(x)$ & the rare weight \eqref{shr:bp:eq:W} & $R_n(x)$ plays this role\\
$Y_i$; $Z_{i,n}$, $V_n$, $V$ & $(2X_i-n)/\sqrt n$ for the mixed $X_i$; the components of \eqref{shr:bp:eq:scaling} & $Y_i$ the same formula for independent $X_i$; $Z_i$ standardized counts (Section~\ref{shr:sec:GUE})\\
$\eps$; $\mathcal E_\eps$ & $n^{-1/2}$; the formal moment functional \eqref{shr:bp:eq:Qr} & $\eps_j$ random signs; $\mathcal E_z$ \eqref{shr:eq:Ez}\\
$\Psi_r$; $\mathcal Q_r(z)$; $Q_m(n)$ & the interaction factor \eqref{shr:bp:eq:Psidef}; its normalized expectation (written $Q_r(0)=1$ once in \eqref{shr:bp:eq:Qr}); the common-phase count $T_{m,B^mD^m}(n)$ \eqref{shr:bp:eq:two-sector} & ---\\
$A_m(z)$; $\mathcal A_{5,L}$, $\mathcal L_5$ & the Stirling factor \eqref{shr:bp:eq:outer}; truncations and leading term of \eqref{shr:bp:eq:five} & $A=\binom m3$, the event $\mathcal A$; $A_{5,L}$, $L_5$ for $T_5$\\
$S_1$, $S_2$; $S_h(t)$; $\mathcal S$ & $\sum_iy_i$ and $\sum_iy_i^2$ (Section~\ref{shr:bp:sec:coefficients}); bounded-height Dyck series; the event that row $m$ starts after row~1 ends & $S_1=h_1(q)=\sum_{i<j}(Y_i+Y_j)^2$\\
$p_i$, $q_i$; $p$ & rank marks \eqref{shr:bp:eq:marks}; a fixed success probability in \eqref{shr:bp:eq:binomialVdm} & $p$ the threshold, $p_r$ power sums, $q_{ij}$\\
$P_j$; $P_1(y)$ & monic orthogonal polynomials of $\operatorname{Bin}(n,p)$ in $x$; the first coefficient of $\Psi_r$ \eqref{shr:bp:eq:P1} & monic orthogonal $P_j(y)$ of the binomial Vandermonde ensemble (Section~6.2)\\
$\lambda_m$, $\beta_m$, $a$ & $m\log m-\log4$, $(m-2)^2/2$, $\log(y/\mathcal K_m)$ \eqref{shr:bp:eq:inverseparams} & $\lambda=m\log m$, $\beta_j$ recurrence coefficients, $L=\log(y/K)$\\
$L$ & a truncation order; $|v|-|u|$ in Lemma~\ref{shr:bp:lem:kernel} & a truncation order; $\log(y/K)$\\
$h(w)$ & the maximal Dyck height of a word & $h_r(q)$, $h_j$\\
$C_a(n,p)$, $c_k$, $N_k$ & centred binomial moments and series coefficients (Appendix~\ref{shr:bp:app:algorithm}) & $c_{m,j}$, $C(z)$, $C_{n-1}$, $C_{m,L}$\\
$\M$, $\Sym_r$, $\Diag$, $\ell_s$ & signed multinomial \eqref{shr:bp:eq:M}, the symmetric group, the diagonal, event labels \eqref{shr:bp:eq:rank} & ---\\
\bottomrule
\end{longtable}}
\end{writenote}

\section{The question and the scope of the answer}
\label{shr:bp:sec:scope}

\subsection{The OEIS family}
For $m,n\ge1$, let $T_m(n)$ be the number of arrays
$A=(a_{i,j})_{1\le i\le m,\,1\le j\le n}$ containing each of
$1,\ldots,mn$ once, with increasing rows, columns, diagonals, and downward
antidiagonals. Equivalently, the following comparisons suffice whenever
the cells exist:
\begin{equation}
 a_{i,j}<a_{i,j+1},\qquad
 a_{i,j}<a_{i+1,j},\qquad
 a_{i,j+1}<a_{i+1,j}.
 \label{shr:bp:eq:poset}
\end{equation}
The diagonal comparison follows from row and column increase.
We set $T_m(0)=1$. This is the family A181196; the height-three, height-four,
and height-five rows are A181197, A181198, and A181199
\cite{oeis196,oeis197,oeis198,oeis199}.
The direction of the limit matters: \emph{$m$ is fixed and $n$ grows}.
A fixed-width strip theorem with a growing number of rows is a different
statement.

Kauers and Koutschan discuss the height-four and height-five problems in
Section~6.4 of their 2023 paper \cite{KK}. The OEIS entries inspected for
this report still label their displayed recurrences as conjectural:
order two and degree nine for A181198, and order three and degree twenty-four
for A181199. The same paper gives related finite-sum conjectures. These
specific formulas are not used anywhere in our proofs.

\subsection{What the preceding repository report already established}
The report \emph{Fixed-height shifted rectangles} in ProveIt
(Part~\ref{shr:part:one} of the present report) obtains an all-order expansion for $T_m(n)$, including
\begin{equation}
 T_m(n)\sim K_m m^{mn}n^{-\alpha_m},\qquad
 \alpha_m=\frac{m^2-1}{2},\qquad
 K_m=\frac{\sqrt m\prod_{j=0}^{m-1}j!}
 {4^{\binom m2}(2\pi)^{(m-1)/2}}.
 \label{shr:bp:eq:priorleading}
\end{equation}
It proves the leading asymptotic conjecture for A181199 without assuming
its recurrence, and proves nonalgebraicity at every height $m\ge3$.
Its Question~7 asks whether the generating function is D-finite for every
fixed height. Its coarse exponentially small boundary bound is explicitly
not claimed sharp. These are the two gaps addressed here.

The repository snapshot used for the comparison is
\begin{center}\small\ttfamily
db20eb37982abc0f163c6d308d392f5ac4c9bc62.
\end{center}
The source/status file in the package records the exact report path,
inspected literature, and limited identifier searches. We do not infer
historical novelty merely from an OEIS conjecture label or a repository
question.

\subsection{New statements developed in this manuscript}
Our exact-enumeration result is stronger than existence of an unspecified
recurrence.
\begin{theorem}[Rational diagonality at every fixed height]
\label{shr:bp:thm:diagonal-main}
For every fixed integer $m\ge1$,
\[
 F_m(z)=\sum_{n\ge0}T_m(n)z^n
\]
is the diagonal of a rational power series over $\Q$. In particular it is
D-finite, and $(T_m(n))_{n\ge0}$ is P-recursive. More generally, each fixed
ordering of row births and row completions has a rational-diagonal generating
function, after assigning arbitrary rational values at $n=0,1$.
\end{theorem}

The proof, given in Sections~\ref{shr:bp:sec:events}--\ref{shr:bp:sec:binomial}, is a
finite construction from chamber kernels. It provides at most $m(m-1)$
free integer summation indices for each event ordering. This count concerns
summation variables, not the number of variables in a minimal rational
function or the order of a minimal differential equation.

For $m\ge2$, write
\begin{equation}
 B_i=a_{i,1},\qquad D_i=a_{i,n},\qquad
 R_m(n)=\#\{A:D_1<B_m\},\quad n\ge1,
 \label{shr:bp:eq:raredef}
\end{equation}
and set $R_m(0)=0$. Thus $R_m(n)$ counts arrays for which the first row
finishes before the last row starts. When $n\ge2$, its complement has all
rows active simultaneously at some intermediate stage.
Put
\begin{equation}
 r=m-2,\qquad d_{\rm R}=\binom r2,\qquad J_r=\prod_{j=0}^{r-1}j!,
 \qquad J_0=1.
 \label{shr:bp:eq:rnotation}
\end{equation}
Empty products are one, throughout. In the chamber formulas of
Sections~\ref{shr:bp:sec:events}--\ref{shr:bp:sec:words}, $r$ may instead denote a general
active-block size; from Section~\ref{shr:bp:sec:cut} onward it is $m-2$.

\begin{theorem}[Sharp rare-boundary expansion]
\label{shr:bp:thm:rare-main}
For each fixed $m\ge2$ and each fixed $L\ge0$,
\begin{equation}
 R_m(n)=\mathcal K_m\left(\frac{m^m}{4}\right)^n n^{-r^2/2}
 \left(\sum_{j=0}^{L}\frac{b_{m,j}^{\rm R}}{n^j}
       +O_{m,L}(n^{-L-1})\right),
 \label{shr:bp:eq:rareall}
\end{equation}
where
\begin{equation}
 \mathcal K_m=
 \sqrt{\frac m2}\,(2\pi)^{-r/2}4^{-d_{\rm R}}9^{-r}J_r,
 \qquad b_{m,0}^{\rm R}=1,
 \label{shr:bp:eq:rareK}
\end{equation}
and all $b_{m,j}^{\rm R}$ are effectively computable rational numbers.
The first correction is
\begin{equation}
 b_{m,1}^{\rm R}
 =\frac{r(6r^3-52r^2-268r-283)}{72(r+2)}.
 \label{shr:bp:eq:rare-b1}
\end{equation}
At $m=2$, the assertion reduces to the exact identity $R_2(n)=1$.
\end{theorem}

The proof occupies Sections~\ref{shr:bp:sec:cut}--\ref{shr:bp:sec:coefficients}.
It is independent of \eqref{shr:bp:eq:priorleading}. The latter is used only when
we subsequently divide $R_m(n)$ by $T_m(n)$.

\begin{resultbox}
\textbf{What has and has not been resolved.}
The fixed-height D-finiteness question is answered affirmatively by an
exact formula, and a sharp rare-event sector is obtained at every height.
The published candidate low-order recurrences and finite sums remain
unproved in this manuscript. Our all-order series are Poincar\'e expansions
with fixed height and fixed truncation order. No Borel summability,
Stokes constants, or complete exponentially improved transseries is claimed.
\end{resultbox}

\begin{remark}[\textbf{[write]} Research question~7 of Part~I is answered in part]
\label{shr:bp:rem:q7}
Added 5 October 2026, batch~98. The abstract printed at the head of this
Part says that Theorem~\ref{shr:bp:thm:diagonal-main} resolves ``an
explicit question in the preceding ProveIt report'', and the delivered
README says ``This resolves Question~7''. That overstates it. Research
question~7 of Part~\ref{shr:part:one} asks two things: whether $F_m$ is
D-finite for every fixed $m$, \emph{and} ``how does the minimal annihilating
order grow with $m$?''. Theorem~\ref{shr:bp:thm:diagonal-main} answers the
first affirmatively. It says nothing about the second: the construction
gives at most $m(m-1)$ summation indices, which the text above states is
not a bound on the order of a minimal differential equation,
Remark~\ref{shr:bp:rem:constructive} says that no minimal operator is
produced, and the manuscript's own fourth research question (here Research
question~11) asks for exactly this growth. The result box above (``The
fixed-height D-finiteness question is answered affirmatively'') is
accurate. The second half of Research question~7 stays open; a dated note
at that question says so.
\end{remark}

\begin{remark}[\textbf{[write]} Research question~4 of Part~I is not settled by $R_m$]
\label{shr:bp:rem:q4}
Added 5 October 2026, batch~98. Section~\ref{shr:bp:sec:scope} lists the
non-sharp boundary bound of Part~\ref{shr:part:one} among ``the two gaps
addressed here''. The sharp question there, Research question~4, asks for
the asymptotic of the boundary error $E_{m,n}(1/2)$ of
Theorem~\ref{shr:thm:cut}. Part~II does not determine it. Its rare count
is a different and much smaller sector. In the order-statistics model of
Section~4.1, the rare arrays have every row-$1$ sample below every row-$m$
sample, so at the cut $p=1/2$ either $X_1=n$ or $X_m=0$; hence the rare
event lies inside the boundary event and
\[
 R_m(n)\le M_m(n)\,E_{m,n}(1/2)
\]
(Proposition~\ref{shr:bp:prop:boundary}(c), which also gives the growth
of the ratio). Exact values, computed at intake from the definitions
(row-state dynamic program for $T_m(n)$ and $R_m(n)$, the rational sum
\eqref{shr:eq:interiorsum} for $I_m(n)$, and
$M_m(n)E_{m,n}(1/2)=T_m(n)-M_m(n)I_m(n)$):
\begin{center}\small
\begin{tabular}{@{}rrrrr@{}}
\toprule
$m$ & $n$ & $M_m(n)E_{m,n}(1/2)/R_m(n)$ & $2^nE_{m,n}(1/2)$ &
$E_{m,n}(1/2)/\bigl(2^{1-n}3^{1-m}J_{m-1}(4n)^{-\binom{m-1}2}\bigr)$\\
\midrule
3 & 8 & $4.892$ & $7.165\cdot10^{-3}$ & $1.0317$\\
3 & 12 & $37.20$ & $4.835\cdot10^{-3}$ & $1.0444$\\
3 & 16 & $357.5$ & $3.594\cdot10^{-3}$ & $1.0351$\\
3 & 20 & $3907$ & $2.855\cdot10^{-3}$ & $1.0276$\\
3 & 30 & $2.053\cdot10^{6}$ & $1.885\cdot10^{-3}$ & $1.0178$\\
3 & 40 & $1.328\cdot10^{9}$ & $1.407\cdot10^{-3}$ & $1.0131$\\
2 & 20 & $8.962\cdot10^{4}$ & $0.6817$ & $1.0226$\\
2 & 80 & $5.104\cdot10^{22}$ & $0.6704$ & $1.0056$\\
4 & 12 & $7.183$ & $1.381\cdot10^{-6}$ & $1.0311$\\
4 & 20 & $374.2$ & $2.995\cdot10^{-7}$ & $1.0350$\\
\bottomrule
\end{tabular}
\end{center}
At $m=3$ the boundary error is $4.9$, $37$, $357$ and $3907$ times the rare
count at $n=8$, $12$, $16$, $20$, and the factor keeps growing roughly like
$2^n$. As a probability the boundary error is $2^{-n}$ times a power of
$n$ (second column), while $R_m(n)/M_m(n)\le\binom{2n}n^{-1}$ is at most
of order $4^{-n}$ by \eqref{shr:bp:eq:rare-upper}. So Theorem~\ref{shr:bp:thm:rare-main} is a
sharp theorem about a sub-sector of the boundary event, not the sharp
boundary asymptotic; Research question~4 stays open, re-scoped in a dated
note there. The write's Proposition~\ref{shr:bp:prop:boundary} supplies
its leading term (last column). The program used for the table is the
intake's and is not shipped; the same numbers follow from the shipped
enumerator \texttt{count} of \texttt{code/\allowbreak 02-boundary-phases-tableaux.py}
and the sum \eqref{shr:eq:interiorsum}.
\end{remark}

\section{Boundary events and the exact walk model}
\label{shr:bp:sec:events}

\subsection{Prefix states, including both exceptional boundaries}
Insert the labels $1,2,\ldots,mn$ in order, and let $x_i$ be the number of
filled cells in row $i$. Row increase makes the filled cells in each row
a prefix. The state space is exactly
\begin{equation}
 n\ge x_1\ge x_2\ge\cdots\ge x_m\ge0,
 \qquad x_i=x_{i+1}\Longrightarrow x_i\in\{0,n\}.
 \label{shr:bp:eq:states}
\end{equation}
Indeed, if a lower row has $j<n$ filled cells, its last filled cell forces
cell $j+1$ in the row above to be filled. If the lower row is full, column
increase forces the upper row to be full. Conversely these conditions
make the filled set an order ideal of \eqref{shr:bp:eq:poset}. Consequently paths
from $(0,\ldots,0)$ to $(n,\ldots,n)$ through \eqref{shr:bp:eq:states}, using unit
steps $+e_i$, are in bijection with the arrays.

A row is \emph{unborn}, \emph{active}, or \emph{completed} according as
its coordinate is $0$, belongs to $\{1,\ldots,n-1\}$, or is $n$.
Completed rows form a prefix and unborn rows a suffix. Thus the active rows
always form one consecutive block, and their coordinates are strictly
decreasing. The ties at zero and at $n$ cannot be discarded: forbidding
either kind changes the enumeration.

For $n\ge2$, birth and completion are different unit steps. Births occur
in row order $1,\ldots,m$, and completions also occur in that order.
Record a letter $B$ for a birth and $D$ for a completion. Every array gives
a word $w$ of length $2m$ with $m$ letters of each kind and with at least
as many $B$'s as $D$'s in every prefix. Such a word is a Dyck word, with
$B$ interpreted as an up-step. Conversely the $i$th $B$ means the birth
of row $i$, and the $i$th $D$ means the completion of row $i$; no further
row-labeling data are needed.

\subsection{The strict chamber between two events}
For $r\ge0$ define
\begin{equation}
 \mathcal C_r(n)=\{(u_1,\ldots,u_r):n-1\ge u_1>\cdots>u_r\ge1\},
 \qquad \mathcal C_0(n)=\{()\}.
 \label{shr:bp:eq:chamber}
\end{equation}
During an interval with no birth or completion, its $r$ active coordinates
move by $+e_i$ in this chamber. If $u,v\in\mathcal C_r(n)$, write
$K_r(u,v)$ for the number of such walks from $u$ to $v$, including a
zero-length walk when $u=v$.

For an integer vector $a=(a_1,\ldots,a_r)$ define the everywhere-defined
integer multinomial
\begin{equation}
 \M(a)=
 \begin{cases}
 (a_1+\cdots+a_r)!/(a_1!\cdots a_r!),&a_i\ge0\ \text{for all }i,\\
 0,&\text{otherwise},
 \end{cases}
 \qquad \M(())=1.
 \label{shr:bp:eq:M}
\end{equation}

\begin{lemma}[Reflection kernel]
\label{shr:bp:lem:kernel}
For strict endpoints $u,v\in\mathcal C_r(n)$,
\begin{equation}
 K_r(u,v)=\sum_{\sigma\in\Sym_r}\sgn(\sigma)
 \M(v_1-u_{\sigma(1)},\ldots,v_r-u_{\sigma(r)}).
 \label{shr:bp:eq:kernel}
\end{equation}
Equivalently, when $L=|v|-|u|\ge0$,
\begin{equation}
 K_r(u,v)=L!\det_{1\le i,j\le r}\frac1{(v_i-u_j)!},
 \label{shr:bp:eq:detkernel}
\end{equation}
where $1/k!=0$ for negative integers $k$. For $r=0$, the value is one.
\end{lemma}

\begin{proof}
Consider the signed family of all unrestricted $+e_i$ walks from a
permutation of $u$ to $v$. The number starting at
$(u_{\sigma(1)},\ldots,u_{\sigma(r)})$ is the multinomial in
\eqref{shr:bp:eq:kernel}. Cancel every walk that has a collision of two
coordinates: choose its earliest collision time, breaking any tie by a
fixed ordering of coordinate pairs, and exchange the histories of the
colliding coordinates up to that time. The starting permutation is changed
by a transposition, the endpoint and step set are unchanged, and the
operation is an involution. Before the chosen time, all equality relations
are preserved under this exchange, so the same collision is selected on
the second application.

A walk remaining after cancellation has no collisions. Its coordinate
order cannot change, since a unit step cannot cross an integer coordinate
without meeting it. Its strictly decreasing endpoint forces its starting
permutation to be the identity, and it stays in the decreasing chamber.
All coordinates are between $1$ and $n-1$: along an uncancelled walk each
coordinate is nondecreasing and lies between its own endpoints. Thus no
extra wall reflection is needed for the finite upper bound. The surviving
walks are exactly those counted by $K_r(u,v)$.
\end{proof}

The proof is the classical reflection argument, included to specify
precisely which walls are present. Formula~\eqref{shr:bp:eq:kernel}, rather than
an unqualified factorial quotient, is the appropriate form for subsequent
binomial-sum manipulations: it is defined even when an individual
displacement is negative.

\subsection{Birth and completion maps}
If $v\in\mathcal C_r(n)$, a birth changes the active vector by
\begin{equation}
 \mathsf B(v)=(v_1,\ldots,v_r,1),
 \quad\text{allowed if }r=0\text{ or }v_r\ge2.
 \label{shr:bp:eq:Bmap}
\end{equation}
A completion changes it by
\begin{equation}
 \mathsf D(v)=(v_2,\ldots,v_r),
 \quad\text{allowed if }r\ge1\text{ and }v_1=n-1.
 \label{shr:bp:eq:Dmap}
\end{equation}
The completion step increases the old first active coordinate from $n-1$
to $n$ before removing it. All previously completed coordinates are already
$n$. These two maps account exactly for the exceptional boundary ties.

\section{Finite binomial sums and rational diagonals}
\label{shr:bp:sec:binomial}

\subsection{The event-skeleton formula}
Fix a Dyck word $w=w_1\cdots w_{2m}$. Let $b_s,d_s$ be the numbers of
births and completions in its first $s$ letters, and set
$r_s=b_s-d_s$, for $0\le s\le2m$. Let $u_s$ be the active vector
immediately after event $s$, and $v_s$ the active vector immediately
before event $s+1$. The initial vector is $u_0=()$.
A skeleton consists of vectors $v_s\in\mathcal C_{r_s}(n)$ with
\begin{equation}
 u_{s+1}=
 \begin{cases}
 \mathsf B(v_s),&w_{s+1}=B,\\
 \mathsf D(v_s),&w_{s+1}=D,
 \end{cases}
 \qquad 0\le s<2m,
 \label{shr:bp:eq:skeletonmap}
\end{equation}
subject to the corresponding legality conditions. All $u_s$ are thereby
affine expressions in earlier endpoint coordinates and the constants
$1,n$. Define $T_{m,w}(n)$ to count arrays with this event word.

\begin{theorem}[Exact finite-phase enumeration]
\label{shr:bp:thm:skeleton}
For $n\ge2$,
\begin{equation}
 T_{m,w}(n)=
 \sum_{\text{legal skeletons for }w}
 \prod_{s=0}^{2m-1}K_{r_s}(u_s,v_s),
 \qquad
 T_m(n)=\sum_{w\in\mathcal D_m}T_{m,w}(n),
 \label{shr:bp:eq:skeleton}
\end{equation}
where $\mathcal D_m$ denotes the Dyck words with $m$ births.
Each first sum can be written with at most $m(m-1)$ free integer indices,
each ranging between $0$ and $n$, and with only affine equality and
inequality constraints.
\end{theorem}

\begin{proof}
Cut a valid prefix path immediately before and after every boundary event.
The intervening pieces are chamber walks with the stated endpoints.
Conversely, choosing a legal skeleton and one chamber walk in each interval
reconstructs a unique full prefix path. The events each contribute a
single unit step, and the intervals concatenate in a prescribed order;
there is no additional interleaving factor. This proves the product and sum.

There are initially $\sum_{s=0}^{2m-1}r_s$ endpoint slots. This sum is at
most $m^2$: at each position $s$ the height is at most
$\min(s,2m-s)$, with equality throughout for $B^mD^m$. Each of the $m$
completion events fixes one different endpoint coordinate to $n-1$.
Eliminating these coordinates leaves at most $m^2-m$ free slots. The
chamber bounds, strict gaps, birth conditions, and all endpoint
substitutions are affine in $n$ and the remaining coordinates.
\end{proof}

As a consistency check, every skeleton with nonzero weight satisfies
\begin{equation}
 \sum_{s=0}^{2m-1}(|v_s|-|u_s|)=m(n-2).
 \label{shr:bp:eq:unmarkedsteps}
\end{equation}
Each of the $m$ rows has precisely $n-2$ non-event steps. Equivalently,
telescoping the active sums adds one at a birth and removes $n-1$ at a
completion, which gives the same identity.

\subsection{An explicit binomial encoding of every constraint}
For integers $a,b$, use
$\binom ab=[t^b](1+t)^a$ when $b\ge0$, and zero when $b<0$.
In particular,
\begin{equation}
 H(a)=\binom aa=\ind_{\{a\ge0\}},\qquad
 \delta(a)=H(a)H(-a)=\ind_{\{a=0\}}.
 \label{shr:bp:eq:Hdelta}
\end{equation}
Every affine inequality or equality in \eqref{shr:bp:eq:skeleton} is therefore
a product of such factors. For example, strictness $v_i>v_{i+1}$ is
$H(v_i-v_{i+1}-1)$, and a completion uses $\delta(v_1-n+1)$.
Moreover, for $r\ge1$,
\begin{equation}
 \M(a_1,\ldots,a_r)=
 \left(\prod_{i=1}^r H(a_i)\right)
 \prod_{i=1}^{r-1}\binom{a_i+\cdots+a_r}{a_i}.
 \label{shr:bp:eq:binomialM}
\end{equation}
If a displacement is negative the leading indicator makes the term zero;
all other factors remain defined. If all displacements are nonnegative,
the binomial product telescopes to the multinomial.

Expand the fixed-size permutation sums in \eqref{shr:bp:eq:kernel}, replace each
multinomial by \eqref{shr:bp:eq:binomialM}, sum each free coordinate from $0$ to
$n$, and insert the legality factors. The result is a finite linear
combination of bounded multiple sums of products of binomial coefficients
with affine arguments. There are no reciprocal binomial coefficients,
no factorial poles, and no hidden variable number of summation indices.

\subsection{The imported diagonal theorem and its application}
We use the following result of Bostan, Lairez, and Salvy
\cite[Theorem~3.5 and Corollary~3.6]{BLS}: a univariate multiple binomial
sum has a generating function that is the diagonal of a rational power
series; such sequences are P-recursive. Their class is closed under affine
index substitutions, products, sums, and finite partial summation.

For clarity, consider a rational function regular at the origin, with
expansion
\[
 Q(x_1,\ldots,x_k)=\sum q_{n_1,\ldots,n_k}x_1^{n_1}\cdots x_k^{n_k}.
\]
Its diagonal is
\[
 \Diag Q(z)=\sum_{n\ge0}q_{n,\ldots,n}z^n.
\]
A power series is D-finite if it satisfies a nonzero linear differential
equation with polynomial coefficients. Equivalently its coefficient
sequence satisfies a linear recurrence with polynomial coefficients.

\begin{proof}[Proof of Theorem~\ref{shr:bp:thm:diagonal-main}]
The preceding encoding exhibits $T_{m,w}(n)$, for $n\ge2$, as a multiple
binomial sum. To define a sequence on all integers, multiply by $H(n-2)$
and assign zero below two. The sums from $0$ to $n$ can first be interpreted
as the extended partial sums of the binomial-sum algebra; the multiplier
removes their values outside the intended range. Adding any finitely
supported correction preserves that algebra, since shifted Kronecker
deltas are allowed. The imported theorem gives a rational diagonal for
each event-word sequence. Summing the finitely many words and adding the
terms at $n=0,1$ gives $F_m$.

This argument proves rational diagonality, not merely a heuristic closure
claim about hypergeometric expressions. D-finiteness and P-recursiveness
then follow from the same theorem. At $m=1$, one may also observe directly
that $F_1(z)=1/(1-z)$.
\end{proof}

\begin{remark}[Constructive does not mean a compact telescoper is supplied]
\label{shr:bp:rem:constructive}
The skeleton and binomial formulas are explicit finite input to the
rational-integral and diagonal constructions of \cite{BLS}. This is an
effective existence proof. We have not carried out that elimination to
produce a small rational function or a minimal differential operator for
general $m$. In particular, it does not certify that either operator
currently displayed in OEIS is the correct one.
\end{remark}

\begin{writenote}[Added 5 October 2026, batch~101]
The last sentence of the remark concerns this existence proof only.
Parts~\ref{shr:h4:part} and~\ref{shr:h5:part} certify both displayed
operators by another method (Theorems~\ref{shr:h4:thm:oeis}
and~\ref{shr:h5:thm:operator}), from Pfaffian formulas for the transfer
integral rather than from the event sum. Report~91, printed in
Part~\ref{shr:sa:part}, proves Theorem~\ref{shr:bp:thm:diagonal-main} by the
same route as this section (Theorem~\ref{shr:sd:thm:diagonal}).
\end{writenote}

Combined with the previous report's algebraicity classification
(Part~\ref{shr:part:one}, Theorem~\ref{shr:thm:algebraic}), the result places the entire family in a precise hierarchy:
$F_m$ is algebraic for $m\le2$, and transcendental but D-finite for
$m\ge3$. The nonalgebraicity part of that statement is inherited from the
previous report, not reproved or claimed as new here.

\section{Which event words occur, and which statistics remain D-finite?}
\label{shr:bp:sec:words}

\subsection{A complete realizability criterion}
Let $h(w)$ be the maximum Dyck height of $w$.
\begin{theorem}[Event-word realizability]
\label{shr:bp:thm:realizability}
For $m\ge1$, $n\ge2$, and $w\in\mathcal D_m$,
\begin{equation}
 T_{m,w}(n)>0\quad\Longleftrightarrow\quad h(w)\le n-1.
 \label{shr:bp:eq:realizability}
\end{equation}
Consequently every Dyck word of semilength $m$ occurs as soon as $n\ge m+1$.
\end{theorem}
\begin{proof}
At height $r$ there must be $r$ distinct active coordinates in
$\{1,\ldots,n-1\}$, so $r\le n-1$ is necessary.

For sufficiency, follow the prescribed word. Before a birth at active
height $r$, the height bound implies $r\le n-2$. Move the existing
coordinates, in top-to-bottom order, to
$(n-1,n-2,\ldots,n-r)$. Each target dominates the present coordinate, since
any decreasing vector in $\mathcal C_r(n)$ has $u_i\le n-i$. The moves
are legal because the coordinate above has already reached its target.
The last target is at least two, so a new coordinate one can be appended.
Before a completion, simply move the top active coordinate to $n-1$ and
complete that row. Neither operation creates any additional boundary
event. Continuing to the end constructs a valid path with exactly $w$.
\end{proof}

The number of event words available at width $n$ is thus a classical
bounded-height Dyck number. If $S_h(t)$ is its generating function by
semilength with height at most $h$, the first-return decomposition gives
\begin{equation}
 S_0(t)=1,\qquad S_h(t)=\frac1{1-tS_{h-1}(t)}\quad(h\ge1).
 \label{shr:bp:eq:boundedDyck}
\end{equation}
This counts possible event types, not arrays: the sector weights are
generally far from equal.

At $n=2$ there is only the alternating word $(BD)^m$, and $T_m(2)=1$.
At $n=3$, height is at most two. A one-row active phase has only one
coordinate to move, while a two-row active phase has the unique vector
$(2,1)$. Thus each permitted word has weight one. Formula
\eqref{shr:bp:eq:boundedDyck} gives $T_m(3)=2^{m-1}$ for $m\ge1$, recovering a
known small-width formula \cite{Sun}, not a new one.

\subsection{Small worked examples}
For two rows and $n\ge2$, the alternating word has weight one. For the
other word,
\begin{align}
 T_{2,BBDD}(n)
 &=\sum_{a=2}^{n-1}\sum_{b=1}^{n-2}
 K_2\bigl((a,1),(n-1,b)\bigr)\notag\\
 &=\sum_{a=2}^{n-1}\sum_{b=1}^{n-2}
 \bigl(\M(n-1-a,b-1)-\M(n-2,b-a)\bigr).
 \label{shr:bp:eq:two-row}
\end{align}
This is an explicit binomial sum without exceptional denominator values.
Adding one recovers $T_2(n)=C_{n-1}$, the usual Catalan control.

At $m=3,n=4$, the sector weights are
\begin{center}
\begin{tabular}{@{}lrrrrr@{}}
\toprule
Event word & $BBBDDD$ & $BBDBDD$ & $BBDDBD$ & $BDBBDD$ & $BDBDBD$\\
\midrule
Number of arrays & $4$ & $16$ & $4$ & $4$ & $1$\\
\bottomrule
\end{tabular}
\end{center}
Their sum is $29$, and the rare subcount is $25$. This also illustrates
that having a full-height phase is possible but not yet typical at small
width.

\subsection{Birth and completion ranks}
In a skeleton, immediately before event $s+1$ there are $d_s$ completed
rows and $|v_s|$ cells in active rows. Hence its label is the affine form
\begin{equation}
 \ell_{s+1}=d_s n+|v_s|+1.
 \label{shr:bp:eq:rank}
\end{equation}
Let $p_1,\ldots,p_m,q_1,\ldots,q_m$ be independent indeterminates and define
\begin{equation}
 T_m(n;\mathbf p,\mathbf q)
 =\sum_A\prod_{i=1}^m p_i^{B_i}q_i^{D_i}.
 \label{shr:bp:eq:marks}
\end{equation}
The birth and completion ranks are labels, not column positions.

\begin{corollary}[Marked enumeration and moment totals]
\label{shr:bp:cor:marks}
For fixed $m$, the generating function in $n$ of
$T_m(n;\mathbf p,\mathbf q)$ is a rational diagonal over
$\Q(\mathbf p,\mathbf q)$. For every fixed polynomial $P$ with rational coefficients in the
$2m$ event ranks,
\[
 n\longmapsto\sum_A P(B_1,\ldots,B_m,D_1,\ldots,D_m)
\]
is a multiple binomial sum and is P-recursive. The same statements hold
within each fixed event sector.
\end{corollary}
\begin{proof}
Insert the weights from \eqref{shr:bp:eq:rank} into \eqref{shr:bp:eq:skeleton}.
They are geometric sequences with affine exponents over the indicated
coefficient field. Polynomial weights are also in the binomial-sum
algebra, since an integer index is $\binom{k}{1}$. The same closure proof
applies. When $n=1$, use the appropriate finitely supported marked term;
birth and completion coincide within a row only in this exceptional case.
\end{proof}

This proves the assertion for \emph{unnormalized moment totals}.
An expectation divides by $T_m(n)$; P-recursiveness of such quotients is
not asserted.
Rotation of the array by $180$ degrees followed by replacement of each
label $a$ by $mn+1-a$ gives the useful exact symmetry
\begin{equation}
 B_i\longleftrightarrow mn+1-D_{m+1-i},\qquad
 T_{m,w}(n)=T_{m,\overline{w^{\rm rev}}}(n),
 \label{shr:bp:eq:dual}
\end{equation}
where the bar interchanges $B$ and $D$.

\section{A first-completion cut for the rare sector}
\label{shr:bp:sec:cut}

\subsection{The dominant rare configuration}
For $m\ge2$ and $n\ge2$, define the core subcount
\begin{equation}
 H_m(n)=\#\{A:B_{m-1}<D_1<B_m\}.
 \label{shr:bp:eq:coredef}
\end{equation}
For $m=2$ this equals $R_2(n)=1$. For $m\ge3$, the remainder has
$D_1<B_{m-1}$: the first row finishes while at least the last two rows
are still unborn. This extra separation will give an exponentially
smaller rate.

The classical shifted hook product for a strict positive partition
$\lambda=(\lambda_1>\cdots>\lambda_k>0)$ is
\begin{equation}
 g^\lambda=
 \frac{|\lambda|!}{\prod_i\lambda_i!}
 \prod_{i<j}\frac{\lambda_i-\lambda_j}{\lambda_i+\lambda_j}.
 \label{shr:bp:eq:hook}
\end{equation}
It counts chains of strict partitions obtained by adding one cell at a
time, equivalently standard shifted tableaux; see \cite{Sun}. This is an
imported classical enumeration formula, not a new hook theorem.

\begin{proposition}[Exact core sum]
\label{shr:bp:prop:core}
With $r=m-2$, for all $m,n\ge2$,
\begin{equation}
 H_m(n)=
 \sum_{n-2\ge x_1>\cdots>x_r\ge1}
 g^{(n-1,x_1,\ldots,x_r)}
 g^{(n,n-x_r,\ldots,n-x_1)}.
 \label{shr:bp:eq:hookcore}
\end{equation}
For $r=0$ the sum has one term, equal to one.
\end{proposition}
\begin{proof}
Immediately before the first completion, row $1$ has length $n-1$ and
the last row is empty. In the core sector every intervening row is
active, so their lengths are $n-2\ge x_1>\cdots>x_r\ge1$.
Before this point no row has completed. Removing zero rows, the prefix
path is a chain of strict partitions ending at $(n-1,x_1,\ldots,x_r)$,
so it has the first hook count in \eqref{shr:bp:eq:hookcore}.

The next step completes row $1$. Reverse all remaining steps and rotate
rows; the complementary lengths are $(n,n-x_r,\ldots,n-x_1)$, a strict
positive partition. Every reversed prefix stays in the strict-partition
region until the endpoint. In particular, a coordinate equal to $n$ at
the endpoint causes no repeated completed-row boundary, since all other
parts are strictly smaller. Thus the suffix has the second hook count.
Their label sets are fixed by the size of the prefix, and the intervening
completion step is unique. The two chains concatenate bijectively to the
core arrays.
\end{proof}

\subsection{Two useful nonasymptotic error bounds}
Write
\begin{equation}
 M_m(n)=\frac{(mn)!}{(n!)^m}.
 \label{shr:bp:eq:Mtotal}
\end{equation}
This counts all interleavings of $m$ individually sorted rows, before
imposing cross-row comparisons.

\begin{proposition}[Separation bounds]
\label{shr:bp:prop:separation}
For $m\ge2$ and $n\ge1$,
\begin{equation}
 R_m(n)\le\frac{M_m(n)}{\binom{2n}{n}}.
 \label{shr:bp:eq:rare-upper}
\end{equation}
For $m\ge3$ and $n\ge2$,
\begin{equation}
 0\le R_m(n)-H_m(n)\le\frac{M_m(n)}{\binom{3n}{n}}.
 \label{shr:bp:eq:core-error}
\end{equation}
For $m=2$ the difference is zero.
\end{proposition}
\begin{proof}
In a uniform row interleaving, restrict attention to rows $1$ and $m$.
Of their $\binom{2n}{n}$ possible relative words, exactly one places all
of row $1$ before all of row $m$. Dropping other comparisons bounds the
rare sector and proves \eqref{shr:bp:eq:rare-upper}.

For the remainder, rows $m-1$ and $m$ must both start after row $1$ has
finished. Among the relative interleavings of these three rows, the
probability that the first $n$ letters all come from row $1$ is
\[
 \frac{(2n)!/(n!)^2}{(3n)!/(n!)^3}=\frac1{\binom{3n}{n}}.
\]
Again dropping all other constraints gives an upper bound.
\end{proof}

These estimates have explicit constants and hold at finite $n$. Once
the core asymptotic is known, \eqref{shr:bp:eq:core-error} is smaller than the
core by a polynomial factor times $(16/27)^n$. It is therefore negligible
at every inverse-power order on the rare scale. We do not claim that
\eqref{shr:bp:eq:core-error} gives the sharp next exponential sector.

\subsection{[write] The boundary error of Part~I's central cut}
\label{shr:bp:sec:boundary}

\begin{writenote}[Added 5 October 2026, batch~98]
This subsection is the write's, not the manuscript's. It gives the leading
asymptotic of the boundary error asked for in Research question~4 of
Part~\ref{shr:part:one}, and places the rare sector of this Part inside
that error (Remark~\ref{shr:bp:rem:q4}). It uses only
Part~\ref{shr:part:one}'s model and estimates and the hook argument of
Proposition~\ref{shr:bp:prop:core}. Like the rest of the report it is
unrefereed; the exact identity (a) and the asymptotic (b) were checked
against exact values for $m=2,3,4$ (table of Remark~\ref{shr:bp:rem:q4}).

\emph{Dated note (5~October 2026).} An independent adversarial check of
the results added at this write, made by the intake after the write, found
all of them valid: Proposition~\ref{shr:bp:prop:boundary}, parts (a), (b)
and~(c); the table of Remark~\ref{shr:bp:rem:q4}, all ten rows to every
printed digit; and the completed arguments for \eqref{shr:bp:eq:CLT} and
\eqref{shr:bp:eq:tails} in the notes after them. The check was a careful
reading with independent exact computations: $T_m(n)$ by a new prefix-state
recursion, cross-checked against a direct linear-extension count of the
poset for $mn\le22$; $E_{m,n}(1/2)$ in exact rationals from conditional
path counts, without the hook formula; $R_m(n)$ recomputed; 116 exact rows
($m=2,\ldots,5$, $n$ up to $80$ at $m=2$), with floating extensions to
$n=1600$ at $m=3$; and the bound \eqref{shr:bp:eq:tails} tested at every
atom of the beta--binomial law for $n\le400$. In those rows $D_{m,n}$ lies far
inside its bound. It is not a formal verification or an external review.
No fix was needed; the observed rate in (b) is recorded after the proof.
\end{writenote}

\begin{proposition}[\textbf{[write]} Leading boundary error at the central cut]
\label{shr:bp:prop:boundary}
Let $m\ge1$, $n\ge2$, and let $E_{m,n}(1/2)$ be the boundary error of
Theorem~\ref{shr:thm:cut} at $p=1/2$. Let $X'_1,\ldots,X'_{m-1}$ be
independent $\operatorname{Bin}(n,1/2)$ variables,
$\mathcal I'=\{1\le X'_i\le n-1\text{ for all }i\}$, and
\[
 W'_n(x)=\prod_{1\le i<j\le m-1}
 \frac{(x_i-x_j)^2}{(x_i+x_j)(2n-x_i-x_j)}
 \prod_{i=1}^{m-1}\frac{x_i}{2n-x_i}.
\]
\begin{enumerate}[label=(\alph*)]
\item Exactly,
\begin{equation}
 E_{m,n}(1/2)=\frac{2^{1-n}}{(m-1)!}\,
 \E\bigl[\ind_{\mathcal I'}W'_n(X')\bigr]+D_{m,n},
 \qquad 0\le D_{m,n}\le 2m(m-1)\,4^{-n}.
 \label{shr:bp:eq:boundary-exact}
\end{equation}
\item As $n\to\infty$,
\begin{equation}
 E_{m,n}(1/2)\sim 2^{1-n}\,3^{1-m}J_{m-1}\,(4n)^{-\binom{m-1}2},
 \qquad J_{m-1}=\prod_{j=0}^{m-2}j!.
 \label{shr:bp:eq:boundary-lead}
\end{equation}
\item Consequently
$M_m(n)E_{m,n}(1/2)/T_m(n)\sim2^{1-n}3^{1-m}(4n)^{m-1}/(m-1)!$, and, for
$m\ge2$, $R_m(n)\le M_m(n)E_{m,n}(1/2)$ for every $n\ge2$, with
\[
 \frac{M_m(n)E_{m,n}(1/2)}{R_m(n)}\sim
 \frac{2\cdot3^{m-3}(m-2)!}{\sqrt\pi\,4^{m-2}}\;2^n\,n^{-(2m-3)/2}.
\]
\end{enumerate}
\end{proposition}

\begin{proof}
Use the order-statistics model of Section~4.1 at $p=1/2$: $X_i$ counts the
row-$i$ samples below $1/2$, the $X_i$ are independent
$\operatorname{Bin}(n,1/2)$, and the proof of Theorem~\ref{shr:thm:cut}
identifies $E_{m,n}(1/2)$ with $\Pp(\mathcal A\cap\mathcal I^c)$.

(a) On $\mathcal A$ the cells with values below $1/2$ form an order ideal,
so $X$ is a prefix state \eqref{shr:eq:states}: $X_i=0$ forces
$X_{i+1}=\cdots=X_m=0$, and $X_i=n$ forces $X_1=\cdots=X_i=n$. Hence on
$\mathcal A\cap\mathcal I^c$ either exactly one coordinate lies in
$\{0,n\}$, and then it is $X_m=0$ or $X_1=n$ with every other coordinate
in $[1,n-1]$, or at least two coordinates lie in $\{0,n\}$. The second
event, whose probability on $\mathcal A$ we call $D_{m,n}$, has
probability at most $\binom m2(2\cdot2^{-n})^2=2m(m-1)4^{-n}$.

Let $x=(x_1,\ldots,x_{m-1},0)$ with $n-1\ge x_1>\cdots>x_{m-1}\ge1$. As in
Section~4.1, $\Pp(\mathcal A\mid X=x)$ is the product of a prefix and a
suffix probability. Every prefix path to $x$ has last coordinate $0$ and
first coordinate below $n$, so its states are strict partitions and the
prefix count is $g^{(x_1,\ldots,x_{m-1})}$. The rotated complement
\eqref{shr:eq:complement} of $x$ is $(n,n-x_{m-1},\ldots,n-x_1)$, a strict
positive partition in which only the first part can reach $n$; as in the
proof of Proposition~\ref{shr:bp:prop:core}, the suffix count is
$g^{(n,n-x_{m-1},\ldots,n-x_1)}$. Dividing both counts by the multinomial
numbers of interleavings and using \eqref{shr:eq:shiftedhook}, the pairs
among $x_1,\ldots,x_{m-1}$ contribute
$\frac{x_i-x_j}{x_i+x_j}\cdot\frac{x_i-x_j}{2n-x_i-x_j}$ and the pairs of
the part $n$ with the parts $n-x_i$ contribute $\frac{x_i}{2n-x_i}$, so
$\Pp(\mathcal A\mid X=x)=W'_n(x_1,\ldots,x_{m-1})$. The weight $W'_n$ is
symmetric and vanishes at ties; summing over decreasing vectors and
multiplying by $\Pp(X_m=0)=2^{-n}$, the stratum $\{X_m=0\}$ contributes
$2^{-n}\E[\ind_{\mathcal I'}W'_n(X')]/(m-1)!$. Reflecting every sample
$u\mapsto1-u$ and reversing the rows preserves the model and $\mathcal A$
and maps $X$ to $(n-X_m,\ldots,n-X_1)$, so the stratum $\{X_1=n\}$
contributes the same. This proves \eqref{shr:bp:eq:boundary-exact}. (For
$m=1$ there are no pairs, $W'_n=1$, $D_{1,n}=0$, and
$E_{1,n}(1/2)=2^{1-n}$.)

(b) Put $c=\binom{m-1}2$, $Y_i=(2X'_i-n)/\sqrt n$ and
$\mathcal C'_n=\{|Y_i|\le\sqrt n/2\text{ for all }i\}$; then
$\mathcal C'_n\subset\mathcal I'$. On $\mathcal C'_n$, exactly as in
\eqref{shr:eq:scaledweight},
\[
 (4n)^cW'_n(X')=\Delta(Y)^2
 \prod_{i<j}\Bigl(1-\frac{(Y_i+Y_j)^2}{4n}\Bigr)^{-1}
 \prod_{i}\frac{1+Y_i/\sqrt n}{3-Y_i/\sqrt n},
\]
where every pair factor lies in $[1,4/3]$ and every single factor in
$[1/7,3/5]$. The right side is at most $(4/3)^c\Delta(Y)^2$, a uniformly
integrable family because the moments of each fixed degree of the $Y_i$
are bounded uniformly in $n$ (Section~5.4), and it converges to
$3^{1-m}\Delta(y)^2$ uniformly on compact sets. Since $Y$ converges in
distribution to a vector $G$ of independent standard normals,
$(4n)^c\E[\ind_{\mathcal C'_n}W'_n(X')]\to3^{1-m}\E\Delta(G)^2
=3^{1-m}(m-1)!\,J_{m-1}$ by \eqref{shr:eq:gaussZ} in $m-1$ variables.
Off $\mathcal C'_n$, $0\le W'_n\le1$ (Remark~\ref{shr:rem:Rbound} and
$x_i\le n$), and $\Pp(\mathcal C_n'^{\,c})\le2(m-1)e^{-n/8}$ by
\eqref{shr:eq:central-tail}, which is $o(n^{-c})$. Inserting this into
\eqref{shr:bp:eq:boundary-exact}, where $D_{m,n}=O(4^{-n})$ is
$o(2^{-n}n^{-c})$, gives \eqref{shr:bp:eq:boundary-lead}.

(c) The first relation is \eqref{shr:bp:eq:boundary-lead} divided by
$T_m(n)/M_m(n)\sim J_m(4n)^{-\binom m2}$ \eqref{shr:eq:mainM}, since
$J_m=(m-1)!\,J_{m-1}$ and $\binom m2-\binom{m-1}2=m-1$. In the same model
the rare event $\{D_1<B_m\}$ is the event that every row-$1$ sample is
below every row-$m$ sample. If the largest row-$1$ sample is below $1/2$,
then $X_1=n$; otherwise the smallest row-$m$ sample exceeds $1/2$ and
$X_m=0$. Either way $X\notin\mathcal I$, so
$R_m(n)/M_m(n)=\Pp(\mathcal A\cap\{D_1<B_m\})
\le\Pp(\mathcal A\cap\mathcal I^c)=E_{m,n}(1/2)$. The ratio follows from
\eqref{shr:bp:eq:boundary-lead}, Theorem~\ref{shr:bp:thm:rare-main} and
$M_m(n)\sim\sqrt m(2\pi n)^{(1-m)/2}m^{mn}$ (Section~5.5): the powers of
$m^m$ cancel, $\sqrt m/\sqrt{m/2}=\sqrt2$,
$(2\pi)^{(1-m)/2+(m-2)/2}=(2\pi)^{-1/2}$, $3^{1-m}9^{m-2}=3^{m-3}$,
$J_{m-1}/J_{m-2}=(m-2)!$, $4^{-\binom{m-1}2+\binom{m-2}2}=4^{2-m}$, and the
power of $n$ is $-\binom{m-1}2+\frac{1-m}2+\frac{(m-2)^2}2=-\frac{2m-3}2$.
\end{proof}

At $m=2$, (c) reads $\binom{2n}nE_{2,n}(1/2)\sim\frac{2}{3\sqrt\pi}2^n
n^{-1/2}$, since $R_2(n)=1$.

\emph{The rate in (b): an observation (dated 5~October 2026).}
Part~(b) states no rate. The independent check recorded at the head of
this subsection observed numerically that the ratio of $E_{m,n}(1/2)$ to
the right side of \eqref{shr:bp:eq:boundary-lead} behaves like
$1+c_m/n+O(n^{-2})$, with no
$n^{-1/2}$ term (consistent with the symmetry $Y\mapsto-Y$):
$n$ times the excess tends to about $0.500$ at $m=3$ (to $n=1600$), to
about $0.667$ at $m=4$ (to $n=400$), and is still settling near $1.42$ at
$m=5$ (to $n=100$). For $m=2$ the constant is $c_2=4/9$, by a short
calculation. There are no pairs and $J_1=1$, so by
\eqref{shr:bp:eq:boundary-exact} the ratio is
$3\E[\ind_{\mathcal I'}X'_1/(2n-X'_1)]+O(2^{-n})$. The excluded atom
$X'_1=0$ has weight~$0$ and $X'_1=n$ has probability $2^{-n}$. With
$t=(2X'_1-n)/n\in[-1,1]$,
\[
 \frac{X'_1}{2n-X'_1}=\frac{1+t}{3-t}
 =\frac13+\frac{4t}9+\frac{4t^2}{27}+\frac{4t^3}{81}+O(t^4)
\]
uniformly, and $\E t=\E t^3=0$, $\E t^2=1/n$, $\E t^4=O(n^{-2})$, so the
ratio is $1+\frac4{9n}+O(n^{-2})$. The values $c_3\approx\frac12$ and
$c_4\approx\frac23$ are numerical observations only; the expansion to all
orders is item~(F2) of Section~\ref{shr:bp:sec:further}.

Part~(a) is also the first step of an
all-order expansion: the expectation in
\eqref{shr:bp:eq:boundary-exact} is an interior sum of Part~I's type with
$m-1$ rows and one extra rational factor per row, so the method of
Section~5 should expand it to every order in $1/n$. That expansion is not
carried out here (Section~\ref{shr:bp:sec:further}).

\section{A beta--binomial representation and the leading constant}
\label{shr:bp:sec:beta}

\subsection{A random cut at the first completion}
Take $m$ independent rows of $n$ independent uniform random variables on
$(0,1)$ and sort within each row. Their global ranks are uniform among
the $M_m(n)$ row interleavings. Let $U$ be the maximum of the top row,
and let $\mathcal S$ be the event that every entry in the bottom row
exceeds $U$. Its probability is $\binom{2n}{n}^{-1}$.
The joint density of $U$ and $\mathcal S$ is
$n u^{n-1}(1-u)^n$. Consequently
\begin{equation}
 U\mid\mathcal S\sim\operatorname{Beta}(n,n+1).
 \label{shr:bp:eq:betaU}
\end{equation}
Conditional on $U=u$ and $\mathcal S$, the numbers of middle-row samples
below $u$ are independent variables
\begin{equation}
 X_i\mid U=u\sim\operatorname{Bin}(n,u),\qquad 1\le i\le r.
 \label{shr:bp:eq:X}
\end{equation}
The $X_i$ are exchangeable but \emph{not} independent after integrating
out $U$.

Set $\mathcal I=\{1\le X_i\le n-2\text{ for all }i\}$ and define
\begin{align}
 W_n(x)&=
 \prod_{i<j}\frac{(x_i-x_j)^2}{(x_i+x_j)(2n-x_i-x_j)}
 \prod_{i=1}^{r}
 \frac{(n-1-x_i)x_i}{(n-1+x_i)(2n-x_i)}.
 \label{shr:bp:eq:W}
\end{align}
Denominators are positive on $\mathcal I$; equal coordinates give zero.
Every pair factor and every single-coordinate factor belongs to $[0,1]$,
so $0\le W_n\le1$ there.

\begin{lemma}[Exact mixture identity]
\label{shr:bp:lem:mixture}
For all $m,n\ge2$,
\begin{equation}
 \frac{H_m(n)}{M_m(n)}=
 \frac{1}{r!\binom{2n}{n}}
 \E\bigl[\ind_{\mathcal I}W_n(X)\bigr],
 \label{shr:bp:eq:mixture}
\end{equation}
with the beta--binomial law \eqref{shr:bp:eq:betaU}--\eqref{shr:bp:eq:X}.
\end{lemma}
\begin{proof}
Given the counts $x$, the samples below $U$ in the top and middle rows
are independent sorted uniforms. Relative to their unconstrained
multinomial interleavings, the first hook count contributes
\[
 \prod_i\frac{n-1-x_i}{n-1+x_i}
 \prod_{i<j}\frac{x_i-x_j}{x_i+x_j}
\]
when $x_1>\cdots>x_r$. Above $U$, the complementary hook count similarly
contributes
\[
 \prod_i\frac{x_i}{2n-x_i}
 \prod_{i<j}\frac{x_i-x_j}{2n-x_i-x_j}.
\]
All cross-cut comparisons are automatic for the legal prefix state.
The product is $W_n(x)$. Distinct middle counts have exactly one decreasing
ordering; the mixture law and the weight are symmetric in those counts.
Thus summing over decreasing counts is $1/r!$ times the unrestricted
expectation on $\mathcal I$. Multiplying by $\PP(\mathcal S)$ proves the
identity. This is also a probabilistic derivation of \eqref{shr:bp:eq:hookcore}.
\end{proof}

\subsection{Scaling and the common Gaussian mode}
Let
\begin{equation}
 \eps=n^{-1/2},\qquad Y_i=\frac{2X_i-n}{\sqrt n},\qquad
 V_n=2\sqrt n\,(U-\tfrac12),\qquad
 Z_{i,n}=\frac{2(X_i-nU)}{\sqrt n}.
 \label{shr:bp:eq:scaling}
\end{equation}
Then $Y_i=Z_{i,n}+V_n$. A direct Laplace expansion of the beta density gives
$V_n\Rightarrow V\sim N(0,1/2)$. Conditional binomial characteristic
functions, uniformly for $U$ near $1/2$, give joint convergence
\begin{equation}
 (Z_{1,n},\ldots,Z_{r,n},V_n)
 \ \Rightarrow\ (Z_1,\ldots,Z_r,V),
 \label{shr:bp:eq:CLT}
\end{equation}
where the $Z_i$ are independent $N(0,1)$ variables and are independent of
$V$. For example, expanding
$(1-U+Ue^{2it/\sqrt n})^n e^{-2itnU/\sqrt n}$ gives the limiting factor
$e^{-t^2/2}$ on $U=1/2+o(1)$; the beta mass outside that neighborhood tends
to zero. The same argument applies to a fixed vector of $t_i$.

\begin{writenote}[Added 5 October 2026, batch~98: the convergence \eqref{shr:bp:eq:CLT} in detail]
The two sentences above are a sketch; here is the argument with its
uniformity made explicit. The density of $U$ given $\mathcal S$ is
$f(u)=n\binom{2n}n u^{n-1}(1-u)^n$, so for $|v|<\sqrt n$ the density of
$V_n$ is
\[
 \sqrt n\,\binom{2n}n4^{-n}\Bigl(1-\frac v{\sqrt n}\Bigr)
 \Bigl(1-\frac{v^2}n\Bigr)^{n-1}\longrightarrow\frac{e^{-v^2}}{\sqrt\pi},
\]
the $N(0,1/2)$ density, and Scheff\'e's lemma gives $V_n\Rightarrow V$.
Conditionally on $U=u$ the $Z_{i,n}$ are independent with characteristic
function $\varphi_n(t;u)=\bigl(\E e^{ih(\xi-u)}\bigr)^n$, where
$h=2t/\sqrt n$ and $\xi$ is Bernoulli$(u)$. The cumulant function
$\kappa(h)=\log\E e^{ih(\xi-u)}$ has $\kappa(0)=\kappa'(0)=0$,
$\kappa''(0)=-u(1-u)$ and, since $|\E e^{ih(\xi-u)}|\ge1-h^2/8$, a third
derivative bounded uniformly in $u\in[0,1]$ for $|h|\le1/2$. Hence
$\log\varphi_n(t;u)=-2u(1-u)t^2+O(|t|^3/\sqrt n)$ uniformly in $u$, for
$|t|\le\sqrt n/4$. As $|\varphi_n|\le1$,
\[
 \E\exp\Bigl(isV_n+i\sum_it_iZ_{i,n}\Bigr)
 =\E\Bigl[e^{isV_n}e^{-2U(1-U)\sum_it_i^2}\Bigr]+o(1)
 \longrightarrow\E e^{isV}\,e^{-\sum_it_i^2/2},
\]
because $2U(1-U)=\frac12-V_n^2/(2n)$ and $V_n\Rightarrow V$. This is
\eqref{shr:bp:eq:CLT}.
\end{writenote}

Let $\Delta(y)=\prod_{i<j}(y_i-y_j)$. Algebraic substitution into
\eqref{shr:bp:eq:W} gives the exact factorization
\begin{equation}
 W_n(X)=(4n)^{-d_{\rm R}}9^{-r}\Delta(Y)^2\Psi_r(\eps,Y)
 \quad\text{on }\mathcal I,
 \label{shr:bp:eq:Psidef0}
\end{equation}
where
\begin{align}
 \Psi_r(\eps,y)
 &=\prod_{i<j}\left(1-\frac{\eps^2(y_i+y_j)^2}{4}\right)^{-1}\notag\\
 &\quad\times\prod_{i=1}^{r}
 \frac{9(1-2\eps^2-\eps y_i)(1+\eps y_i)}
 {(3-2\eps^2+\eps y_i)(3-\eps y_i)}.
 \label{shr:bp:eq:Psidef}
\end{align}
The factor $9^{-r}$ has a simple origin: at $x_i\sim n/2$, the interaction
with the completed top row and the unborn bottom row tends to
$(1/3)(1/3)=1/9$ for each middle row.

\subsection{Gaussian normalization}
For independent standard normal $Z_i$,
\begin{equation}
 \E\Delta(Z)^2=r!J_r.
 \label{shr:bp:eq:gaussianZ}
\end{equation}
One elementary proof writes the Vandermonde determinant in the basis of
monic Hermite polynomials. Their squared norms for the standard normal
measure are $j!$; determinant expansion and orthogonality give
$r!\prod_{j=0}^{r-1}j!$. The Hermite norms themselves follow from
\[
 \E\exp(tZ-t^2/2)\exp(sZ-s^2/2)=e^{ts}.
\]
Since adding a common scalar does not change a Vandermonde,
\begin{equation}
 \Delta(Z_1+V,\ldots,Z_r+V)=\Delta(Z).
 \label{shr:bp:eq:translation}
\end{equation}
Thus the extra beta fluctuation contributes a common Gaussian shift, but
does not alter the leading Vandermonde normalization.

For completeness, moment convergence is justified by uniform tails.
The beta density is proportional to
$[u(1-u)]^{n-1}(1-u)$, and
$u(1-u)=1/4-(u-1/2)^2$. It follows, by comparing the normalizing integral
with its central part, that
$\PP(|V_n|>t)\le C e^{-ct^2}$ in the relevant range, with harmless changes
of constants near the endpoints. Conditional binomial exponential bounds
give the same form for $Z_{i,n}$. In particular, all fixed moments of the
$Y_i$ are uniformly bounded, and for any fixed $0<\eta<1/2$,
\begin{equation}
 \PP\{\max_i|Y_i|>n^\eta\}\le C_r e^{-c_r n^{2\eta}}.
 \label{shr:bp:eq:tails}
\end{equation}
These facts also follow from the exact moment formula in the next section.

\begin{writenote}[Added 5 October 2026, batch~98: explicit constants in \eqref{shr:bp:eq:tails}]
The phrase ``with harmless changes of constants near the endpoints'' can be
replaced by explicit constants. With $\delta=u-\frac12$, $1-u\le1$ and the
standard bound $\binom{2n}n\le4^n/\sqrt{\pi n}$,
\[
 f(u)\le n\binom{2n}n\bigl(\tfrac14-\delta^2\bigr)^{n-1}
 \le4\sqrt{n/\pi}\,e^{-4(n-1)\delta^2}.
\]
Integrating over $|\delta|>s$ and using $(s+v)^2\ge s^2+v^2$ gives
$\PP(|U-\frac12|>s)\le2\sqrt{n/(n-1)}\,e^{-4(n-1)s^2}$; with
$s=t/(2\sqrt n)$ and $n\ge2$, $\PP(|V_n|>t)\le2\sqrt2\,e^{-t^2/2}$ for all
$t\ge0$. Given $U=u$, Hoeffding's inequality for the $n$ Bernoulli
variables of row~$i$ gives $\PP(|Z_{i,n}|>t\mid U=u)\le2e^{-t^2/2}$,
uniformly in $u$. Since $Y_i=Z_{i,n}+V_n$, it follows that
$\PP(|Y_i|>t)\le5e^{-t^2/8}$ for all $n\ge2$ and $t\ge0$. This gives the
uniform moment bounds, and \eqref{shr:bp:eq:tails} with $C_r=5r$ and
$c_r=1/8$, with no exception near the endpoints.
\end{writenote}

On the central event in \eqref{shr:bp:eq:tails}, $\Psi_r\to1$ with denominators
uniformly separated from zero. Outside it, $W_n\le1$ makes its contribution
to \eqref{shr:bp:eq:mixture} smaller than any power of $n^{-1}$, even after
multiplication by $n^{d_{\rm R}}$. Therefore
\begin{equation}
 \frac{H_m(n)}{M_m(n)}\sim
 \frac{J_r}{\binom{2n}{n}(4n)^{d_{\rm R}}9^r}.
 \label{shr:bp:eq:core-leading}
\end{equation}
Using Stirling's formula for $M_m(n)$ and $\binom{2n}{n}$ gives
\eqref{shr:bp:eq:rareK} and the leading term of Theorem~\ref{shr:bp:thm:rare-main}.
Finally \eqref{shr:bp:eq:core-error} transfers it to $R_m(n)$.
The power is
\[
 \frac{m-1}{2}-\frac12+d_{\rm R}=\frac r2+\binom r2=\frac{r^2}{2},
\]
which explains the drop from the total-count exponent.

\section{All-order coefficients and the universal first correction}
\label{shr:bp:sec:coefficients}

\subsection{Exact factorial moments of the mixture}
For nonnegative integers $k_1,\ldots,k_r$, let $K=k_1+\cdots+k_r$.
Conditional binomial factorial moments and the beta integral give
\begin{equation}
 \E\prod_{i=1}^r\fall{X_i}{k_i}
 =\left(\prod_{i=1}^r\fall{n}{k_i}\right)
 \frac{\rise{n}{K}}{\rise{2n+1}{K}}.
 \label{shr:bp:eq:factorialmoment}
\end{equation}
Here falling and rising factorials are distinguished typographically.
Every joint ordinary moment of $X$ is a finite rational function of $n$,
by the Stirling-number identity
$x^a=\sum_k \Stirling{a}{k}\fall{x}{k}$. Therefore every joint moment of
$Y$ has an effective Laurent expansion in $\eps=n^{-1/2}$. For a monomial
of total degree $A$, only powers of $\eps$ of parity $A$ occur.
The bounded-moment argument above ensures that no negative powers survive.

Let $\mathcal E_\eps$ be the linear functional on polynomials in
$y_1,\ldots,y_r$ obtained by replacing each monomial by this formal moment
expansion. Set
\begin{equation}
 \mathcal Q_r(z)=\frac1{r!J_r}
 \mathcal E_\eps\!\left[\Delta(y)^2\Psi_r(\eps,y)\right]
 \bigg|_{\eps^2=z},
 \qquad Q_r(0)=1.
 \label{shr:bp:eq:Qr}
\end{equation}
This is interpreted coefficient by coefficient, not as convergence of a
formal infinite series. The symmetry
$\Psi_r(-\eps,-y)=\Psi_r(\eps,y)$ shows that its coefficient of $\eps^k$
has polynomial parity $k$. Since $\Delta^2$ has even total degree, the
moment parity rule shows that the right-hand side of \eqref{shr:bp:eq:Qr}
contains only even powers of $\eps$. Thus $\mathcal Q_r$ is a well-defined series
in $z$ with rational coefficients.

The remaining Stirling factor is
\begin{equation}
 A_m(z)=\exp\left\{
 \sum_{k\ge1}\frac{B_{2k}}{2k(2k-1)}
 \left(m^{1-2k}-m+2-2^{1-2k}\right)z^{2k-1}
 \right\},
 \label{shr:bp:eq:outer}
\end{equation}
where $B_{2k}$ are Bernoulli numbers. We obtain the explicit all-order rule
\begin{equation}
 \boxed{\displaystyle
 b_{m,j}^{\rm R}=[z^j]A_m(z)\mathcal Q_{m-2}(z).}
 \label{shr:bp:eq:coefficient-rule}
\end{equation}
For a prescribed $j$, all operations are finite: expand $\Psi$ through
$\eps^{2j}$, take finitely many moments from
\eqref{shr:bp:eq:factorialmoment}, and truncate \eqref{shr:bp:eq:outer} at $z^j$.

\subsection{Why the formal calculation is an asymptotic theorem}
\begin{proof}[Completion of the all-order part of Theorem~\ref{shr:bp:thm:rare-main}]
Fix $L$. On $\max|Y_i|\le n^\eta$, choose any $0<\eta<1/2$.
Each rational factor of \eqref{shr:bp:eq:Psidef} has denominator bounded away
from zero for large $n$. Taylor expansion through order $2L+1$ has a
remainder bounded by
\[
 C_{r,L}\eps^{2L+2}(1+\|Y\|)^{N_{r,L}}
\]
there, for some fixed exponent $N_{r,L}$. This follows either by finite
geometric-series remainders in the displayed denominators, or by Taylor's
formula on the line segment from $0$ to $\eps$; throughout that segment
$|\eps Y_i|$ tends uniformly to zero. After multiplication by $\Delta^2$,
uniform bounded moments bound the expected remainder by
$O_{r,L}(n^{-L-1})$.

The complement of the central event contributes exponentially little to
the original weighted expectation, by $0\le W_n\le1$ and
\eqref{shr:bp:eq:tails}. It also contributes exponentially little to each
polynomial Taylor term, by the same tail bound and higher bounded
moments. We may consequently replace each truncated polynomial expectation
by its full exact expectation. Formula~\eqref{shr:bp:eq:factorialmoment} expands
these rational functions to the desired order with an ordinary rational
remainder estimate. The parity argument removes half-integral orders.
This proves the Poincar\'e expansion specified by $\mathcal Q_r$.

Stirling's formula with its fixed-order remainder yields $A_m$ and
\eqref{shr:bp:eq:coefficient-rule}. Finally
\[
 \frac{M_m(n)/\binom{3n}{n}}
 {M_m(n)/\{\binom{2n}{n}n^{d_{\rm R}}\}}
 =n^{d_{\rm R}}\frac{\binom{2n}{n}}{\binom{3n}{n}}
 =O_r\!\left(n^{d_{\rm R}}(16/27)^n\right).
\]
Thus the difference between $R_m$ and $H_m$ affects no coefficient of the
rare-scale inverse-power expansion. All coefficients are rational by
construction, and the case $r=0$ is the exact identity already noted.
\end{proof}

The theorem asserts effective coefficients and existence of fixed-order
error constants. The package does not supply optimized numerical values
for these constants or a rigorously certified onset of the asymptotic
regime.

\subsection{An exact Vandermonde identity}
It is useful to derive the first correction without a large polynomial
expansion. For $X\sim\operatorname{Bin}(n,p)$ consider
\begin{equation}
 F_x(t)=(1+(1-p)t)^x(1-pt)^{n-x}.
\end{equation}
The polynomial $P_j(x)=j![t^j]F_x(t)$ is monic of degree $j$. Direct
binomial expectation gives
\begin{equation}
 \E F_X(t)F_X(s)=(1+p(1-p)ts)^n.
\end{equation}
Hence the $P_j$ are orthogonal and have squared norms
$j!\fall{n}{j}[p(1-p)]^j$. Replacing Vandermonde columns by these monic
polynomials and expanding the determinant gives, for independent
binomial variables conditional on $U=p$,
\begin{equation}
 \E[\Delta(Y)^2\mid U=p]
 =r!J_r\left(\prod_{j=0}^{r-1}\frac{\fall nj}{n^j}\right)
 [4p(1-p)]^{d_{\rm R}}.
 \label{shr:bp:eq:binomialVdm}
\end{equation}
Integrating the beta law,
\begin{equation}
 \E[U^{d_{\rm R}}(1-U)^{d_{\rm R}}]=
 \frac{\rise n{d_{\rm R}}\rise{n+1}{d_{\rm R}}}{\rise{2n+1}{2d_{\rm R}}}.
 \label{shr:bp:eq:betad}
\end{equation}
It follows by expanding these finite products that
\begin{equation}
 \frac{\E\Delta(Y)^2}{r!J_r}
 =1-\frac{r(r-1)(2r-1)}{12n}+O_r(n^{-2}).
 \label{shr:bp:eq:vdm-correction}
\end{equation}
This identity isolates the correction caused by finite binomial and beta
moments before inserting the rational interaction factors.

\subsection{Evaluation of the interaction correction}
Write $S_1=\sum_i y_i$ and $S_2=\sum_i y_i^2$. Expanding
\eqref{shr:bp:eq:Psidef} begins with
\begin{equation}
 \Psi_r(\eps,y)=1+\eps^2P_1(y)+O(\eps^3),\qquad
 P_1(y)=\frac{9r-50}{36}S_2+\frac14S_1^2-\frac{4r}{3}.
 \label{shr:bp:eq:P1}
\end{equation}
The $O(\eps^3)$ notation here denotes a formal polynomial expansion;
the previous proof supplies its expectation-level meaning.

Under the Gaussian law tilted by $\Delta(Z)^2$, two elementary identities
are
\begin{equation}
 \E_*\sum_i Z_i^2=r^2,\qquad
 \E_*\left(\sum_iZ_i\right)^2=r.
 \label{shr:bp:eq:gue-moments}
\end{equation}
For the first, scale the Gaussian integral: the degree of $\Delta^2$ is
$r(r-1)$ and the volume dimension is $r$, so differentiation in the
Gaussian variance gives $r^2$. For the second, separate the center of
mass from its orthogonal complement. The Vandermonde is independent of
the center of mass, whose sum coordinate is $N(0,r)$.
Since $Y_i=Z_i+V$ and $\Var V=1/2$, with the tilt still independent of $V$,
\begin{equation}
 \E_*S_2=r^2+r/2,\qquad
 \E_*S_1^2=r+r^2/2.
 \label{shr:bp:eq:shifted-moments}
\end{equation}
Substitution into \eqref{shr:bp:eq:P1} yields
\begin{equation}
 \E_*P_1=\frac{r(9r^2-41r-64)}{36}.
\end{equation}
Combining with \eqref{shr:bp:eq:vdm-correction}, the first coefficient of $\mathcal Q_r$ is
\begin{equation}
 [z]\mathcal Q_r(z)=\frac{r(3r^2-32r-67)}{36}.
 \label{shr:bp:eq:q1}
\end{equation}
The first coefficient of $A_m$ is
\begin{equation}
 [z]A_m(z)=\frac18+\frac{1/m-m}{12}.
 \label{shr:bp:eq:a1}
\end{equation}
Adding \eqref{shr:bp:eq:q1} and \eqref{shr:bp:eq:a1}, with $m=r+2$, gives
\eqref{shr:bp:eq:rare-b1}. This proves the universal first correction independently
of the higher-order implementation.

\section{Height five, conditional shapes, and exponential separation}
\label{shr:bp:sec:five}

\subsection{Three corrections for the A181199 rare subcount}
Applying \eqref{shr:bp:eq:coefficient-rule} at $m=5$ gives
\begin{corollary}[Height-five rare expansion]
\label{shr:bp:cor:five}
For $R_5(n)=\#\{A:a_{1,n}<a_{5,1}\}$,
\begin{align}
 R_5(n)=&\frac{\sqrt5}{93312\pi^{3/2}}
 \left(\frac{3125}{4}\right)^n n^{-9/2}\notag\\
 &\times\left(1-\frac{1393}{120n}
 +\frac{656449}{28800n^2}
 +\frac{9833844689}{155520000n^3}
 +O(n^{-4})\right).
 \label{shr:bp:eq:five}
\end{align}
\end{corollary}
The rational coefficients are finite evaluations of the proven formula,
not a fit to sequence values. Appendix~\ref{shr:bp:app:algorithm} describes the
implementation-independent arithmetic used to reproduce them.
For comparison, the first three corrections at adjacent heights are
\begin{center}\small
\begin{tabular}{@{}crrr@{}}
\toprule
$m$ & $b_{m,1}^{\rm R}$ & $b_{m,2}^{\rm R}$ & $b_{m,3}^{\rm R}$\\
\midrule
$2$ & $0$ & $0$ & $0$\\[2pt]
$3$ & $-199/72$ & $-6479/10368$ & $-491815/2239488$\\[2pt]
$4$ & $-979/144$ & $74083/13824$ & $25693405/1990656$\\[2pt]
$5$ & $-1393/120$ & $656449/28800$ & $9833844689/155520000$\\
\bottomrule
\end{tabular}
\end{center}
Here $R_5$ is a newly considered refinement of the A181199 counting
problem; it is not assigned a new OEIS identifier in this report.

\subsection{How rare is the absence of a common active phase?}
Let $Q_m(n)=T_{m,B^mD^m}(n)$ for $n\ge2$. Then exactly
\begin{equation}
 T_m(n)=Q_m(n)+R_m(n).
 \label{shr:bp:eq:two-sector}
\end{equation}
Both $Q_m$ and $R_m$ separately have rational-diagonal generating
functions: each is a fixed finite sum of the event sectors in
Theorem~\ref{shr:bp:thm:diagonal-main}, up to the harmless terms at $n=0,1$.
Using \eqref{shr:bp:eq:priorleading} from the preceding report gives
\begin{corollary}[Rare-event probability]
\label{shr:bp:cor:probability}
For fixed $m\ge2$, a uniformly chosen shifted rectangle satisfies
\begin{equation}
 \PP(D_1<B_m)=\frac{R_m(n)}{T_m(n)}
 \sim\frac{\sqrt\pi\,4^{2m-3}}
 {9^{m-2}(m-2)!(m-1)!}\,n^{2m-5/2}4^{-n}.
 \label{shr:bp:eq:rareprob}
\end{equation}
At height five, the first correction is
\begin{equation}
 \frac{R_5(n)}{T_5(n)}=
 \frac{1024\sqrt\pi}{6561}\,n^{15/2}4^{-n}
 \left(1-\frac{509}{24n}+O(n^{-2})\right).
 \label{shr:bp:eq:rareprob5}
\end{equation}
\end{corollary}
\begin{proof}
Divide the leading constants and powers in
\eqref{shr:bp:eq:rareall} and \eqref{shr:bp:eq:priorleading}. The factorial quotient is
$J_{m-2}/J_m=1/((m-2)!(m-1)!)$, and the powers of $4$ give $4^{2m-3}$.
For the correction at height five, the preceding report gives
$b_{5,1}^{\rm total}=48/5$, so subtraction gives
$-1393/120-48/5=-509/24$.
\end{proof}

Thus $Q_m$ has the same dominant-scale inverse-power expansion as $T_m$
to every fixed order, yet their difference has the explicitly identified
exponential scale $(m^m/4)^n$. This is a combinatorially defined,
beyond-all-orders refinement. It does \emph{not} by itself identify a
Stokes multiplier of a chosen analytic interpolation or summation of the
dominant formal series. There may be other exponentially small structures
inside $Q_m$, and the next boundary sector remains to be analyzed.

\subsection{The middle shape conditioned on the rare event}
The same argument gives a probabilistic limit theorem. Use the auxiliary
uniform-sample model of Section~\ref{shr:bp:sec:beta}. Condition on a legal array
and on $D_1<B_m$. Let $X_i$ count middle-row entries below the top-row
maximum $U$, and use the scaling \eqref{shr:bp:eq:scaling}.

\begin{theorem}[Rare conditional shape]
\label{shr:bp:thm:shape}
For fixed $m\ge3$, conditional on the rare event, the probability that all
middle rows are active at the first completion tends to one. On that
core event, the limiting joint law is
\begin{equation}
 (Y_1,\ldots,Y_r,V_n)\Rightarrow
 (\Lambda_1+V,\ldots,\Lambda_r+V,V),
 \label{shr:bp:eq:shape}
\end{equation}
where $V\sim N(0,1/2)$ is independent of
$\Lambda_1>\cdots>\Lambda_r$, whose density on the decreasing chamber is
\begin{equation}
 \frac1{J_r}\Delta(\lambda)^2
 \prod_{i=1}^r\frac{e^{-\lambda_i^2/2}}{\sqrt{2\pi}}.
 \label{shr:bp:eq:orderedGaussian}
\end{equation}
All fixed polynomial moments converge as well.
\end{theorem}
\begin{proof}
The first statement is \eqref{shr:bp:eq:core-error} divided by the asymptotic
for $R_m$. In the mixture representation, conditioning on legality and
the core multiplies the joint law by $W_n(X)$ and restricts the counts to
the decreasing chamber. After normalization,
\eqref{shr:bp:eq:Psidef0}, \eqref{shr:bp:eq:CLT}, and the tail estimates show that the
limiting weight is $\Delta(Z+V)^2=\Delta(Z)^2$. It is independent of $V$,
so the common Gaussian mode is not tilted. Restricting the $Z$ coordinates
to decreasing order changes the normalization from $r!J_r$ to $J_r$.
The tail estimates with additional polynomial factors give moment
convergence. The omitted noncore configurations have exponentially small
conditional probability; their scaled coordinates are bounded by
$O(\sqrt n)$, so they do not affect any fixed moment.
\end{proof}

The auxiliary $U$ is a uniform-order-statistic realization of the
combinatorial cut. The theorem is not a claim that a tableau itself
contains real-valued sample times. In particular, its statement about
$Y_i$, the scaled integer row lengths, is intrinsic; the joint statement
with $U$ uses the specified coupling.

\section{Inverting the rare-count asymptotic profile}
\label{shr:bp:sec:inverse}

For $m\ge3$ put
\begin{equation}
 \lambda_m=m\log m-\log4>0,\qquad
 \beta_m=\frac{(m-2)^2}{2}>0,
 \qquad a=\log(y/\mathcal K_m).
 \label{shr:bp:eq:inverseparams}
\end{equation}
The leading continuous profile is
$y=\mathcal K_m e^{\lambda_m x}x^{-\beta_m}$. Its large solution is
\begin{equation}
 x_0=-\frac{\beta_m}{\lambda_m}
 W_{-1}\!\left(-\frac{\lambda_m}{\beta_m}e^{-a/\beta_m}\right),
 \label{shr:bp:eq:lambert}
\end{equation}
where $W_{-1}$ is the real lower branch. The argument belongs to $(-1/e,0)$
for sufficiently large $y$. The branch is forced by the requirement
$x_0\to\infty$; the principal branch gives the small solution instead.
Lambert-core inversion and inverse-power reversion are standard tools,
also used for the total sequence in Part~\ref{shr:part:one}
(Section~\ref{shr:sec:inverse}). Their application here
uses the different rare-event rate and power.

\begin{proposition}[Two corrections to the inverse profile]
\label{shr:bp:prop:inverse}
Write $b_1=b_{m,1}^{\rm R}$, $b_2=b_{m,2}^{\rm R}$,
$\lambda=\lambda_m$, and $\beta=\beta_m$. The formal large inverse has
\begin{equation}
 x=x_0-\frac{b_1}{\lambda x_0}
 -\frac{1}{x_0^2}
 \left(\frac{b_2-b_1^2/2}{\lambda}
       +\frac{\beta b_1}{\lambda^2}\right)
 +O(x_0^{-3}).
 \label{shr:bp:eq:inverse2}
\end{equation}
Every subsequent coefficient is obtained recursively from
\eqref{shr:bp:eq:coefficient-rule}. For any fixed sufficiently long truncation
of the real asymptotic profile, its genuine large real inverse has the
corresponding finite expansion.
\end{proposition}
\begin{proof}
Take logarithms:
\[
 \lambda x-\beta\log x+\log\left(1+b_1/x+b_2/x^2+\cdots\right)=a.
\]
Substitute $x=x_0+c_1/x_0+c_2/x_0^2+\cdots$ and use
$\lambda x_0-\beta\log x_0=a$. At orders $x_0^{-1}$ and $x_0^{-2}$,
\[
 \lambda c_1+b_1=0,\qquad
 \lambda c_2-\beta c_1+b_2-b_1^2/2=0.
\]
This gives \eqref{shr:bp:eq:inverse2}. At every later order the next unknown
coefficient occurs linearly with nonzero coefficient $\lambda$.
For a fixed real profile truncation, its logarithmic derivative tends to
$\lambda>0$. Taylor's theorem or the real implicit function theorem then
turns the recursive formal truncation into the stated inverse asymptotic.
\end{proof}

For height five, $b_2-b_1^2/2=-535/12$. Thus, with
$\lambda=\log(3125/4)$ and $\beta=9/2$,
\begin{equation}
 x=x_0+\frac{1393}{120\lambda x_0}
 +\frac1{x_0^2}
 \left(\frac{535}{12\lambda}+\frac{4179}{80\lambda^2}\right)
 +O(x_0^{-3}).
 \label{shr:bp:eq:inverse5}
\end{equation}
The leading logarithmic scale can also be read without $W$:
\[
 x_0=\frac1\lambda\left(a+\beta\log a-\beta\log\lambda
            +O\left(\frac{\log a}{a}\right)\right).
\]

\subsection{A discrete inverse is a staircase}
The sequence $R_m(n)$ is eventually strictly increasing for $m\ge3$,
since its successive ratio tends to $e^{\lambda_m}>1$.
Define $N_m(y)=\min\{n\ge1:R_m(n)\ge y\}$ for large $y$.
Let $\xi_L(y)$ be the large inverse of the truncated continuous profile
\[
 \mathcal K_m e^{\lambda_m x}x^{-\beta_m}
 \sum_{j=0}^L b_{m,j}^{\rm R}x^{-j}.
\]
The fixed-order remainder and a logarithmic derivative bounded away from
zero imply that, for some $C_{m,L}$ and all sufficiently large $y$,
\begin{equation}
 \left\lceil\xi_L(y)-C_{m,L}\xi_L(y)^{-L-1}\right\rceil
 \le N_m(y)\le
 \left\lceil\xi_L(y)+C_{m,L}\xi_L(y)^{-L-1}\right\rceil.
 \label{shr:bp:eq:staircase}
\end{equation}
Indeed, the asymptotic relative error is $O(n^{-L-1})$, so the two smooth
comparison profiles cross the target within that horizontal distance.
When the two displayed ceilings coincide, their common value is the
threshold. No such agreement can be asserted uniformly near integer
crossings without explicit error constants and distance information.
The package's numerical inverse checks are therefore profile diagnostics,
not integer-threshold certificates.

\begin{writenote}[Added 5 October 2026, batch~98: the inversion is an instance]
As for Part~\ref{shr:part:one}'s \eqref{shr:eq:inverseW} (note at the end
of Section~\ref{shr:sec:inverse}), \eqref{shr:bp:eq:lambert} is a case of
Theorem~J.23 (\texttt{p0:thm:lambert-core}) of the transseries volume
\texttt{Analysis/\allowbreak Transseries/\allowbreak docs/\allowbreak series-\allowbreak and-\allowbreak transseries/\allowbreak Transseries\_\allowbreak And\_\allowbreak Inversion/},
with $a=\lambda_m>0$ and $b=-\beta_m<0$, whose branch rule for $b<0$
selects $W_{-1}$; the caution about integer thresholds and
\eqref{shr:bp:eq:staircase} are the separation condition of that volume's
staircase theorem J.43 (\texttt{p0:thm:staircase}). The manuscript itself
calls these tools standard. Two steps are argued in one sentence each and
are recorded in Section~\ref{shr:bp:sec:further}: the last sentence of
Proposition~\ref{shr:bp:prop:inverse} (that the genuine real inverse has
the formal expansion) and the staircase \eqref{shr:bp:eq:staircase}, whose
constant $C_{m,L}$ and range of validity are existential.
\end{writenote}

\section{Exact checks, numerical diagnostics, and reproducibility}
\label{shr:bp:sec:computation}

\subsection{Independent computations}
The package includes three genuinely different finite enumerations:
unit-step prefix-state dynamic programming; an independent bitmask
order-ideal enumerator built directly from the cell comparisons; and an
event-skeleton evaluator using the signed multinomial kernel. The hook
core sum and the symbolic beta-moment expansion are separate again.
No candidate OEIS recurrence is used to generate a test value.

The verification script completed \textbf{1,954 exact checks}. The
categories and counts are recorded in \texttt{data/verification.json}:
\begin{center}\small
\begin{tabular}{@{}lr@{}}
\toprule
Check & Number\\
\midrule
Independent poset enumeration versus prefix states & 28\\
Reflection kernel versus independent chamber walks & 887\\
Event decomposition versus prefix states & 25\\
Event-word realizability and rotation symmetry & 640\\
Rare event-word sum and the worked example & 21\\
OEIS reference values (heights four and five) & 20\\
Core lower bound, separation bound, and core error bound & 228\\
Classical Catalan and small-width controls & 64\\
Universal first correction and saved higher coefficients & 7\\
Common-mode Gaussian moments and exact beta moments & 34\\
\midrule
Total & 1954\\
\bottomrule
\end{tabular}
\end{center}
The reference terms are small data excerpts from A181198 and A181199,
with source URLs and access date in the audit. The tests regenerate these
values rather than assuming their recurrence. Finite checks corroborate
the arguments; they are not proofs of an infinite identity or a substitute
for peer review.

\subsection{Accuracy on the rare exponential scale}
Let $\mathcal L_5(n)$ denote the leading term in \eqref{shr:bp:eq:five}, and put
$\mathcal A_{5,L}(n)=\mathcal L_5(n)\sum_{j=0}^L b_{5,j}^{\rm R}n^{-j}$.
The following entries are relative errors $\mathcal A_{5,L}(n)/R_5(n)-1$
computed against exact integer rare counts.
\begin{center}\small
\begin{tabular}{@{}rrrrr@{}}
\toprule
$n$ & $L=0$ & $L=1$ & $L=2$ & $L=3$\\
\midrule
$10$ & $4.46198$ & $-1.87847$ & $-0.633499$ & $-0.288127$\\
$20$ & $1.05994$ & $-0.135683$ & $-0.0183006$ & $-0.00201881$\\
$40$ & $0.379221$ & $-0.0210405$ & $-0.00139232$ & $-0.0000296524$\\
$60$ & $0.229776$ & $-0.00815163$ & $-0.000365317$ & $-0.00000531141$\\
\bottomrule
\end{tabular}
\end{center}
These values illustrate both the usefulness of the corrections and the
slow onset of the leading asymptotic at height five. The negative
first-order approximation at $n=10$ is not a contradiction: fixed-order
asymptotic expansions are not guaranteed positive or accurate at small
$n$.

The core already captures the rare count particularly well at moderate
width:
\begin{center}\small
\begin{tabular}{@{}rll@{}}
\toprule
$n$ & $R_5(n)/\mathcal L_5(n)$ & $H_5(n)/R_5(n)$\\
\midrule
$10$ & $0.183083833013$ & $0.862684900894$\\
$20$ & $0.485450790529$ & $0.999319561936$\\
$40$ & $0.725047021487$ & $0.999999971834$\\
$60$ & $0.813156328776$ & $0.999999999998899$\\
\bottomrule
\end{tabular}
\end{center}
The underlying integer values and 70-digit calculations are provided in
\texttt{data/numerics.csv}. Displayed decimals are diagnostics, not
interval enclosures.

\subsection{Reproducing the package}
From the package root, a Python environment with SymPy and mpmath can run
\begin{verbatim}
python code/verify.py
python code/coefficients.py --height 5 --order 3
python code/numerics.py
python code/inverse_profile.py --index 60 --height 5
\end{verbatim}
The first command reruns the exact tests and writes their log. The second
regenerates the three height-five rational corrections. The third
regenerates the decimal comparison tables from exact rare counts.
The fourth checks the inverse profile at a target obtained from an exact
rare count; it is not a threshold-certification procedure.
The tested environment is Python 3.13.5, SymPy 1.14.0, and mpmath 1.3.0.
The standard-library enumerators themselves do not require either
symbolic or numerical package. No Maple diagonal-elimination computation,
Lean proof, or Rocq proof is claimed.

\begin{writenote}[Added 5 October 2026, batch~98: the shipped layout]
``The package'' in this section is the delivered archive. In ProveIt its
files carry the prefix \texttt{02-boundary-phases-}:
\begin{center}\footnotesize
\begin{tabular}{@{}ll@{}}
\toprule
Delivered as & Shipped as\\
\midrule
\texttt{code/}\emph{name}\texttt{.py} (five scripts) & \texttt{code/02-boundary-phases-}\emph{name}\texttt{.py}\\
\texttt{build.sh} & \texttt{code/02-boundary-phases-build.sh}\\
\texttt{data/}\emph{name} (nine records) & \texttt{data/02-boundary-phases-}\emph{name}\\
\texttt{requirements.txt} & \texttt{data/02-boundary-phases-requirements.txt}\\
\texttt{notes/PROPOSED\_OEIS\_NOTE.md} & \texttt{data/02-boundary-phases-notes-PROPOSED\_OEIS\_NOTE.md}\\
\texttt{notes/SOURCES\_AND\_STATUS.md} & \texttt{02-boundary-phases-notes-SOURCES\_AND\_STATUS.md}\\
\bottomrule
\end{tabular}
\end{center}
The scripts import one another and read and write \texttt{data/} under
their delivered names (\texttt{verify.py} imports \texttt{tableaux}, and
every script sets its root to the parent of \texttt{code/}), and
\texttt{build.sh} changes to its own directory; so they run only on a copy
with the delivered names restored, never in place, and the README of the
report gives the commands. At intake (5~October 2026, on such a copy;
Python~3.13.5, SymPy~1.14.0, mpmath~1.3.0) \texttt{verify.py} passed all
1{,}954 checks in 66~s (the package recorded 23~s), \texttt{numerics.py},
\texttt{coefficients.py --height 5 --order 3} and
\texttt{inverse\_profile.py --index 60 --height 5} reproduced the shipped
records after removal of carriage returns, and the only differences were
the timing fields of \texttt{verification.json}, \texttt{verification.txt}
and \texttt{numerics.json}. The check table above merges some of the
script's sixteen categories (its $640=320+320$, $21=20+1$, $228=3\cdot76$,
$7=4+3$ and $34=6+28$); the total agrees.

\texttt{code/\allowbreak 02-boundary-phases-verify.py} embeds the first ten
terms of A181198 and A181199 as reference values. They are OEIS content,
published by The OEIS Foundation Inc.\ under the Creative Commons
Attribution-ShareAlike 4.0 licence, so those values are third-party data
under that licence, not MIT-0. The draft
\texttt{data/\allowbreak 02-boundary-phases-notes-PROPOSED\_OEIS\_NOTE.md}
is headed ``NOT SUBMITTED'' and stays so: nothing was submitted to OEIS by
the package's author or by the intake, and its instruction not to remove
the conjectural status of the displayed recurrences stands.
\end{writenote}

\begin{writenote}[Added 5 October 2026, batch~101: the OEIS draft]
The draft's instruction ``Do not remove the conjectural status of the
displayed specific recurrences or finite-sum formulas on the basis of this
result alone'' was right for this manuscript, which proves no recurrence,
and the shipped file stays byte-identical. Since 5~October 2026
Parts~\ref{shr:h4:part} and~\ref{shr:h5:part} prove those recurrences and
finite sums. The draft has not been updated, nothing has been submitted to
OEIS, and any change to the OEIS entries is for a human editor after
review.
\end{writenote}

\section{Further research questions and concrete next steps}
\label{shr:bp:sec:future}

\begin{question}[Certify the particular low-order recurrences]
Convert the event sum at heights four and five to a reduced rational
integral, compute a telescoper, and compare it with the conjectured OEIS
operators. A successful certificate must also control singular initial
indices. This would settle a stronger statement than the existence theorem
proved here and could resolve Conjectures~18 and~19 of \cite{KK}.
\end{question}

\begin{writenote}[Added 5 October 2026, batch~101: Research question~8 answered]
Conjectures~18 and~19 of \cite{KK} and both OEIS operators are proved in
Parts~\ref{shr:h4:part} and~\ref{shr:h5:part}, with the singular initial
indices controlled: the operators hold for $n\ge1$, and their residuals at
$n=0$ under $a_0=1$ are $-2160$ and $5621993879040000$. The route is not the
one proposed here: no reduced rational integral of the event sum and no
telescoper, but Pfaffian identities for the transfer integral, one or two
polynomial periods, and explicit rational certificates (a divergence
certificate at height four, contiguity and a moment closure at height
five). A compact certificate from the event sum or the order polytope
remains of interest (Research question~17).
\end{writenote}

\begin{question}[Reduce the dimension of the exact formula]
The bound $m(m-1)$ is deliberately structural rather than optimal. Can
adjacent chamber kernels be convolved before summation, or can a
Pfaffian/determinantal identity absorb the ordered births and deaths?
A substantially smaller representation would make the constructive
D-finiteness theorem computationally practical at larger heights.
\end{question}

\begin{question}[Determine deeper boundary sectors]
The error in \eqref{shr:bp:eq:core-error} contains arrays for which two or more
bottom rows are unborn at the first completion. Determine its sharp
rate, constant, and all-order expansion, then stratify by the number of
unborn rows and early completions. The crude multinomial bound supplies
an upper rate, not the next asymptotic term.
\end{question}

\begin{question}[Find minimal differential complexity]
How do the minimal recurrence order, polynomial degree, and minimal
rational-diagonal dimension grow with $m$? Does the event decomposition
impose a useful uniform bound much smaller than a general elimination
bound? The conjectured small orders at heights four and five suggest
substantial cancellation that the existence proof does not explain.
\end{question}

\begin{writenote}[Added 5 October 2026, batch~101: Research question~11 advanced]
At $m=4$ the minimal order of a recurrence is two, and among recurrences
of order two the minimal degree is nine
(Proposition~\ref{shr:h4:prop:minimal}); at $m=5$ the minimal order is two
or three (Proposition~\ref{shr:h5:prop:nothyper}). The small orders are
explained by the Pfaffian reductions of Parts~\ref{shr:h4:part}
and~\ref{shr:h5:part}, which leave one period at height four and two at
height five. Open: the degrees over all orders, order two against three at
$m=5$ (items~(F5) and~(F9)), and the growth with $m$. (Added 5~October
2026, after the intake's independent check: at $m=4$ the smallest degree
is eight at order three and seven at orders four to eight,
Proposition~\ref{shr:h4:prop:higher}; at $m=5$ there are recurrences of
order four and degree $18$ and of order five and degree $16$,
Proposition~\ref{shr:h5:prop:higher}. The degrees at higher orders stay
open.)
\end{writenote}

\begin{question}[Describe event-rank limit laws]
Corollary~\ref{shr:bp:cor:marks} gives exact D-finite moment totals. Use these
or the beta representation to derive joint laws for $B_i,D_i$, both
unconditionally and under $D_1<B_m$. Determine whether the one-time
conditional shape theorem extends to a process of interacting paths near
the first completion.
\end{question}

\begin{question}[Make height growth uniform]
The proofs fix $m$. Determine ranges $m=m(n)$ in which the Gaussian
reduction and rare-sector asymptotics remain uniform. The dimensions of
the Vandermonde and the interaction product grow quadratically, so fixed
height error estimates cannot simply be reused. Identify the crossover
between a probable common active phase and a boundary-dominated regime.
\end{question}

\begin{question}[Obtain certified finite-width bounds and inversion]
Replace the existential remainders in \eqref{shr:bp:eq:rareall} by computable,
useful bounds with explicit onset. Combine them with
\eqref{shr:bp:eq:staircase} to certify integer thresholds, including inputs very
close to an integer crossing. This requires more than high-precision
floating-point evaluation of the asymptotic series.
\end{question}

\begin{question}[Generalize and formalize the finite-event method]
Investigate unequal row lengths and other fixed-offset antidiagonal
constraints whose active states may still be described by a finite
collection of reflection chambers. Separately, formalize the prefix-state
bijection, reflection involution, and binomial-sum encoding in a proof
assistant. Each is a finite combinatorial component with a clear interface;
a formal proof of the analytic rare-event expansion is a larger second
stage.
\end{question}

\begin{writenote}[Added 5 October 2026, batch~98: the questions of Part~II and Part~I]
The eight questions above are the manuscript's Research questions~1--8,
numbered~8--15 here. Question~8 refines Part~\ref{shr:part:one}'s
Research question~1 (a telescoper route from the event sum, and
Conjectures~18--19 of \cite{KK}); Question~11 is the open second half of
Part~\ref{shr:part:one}'s Research question~7; Questions~12, 13 and~14
continue Research questions~5, 3 and~6 of Part~\ref{shr:part:one}.
\end{writenote}

\subsection{Further questions and research recorded at the write}
\label{shr:bp:sec:further}

Added 5 October 2026, batch~98, under Vladimir Reshetnikov's standing rule
that unproved claims are recorded as open questions rather than dropped.
Each item names its source, the argument it came with, and what is
missing.
\begin{enumerate}[label=(F\arabic*),leftmargin=*]
\item \emph{Other exponentially small structures} (the manuscript, after
Corollary~\ref{shr:bp:cor:probability}). The text says there ``may be
other exponentially small structures inside $Q_m$, and the next boundary
sector remains to be analyzed''. No argument is offered either way;
Proposition~\ref{shr:bp:prop:boundary} shows that the boundary error of
Part~\ref{shr:part:one}'s central cut has rate $2^{-n}$, far above the
rare rate, but that error is an artefact of the cut, not a sector of
$T_m$. Missing: a classification of the event sectors of
Theorem~\ref{shr:bp:thm:skeleton} by exponential rate.
\item \emph{All orders for the boundary error} (the write;
Research question~4 of Part~\ref{shr:part:one}, re-scoped).
Proposition~\ref{shr:bp:prop:boundary} gives the leading term of
$E_{m,n}(1/2)$. Sketch: by \eqref{shr:bp:eq:boundary-exact} the leading
stratum is an interior sum with $m-1$ independent binomial rows and the
extra factor $\prod_i(1+\eps Y_i)/(3-\eps Y_i)$, whose odd powers of
$\eps$ cancel in expectation by the symmetry $Y\mapsto-Y$; Section~5's
moment polynomials and Taylor control should give an expansion in powers
of $1/n$. Missing: writing that argument, computing its coefficients, and
the analogous analysis of the $O(4^{-n})$ strata $D_{m,n}$ and of how the
interior sum and the boundary terms combine.
\item \emph{Sketch-level steps of Sections~\ref{shr:bp:sec:beta},
\ref{shr:bp:sec:five} and~\ref{shr:bp:sec:inverse}.} The joint convergence
\eqref{shr:bp:eq:CLT} and the tail bound \eqref{shr:bp:eq:tails} were
sketches; the write's notes after them supply complete arguments. Still
argued only in outline: the moment convergence in
Theorem~\ref{shr:bp:thm:shape} (``the tail estimates with additional
polynomial factors give moment convergence''; with the explicit tails it
reduces to uniform integrability of $|Y|^kW_n(X)/\E W_n(X)$, which needs
$W_n\le Cn^{-d_{\rm R}}\Delta(Y)^2$ on the central event and the
exponential tails off it); the last sentence of
Proposition~\ref{shr:bp:prop:inverse} (the genuine real inverse has the
formal expansion: a real implicit-function argument, as in Theorem~J.23 of
the transseries volume); and the staircase \eqref{shr:bp:eq:staircase},
whose constant $C_{m,L}$ and threshold for ``sufficiently large $y$'' are
existential (Research question~14 asks for explicit ones).
\item \emph{Priority of Theorem~\ref{shr:bp:thm:diagonal-main}} (the
manuscript and its source audit). The audit searched the OEIS entries,
Kauers--Koutschan, Bostan--Lairez--Salvy, Sun and the web for the sequence
identifiers with ``proved'', ``D-finite'', ``shifted rectangles'' and
``diagonal'', and found no earlier proof; it did not search the general
literature on D-finiteness of linear-extension counts of fixed-height
posets or of walks in Weyl chambers, where the theorem might be a known
special case. Missing: that search.
\end{enumerate}

\begin{writenote}[Added 5 October 2026, batch~101: item~(F4)]
Report~91, dated 2~October 2026 and printed in Part~\ref{shr:sa:part},
proves Theorem~\ref{shr:bp:thm:diagonal-main} by the same route
(Theorem~\ref{shr:sd:thm:diagonal}): Dyck words of births and saturations,
reflection determinants, everywhere-defined binomial products, and
Theorem~3.5 and Corollary~3.6 of Bostan, Lairez and Salvy. Neither
manuscript knew the other. Within this report the earlier proof by date
line is therefore Report~91's (Section~\ref{shr:sa:sec:chronology}); the
search of the general literature that (F4) asks for is still missing.
\end{writenote}

\section{Conclusion}
The obstruction in the shifted-rectangle model is not the strict
interior walk but the coexistence of births, completions, and two kinds
of boundary ties. Recording those events turns the exact enumeration into
a finite binomial construction. That proves rational diagonality and
D-finiteness for every fixed height without a guessed recurrence.

The same boundaries also carry information invisible to the dominant
inverse-power expansion. The absence of a common active phase has rate
$4^{-n}$ relative to the main sequence, an explicit polynomial multiplier,
and its own all-order expansion. A random first-completion cut exposes
both the factor $9^{-r}$ and a Gaussian common mode independent of the
Vandermonde repulsion. These statements are precise additions to the
preceding asymptotic report. They leave clear next targets: compact
operators, the specific OEIS recurrence certificates, and deeper
exponential boundary strata.

\clearpage
\part{A proof of the height four shifted rectangle conjecture}
\label{shr:h4:part}

\noindent{\Large Exact enumeration for A181198 and Kauers and Koutschan
Conjecture~18\par}
\smallskip
\noindent{\large Research report 231, prepared for Vladimir Reshetnikov,
October 5, 2026\par}
\medskip

\paragraph{Main conclusion.}
The full height-four array count satisfies the specific second-order
recurrence recorded in OEIS A181198 and the finite sum printed as
Conjecture 18 of Kauers and Koutschan (2023). The proof starts with the
literal array inequalities, reduces their order-polytope volume to a
weighted triangle integral, and supplies an explicit polynomial divergence
certificate with all boundary terms.

\paragraph{Abstract.}
Column interlacing gives an exterior Volterra transfer. An elementary
four-variable Pfaffian calculation reduces the entire count to two kernel
traces. A rank-two perturbation and integrations by parts express those
traces in one weighted period. A six-coefficient polynomial vector field
then proves a first-order inhomogeneous recurrence. Exact elimination gives
the prescribed OEIS operator; variation of constants gives the actual
printed finite sum, including its initial value. We also derive rational
all-fixed-order asymptotic coefficients directly from the now-proved
recurrence, with effective contraction remainder bounds. Finite
Lambert--$W_{-1}$ models yield rigorously qualified inverse-index envelopes.
The leading asymptotic and existence of all-order expansions were already
established in the earlier fixed-height report and are credited as such.

\paragraph{Status and limits.}
AI-assisted research manuscript, not a peer-reviewed or proof-assistant
formalization. Exact verification code is supplied. The proof concerns
height four only; Conjecture 19 for height five remains open here.
Historical priority is not certified. No canonical exponentially small
transseries is asserted.

\begin{writenote}[Added 5 October 2026, batch~101: provenance of Part~III]
Part~III is a single manuscript, printed in full: Research Report~231 of
the session bundle that arrived in commit \texttt{60f54ea06} and was placed
in commit \texttt{f7c612c72} (batch~101 of ProveIt's incoming-reports
intake) as an addition to this report. Its title is \emph{A proof of the
height four shifted rectangle conjecture}, subtitled ``Exact enumeration for
A181198 and Kauers and Koutschan Conjecture 18'', dated October~5, 2026 (a
25-page PDF; a main file and eight section files, 1{,}518 lines with 108
labels). It arrived as the archive \texttt{Report231.zip} (514{,}097 bytes,
20 files under the doubled wrapper directory
\texttt{Report231/\allowbreak Report231/}). Its title page reads ``Prepared
for Vladimir Reshetnikov'' and its PDF metadata name ``Research report
prepared for Vladimir Reshetnikov''; no human author and no tool are named.
The ``Status and limits'' paragraph above (``Conjecture 19 for height five
remains open here'') describes the manuscript; Part~\ref{shr:h5:part} proves
Conjecture~19.

\emph{Pin.} The manuscript names no ProveIt commit. It pins this report's
source by its Git blob \texttt{2b40ded9\allowbreak daa6549f\allowbreak 786a52b8\allowbreak 80d5477d\allowbreak e3b5adfc},
which is \texttt{article.tex} as committed in \texttt{6fef5383b}
(3~October 2026, the batch-85 write) and unchanged until
\texttt{b71545625} added Part~\ref{shr:bp:part} on 5~October 2026. So it
read Part~\ref{shr:part:one} alone, with the numbers printed here. It also
inspected a 24-page ``user-held copy'' of the manuscript that is now
Part~\ref{shr:bp:part}, which it calls ``this private companion'' and cites
for overlap disclosure only, not as a mathematical dependency; that
manuscript is public as Part~\ref{shr:bp:part}.

\emph{Files.} Its four programs, its receipt, its code README and its
source record are shipped with the prefix \texttt{03-height-four-}, in
\texttt{code/}, in \texttt{data/} and at the report root; the note at the
end of Section~33.3 lists them. Its
\LaTeX\ source, delivered README, PDF and \texttt{MANIFEST.sha256} (19
entries, all verified at placement) are not shipped and survive in the
arrival commit's archive. The manuscript is unrefereed and AI-assisted; its
proofs are ordinary mathematical proofs, nothing in it is formalized in
Lean or Rocq, and its place in this report confers no formal status.

\emph{What it answers.} Theorems~\ref{shr:h4:thm:main},
\ref{shr:h4:thm:oeis} and~\ref{shr:h4:thm:sum} prove, from the array
definition and for every $n\ge1$ (the finite sum for $n>1$), the
order-two, degree-nine recurrence displayed as conjectural in OEIS A181198
and the finite sum printed as Conjecture~18 of Kauers and Koutschan
\cite{KK}. This answers the height-four half of Research question~1 of
Part~\ref{shr:part:one} and of Research question~8 of
Part~\ref{shr:bp:part}; Part~\ref{shr:h5:part} answers the height-five
half. Dated notes at both questions, after Remark~8.1, at Research
questions~6, 7 and~11, and after Remark~\ref{shr:bp:rem:constructive} say
so. The manuscript also shows that the operator does not extend to $n=0$
under $a_0=1$ (residual $-2160$). The write adds Section~\ref{shr:h4:sec:minimal}:
the operator has minimal order, and every order-two recurrence for A181198
is a polynomial multiple of it.

\emph{Overlap, credited.} Section~\ref{shr:h4:sec:asymptotics} recovers,
from the recurrence, the case $m=4$ of Part~\ref{shr:part:one}'s
Theorem~\ref{shr:thm:main}, as the manuscript says: its $c_j$ are
$\tfrac3{1024}c_{4,j}$ ($J_4/4^6=3/1024$), its $d_j$ are $b_{4,j}$, and
$d_1,\ldots,d_5$ equal the entries of \texttt{data/\allowbreak
coefficients\_\allowbreak m1\_\allowbreak to\_\allowbreak m6.json}. The sixth coefficient
$d_6=-79147528275/268435456$ is new to the repository; at placement,
Part~\ref{shr:part:one}'s engine \texttt{code/\allowbreak coefficients.py} run at order
six on a copy gave the same $b_{4,6}$, so two independent routes agree.
Its $C_*$, $\beta_*$ and $\lambda$ are $K_4$, $\alpha_4$ and $4\log4$. The
effective error constants of Sections~\ref{shr:h4:sec:asymptotics}
and~\ref{shr:h4:sec:inverse} are new; the finite Lambert model and the
two-ceiling envelope use the apparatus of Section~\ref{shr:sec:inverse},
Section~\ref{shr:bp:sec:inverse} and Theorems~J.23 and~J.43 of the
transseries volume (see the note at the end of Section~10.2). The transfer,
the rank-two identity and the scalar integrals of
Sections~\ref{shr:h4:sec:transfer}--\ref{shr:h4:sec:rank} are re-derived
independently in Part~\ref{shr:h5:part}, and so is the period recurrence,
Theorem~\ref{shr:h4:thm:Z} (Remark~\ref{shr:h5:rem:sameZ}). The text shares
0.0\,\% and 0.6\,\% of its word 8-grams with Parts~\ref{shr:part:one}
and~\ref{shr:bp:part}, and 5.7\,\% with Part~\ref{shr:h5:part}.

\emph{Source checks.} The manuscript transcribes Conjecture~18 from
arXiv:2303.02793v2 (24~April 2023, Section~6.4, printed page~33) and the
operator from the A181198 internal record, retrieved 5~October 2026
(revision~39 of 1~January 2024, still labelled conjectural); its bounded
priority search, which included arXiv:2607.24832, found no matching proof
(Section~\ref{shr:h4:sec:scope}). At placement the intake did not
re-inspect the printed page of \cite{KK}; it checked that $p_0,p_1,p_2$
equal the polynomials of the live A181198 entry, checked the certificate
\eqref{shr:h4:eq:polycert}, the edge identities and \eqref{shr:h4:eq:edgeflux}
in SymPy (residual zero), and checked Theorems~\ref{shr:h4:thm:main},
\ref{shr:h4:thm:oeis} and~\ref{shr:h4:thm:sum}, \eqref{shr:h4:eq:affine}
and \eqref{shr:h4:eq:Zrec} for $n\le30$ against an independent row-state
count. Its verifier, rerun at the write on a copy (Section~33.3), printed
its receipt byte for byte.

\emph{Editorial changes.} Section~$k$ of the manuscript is
Section~$k+25$ here; theorem and equation numbers follow their sections.
Every label carries the prefix \texttt{shr:h4:}. Its title, subtitle and
title-page paragraphs open this Part; its table of contents is dropped. Its
two citations of its predecessor now point to Part~\ref{shr:part:one} and
its citation of the companion to Part~\ref{shr:bp:part}. The bibliographies
are merged: Kauers--Koutschan, Sun and DLMF~\S5.11 are
Part~\ref{shr:part:one}'s entries (the manuscript retrieved DLMF~5.11.1 and
\S5.11(ii) on 5~October 2026 for the positive-real Stirling remainder
bound), and its OEIS entry is \cite{oeis198int}. Its bibliography said of
Part~\ref{shr:part:one}, pinned by the blob above, that ``its all-height
asymptotics are prior results relative to the present report; its scope
leaves the particular recurrences unproved'', and of
Part~\ref{shr:bp:part} that ``the inspected copy proves rational
diagonality and D-finiteness and records the particular recurrence problem
as open''. Remarks take this report's theorem style, and the four
command-line options of Section~33.3 are typeset \texttt{-{}-}\emph{option},
because this report's typewriter font would join their double hyphens into
a dash. The \textbf{[write]}
notes, Section~\ref{shr:h4:sec:minimal} and
Section~\ref{shr:h4:sec:further} are new. No other statement, proof or
number of the manuscript was changed.
\end{writenote}

\begin{writenote}[Added 5 October 2026, batch~101: notation of Part~III]
Part~III keeps the manuscript's symbols. Many clash with those of the
other Parts; within Part~III each has the meaning below, and the right
column gives the tempting false reading from elsewhere in this report.
The count is $a_n=T_4(n)$ and the array entries are $x_{i,j}$. Shared
without change: $T_m(n)$, $\Delta$, the Lambert branch $W_{-1}$ and the
Bernoulli numbers $B_{2j}$.
{\footnotesize
\begin{longtable}{@{}>{\raggedright\arraybackslash}p{0.17\linewidth}>{\raggedright\arraybackslash}p{0.43\linewidth}>{\raggedright\arraybackslash}p{0.34\linewidth}@{}}
\toprule
symbol & meaning in Part~III & \emph{not}: meaning elsewhere\\
\midrule\endhead
$a_n$; $M_n$; $b_n$ & $T_4(n)$; $M_4(n)=(4n)!/(n!)^4$; $a_n/M_n$ \eqref{shr:h4:eq:normalizers} & array entries $a_{i,j}$; the coefficients $b_{m,j}$ of \eqref{shr:eq:mainK}\\
$A_n$, $B_n$, $D_n$ & coefficients of \eqref{shr:h4:eq:inhom} & $A=\binom m3$, $B=m(m-1)^2(m-2)/12$ (Section~6.1); Part~II's ranks $B_i,D_i$\\
$V$, $K$, $W$ & the Volterra kernel, $K=V^{n-1}$ \eqref{shr:h4:eq:K}, $W=V^n(V^*)^n$ \eqref{shr:h4:eq:W} & $K_m$ \eqref{shr:eq:Km}; Part~II's $K_r(u,v)$ and $W_n(x)$\\
$E$; $E_K$, $E'$ & the sign kernel \eqref{shr:h4:eq:E}; error constants (Sections~\ref{shr:h4:sec:asymptotics}--\ref{shr:h4:sec:inverse}) & the boundary error $E_{m,n}(p)$; $\E$\\
$J$; $J_0$; $J_n$ & $\one\otimes\one$; a scalar \eqref{shr:h4:eq:scalars}; $D_nM_n$ in Section~30.2 & $J_m=\prod_{j<m}j!$\\
$\mathcal A$; $B$ & the kernels $KEK^*$ and $E\mathcal A$ \eqref{shr:h4:eq:ABkernels} & the event $\mathcal A$ of Section~4\\
$\alpha(t)$, $\beta(t)$; $\beta_*$ & $t^n/n!$ and $t^{n-1}/(n-1)!$ \eqref{shr:h4:eq:alphabeta}; $15/2=\alpha_4$ & $\alpha_m=(m^2-1)/2$; recurrence coefficients $\beta_j$\\
$h$, $e$, $w$; $c$, $H$, $U$, $I$, $K_0$, $L$ & functions and scalars of Section~\ref{shr:h4:sec:rank} & $h_r(q)$, $h_j$; $H(\lambda)$; $L$ a truncation order\\
$N$ & $n!$ (Sections~\ref{shr:h4:sec:rank}--\ref{shr:h4:sec:certificate}); an integer onset (Section~\ref{shr:h4:sec:asymptotics}); the threshold $N(y)$ \eqref{shr:h4:eq:threshold} & ---\\
$Z_n$ & the weighted triangle period \eqref{shr:h4:eq:Zdef}, \eqref{shr:h4:eq:Ztriangle}, the same as Part~IV's $Z_n$ & $Z_m(n)$ \eqref{shr:eq:binomZ}; $Z_i$ (Section~9)\\
$\mathcal C_m$; $\mathcal D$, $f$; $I_j$ & the ordered simplex; the triangle and $uv(u+v-1)$ \eqref{shr:h4:eq:triangle}; triangle integrals (Section~\ref{shr:h4:sec:certificate}) & the central event $\mathcal C_n$; Part~II's chamber $\mathcal C_r(n)$ and Dyck words $\mathcal D_m$\\
$P$, $Q$; $p_{ij}$; $\gamma_n$ & the certificate field \eqref{shr:h4:eq:PQ}; its coefficients \eqref{shr:h4:eq:pcoeff} & Part~I's $P_j$; Part~II's $Q_m(n)$, $p_i$; Part~IV's $P(n)$, $Q(n)$, $\gamma_n$\\
$g_n$, $r_n$, $m_n$, $q_n$, $c_n$ & Section~\ref{shr:h4:sec:finite} & $c_{m,j}$; $r=m-2$ (Part~II)\\
$p_0$, $p_1$, $p_2$ & the A181198 operator \eqref{shr:h4:eq:poperator} & power sums $p_r$\\
$F_n$, $s_n$, $t_k$ & variation of constants \eqref{shr:h4:eq:F}, \eqref{shr:h4:eq:t} & the generating function $F_m(z)$; Part~IV's $F_n=a_n/M_n$\\
$(x)_j$ & the rising factorial (Section~30.4) & Part~I's falling factorial $n^{\underline k}$; Part~II and Part~IV write rising factorials $(x)^{\overline k}$\\
$R(n)$, $S(n)$ & the coefficients \eqref{shr:h4:eq:R}, \eqref{shr:h4:eq:S} & $R_n$, $R_m$, $R(z)$, $R_m(n)$; Part~IV's ratios $R(n)$, $S(n)$\\
$C(t)$, $c_j$; $d_j$ & the formal series of $n^6b_n$, with $c_j=\tfrac3{1024}c_{4,j}$; $d_j=b_{4,j}$ & the Catalan series $C(z)$; $c_{m,j}$; $d=\binom m2$\\
$C_*$, $\lambda$, $\ell_j$, $s_j$ & $K_4$; $\log256$; logarithmic and inverse coefficients & $\lambda=m\log m$ (Section~10)\\
$\Psi_K$, $\xi_K$, $w(y)$, $\Delta_K$, $G_K$, $H_K$, $W_K$, $X_K$, $C_K$ & the finite inverse model of Section~\ref{shr:h4:sec:inverse} & the Vandermonde $\Delta$; $W_{-1}$; Part~II's $C_{m,L}$\\
\bottomrule
\end{longtable}}
\end{writenote}

\section{The theorem and its scope}
\label{shr:h4:sec:scope}
For integers $m,n\ge1$, let $T_m(n)$ count the $m\times n$ arrays
$(x_{i,j})$ containing the integers $1,\ldots,mn$ exactly once and satisfying
\begin{align}
 x_{i,j}&<x_{i,j+1},& x_{i,j}&<x_{i+1,j},\notag\\
 x_{i,j}&<x_{i+1,j+1},&x_{i,j+1}&<x_{i+1,j},
 \label{shr:h4:eq:literal}
\end{align}
whenever the cells exist. Rows increase from left to right and columns from
top to bottom. The third inequality follows from the first two. The fourth,
for a downward antidiagonal, is essential. Throughout this report
\emph{height is four and width is $n$}. Write
\begin{equation}
 a_n=T_4(n),\qquad M_n=\frac{(4n)!}{(n!)^4},\qquad b_n=\frac{a_n}{M_n}.
 \label{shr:h4:eq:normalizers}
\end{equation}
The initial values are
\[
 a_1=a_2=1,\qquad a_3=8,\quad a_4=169,\quad a_5=6392.
\]
Width two gives a single interlaced chain, explaining $a_2=1$ without
recourse to a recurrence.

\begin{theorem}[The exact inhomogeneous identity]
\label{shr:h4:thm:main}
For every integer $n\ge1$,
\begin{equation}
 A_n a_{n+1}+B_n a_n=D_n M_n,
 \label{shr:h4:eq:inhom}
\end{equation}
where
\begin{align}
 A_n&=3(n-1)(n+1)(3n+1)(3n+2),\notag\\
 B_n&=8n(2n+1)(4n-1)(4n+1),\label{shr:h4:eq:ABD}\\
 D_n&=\frac{6(4n+1)(7n^2-1)}{(n+1)^2(2n-1)^2(2n+1)}.\notag
\end{align}
For $n\ge2$ this identity and $a_2=1$ determine the whole sequence.
It implies the exact order-two operator in Theorem~\ref{shr:h4:thm:oeis}
and the printed finite sum in Theorem~\ref{shr:h4:thm:sum}.
\end{theorem}

The first-order coefficient $A_1$ is zero, so the identity at $n=1$ fixes
$a_1$ rather than $a_2$. We keep this degeneracy explicit. The resulting
second-order recurrence is valid for $n\ge1$. It does \emph{not} hold at
$n=0$ if one sets the empty-array value $a_0=1$.

\subsection{What is resolved and what is reused}
Kauers and Koutschan~\cite{KK} give the height-four formula as
Conjecture~18 in Section~6.4, printed page~33 of arXiv version~2.
Their statement is a finite sum, rather than merely an unnamed
D-finiteness assertion. The corresponding order-two recurrence is printed
in OEIS A181198~\cite{oeis198int}. Sections~\ref{shr:h4:sec:transfer}--\ref{shr:h4:sec:finite}
prove both exact statements from the full array definition.

The earlier report \emph{Fixed-height shifted rectangles}
(Part~\ref{shr:part:one}) already proves the leading asymptotic and all-fixed-order
expansions for every fixed height. Its scope expressly leaves the specific
height-four and height-five recurrences unproved. A later companion report
on rational diagonals and rare boundary phases (Part~\ref{shr:bp:part}) establishes
D-finiteness by a different route and likewise leaves these particular
recurrences unresolved. We do not reclassify those earlier asymptotic or
D-finiteness results as new. The additional exact result here is the
counting-to-operator certificate for height four, together with the
simpler coefficient and remainder machinery that this recurrence permits.

The proof uses classical order-polytope, determinant-integration, and
Pfaffian ideas. The finite-dimensional integral identities needed here are
proved directly, so no unstated Fredholm determinant theorem or
symbolic-integration oracle is a mathematical dependency. Sun's work on
shifted strips~\cite{Sun} supplies relevant background; the orientation
with fixed width and growing height is a different enumeration problem.

\subsection{Source status and limitations}
The source check used the primary 2023 article, the A181198 entry retrieved
on October~5, 2026, and the existing report source. The retrieved OEIS
internal snapshot still calls the recurrence and formula conjectural; its
revision record is revision~39, January~1, 2024. That label is evidence of
the inspected entry's status, not proof of global novelty. Targeted
identifier and exact-coefficient searches, including inspection of the
2026 paper arXiv:2607.24832, did not locate a matching proof. These are
bounded search results. The report makes no first-ever or exhaustive
priority claim.

The height-five operator and Conjecture~19 are outside the theorem.
Splitting the count into asymptotic event sectors does not separately
annihilate those sectors with the whole-sequence operator; the proof below
avoids that issue by retaining the entire order polytope. The expansions
in Section~\ref{shr:h4:sec:asymptotics} are Poincar\'e expansions at every
\emph{fixed} order. They are not assertions about convergence, optimal
truncation, Borel summation, or a complete exponentially improved
transseries.

\subsection{Outline of the proof}
The exact argument has four stages:
\begin{enumerate}
\item Compose the interlacing indicator along consecutive columns, with
      ordered-simplex normalization retained exactly.
\item Evaluate the resulting four-variable determinant integral in terms
      of two traces, and reduce those traces to one scalar period $Z_n$.
\item Prove a first-order recurrence for $Z_n$ by a polynomial divergence
      identity on a triangle, evaluating every edge.
\item Substitute back into the full count, eliminate the hypergeometric
      forcing term, and solve by variation of constants.
\end{enumerate}
Every identity is uniform in the integer parameter. Finite counting tests
are separate transcription and implementation checks, never replacements
for these arguments.

\section{Column transfer and the Pfaffian identity}
\label{shr:h4:sec:transfer}
\subsection{The order polytope and interlacing determinant}
Let
\[
 \mathcal C_m=\{(x_1,\ldots,x_m):0<x_1<\cdots<x_m<1\}.
\]
Give each array cell a real coordinate in $(0,1)$ and impose the same
strict comparisons as in \eqref{shr:h4:eq:literal}. Coordinates that coincide
form a null set. Every admissible total order of the $mn$ coordinates
contributes an open simplex of volume $1/(mn)!$. Hence the volume of this
order polytope is $T_m(n)/(mn)!$.

For two consecutive columns $x,y\in\mathcal C_m$, the comparisons are
exactly
\begin{equation}
 x_1<y_1<x_2<y_2<\cdots<x_m<y_m.
 \label{shr:h4:eq:interlace}
\end{equation}
Rows give $x_i<y_i$ and antidiagonals give $y_i<x_{i+1}$. These comparisons
also imply all column and downward-diagonal comparisons involving these
columns. For $n\ge2$ the adjacent-column conditions therefore include all
constraints, with no restriction on whether a row has begun or finished
in a label-prefix description.

Define the Volterra kernel and its operator by
\[
 V(t,s)=\one_{\{s<t\}},\qquad (Vg)(t)=\int_0^t g(s)\dd s.
\]
For ordered distinct $x,y$, we have
\begin{equation}
 \det[V(y_i,x_j)]_{i,j=1}^m
 =\one_{\{x_1<y_1<\cdots<x_m<y_m\}}.
 \label{shr:h4:eq:indicator}
\end{equation}
Indeed, let $k_i$ be the number of $x_j$ below $y_i$. Row $i$ consists of
$k_i$ ones followed by zeros, with $0\le k_1\le\cdots\le k_m\le m$.
A nonzero determinant requires distinct nonzero rows, hence $k_i=i$.
In that case the matrix is unit lower triangular. This proves both the
indicator and its positive sign.

\begin{lemma}[Composition on the ordered simplex]
\label{shr:h4:lem:andreief}
For bounded kernels $F,G$ on $[0,1]^2$,
\begin{equation}
 \int_{\mathcal C_m}
 \det[F(y_i,z_j)]\det[G(z_i,x_j)]\dd z
 =\det\left[\int_0^1F(y_i,t)G(t,x_j)\dd t\right].
 \label{shr:h4:eq:andreief}
\end{equation}
\end{lemma}
\begin{proof}
The product of determinants is symmetric in the $z$ variables, because
permuting them changes the sign of each factor. Its cube integral is thus
$m!$ times its ordered-simplex integral. Expand both determinants in the
cube integral and integrate each product by Fubini. For every relative
permutation contributing to the determinant on the right there are $m!$
choices of the first permutation, all with the same resulting sign.
That factor cancels the ordered-simplex factor. No factorial remains.
\end{proof}

Composing the $n-1$ adjacent-column indicators now gives, for $n\ge2$,
\begin{equation}
 \frac{T_m(n)}{(mn)!}
 =\int_{\mathcal C_m}\int_{\mathcal C_m}
       \det[K(y_i,x_j)]\dd x\dd y,
 \quad K=V^{n-1},
 \label{shr:h4:eq:transfer}
\end{equation}
where
\begin{equation}
 K(t,s)=\frac{(t-s)_+^{n-2}}{(n-2)!}.
 \label{shr:h4:eq:K}
\end{equation}
For $n=2$ the exponent-zero notation denotes the indicator $s<t$;
its value on equality is immaterial. Formula~\eqref{shr:h4:eq:K} follows by
integrating the simplex $s<t_1<\cdots<t_{n-2}<t$. The case $n=1$ is
handled by its direct count and never requires an identity distribution
as an integral kernel.

\subsection{An elementary four-variable de Bruijn calculation}
Set
\begin{equation}
 E(t,s)=\sgn(s-t),\qquad J=\one\otimes\one,
 \qquad E=J-2V=2V^*-J.
 \label{shr:h4:eq:E}
\end{equation}
Values on the diagonal may be assigned arbitrarily. For an antisymmetric
$4\times4$ matrix $C$, our Pfaffian convention is
\[
 \Pf C=C_{12}C_{34}-C_{13}C_{24}+C_{14}C_{23}.
\]
In particular, for $x\in\mathcal C_4$, every upper-triangular entry of
$[E(x_i,x_j)]$ is $1$, so its Pfaffian is $1-1+1=1$.

Let $f_i$ be bounded functions and put
\[
 A_{ij}=\int_0^1\int_0^1 f_i(s)E(s,t)f_j(t)\dd s\dd t.
\]
Then
\begin{equation}
 \int_{\mathcal C_4}\det[f_i(x_j)]\dd x=\Pf A.
 \label{shr:h4:eq:debruijn}
\end{equation}
Here is a direct derivation. If $D(x)=\det[f_i(x_j)]$, alternation gives
\[
 \Pf A=\frac18\int_{[0,1]^4}D(x)E(x_1,x_2)E(x_3,x_4)\dd x.
\]
This follows by expanding the standard $1/(2^2\,2!)$ permutation formula
for the Pfaffian. Relabeling variables shows that the integral of
$D(x)(-E_{13}E_{24})$ and that of $D(x)E_{14}E_{23}$ each equal the
integral of $D(x)E_{12}E_{34}$. Thus the right side is
\[
 \frac1{24}\int_{[0,1]^4}D(x)\Pf[E(x_i,x_j)]\dd x.
\]
The integrand is symmetric, and on the ordered simplex its Pfaffian factor
is $1$. This is exactly the left side of \eqref{shr:h4:eq:debruijn}. The
orientation can also be checked in size two: the corresponding entry is
\[
 \int_{s<t}\bigl(f_i(s)f_j(t)-f_i(t)f_j(s)\bigr)\dd s\dd t.
\]

Apply \eqref{shr:h4:eq:debruijn} to $f_i(s)=K(y_i,s)$ and define
\begin{equation}
 \mathcal A=KEK^*,\qquad B=E\mathcal A=EKEK^*.
 \label{shr:h4:eq:ABkernels}
\end{equation}
The kernel $\mathcal A$ is antisymmetric. The inner integral in
\eqref{shr:h4:eq:transfer} is $\Pf[\mathcal A(y_i,y_j)]$.

\subsection{The two traces and their exact constants}
The product $\Pf[\mathcal A(y_i,y_j)]\Pf[E(y_i,y_j)]$ is symmetric, so
\begin{equation}
 \frac{a_n}{(4n)!}
 =\frac1{24}\int_{[0,1]^4}\Pf[\mathcal A_{ij}]\Pf[E_{ij}]\dd y.
 \label{shr:h4:eq:doublepf}
\end{equation}
In this subsection, a trace denotes its finite kernel integral. Thus
\[
 \tr B=\int_0^1B(t,t)\dd t,\qquad
 \tr(B^2)=\int_0^1\int_0^1B(t,s)B(s,t)\dd s\dd t.
\]
Expanding the two three-term Pfaffians gives nine products. The three
matching pairings each integrate to $(\tr B)^2$. For example,
$\int E_{12}\mathcal A_{12}=-\tr B$, and its square has the positive
sign. The six signed crossed pairings each integrate to $-\tr(B^2)$.
One representative follows from
\[
 \tr(B^2)=\int E_{12}\mathcal A_{23}E_{34}\mathcal A_{41}\dd y
 =-\int E_{12}E_{34}\mathcal A_{14}\mathcal A_{23}\dd y;
\]
the others follow by relabeling the four variables and retaining the
Pfaffian signs. Consequently
\begin{equation}
 \boxed{\displaystyle
 \frac{a_n}{(4n)!}=\frac{(\tr B)^2}{8}-\frac{\tr(B^2)}4,\qquad n\ge2.}
 \label{shr:h4:eq:traceformula}
\end{equation}
The constants are exactly $3/24$ and $6/24$. All the kernels used here
are bounded, piecewise polynomial functions. Products are absolutely
integrable on compact cubes, so Fubini and the relabelings are justified.
This is a finite integral calculation, not a formal infinite-dimensional
Fredholm series argument.

\section{Rank reduction to one weighted period}
\label{shr:h4:sec:rank}
Throughout this section $n\ge2$ and $N=n!$. To keep the count $a_n$
distinct from scalar functions, use
\begin{equation}
 \alpha(t)=\frac{t^n}{N},\qquad
 \beta(t)=\frac{t^{n-1}}{(n-1)!},\qquad \alpha'=\beta.
 \label{shr:h4:eq:alphabeta}
\end{equation}
Write $p\otimes q$ for the rank-one kernel $p(t)q(s)$, and let
$\ip pq=\int_0^1p(t)q(t)\dd t$. Define
\begin{align}
 W&=V^n(V^*)^n,\notag\\
 W(t,s)&=\frac1{((n-1)!)^2}
  \int_0^{\min(t,s)}(t-r)^{n-1}(s-r)^{n-1}\dd r,\label{shr:h4:eq:W}\\
 h(t)&=W(t,1),\qquad e(t)=2h(t)-\frac{\beta(t)}N.\notag
\end{align}
The kernel $W$ is symmetric and continuous, with polynomial expressions
on the two triangles of the square.

\subsection{An exact rank-two perturbation}
Use both expressions for $E$ in \eqref{shr:h4:eq:E}, preserving factor order:
\begin{align*}
 B&=(JV^{n-1}-2V^n)
       \bigl(2(V^*)^n-J(V^*)^{n-1}\bigr)\\
  &=-4W+\one\otimes e+2\alpha\otimes\beta.
\end{align*}
Thus
\begin{equation}
 B=-4W+\one\otimes e+2\alpha\otimes\beta.
 \label{shr:h4:eq:ranktwo}
\end{equation}
For clarity, the four product reductions are
\begin{align*}
 JV^{n-1}(V^*)^n&=\one\otimes h,&
 JV^{n-1}J(V^*)^{n-1}&=\frac{\one\otimes\beta}N,\\
 V^n(V^*)^n&=W,&
 V^nJ(V^*)^{n-1}&=\alpha\otimes\beta.
\end{align*}
For the first, integrate the first argument of $V^{n-1}$ to obtain
$(1-r)^{n-1}/(n-1)!$, leaving exactly $h$. For the second, the intervening
scalar integral is $1/N$. The remaining two follow directly from
$V^n\one=\alpha$ and $V^{n-1}\one=\beta$.

The kernel identities needed below are
\begin{align}
 W(t,0)&=0,& (\partial_t+\partial_s)W(t,s)&=\beta(t)\beta(s),
 \label{shr:h4:eq:Wdiff}\\
 W(rt,rs)&=r^{2n-1}W(t,s),
 \label{shr:h4:eq:Whom}
\end{align}
and the scalar integrals are
\begin{align}
 \int_0^1h&=\frac1{2N^2},&
 \int_0^1\beta&=\frac1N,&
 c:=\ip\alpha\beta&=\frac1{2N^2},\notag\\
 \int_0^1e&=0,&
 \tr W&=\frac1{2n(2n-1)((n-1)!)^2}.
 \label{shr:h4:eq:elementary}
\end{align}
To prove the differential identity, the sum of the two outer derivatives
of the integrand of \eqref{shr:h4:eq:W} is minus its derivative in $r$. The lower
endpoint produces $\beta(t)\beta(s)$. The upper endpoint vanishes for
$n\ge2$. The identities hold almost everywhere, which is sufficient for
the following integrations. When the square is split at $t=s$, its inner
boundary contributions cancel because the kernels are continuous.
Homogeneity follows by the substitution $r\mapsto r\rho$ with
$\rho>0$. The integral of $h$ can be checked explicitly by Fubini:
\[
 \int_0^1h(t)\dd t
 =\frac1{n((n-1)!)^2}\int_0^1(1-r)^{2n-1}\dd r
 =\frac1{2N^2}.
\]
The remaining integrals are monomial integrals.

Taking the trace in \eqref{shr:h4:eq:ranktwo} now gives
\begin{equation}
 \tr B=-4\tr W+\int e+2c
       =-\frac1{(2n-1)N^2}.
 \label{shr:h4:eq:trB}
\end{equation}

\subsection{Every integration by parts for the square trace}
Introduce
\begin{equation}
 \begin{gathered}
 H=\ip\alpha h,\qquad U=\tr(W^2),\qquad w=W\one,\\
 J_0=\ip h w,\qquad K_0=\ip\beta w,\qquad I=\ip\beta{W\alpha}.
 \end{gathered}
 \label{shr:h4:eq:scalars}
\end{equation}
Integrate the sum of the two derivatives of
$\alpha(t)\alpha(s)W(t,s)$ over $[0,1]^2$. The lower edges vanish, because
$\alpha(0)=0$ and $W(t,0)=0$. Each upper edge contributes $H/N$.
The differentiated integrand integrates to $2I+c^2$, using symmetry of
$W$ and \eqref{shr:h4:eq:Wdiff}. Hence
\begin{equation}
 2I+c^2=\frac{2H}N,\qquad I=\frac HN-\frac{c^2}2.
 \label{shr:h4:eq:I}
\end{equation}
Next apply the same operation to $\alpha(t)W(t,s)$. Its two upper edges
are $(\int h)/N$ and $H$, and its two lower edges vanish. Its interior
integral is $K_0+c/N$. By \eqref{shr:h4:eq:elementary},
\begin{equation}
 K_0+\frac cN=H+\frac1N\int h,
 \qquad K_0=H.
 \label{shr:h4:eq:K0}
\end{equation}

For a rank-one kernel, $\tr(p\otimes q)=\ip qp$ and
$(p\otimes q)(r\otimes s)=\ip qr\,p\otimes s$.
Squaring \eqref{shr:h4:eq:ranktwo} therefore gives exactly
\begin{equation}
 \tr(B^2)=16U-8\ip e w-16I+4c^2+\frac4N\ip e\alpha.
 \label{shr:h4:eq:trsquareexpand}
\end{equation}
The term $(\int e)^2$ has vanished; the cross term of the two rank-one
kernels supplies the last term, including its factor $4$. Since
\[
 \ip e w=2J_0-\frac HN,\qquad
 \ip e\alpha=2H-\frac cN,
\]
substitution of \eqref{shr:h4:eq:I} reduces \eqref{shr:h4:eq:trsquareexpand} to
\begin{equation}
 \tr(B^2)=16U-16J_0+\frac1{N^4}.
 \label{shr:h4:eq:trsquareJ}
\end{equation}
All occurrences of $H/N$ in this substitution cancel.

Integrating \eqref{shr:h4:eq:Wdiff} with respect to $s$ gives
\[
 w'(t)+h(t)=\frac{\beta(t)}N,
 \qquad w(0)=0,\quad w(1)=\frac1{2N^2}.
\]
Multiply by $w$ and integrate. The only integration-by-parts boundary
is $[w(t)^2/2]_0^1=1/(8N^4)$. By \eqref{shr:h4:eq:K0},
\begin{equation}
 J_0=\frac{K_0}N-\frac1{8N^4}
     =\frac HN-\frac1{8N^4}.
 \label{shr:h4:eq:J0}
\end{equation}
Thus
\begin{equation}
 \tr(B^2)=16U-\frac{16H}N+\frac3{N^4}.
 \label{shr:h4:eq:trB2}
\end{equation}
This calculation has retained the full trace, including its boundary
terms. No asymptotic replacement has been made.

\subsection{Reducing the remaining square integral}
Symmetry and \eqref{shr:h4:eq:Whom} give, by $t=rs$ on $0<t<s<1$,
\begin{align}
 U&=2\int_0^1\int_0^sW(t,s)^2\dd t\dd s\notag\\
  &=2\int_0^1\int_0^1s^{4n-1}h(r)^2\dd r\dd s
   =\frac1{2n}\int_0^1h(t)^2\dd t.
 \label{shr:h4:eq:Uhom}
\end{align}
The Jacobian is $s$, which is responsible for the exponent $4n-1$.
The defining integral for $h$ obeys
\begin{equation}
 (1-t)h'(t)+(2n-1)h(t)
 =\frac{t^{n-1}}{((n-1)!)^2}.
 \label{shr:h4:eq:hODE}
\end{equation}
For $n\ge2$, integrate minus the $r$ derivative of
$(t-r)^{n-1}(1-r)^n$. The negative derivative is
\[
 (n-1)(1-t)(t-r)^{n-2}(1-r)^{n-1}
 +(2n-1)(t-r)^{n-1}(1-r)^{n-1}.
\]
The lower endpoint contributes $t^{n-1}$ and the upper endpoint vanishes.
This proves \eqref{shr:h4:eq:hODE}. At $n=1$, the same equation is checked
directly from $h(t)=t$.

Let $L=\int_0^1 t^{n-1}h(t)\dd t$. Multiplying \eqref{shr:h4:eq:hODE} by $2h$
and integrating yields
\begin{equation}
 (4n-1)\int_0^1h^2=\frac{2L}{((n-1)!)^2}.
 \label{shr:h4:eq:h2}
\end{equation}
Indeed, the first term integrates to
$[(1-t)h(t)^2]_0^1+\int h^2=\int h^2$; both endpoints are zero, by
$h(0)=0$ and the factor $1-t$.
Multiplication instead by $t^n$ gives
\[
 -nL+3n\int_0^1t^nh(t)\dd t=\frac1{2n((n-1)!)^2}.
\]
Here $[t^n(1-t)h(t)]_0^1=0$, and differentiating $t^n(1-t)$ produces
$n t^{n-1}-(n+1)t^n$. Since $\int t^nh=NH$, this says
\begin{equation}
 L=3NH-\frac1{2N^2}.
 \label{shr:h4:eq:L}
\end{equation}
Combining \eqref{shr:h4:eq:Uhom}--\eqref{shr:h4:eq:L} gives
\begin{equation}
 U=\frac{3nH}{(4n-1)N}-\frac{n}{2(4n-1)N^4}.
 \label{shr:h4:eq:UH}
\end{equation}

\subsection{The exact full-count formula}
Substitute \eqref{shr:h4:eq:trB}, \eqref{shr:h4:eq:trB2}, and \eqref{shr:h4:eq:UH} into
\eqref{shr:h4:eq:traceformula}. Define
\begin{equation}
 Z_n=N^3H=N^2\int_0^1t^nh(t)\dd t.
 \label{shr:h4:eq:Zdef}
\end{equation}
After collecting rational terms, the result is
\begin{equation}
 \frac{a_n}{M_n}
 =\frac{4(n-1)}{4n-1}Z_n
 -\frac{32n^3-56n^2+28n-5}{8(2n-1)^2(4n-1)}.
 \label{shr:h4:eq:affine}
\end{equation}
For example, before the final simplification the right side is
\[
 \frac1{8(2n-1)^2}-4N^4U+4N^3H-\frac34,
\]
which provides an explicit check on the normalization.

In the double integral defining \eqref{shr:h4:eq:Zdef}, set $u=t$ and $v=1-r$.
The domain $0<r<t<1$ becomes the triangle
\begin{equation}
 \mathcal D=\{(u,v):0<u<1,\ 0<v<1,\ u+v>1\},
 \qquad f(u,v)=uv(u+v-1).
 \label{shr:h4:eq:triangle}
\end{equation}
The Jacobian has absolute value $1$, and the factorial factor is
$N^2/((n-1)!)^2=n^2$. Thus
\begin{equation}
 Z_n=n^2\int_{\mathcal D}u f^{n-1}\dd u\dd v
     =\frac{n^2}2\int_{\mathcal D}(u+v)f^{n-1}\dd u\dd v.
 \label{shr:h4:eq:Ztriangle}
\end{equation}
The second equality uses the symmetry exchanging $u$ and $v$.
Formula~\eqref{shr:h4:eq:affine} was derived for $n\ge2$. It also holds at $n=1$
by direct checking: its coefficient of $Z_1$ is zero, and its remaining
value is $1/24=a_1/M_1$. Independently,
\eqref{shr:h4:eq:Ztriangle} gives $Z_1=1/3$. This extension uses no singular
kernel calculation.

\section{A polynomial certificate on the triangle}
\label{shr:h4:sec:certificate}
Set
\[
 I_j=\int_{\mathcal D}(u+v)f(u,v)^j\dd u\dd v,\qquad j\ge0,
\]
so $Z_n=n^2I_{n-1}/2$. We now prove the exact period recurrence.

\begin{theorem}[The weighted triangle recurrence]
\label{shr:h4:thm:Z}
For every integer $n\ge1$,
\begin{equation}
 3(3n+1)(3n+2)Z_{n+1}+(n+1)^2Z_n
 =\frac{28n^3+50n^2+29n+5}{2(2n+1)}.
 \label{shr:h4:eq:Zrec}
\end{equation}
\end{theorem}

\subsection{The six coefficients}
For a fixed integer $n\ge1$, define
\begin{align}
 \gamma_n&=\frac{3(3n+1)(3n+2)}{n^2},\notag\\
 p_{21}&=\frac{(n+3)(3n+1)(3n+2)}{n^2(n+1)^2},&
 p_{20}&=-\frac{3n+1}{n(n+1)},\notag\\
 p_{12}&=\frac{3(3n+1)(3n+2)}{n^2(n+1)},&
 p_{11}&=-\frac{(3n+1)(7n^2+14n+6)}{n^2(n+1)^2},\label{shr:h4:eq:pcoeff}\\
 p_{10}&=\frac{n-1}{n(n+1)},&p_{00}&=\frac1n.\notag
\end{align}
Use the polynomial vector field
\begin{align}
 P(u,v)&=uv\bigl(p_{21}u^2v+p_{20}u^2+p_{12}uv^2
                      +p_{11}uv+p_{10}u+p_{00}\bigr),\notag\\
 Q(u,v)&=P(v,u).
 \label{shr:h4:eq:PQ}
\end{align}
Its certificate identity is
\begin{equation}
 f(P_u+Q_v)+(n-1)(Pf_u+Qf_v)
       =f(\gamma_n f+1)(u+v).
 \label{shr:h4:eq:polycert}
\end{equation}
This is a polynomial identity in $u,v$, over rational functions of $n$.
It can be verified simply by substituting \eqref{shr:h4:eq:pcoeff} and comparing
coefficients. More concretely, multiply it by $n^2(n+1)^2$ to obtain an
identity in $\Q[n,u,v]$ with no rational-function arithmetic left. The
packaged verifier performs that coefficient comparison with integer and
rational arithmetic, and the residual polynomial is identically zero.
This displayed certificate, rather than the procedure that discovered it,
is the proof object.

For $n\ge2$, multiplication of \eqref{shr:h4:eq:polycert} by $f^{n-2}$ gives
\begin{equation}
 \diver\bigl(f^{n-1}P,f^{n-1}Q\bigr)
   =f^{n-1}(\gamma_n f+1)(u+v).
 \label{shr:h4:eq:divcert}
\end{equation}
For $n=1$, cancel the common polynomial factor $f$ in
\eqref{shr:h4:eq:polycert} first. This gives $P_u+Q_v=(\gamma_1f+1)(u+v)$,
which is precisely \eqref{shr:h4:eq:divcert}. In particular there is no division
by a vanishing boundary function. For all the stated indices, both sides
of \eqref{shr:h4:eq:divcert} are polynomials on the closed triangle.

\subsection{The slanted edge including the exceptional exponent}
The lower slanted edge is $v=1-u$, with outward normal parallel to
$(-1,-1)$. Its normal flux, in the parameter $u$, is
$-f^{n-1}(P+Q)\dd u$. The stronger polynomial identity
\begin{equation}
 P+Q=f\left((p_{21}+p_{12})uv+p_{20}(u+v)-\frac2n\right)
 \label{shr:h4:eq:bottom}
\end{equation}
holds by substitution. Therefore $P(u,1-u)+Q(u,1-u)=0$ identically.
The flux vanishes even at $n=1$, when $f^{n-1}=1$ and vanishing of that
power alone would not settle the boundary. The vertices require no
separate contribution in the planar divergence theorem.

\subsection{Both outer edges and their exact integral}
At the vertical edge $u=1$, the outward flux is
$f(1,v)^{n-1}P(1,v)\dd v$, where
\begin{align}
 f(1,v)&=v^2,\notag\\
 P(1,v)&=p_{12}v^3
 -\frac{(3n+1)(4n+3)}{n(n+1)^2}v^2-\frac vn.
 \label{shr:h4:eq:edge}
\end{align}
At the horizontal edge $v=1$, symmetry gives the same integral.
Consequently their sum is
\begin{align}
 2\int_0^1v^{2n-2}P(1,v)\dd v
 &=2\left(\frac{p_{12}}{2n+2}
 -\frac{(3n+1)(4n+3)}{n(n+1)^2(2n+1)}-\frac1{2n^2}\right)\notag\\
 &=\frac{28n^3+50n^2+29n+5}{n^2(n+1)^2(2n+1)}.
 \label{shr:h4:eq:edgeflux}
\end{align}
Each exponent being integrated is nonnegative for $n\ge1$. All integrals
are ordinary proper integrals of polynomials; no convergence or limiting
boundary prescription is needed.

\begin{proof}[Proof of Theorem~\ref{shr:h4:thm:Z}]
Integrate \eqref{shr:h4:eq:divcert} and apply the divergence theorem on the
closed triangle. Equations~\eqref{shr:h4:eq:bottom} and \eqref{shr:h4:eq:edgeflux}
account for all three edges. Thus
\[
 \gamma_n I_n+I_{n-1}
 =\frac{28n^3+50n^2+29n+5}{n^2(n+1)^2(2n+1)}.
\]
Multiplying by $n^2(n+1)^2/2$ and using $Z_j=j^2I_{j-1}/2$ proves
\eqref{shr:h4:eq:Zrec} for every $n\ge1$.
\end{proof}

Direct period integration gives
\[
 Z_1=\frac13,\qquad Z_2=\frac{13}{45},\qquad
 Z_3=\frac{461}{1680}.
\]
These values provide useful checks on orientation and normalization; the
uniform recurrence follows from the certificate independently of them.

\section{The exact recurrence and printed finite sum}
\label{shr:h4:sec:finite}
\subsection{Returning from the period to the array count}
Write the full-count formula \eqref{shr:h4:eq:affine} as
\[
 b_n=g_nZ_n+r_n,\qquad
 g_n=\frac{4(n-1)}{4n-1},\quad
 r_n=-\frac{32n^3-56n^2+28n-5}{8(2n-1)^2(4n-1)}.
\]
The multinomial normalizer has ratio
\begin{equation}
 m_n:=\frac{M_{n+1}}{M_n}
 =\frac{8(2n+1)(4n+1)(4n+3)}{(n+1)^3}.
 \label{shr:h4:eq:Mratio}
\end{equation}
For $n\ge2$, solve \eqref{shr:h4:eq:Zrec} for $Z_{n+1}$ and substitute in
\[
 A_nm_n(g_{n+1}Z_{n+1}+r_{n+1})+B_n(g_nZ_n+r_n).
\]
The coefficient of $Z_n$ cancels because
\[
 \frac{A_nm_ng_{n+1}(n+1)^2}{3(3n+1)(3n+2)}=B_ng_n.
\]
The remaining rational expression is
\begin{equation}
 A_nm_n\left(g_{n+1}
 \frac{28n^3+50n^2+29n+5}{6(2n+1)(3n+1)(3n+2)}+r_{n+1}\right)
 +B_nr_n=D_n.
 \label{shr:h4:eq:rationalreturn}
\end{equation}
This equality follows by a common-denominator expansion and is one of the
exact polynomial checks supplied with the report. Thus
$A_nm_nb_{n+1}+B_nb_n=D_n$, proving \eqref{shr:h4:eq:inhom} for $n\ge2$.

For $n=1$, $A_1=0$, $B_1=360$, and $D_1M_1=360$, so the identity holds
by $a_1=1$. The width-two count is independently $a_2=1$. All factors of
$A_n$ are nonzero for $n\ge2$. This proves Theorem~\ref{shr:h4:thm:main},
including its precise initial-index statement.

\subsection{Eliminating the hypergeometric forcing term}
Let $J_n=D_nM_n$. Direct cancellation gives
\begin{equation}
 q_n:=\frac{J_{n+1}}{J_n}
 =\frac{8(2n-1)^2(4n+3)(4n+5)(7n^2+14n+6)}
 {(n+1)(n+2)^2(2n+3)(7n^2-1)}.
 \label{shr:h4:eq:qforcing}
\end{equation}
Subtract $q_n$ times \eqref{shr:h4:eq:inhom} at $n$ from that identity at $n+1$.
Multiply the result by
\begin{equation}
 c_n=\frac{(n+1)(n+2)^2(2n+3)(7n^2-1)}n.
 \label{shr:h4:eq:elimmult}
\end{equation}
The denominator is nonzero for the stated range $n\ge1$.

\begin{theorem}[The A181198 operator]
\label{shr:h4:thm:oeis}
For every integer $n\ge1$,
\begin{equation}
 p_2(n)a_{n+2}+p_1(n)a_{n+1}+p_0(n)a_n=0,
 \label{shr:h4:eq:oeis}
\end{equation}
where
\begin{align}
 p_2(n)&=3(n+1)(2n+3)(3n+4)(3n+5)(7n^2-1)(n+2)^3,\notag\\
 p_1(n)&=-8(n+1)(2n+1)(4n+3)(4n+5)\notag\\[-3pt]
 &\hspace{14mm}\cdot
 (364n^5+84n^4-1025n^3-534n^2+157n+54),\label{shr:h4:eq:poperator}\\
 p_0(n)&=-64(2n-1)^2(2n+1)(4n-1)(4n+1)(4n+3)(4n+5)\notag\\[-3pt]
 &\hspace{14mm}\cdot(7n^2+14n+6).\notag
\end{align}
These are the coefficient polynomials recorded for OEIS A181198.
\end{theorem}
\begin{proof}
The three coefficient identities are
\[
 p_2=c_nA_{n+1},\qquad
 p_1=c_n(B_{n+1}-q_nA_n),\qquad p_0=-c_nq_nB_n.
\]
They follow by direct factorization, or by clearing denominators and
comparing polynomial coefficients. The elimination just performed thus
proves \eqref{shr:h4:eq:oeis}. Its leading coefficient is nonzero for every
positive integer $n$, so $a_1=a_2=1$ uniquely specifies its solution.
\end{proof}

The natural empty-array convention would give $a_0=1$. Substitution of
$a_0=a_1=a_2=1$ in the polynomial left side at $n=0$ gives $-2160$, not
zero. The pole in \eqref{shr:h4:eq:elimmult} is therefore substantive. Neither
the proof nor the theorem extends this operator to $n=0$.

\subsection{Solving the first-order recurrence}
For $n\ge2$, put
\begin{equation}
 F_n=\frac{(-1)^{n+1}n(n-1)(4n-2)!}{n!(3n)!},
 \label{shr:h4:eq:F}
\end{equation}
and for $k\ge2$ put
\begin{equation}
 t_k=\frac{3(-1)^k(7k^2-1)(3k)!}
 {(k!)^3(k-1)k(k+1)^2(2k-1)^2(2k+1)^2}.
 \label{shr:h4:eq:t}
\end{equation}
Factorial cancellation proves the two exact relations
\begin{equation}
 \frac{F_{n+1}}{F_n}=-\frac{B_n}{A_n},\qquad
 A_nF_{n+1}t_n=D_nM_n\qquad(n\ge2).
 \label{shr:h4:eq:variation}
\end{equation}
Writing $a_n=F_ns_n$ therefore turns \eqref{shr:h4:eq:inhom} into
$s_{n+1}-s_n=t_n$. Since $F_2=-1$ and $a_2=1$, $s_2=-1$, and
\begin{equation}
 a_n=F_n\left(-1+\sum_{k=2}^{n-1}t_k\right),\qquad n\ge2.
 \label{shr:h4:eq:factorialsum}
\end{equation}
The sum is empty for $n=2$, giving its correct value without a separate
limit convention.

\subsection{Literal correspondence with Conjecture 18}
For a nonnegative integer $j$, write $(x)_j=x(x+1)\cdots(x+j-1)$ for
the rising factorial, with $(x)_0=1$. Generalized binomial coefficients
have the usual polynomial definition.

\begin{theorem}[Kauers and Koutschan Conjecture 18]
\label{shr:h4:thm:sum}
For each integer $n>1$,
\begin{align}
 a_n={}&\frac{(-64)^n(n-1)(-\tfrac12)_{2n}(\tfrac12)_n}{4(3n)!}
 \left(-1+3\sum_{k=2}^{n-1}
 \frac{(-4)^k(7k^2-1)\binom{3k}{2k}\binom{k+1/2}{k}}
 {(k-1)k(k+1)^2(2k-1)^2(2k+1)^3}\right).
 \label{shr:h4:eq:printed}
\end{align}
\end{theorem}
\begin{proof}
The elementary products
\begin{equation}
 (-\tfrac12)_{2n}=-\frac{(4n-2)!}{2^{4n-1}(2n-1)!},\qquad
 (\tfrac12)_n=\frac{(2n)!}{4^nn!}
 \label{shr:h4:eq:pochhammer}
\end{equation}
reduce the prefactor in \eqref{shr:h4:eq:printed} exactly to $F_n$ in
\eqref{shr:h4:eq:F}. Also
\[
 \binom{k+1/2}{k}=\frac{(2k+1)!}{4^k(k!)^2}.
\]
Insert this and $\binom{3k}{2k}=(3k)!/((2k)!k!)$ into three times the
summand in \eqref{shr:h4:eq:printed}. The factors $4^k$ cancel, and one factor
$2k+1$ cancels from its denominator, leaving exactly $t_k$.
Thus \eqref{shr:h4:eq:printed} is \eqref{shr:h4:eq:factorialsum}, already proved.
\end{proof}

The rising factorials, summation range, factor $3$, and prefactor have
all been retained. This establishes the statement actually printed in
the primary paper, rather than only a recurrence empirically agreeing
with a few of its values.

\section{Rational asymptotics with effective fixed-order errors}
\label{shr:h4:sec:asymptotics}
The asymptotic conclusions of this section overlap the earlier
all-height theorem (Part~\ref{shr:part:one}). The recurrence gives a short independent
coefficient algorithm and convenient explicit error bounds. In particular,
the leading constant is a recovered result, not a new priority claim.

\subsection{The normalized recurrence is a contraction}
Divide \eqref{shr:h4:eq:inhom} by $A_nM_{n+1}$. For $n\ge2$,
\begin{equation}
 b_{n+1}=R(n)b_n+S(n),
 \label{shr:h4:eq:brec}
\end{equation}
where
\begin{align}
 R(n)&=-\frac{n(n+1)^2(4n-1)}
 {3(n-1)(3n+1)(3n+2)(4n+3)},\label{shr:h4:eq:R}\\
 S(n)&=\frac{7n^2-1}
 {4(n-1)(3n+1)(3n+2)(2n-1)^2(2n+1)^2(4n+3)}.
 \label{shr:h4:eq:S}
\end{align}
Both rational functions are regular for $n\ge2$, and
\begin{equation}
 R(n)\longrightarrow-\frac1{27},\qquad
 S(n)=\frac7{2304}n^{-6}+O(n^{-7}).
 \label{shr:h4:eq:RSlead}
\end{equation}
A convenient global bound is
\begin{equation}
 |R(n)|\le\frac16\qquad(n\ge2).
 \label{shr:h4:eq:contraction}
\end{equation}
Indeed, $n+1\le3n/2$, $4n-1\le4n$, $n-1\ge n/2$, and each of
$3n+1,3n+2,4n+3$ is at least its leading monomial. The numerator is at
most $9n^4$ and the denominator is at least $54n^4$.

Let $t=1/n$, $r(t)=R(1/t)$ and $s(t)=t^{-6}S(1/t)$. These functions are
rational and analytic at zero. Seek
\[
 b_n\sim n^{-6}C(1/n),\qquad C(t)=\sum_{j\ge0}c_jt^j.
\]
The formal equation is
\begin{equation}
 (1+t)^{-6}C\left(\frac{t}{1+t}\right)=r(t)C(t)+s(t).
 \label{shr:h4:eq:formal}
\end{equation}
In the coefficient of $t^j$, the new unknown $c_j$ occurs with coefficient
$1-r(0)=28/27$. Therefore this equation determines one and only one
formal series, with every $c_j\in\Q$. Its first coefficient is
\[
 c_0=\frac{7/2304}{1+1/27}=\frac3{1024}.
\]
Only finitely many rational operations are required for any chosen order.

\begin{theorem}[Every fixed order with a computable error]
\label{shr:h4:thm:bexp}
For every fixed integer $K\ge0$, there are explicitly computable rational
constants $E_K>0$ and integer $N_K\ge2$ such that
\begin{equation}
 \left|b_n-n^{-6}\sum_{j=0}^Kc_jn^{-j}\right|
 \le E_Kn^{-K-7}\qquad(n\ge N_K).
 \label{shr:h4:eq:b-error}
\end{equation}
The coefficients are exactly those of \eqref{shr:h4:eq:formal}. In particular,
\begin{align}
 b_n={}&\frac3{1024}n^{-6}+\frac{15}{1024}n^{-7}
 +\frac{45}{1024}n^{-8}+\frac{99}{1024}n^{-9}\notag\\
 &+\frac{2367}{16384}n^{-10}+\frac{279}{16384}n^{-11}
 -\frac{56751}{65536}n^{-12}+O(n^{-13}).
 \label{shr:h4:eq:bcoefficients}
\end{align}
\end{theorem}

\begin{proof}
Fix $K$, write $p=K+7$ and
$q_K(n)=n^{-6}\sum_{j=0}^Kc_jn^{-j}$. Its exact residual is
\[
 d_K(n)=q_K(n+1)-R(n)q_K(n)-S(n).
\]
The formal construction says $d_K(n)=O(n^{-p})$. It also supplies an
effective rational bound as follows. Express
\begin{equation}
 t^{-p}d_K(1/t)=\frac{P(t)}{H(t)},\qquad H(0)>0,
 \label{shr:h4:eq:PHresidual}
\end{equation}
with rational polynomials $P,H$. Remove any common powers of $t$ before
this expression; the rational function is analytic at zero. Write
$P(t)=\sum P_it^i$, $H(t)=\sum H_it^i$.
Choose an integer $N\ge2$ large enough that
\begin{equation}
 h_N:=H_0-\sum_{i\ge1}|H_i|N^{-i}>0,\qquad
 \rho_N:=\left(\frac N{N+1}\right)^p>\frac16.
 \label{shr:h4:eq:Nchoices}
\end{equation}
Both conditions eventually hold. Then an explicit bound is
\begin{equation}
 C_D=\frac{\sum_i|P_i|N^{-i}}{h_N},\qquad
 |d_K(n)|\le C_Dn^{-p}\quad(n\ge N).
 \label{shr:h4:eq:CD}
\end{equation}
The initial $b_N$ is an exact rational number, obtained from the proved
recurrence or directly from the period formula. Define
\begin{equation}
 E_K=\max\left\{N^p|b_N-q_K(N)|,
                    \frac{C_D}{\rho_N-1/6}\right\}.
 \label{shr:h4:eq:EK}
\end{equation}
If this maximum is zero one may replace it by any positive rational.
Set $e_n=b_n-q_K(n)$. From \eqref{shr:h4:eq:brec},
$e_{n+1}=R(n)e_n-d_K(n)$. If $|e_n|\le E_Kn^{-p}$, then
\[
 |e_{n+1}|\le(E_K/6+C_D)n^{-p}
 \le E_K\rho_Nn^{-p}\le E_K(n+1)^{-p}.
\]
The last inequality uses $(n/(n+1))^p\ge\rho_N$.
The initial inequality follows from \eqref{shr:h4:eq:EK}, proving the bound by
induction. This is a proof for the actual counting solution, not just a
formal-series calculation.
\end{proof}

The constants in \eqref{shr:h4:eq:PHresidual}--\eqref{shr:h4:eq:EK} can be obtained by
finite polynomial arithmetic and finitely many exact recurrence steps.
The bound deliberately favors transparency over numerical sharpness.
It does not impose an upper limit on the fixed order $K$.

\subsection{Restoring the multinomial factor}
The logarithmic Stirling expansion gives
\begin{align}
 \log M_n\sim{}&n\log256-\frac32\log n
             -\log(\sqrt2\,\pi^{3/2})\notag\\
 &+\sum_{j\ge1}\frac{B_{2j}}{2j(2j-1)}
                   \bigl(4^{-(2j-1)}-4\bigr)n^{-(2j-1)},
 \label{shr:h4:eq:logM}
\end{align}
interpreted at each fixed truncation, with $B_{2j}$ the Bernoulli numbers.
In particular, the first logarithmic correction is $-5/(16n)$.
The constant follows directly from
$\sqrt{8\pi n}/(2\pi n)^2=1/(\sqrt2\,\pi^{3/2}n^{3/2})$.

\begin{corollary}[The count expansion]
\label{shr:h4:cor:aexp}
For every fixed $K\ge0$,
\begin{equation}
 a_n=\frac{3}{1024\sqrt2\,\pi^{3/2}}
       256^n n^{-15/2}
       \left(1+\sum_{j=1}^K d_jn^{-j}+O_K(n^{-K-1})\right),
 \label{shr:h4:eq:aexp}
\end{equation}
with rational coefficients. The first six are
\begin{align}
 d_1&=\frac{75}{16},&d_2&=\frac{6905}{512},&
 d_3&=\frac{233985}{8192},\notag\\
 d_4&=\frac{20845307}{524288},&
 d_5&=-\frac{66392235}{8388608},&
 d_6&=-\frac{79147528275}{268435456}.
 \label{shr:h4:eq:dvalues}
\end{align}
\end{corollary}
\begin{proof}
Multiply the expansion of Theorem~\ref{shr:h4:thm:bexp} by the fixed-order
exponential of \eqref{shr:h4:eq:logM}. Its leading coefficient is
$\frac{c_0}{\sqrt2\,\pi^{3/2}}$. Every further coefficient is obtained by finite
addition and multiplication of rational numbers. Stirling's remainder
and \eqref{shr:h4:eq:b-error} justify multiplication at each fixed order.
\end{proof}

The leading term in \eqref{shr:h4:eq:aexp} agrees with the general fixed-height
constant of the prior report at $m=4$. The factor $n^{-6}$ in $b_n$
accounts for the difference between the multinomial exponent $3/2$ and
the count exponent $15/2$.

\subsection{An explicit logarithmic remainder for inversion}
The following construction makes the error term usable without concealing
an unknown constant. Define
\begin{equation}
 C_*:=\frac{3}{1024\sqrt2\,\pi^{3/2}},\qquad
 \lambda=\log256,\qquad \beta_*:=\frac{15}2.
 \label{shr:h4:eq:Cstar}
\end{equation}
For each fixed $K$, there are rational $\ell_1,\ldots,\ell_K$ and
computable $C_K>0,N_K$ such that
\begin{equation}
 \left|\log a_n-\left(\lambda n-\beta_*\log n+\log C_*
                   +\sum_{j=1}^K\ell_jn^{-j}\right)\right|
 \le C_Kn^{-K-1}\quad(n\ge N_K).
 \label{shr:h4:eq:logerror}
\end{equation}
The $\ell_j$ are the coefficients of
$\log(1+\sum_{j\ge1}d_jt^j)$.

Here is one fully explicit construction of $C_K$. Retain an $E=E_K$ and
initial lower bound from Theorem~\ref{shr:h4:thm:bexp}. Put
\[
 u(t)=\sum_{j=1}^K\frac{c_j}{c_0}t^j,
 \quad U=\sum_{j=1}^K\left|\frac{c_j}{c_0}\right|N^{-(j-1)}.
\]
Increase the integer $N$ until $U/N\le1/4$ and
$(E/c_0)N^{-K-1}\le1/4$. Increasing $N$ preserves the already-proved
error bound. For $n\ge N$, both $1+u(1/n)$ and $b_nn^6/c_0$ are at
least $1/2$, so the mean value theorem for $\log$ gives
\[
 \left|\log\frac{b_nn^6}{c_0}-\log(1+u(1/n))\right|
 \le\frac{2E}{c_0}n^{-K-1}.
\]
Let
\[
 L_K(t)=\sum_{h=1}^K\frac{(-1)^{h+1}}h u(t)^h,
 \quad
 V=\sum_{j>K}|[t^j]L_K(t)|N^{-(j-K-1)}.
\]
The omitted tail of the scalar logarithm series is bounded by
$4U^{K+1}n^{-K-1}/(3(K+1))$. The part of the finite polynomial $L_K$
above degree $K$ is bounded by $Vn^{-K-1}$. For $K=0$, take $u=L_K=0$
and $U=V=0$. Thus the $b_n$ part contributes at most
\begin{equation}
 C_b=\frac{2E}{c_0}+V+\frac{4U^{K+1}}{3(K+1)}.
 \label{shr:h4:eq:Cb}
\end{equation}

Let $L=\lfloor(K+1)/2\rfloor$ and $q=2L+1$. The classical positive-real
Stirling remainder after $L$ Bernoulli terms has absolute value at most
$|B_{2L+2}|/((2L+2)(2L+1)x^{2L+1})$; this follows, for example, from the
integral remainder in the Euler--Maclaurin formula for $\log\Gamma$
\cite{DLMFgamma}. Apply it separately at $x=4n$ and at $x=n$. Since $q\ge K+1$,
the multinomial logarithm contributes at most
\begin{equation}
 C_M=\frac{|B_{2L+2}|}{(2L+2)(2L+1)}(4^{-q}+4)N^{K+1-q}
 \label{shr:h4:eq:CM}
\end{equation}
times $n^{-K-1}$. Take $C_K=C_b+C_M$ in \eqref{shr:h4:eq:logerror}.
This provides a finite, effective remainder certificate for every fixed
order, using only the displayed polynomial and rational bounds plus the
standard Stirling inequality.

\subsection{Why this does not specify a full transseries}
The homogeneous equation associated with \eqref{shr:h4:eq:brec} has solution
proportional to $\prod R(j)$, whose successive absolute ratio tends to
$1/27$. The first-order equation therefore explains stability of the
algebraic expansion. It does not make a particular all-orders formal
series into a canonical summed particular solution. Changing an exact
particular solution by a homogeneous solution changes its exponentially
small component without changing any inverse-power coefficient.
Neither the finite-order bounds nor the formal branches alone select
Stokes data or a canonical exponentially small transseries decomposition.
No such conclusion is used here.

\section{Finite Lambert W models and threshold inverses}
\label{shr:h4:sec:inverse}
The sequence is defined at integers. We therefore distinguish a smooth
finite model, which can be inverted as a real function, from the exact
integer threshold of the count. This distinction prevents an unjustified
rounding claim near an integer and avoids differentiating an unspecified
discrete remainder.

\subsection{A monotone finite model}
Fix $K\ge0$ and use the coefficients and error bound from
\eqref{shr:h4:eq:logerror}. Define, for real $x>0$,
\begin{equation}
 \Psi_K(x)=\lambda x-\beta_*\log x+\log C_*
                  +\sum_{j=1}^K\ell_jx^{-j}.
 \label{shr:h4:eq:Psi}
\end{equation}
Choose $X_K$ large enough that
\begin{equation}
 \frac{\beta_*}{X_K}
 +\sum_{j=1}^Kj|\ell_j|X_K^{-j-1}\le\frac\lambda2.
 \label{shr:h4:eq:XK}
\end{equation}
Then $\Psi_K'(x)\ge\lambda/2$ for $x\ge X_K$. Hence for every
$y\ge\exp\Psi_K(X_K)$ there is a unique $\xi_K(y)\ge X_K$ with
\begin{equation}
 \Psi_K(\xi_K(y))=\log y.
 \label{shr:h4:eq:xi}
\end{equation}
This defines an actual finite smooth model and a genuine inverse, rather
than an unspecified interpolation of the sequence.

For $K=0$, inversion is explicit:
\begin{equation}
 w(y):=\xi_0(y)
 =-\frac{\beta_*}{\lambda}
 W_{-1}\left(-\frac{\lambda}{\beta_*}
                   \left(\frac{C_*}{y}\right)^{1/\beta_*}\right).
 \label{shr:h4:eq:Lambert}
\end{equation}
For sufficiently large $y$ the argument is in $(-1/e,0)$. The branch
$W_{-1}$ is selected because it tends to $-\infty$ as the argument tends
to $0$ from below, giving the large positive solution. Indeed, taking
$\beta_*$th roots in $y=C_*e^{\lambda w}w^{-\beta_*}$ gives
$w e^{-\lambda w/\beta_*}=(C_*/y)^{1/\beta_*}$, which proves
\eqref{shr:h4:eq:Lambert}. The principal branch would give the small solution
on the decreasing part of this same leading model.

\subsection{Arbitrary fixed corrections to the Lambert solution}
The finite model admits an explicitly recursive inverse expansion. Put
$t=1/w$ and seek
\begin{equation}
 \widehat\xi_K(y)=w+\Delta_K(1/w),\qquad
 \Delta_K(t)=\sum_{j=1}^Ks_jt^j.
 \label{shr:h4:eq:xihat}
\end{equation}
The coefficients are determined by setting the coefficients of
$t,t^2,\ldots,t^K$ to zero in
\begin{equation}
 G_K(t)=\lambda\Delta_K(t)
 -\beta_*\log(1+t\Delta_K(t))
 +\sum_{j=1}^K\ell_jt^j(1+t\Delta_K(t))^{-j}.
 \label{shr:h4:eq:Ginverse}
\end{equation}
At order $j$, the new $s_j$ occurs only in $\lambda\Delta_K$,
so the recursion is uniquely solvable and
$s_j\in\Q[\lambda^{-1}]$. For example,
\begin{equation}
 \ell_1=\frac{75}{16},\qquad \ell_2=\frac52,
 \quad
 s_1=-\frac{75}{16\lambda},\qquad
 s_2=-\frac5{2\lambda}-\frac{1125}{32\lambda^2}.
 \label{shr:h4:eq:firstinverse}
\end{equation}
Thus corrections begin at $w^{-1}$; the leading Lambert model already
incorporates the logarithmic displacement.

\begin{proposition}[A finite inverse with an effective error]
\label{shr:h4:prop:inverse}
For each fixed $K$, constants $H_K,W_K$ can be computed from the finite
expression \eqref{shr:h4:eq:Ginverse} such that, whenever $w(y)\ge W_K$ and
$y\ge\exp\Psi_K(X_K)$,
\begin{equation}
 |\xi_K(y)-\widehat\xi_K(y)|
 \le\frac{2H_K}{\lambda}\,w(y)^{-K-1}.
 \label{shr:h4:eq:inv-error}
\end{equation}
\end{proposition}
\begin{proof}
Choose $W_K\ge2X_K$ large enough that
$|t\Delta_K(t)|\le1/2$ for $0\le t\le1/W_K$. This can be ensured by
the finite coefficient bound
$\sum_j|s_j|W_K^{-j-1}\le1/2$. Then
$\widehat\xi_K=w(1+t\Delta_K(t))\ge w/2\ge X_K$.
Since $G_K$ has its coefficients through degree $K$ equal to zero,
Taylor's theorem gives
\[
 |G_K(t)|\le H_Kt^{K+1}\quad(0\le t\le1/W_K).
\]
One explicit way to obtain the constant is to differentiate $K+1$ times.
Writing $z(t)=1+t\Delta_K(t)$, the derivative has the form
\[
 G_K^{(K+1)}(t)=\frac{P_K(t)}{z(t)^{2K+1}}
\]
for a computable real polynomial $P_K$ with coefficients in
$\Q[\lambda,\lambda^{-1}]$. For $K=0$ take $P_0=0$.
The claimed form follows by repeated differentiation of the logarithm
and the finitely many negative integer powers, then bringing them to a
common denominator. Thus one may take
\begin{equation}
 H_K=\frac{2^{2K+1}}{(K+1)!}
       \sum_i|[t^i]P_K(t)|W_K^{-i}.
 \label{shr:h4:eq:HK}
\end{equation}
The identity $\Psi_0(w)=\log y$ shows exactly that
$\Psi_K(\widehat\xi_K)-\log y=G_K(1/w)$.
Both $\xi_K$ and $\widehat\xi_K$ lie on the interval where the derivative
is at least $\lambda/2$. The mean value theorem now gives
\eqref{shr:h4:eq:inv-error}.
\end{proof}

All objects in this proof are finite functions. No convergence of the
infinite inverse series is asserted or needed. Interval evaluations of
$\lambda=\log256$ and of the finite polynomials suffice if numerical
certified constants rather than the displayed exact expressions are
wanted.

\subsection{The integer threshold and its rounding envelope}
Define
\begin{equation}
 N(y)=\min\{n\ge1:a_n\ge y\}.
 \label{shr:h4:eq:threshold}
\end{equation}
This is defined for all positive $y$ because $a_n\to\infty$.
The sequence is eventually strictly increasing. One direct effective
proof uses \eqref{shr:h4:eq:logerror} at $K=0$: for large enough $n$,
\begin{align*}
 \log a_{n+1}-\log a_n
 &\ge\lambda-\beta_*\log(1+1/n)
       -C_0\bigl(n^{-1}+(n+1)^{-1}\bigr)\\
 &\ge\lambda-\frac{\beta_*+2C_0}{n}.
\end{align*}
Taking $n\ge2(\beta_*+2C_0)/\lambda$, as well as the onset required by
\eqref{shr:h4:eq:logerror}, makes this at least $\lambda/2$. No derivative of
an error sequence has been taken.

Fix $K$, let $p=K+1$, and choose an integer $N_*$ beyond this monotonicity
onset, beyond the onset of \eqref{shr:h4:eq:logerror}, and beyond $X_K$.
Choose $y$ large enough that
\begin{equation}
 \begin{gathered}
 y>\max_{1\le n<N_*}a_n,\qquad \xi:=\xi_K(y)\ge\max(2N_*+2,4),\\
 \delta:=\frac{2^{p+1}C_K}{\lambda}\xi^{-p}\le1.
 \end{gathered}
 \label{shr:h4:eq:thresholdonset}
\end{equation}
Every condition is verifiable from the finite bounds and initial values.

\begin{theorem}[Qualified threshold envelope]
\label{shr:h4:thm:threshold}
Under \eqref{shr:h4:eq:thresholdonset},
\begin{equation}
 \left\lceil\xi_K(y)-\delta\right\rceil
 \le N(y)\le
 \left\lceil\xi_K(y)+\delta\right\rceil,
 \qquad
 \delta=\frac{2^{K+2}C_K}{\lambda}\xi_K(y)^{-K-1}.
 \label{shr:h4:eq:envelope}
\end{equation}
\end{theorem}
\begin{proof}
Let $m=\lceil\xi+\delta\rceil$. The derivative bound gives
$\Psi_K(m)-\log y\ge\lambda\delta/2$.
Since $m\ge\xi\ge N_*$, its logarithmic error is at most
$C_Km^{-p}\le C_K\xi^{-p}\le\lambda\delta/2$.
Therefore $a_m\ge y$.

Let $j=\lceil\xi-\delta\rceil-1$. Then $j<\xi-\delta$ and
$j\ge\xi-2\ge\xi/2\ge N_*$. Hence
$\Psi_K(j)-\log y<-\lambda\delta/2$, whereas its logarithmic error
is at most $C_Kj^{-p}\le2^pC_K\xi^{-p}=\lambda\delta/2$.
Thus $a_j<y$. Eventual monotonicity and the finite-prefix condition in
\eqref{shr:h4:eq:thresholdonset} rule out an earlier crossing. These two
integer comparisons prove \eqref{shr:h4:eq:envelope}.
\end{proof}

For sufficiently large $y$, the envelope has width less than one in its
real argument. It gives a unique integer whenever its two ceilings
coincide. If an endpoint approaches an integer, the theorem may leave two
adjacent possibilities; checking their exact recurrence values resolves
the threshold. An unconditional instruction to round a real asymptotic
inverse would lose this necessary qualification.

The explicit Lambert correction \eqref{shr:h4:eq:xihat} can replace $\xi_K$.
For example, write $E'=2H_K/\lambda$ and
$D=2^{p+1}C_K/\lambda$. Enlarge the onset until
$|\xi_K-\widehat\xi_K|\le E'w^{-p}$ and $\xi_K\ge w/2$.
Then \eqref{shr:h4:eq:envelope} implies the entirely explicit envelope
\begin{equation}
 \left\lceil\widehat\xi_K-(E'+2^pD)w^{-p}\right\rceil
 \le N(y)\le
 \left\lceil\widehat\xi_K+(E'+2^pD)w^{-p}\right\rceil.
 \label{shr:h4:eq:lambertenvelope}
\end{equation}
Here $w$ is the Lambert expression \eqref{shr:h4:eq:Lambert}, $\widehat\xi_K$
is a finite correction polynomial, and every error constant has already
been specified. The increasing onset is part of the conclusion; it is
not a claim of sharp small-$n$ rounding accuracy.

\section{Exact verification and reproducibility}
\label{shr:h4:sec:verification}
\subsection{Proof certificates and finite checks have different roles}
The proof's universal algebraic objects are \eqref{shr:h4:eq:polycert}, its
boundary identities, the rational substitution
\eqref{shr:h4:eq:rationalreturn}, and the three elimination coefficients.
These are identities in polynomial or rational-function rings. Their
verification does not depend on a finite numerical sample of $n$.
The order-polytope, Pfaffian, integration-by-parts, contraction, and
inverse arguments are the mathematical proofs in the preceding sections.

The accompanying core code uses the Python standard library only. It
implements sparse polynomials over exact rational numbers and rational
functions reduced by polynomial arithmetic. All proof-critical tests
raise explicit exceptions on failure, rather than relying on statements
removed by Python's optimized mode.

The principal files are:
\begin{itemize}
\item \texttt{code/exact\_algebra.py}: rational sparse polynomial and
      rational-function operations;
\item \texttt{code/formal\_series.py}: fixed-order formal arithmetic,
      shifts, Bernoulli numbers, and the Stirling product;
\item \texttt{code/verify\_report231.py}: certificate identities,
      independent exact counting checks, and series checks;
\item \texttt{code/receipt.json}: deterministic output of the verifier;
\item \texttt{build.py}: frozen-inventory verification and isolated
      PDF/archive replay.
\end{itemize}
The package includes the editable \LaTeX\ source and this PDF. It does
not redistribute third-party source papers.

\subsection{Three independent ways to evaluate the finite count}
A legal prefix of an array is described by the number of assigned entries
in each row, $(\lambda_1,\ldots,\lambda_4)$. Rows increase, so each row
prefix is an initial interval. Literal cell comparisons say exactly that
\begin{equation}
 n\ge\lambda_1\ge\lambda_2\ge\lambda_3\ge\lambda_4\ge0,
 \qquad
 \lambda_i=\lambda_{i+1}\ \Longrightarrow\
 \lambda_i\in\{0,n\}.
 \label{shr:h4:eq:prefix}
\end{equation}
To see the strictness clause, an occupied lower-row cell in column $j<n$
requires its upper-row neighbor in column $j+1$ by the antidiagonal
comparison. If the lower row has length $n$, the upper row must also
have length $n$. Empty adjacent rows impose no strictness. Conversely
these conditions include the row, column, and antidiagonal predecessors
of each occupied cell; the diagonal condition is redundant.
Starting from $(0,0,0,0)$, increment one legal coordinate and accumulate
integer path counts until $(n,n,n,n)$. This dynamic program uses no
candidate recurrence or triangle integral.

A second evaluator integrates \eqref{shr:h4:eq:Ztriangle} directly in positive
barycentric coordinates on the triangle. Put
\[
 r=1-u,\qquad s=1-v,\qquad t=1-r-s=u+v-1.
\]
Then $u=s+t$, $v=r+t$, and
\[
 Z_n=n^2\int_{r,s>0,\ r+s<1}(s+t)^n(r+t)^{n-1}t^{n-1}\dd r\dd s.
\]
Expanding the two positive binomial powers and using
\[
 \int_{r,s>0,\ r+s<1}r^a s^b(1-r-s)^c\dd r\dd s
 =\frac{a!b!c!}{(a+b+c+2)!}
\]
gives a finite positive rational sum. It is independent of both the
period recurrence and the prefix-state dynamic program. The displayed
Dirichlet integral follows by two beta integrals, or by integrating
monomials on a simplex.

A third check uses the exact polynomial
\begin{equation}
 h(t)=\sum_{j=0}^{n-1}
 \frac{(-1)^jt^{n+j}}{(n-1-j)!(n+j)!},
 \label{shr:h4:eq:hpoly}
\end{equation}
obtained by expanding $(1-r)^{n-1}$ in the defining integral and applying
a beta integral. It yields
\begin{equation}
 Z_n=(n!)^2\sum_{j=0}^{n-1}
 \frac{(-1)^j}{(n-1-j)!(n+j)!(2n+j+1)}.
 \label{shr:h4:eq:Zfinite}
\end{equation}
The opposite signs make this a less stable floating-point formula, but
exact rational arithmetic has no cancellation error. It is convenient
for independent symbolic checking of the trace reduction.

The package verifies these count evaluations, the printed finite sum,
and both recurrences for every width $1\le n\le20$ where their indices
apply. The first ten values are
\begin{center}
\begin{tabular}{r r@{\qquad}r r}
\toprule
$n$&$a_n$&$n$&$a_n$\\\midrule
1&1&6&352184\\
2&1&7&25097600\\
3&8&8&2152061145\\
4&169&9&212012802584\\
5&6392&10&23263015359672\\\bottomrule
\end{tabular}
\end{center}
These tests help detect transcription errors. They do not establish a
uniform recurrence by induction on the number of successful samples.

\subsection{Replaying the public package}
From the package directory, run
\begin{quote}\ttfamily
python -B code/verify\_report231.py -{}-self-test
\end{quote}
for the standard-library certificate checks. It prints deterministic JSON;
\texttt{-{}-receipt} accepts a fresh destination and refuses an existing
file. To verify the frozen package inventory without rebuilding, use
\begin{quote}\ttfamily
python -B build.py -{}-verify-only
\end{quote}
To rebuild, provide a new directory outside the package:
\begin{quote}\ttfamily
python -B build.py -{}-output /path/to/new-output
\end{quote}
The build repeats exact checks, creates a private TeX format/cache,
compiles the source, compares the PDF with the frozen copy, and writes a
deterministic archive. A TeX Live installation with the packages named in
\texttt{article.tex} is needed for PDF reproduction; the exact checker
itself needs only Python. The build never needs network access.
The same commands may be run with \texttt{-O}; every Python child process
explicitly receives \texttt{-B}, and optimized-mode guards remain active.

The manifest covers every packaged file except itself and rejects missing,
changed, or extra files and unexpected directories. Output paths must be
new, outside the source package, and free of symlink ancestors. Source
inventory is compared before and after the build. The shipped PDF and
archive are pinned by their public digest records. Byte-for-byte PDF
reproduction is expected under the recorded TeX toolchain; a different
font or TeX version can legitimately require a new rendering review.

\begin{writenote}[Added 5 October 2026, batch~101: the shipped layout]
``The package'' in this section is the delivered archive. In ProveIt its
programs and receipt carry the prefix \texttt{03-height-four-}:
\begin{center}\footnotesize
\begin{tabular}{@{}ll@{}}
\toprule
Delivered as & Shipped as\\
\midrule
\texttt{code/}\emph{name}\texttt{.py} (three modules) & \texttt{code/\allowbreak 03-\allowbreak height-\allowbreak four-}\emph{name}\texttt{.py}\\
\texttt{build.py} & \texttt{code/\allowbreak 03-\allowbreak height-\allowbreak four-\allowbreak build.py}\\
\texttt{code/\allowbreak receipt.json} & \texttt{data/\allowbreak 03-\allowbreak height-\allowbreak four-\allowbreak receipt.json}\\
\texttt{code/README.md} & \texttt{03-\allowbreak height-\allowbreak four-\allowbreak code-\allowbreak README.md}\\
\texttt{SOURCES.md} & \texttt{03-\allowbreak height-\allowbreak four-\allowbreak SOURCES.md}\\
\bottomrule
\end{tabular}
\end{center}
The editable source, this PDF and \texttt{MANIFEST.sha256} named above are
not shipped. \texttt{verify\_\allowbreak report231.py} imports \texttt{exact\_algebra}
and \texttt{formal\_series} under their delivered names, and
\texttt{build.py} needs the manifest and the frozen PDF, so neither runs in
place; the verifier runs on a copy with the delivered names restored, and
the README of the report gives the commands. At the write (5~October 2026,
Python~3.14.4 on Windows, on such a copy) \texttt{python -B code/\allowbreak verify\_\allowbreak report231.py
-{}-self-test} printed
\texttt{data/\allowbreak 03-\allowbreak height-\allowbreak four-\allowbreak receipt.json} byte for byte after removal of
carriage returns, with and without \texttt{-O}, in about three seconds.
\texttt{build.py} was not run. The verifier hard-codes the three
coefficient polynomials as displayed in the A181198 entry and checks the
elimination against them; as OEIS content they are published by The OEIS
Foundation Inc.\ under the Creative Commons Attribution-ShareAlike~4.0
licence.
\end{writenote}


\subsection{What remains outside this report}
The height-five Conjecture~19 is not implied by the four-variable
Pfaffian calculation. Higher Pfaffian sizes contain further trace
invariants, so an analogous low-dimensional reduction would require new
work. Likewise, the stable first-order normalized recurrence controls
every fixed algebraic order but does not by itself canonically resolve
exponentially small sectors. Neither extension is claimed here.

The result established in full is the height-four count, its specified
operator, and the primary paper's finite sum. The proof and exact
certificate are available for independent checking without trusting a
recurrence guess, a finite match, or a symbolic algebra package.

\subsection{[write] The order-two operator is minimal}
\label{shr:h4:sec:minimal}

Added 5 October 2026, batch~101. This subsection is the write's, not the
manuscript's, which makes no minimality claim. It shows that A181198 obeys
no first-order recurrence, so the operator of
Theorem~\ref{shr:h4:thm:oeis} has the smallest possible order, and that
every second-order recurrence of the sequence is a polynomial multiple of
that operator. The proof uses only the period recurrence
\eqref{shr:h4:eq:Zrec}, the inhomogeneous identity \eqref{shr:h4:eq:inhom}
and the affine formula \eqref{shr:h4:eq:affine}. In this subsection a
\emph{recurrence of order~$r$} for a sequence $(s_n)$ means polynomials
$q_0,\ldots,q_r\in\Q[n]$ with $q_r\ne0$ and
$\sum_{j=0}^rq_j(n)s_{n+j}=0$ for all sufficiently large integers $n$.

\begin{lemma}[\textbf{[write]} The period is not eventually rational]
\label{shr:h4:lem:Zirr}
There is no rational function $z\in\mathbb C(n)$ with $Z_n=z(n)$ for all
sufficiently large integers $n$.
\end{lemma}

\begin{proof}
Put $\mathsf a(n)=3(3n+1)(3n+2)$, $\mathsf b(n)=(n+1)^2$ and
$\phi(n)=(28n^3+50n^2+29n+5)/(2(2n+1))$. The numerator of $\phi$ equals
$-1/2$ at $n=-1/2$, so $\phi$ has a simple pole at $-1/2$ and no other
pole. Suppose $Z_n=z(n)$ for $n\ge n_1$. By \eqref{shr:h4:eq:Zrec}, the
rational function $\mathsf a(n)z(n+1)+\mathsf b(n)z(n)-\phi(n)$ vanishes at
every integer $n\ge n_1$, hence identically:
\begin{equation}
 \mathsf a(n)z(n+1)+\mathsf b(n)z(n)=\phi(n)\qquad\text{in }\mathbb C(n).
 \label{shr:h4:eq:zrat}
\end{equation}
Suppose $z$ has a pole, and fix a coset $O=\theta+\Z$ of $\mathbb C/\Z$
that contains one. The poles of $z$ in $O$ form a finite nonempty set; let
$\beta$ be its largest and $\gamma$ its smallest element (they differ by a
nonnegative integer).

(i) Near $n=\beta$, the term $\mathsf a(n)z(n+1)$ is regular, because
$\beta+1\in O$ is not a pole. By \eqref{shr:h4:eq:zrat},
$\mathsf b(n)z(n)=\phi(n)-\mathsf a(n)z(n+1)$ is regular at $\beta$ unless
$\phi$ has a pole there. Since $z$ has a pole at $\beta$, either
$\mathsf b(\beta)=0$, that is $\beta=-1$, or $\beta=-1/2$.

(ii) Near $n=\gamma-1$, the term $\mathsf b(n)z(n)$ is regular, because
$\gamma-1\in O$ is not a pole. By \eqref{shr:h4:eq:zrat},
$\mathsf a(n)z(n+1)=\phi(n)-\mathsf b(n)z(n)$ is regular at $\gamma-1$
unless $\gamma-1=-1/2$. Since $z(n+1)$ has a pole at $n=\gamma-1$, either
$\mathsf a(\gamma-1)=0$, that is $\gamma-1\in\{-1/3,-2/3\}$, or
$\gamma-1=-1/2$. Thus $\gamma\in\{1/3,2/3,1/2\}$.

By (i), $O=\Z$ or $O=\tfrac12+\Z$. If $O=\Z$, (ii) is impossible. If
$O=\tfrac12+\Z$, then (i) gives $\beta=-1/2$ and (ii) gives $\gamma=1/2$,
contradicting $\gamma\le\beta$. Hence $z$ has no pole and is a polynomial
(possibly zero). Then the left side of \eqref{shr:h4:eq:zrat} is a
polynomial, while $\phi$ has a pole at $-1/2$. This contradiction proves
the lemma.
\end{proof}

\begin{proposition}[\textbf{[write]} Minimality at height four]
\label{shr:h4:prop:minimal}
The sequence $a_n=T_4(n)$ (A181198) has no recurrence of order zero or
one; in particular it is not hypergeometric. If $q_0,q_1,q_2$ give a
recurrence of order two, then $(q_0,q_1,q_2)=\rho\,(p_0,p_1,p_2)$ for a
polynomial $\rho\in\Q[n]$, with $p_0,p_1,p_2$ as in \eqref{shr:h4:eq:poperator}.
Hence the operator of Theorem~\ref{shr:h4:thm:oeis} has minimal order, and
degree nine is the smallest degree of any recurrence of order two.
\end{proposition}

\begin{proof}
Every $a_n$ with $n\ge1$ is a positive integer (the column-by-column
interlaced filling is one array).

\emph{Order zero.} If $q_0(n)a_n=0$ for all large $n$, then $q_0$ has
infinitely many zeros, so $q_0=0$; this is excluded.

\emph{Order one.} Suppose $q_1(n)a_{n+1}+q_0(n)a_n=0$ for $n\ge n_0$ with
$q_1\ne0$. Enlarging $n_0$, we may assume $n_0\ge2$ and $q_1(n)\ne0$ for
$n\ge n_0$, so $a_{n+1}=\varrho(n)a_n$ there, with
$\varrho=-q_0/q_1\in\Q(n)$. Substituting into \eqref{shr:h4:eq:inhom}
gives
\[
 a_n\bigl(A_n\varrho(n)+B_n\bigr)=D_nM_n\qquad(n\ge n_0).
\]
For $n\ge1$ both $D_n$ and $M_n$ are positive, so
$A_n\varrho(n)+B_n\ne0$ for every $n\ge n_0$; the rational function
$A_n\varrho+B_n$ is therefore not zero, and
\[
 b_n=\frac{a_n}{M_n}=\frac{D_n}{A_n\varrho(n)+B_n}\qquad(n\ge n_0)
\]
is the value of a rational function. By \eqref{shr:h4:eq:affine},
$Z_n=(b_n-r_n)/g_n$ for $n\ge2$, where $g_n=4(n-1)/(4n-1)\ne0$ and $r_n$
are the rational functions of Section~30.1. So $Z_n$ coincides with a
rational function for all $n\ge n_0$, contradicting
Lemma~\ref{shr:h4:lem:Zirr}.

\emph{Order two.} Write $(La)(n)=p_2(n)a_{n+2}+p_1(n)a_{n+1}+p_0(n)a_n$
and $(L'a)(n)=q_2(n)a_{n+2}+q_1(n)a_{n+1}+q_0(n)a_n$. By
Theorem~\ref{shr:h4:thm:oeis} and the hypothesis, both vanish for all large
$n$, and hence so does
\[
 q_2(n)(La)(n)-p_2(n)(L'a)(n)
 =(q_2p_1-p_2q_1)(n)\,a_{n+1}+(q_2p_0-p_2q_0)(n)\,a_n.
\]
By the cases of order zero and one, both polynomial coefficients vanish
identically. With $\rho=q_2/p_2\in\Q(n)$ this says
$(q_0,q_1,q_2)=\rho\,(p_0,p_1,p_2)$. The polynomials $p_0,p_1,p_2$ have no
common factor: the irreducible factors of $p_2$ are $n+1$, $n+2$, $2n+3$,
$3n+4$, $3n+5$ and $7n^2-1$, and none of them divides $p_0$, whose
irreducible factors are $2n\pm1$, $4n-1$, $4n+1$, $4n+3$, $4n+5$ and
$7n^2+14n+6$. Write $\rho=u/v$ in lowest terms. Since $q_j=up_j/v$ is a
polynomial, $v$ divides $up_j$, hence $p_j$, for $j=0,1,2$; so $v$ is a
constant and $\rho$ is a polynomial. Each $p_j$ has degree nine, so
$\deg q_2=9+\deg\rho\ge9$.
\end{proof}

The same proof works over $\mathbb C$ in place of $\Q$. The proposition
says nothing about recurrences of order three or more: a left multiple of
the operator by a rational operator can have polynomial coefficients of
lower degree, and Proposition~\ref{shr:h4:prop:higher} below shows that
this happens here. (Revised 5~October 2026 after the intake's independent
check, recorded at the end of Section~\ref{shr:sa:sec:further}; this
sentence first ended ``and whether this happens here is item~(F5)
below''.) The corresponding statement at height five is
Proposition~\ref{shr:h5:prop:nothyper}.

Proposition~\ref{shr:h4:prop:higher} was found and proved by the intake's
independent check, 5~October 2026; it is neither the manuscript's nor the
write's. Below, $E$ is the forward shift, $Ef(n)=f(n+1)E$, and
$L=p_2(n)E^2+p_1(n)E+p_0(n)$ is the operator of
Theorem~\ref{shr:h4:thm:oeis}; a product $(\sum_kQ_kE^k)L$ with rational
$Q_k$ is taken in the Ore ring $\Q(n)[E]$.

\begin{proposition}[\textbf{[check]} Lower degree at orders three to eight]
\label{shr:h4:prop:higher}
Let $a_n=T_4(n)$ (A181198).
\begin{enumerate}[label=(\alph*),leftmargin=*]
\item The left multiple
\[
 \frac{3E+4(1673n+942)}{7n^2+14n+6}\,L=\sum_{j=0}^{3}q_j(n)E^j
\]
has the polynomial coefficients
\begin{align*}
 q_3&=9(n+2)(n+3)^3(2n+5)(3n+7)(3n+8),\\
 q_2&=12(n+2)(2n+3)(15057n^6+97156n^5+222866n^4+200444n^3\\
    &\hspace{14mm}+33411n^2-21546n+6340),\\
 q_1&=-32(2n+1)(4n+3)(4n+5)(86996n^5-17552n^4-200235n^3\\
    &\hspace{14mm}-68149n^2+40579n+13581),\\
 q_0&=-256(2n-1)^2(2n+1)(4n-1)(4n+1)(4n+3)(4n+5)(1673n+942),
\end{align*}
and $\sum_{j=0}^3q_j(n)a_{n+j}=0$ for every $n\ge1$: A181198 obeys a
recurrence of order three and degree eight.
\item The left multiple
\[
 \Bigl(\frac{E^2}{7n^2+28n+27}
 -\frac{4(7315n^2+21945n+16406)\,E}{(7n^2+14n+6)(7n^2+28n+27)}
 +\frac{239136}{7n^2+14n+6}\Bigr)L=\sum_{j=0}^{4}\tilde q_j(n)E^j
\]
has the polynomial coefficients
\begin{align*}
 \tilde q_4&=3(n+3)(n+4)^3(2n+7)(3n+10)(3n+11),\\
 \tilde q_3&=-4(n+3)(2n+5)(29879n^5+358932n^4+1691327n^3\\
    &\hspace{14mm}+3891444n^2+4346534n+1869504),\\
 \tilde q_2&=32(2n+3)(1069163n^6+8737560n^5+26637818n^4\\
    &\hspace{14mm}+34145688n^3+7805720n^2-19807932n-12702660),\\
 \tilde q_1&=-256(2n+1)(4n+3)(4n+5)(321716n^4-767080n^3\\
    &\hspace{14mm}-1948559n^2-1282655n-449532),\\
 \tilde q_0&=-15304704(2n-1)^2(2n+1)(4n-1)(4n+1)(4n+3)(4n+5),
\end{align*}
and $\sum_{j=0}^4\tilde q_j(n)a_{n+j}=0$ for every $n\ge1$: A181198 obeys
a recurrence of order four and degree seven.
\item Every recurrence of A181198, in the sense of this subsection, holds
for every $n\ge1$. No recurrence of order at most three has degree below
eight, and no recurrence of order at most eight has degree below seven.
Hence the smallest degree of a recurrence of order $r$ is nine for $r=2$,
eight for $r=3$, and seven for $4\le r\le8$.
\end{enumerate}
\end{proposition}

\begin{proof}
(a) and (b). Since $7(n+1)^2-1=7n^2+14n+6$ and $7n^2+28n+27$ is
$7n^2+14n+6$ at $n+1$, the first quadratic divides $p_0(n)$ and
$p_2(n+1)$, and the second divides $p_0(n+1)$ and $p_2(n+2)$. Expanding the products with $Ef(n)=f(n+1)E$,
the coefficient of $E^j$ is $\sum_kQ_k(n)p_{j-k}(n+k)$ ($p_i=0$ for
$i\notin\{0,1,2\}$); clearing the denominators and comparing coefficients
gives the displayed polynomials, a finite identity over $\Q[n]$ checked
exactly. The quadratic $7n^2+14n+6$ has discriminant $28$ and no rational
zero, so neither quadratic vanishes at an integer, and for $n\ge1$
\[
 \sum_jq_j(n)a_{n+j}=\frac{3\,(La)(n+1)+4(1673n+942)\,(La)(n)}{7n^2+14n+6}=0
\]
by Theorem~\ref{shr:h4:thm:oeis}; the same argument with three terms
proves (b).

(c) Let $R=\sum_{j=0}^rq_jE^j$ be a recurrence of order $r$, valid for
$n\ge n_0$; by Proposition~\ref{shr:h4:prop:minimal}, $r\ge2$. Since $p_2\ne0$,
right division by $L$ in $\Q(n)[E]$ gives unique rational functions
$Q_0,\ldots,Q_{r-2}$, $s_0$, $s_1$ with
$R=\bigl(\sum_{k=0}^{r-2}Q_kE^k\bigr)L+s_1E+s_0$. Comparing the
coefficients of $E^{r},E^{r-1},\ldots,E^2$ determines them from the top:
$Q_{r-2}=q_r/p_2(n+r-2)$ and
\[
 Q_k=\frac{q_{k+2}-Q_{k+1}(n)\,p_1(n+k+1)-Q_{k+2}(n)\,p_0(n+k+2)}{p_2(n+k)}
 \qquad(k=r-3,\ldots,0),
\]
with $Q_k=0$ for $k>r-2$; then $s_1=q_1-Q_0p_1-Q_1(n)p_0(n+1)$ and
$s_0=q_0-Q_0p_0$. By induction from the top, the denominator of $Q_k$
divides $\prod_{i=k}^{r-2}p_2(n+i)$, and those of $s_0$, $s_1$ divide
$\prod_{i=0}^{r-2}p_2(n+i)$. The real zeros of $p_2$ are $-2$, $-5/3$,
$-3/2$, $-4/3$, $-1$ and $\pm1/\sqrt7$, all below one, so no
$p_2(n+i)$ with $i\ge0$ vanishes at an integer $n\ge1$, and every $Q_k$,
$s_0$, $s_1$ is finite there. At such $n$ the coefficient identity gives
\[
 (Ra)(n)=\sum_{k=0}^{r-2}Q_k(n)\,(La)(n+k)+s_1(n)a_{n+1}+s_0(n)a_n .
\]
For $n\ge n_0$ the left side and every $(La)(n+k)$ vanish, so
$s_1(n)a_{n+1}+s_0(n)a_n=0$ there. Multiplied by a common denominator,
this is a polynomial relation of order at most one that holds for all
large $n$; by the cases of order zero and one in the proof of
Proposition~\ref{shr:h4:prop:minimal} both of its coefficients vanish
identically, so $s_0=s_1=0$. Hence
$(Ra)(n)=\sum_kQ_k(n)(La)(n+k)=0$ for every $n\ge1$, by
Theorem~\ref{shr:h4:thm:oeis}. This proves the first assertion.

Now let $R$ have order at most $\bar r$ and degree at most $d$, and write
$q_j=\sum_{i=0}^dx_{ji}n^i$ for $0\le j\le\bar r$, with $x_{ji}=0$ for $j$
above the order of $R$. By the first assertion the vector
$(x_{ji})\ne0$ solves the integer linear system
$\sum_{j,i}n^ia_{n+j}x_{ji}=0$, $n=1,\ldots,N$, whose $(\bar r+1)(d+1)$
columns are indexed by $(j,i)$ and whose entries are the exact values of
$a_n$ that Theorem~\ref{shr:h4:thm:oeis} determines from $a_1=a_2=1$. For
$(\bar r,d,N)=(3,7,72)$ the $72\times32$ matrix has rank $32$ over $\Q$,
and for $(\bar r,d,N)=(8,6,103)$ the $103\times63$ matrix has rank $63$
(exact fraction-free computation; also full rank modulo $2^{61}-1$, which
by itself implies full rank over $\Q$; recorded in \texttt{data/\allowbreak check-\allowbreak lower-\allowbreak order-\allowbreak recurrences.\allowbreak json}). So no such $R$
exists. Proposition~\ref{shr:h4:prop:minimal} gives degree nine at order
two, (a) gives degree eight at order three, and (b) and its shifts
$\sum_{j=0}^4\tilde q_j(n+k)a_{n+k+j}=0$, $1\le k\le4$, give degree seven at
orders four to eight.
\end{proof}

Proposition~\ref{shr:h5:prop:higher} is the height-five counterpart. The
coefficients of (a) and (b) and the ranks used in (c) are also in
\texttt{data/\allowbreak check-\allowbreak lower-\allowbreak order-\allowbreak recurrences.\allowbreak json}.

\subsection{Further questions and research recorded at the write}
\label{shr:h4:sec:further}

Added 5 October 2026, batch~101, under Vladimir Reshetnikov's standing rule
that unproved claims are recorded as open questions rather than dropped.
The manuscript states no unproved claim; the items below are what it
constructs without evaluating, what it leaves outside its scope, and the
minimality question that the batch-101 placement assigned to this section.
The labels continue those of Section~\ref{shr:bp:sec:further}.
\begin{enumerate}[label=(F\arabic*),leftmargin=*,start=5]
\item \emph{Minimality of the A181198 recurrence} (the placement; no claim
in the manuscript). Settled for the order, and for the degree among
recurrences of order two, by Proposition~\ref{shr:h4:prop:minimal}.
At orders three to eight it is settled by
Proposition~\ref{shr:h4:prop:higher}: the smallest degree is eight at
order three and seven at each order from four to eight. Still open:
whether some recurrence of order nine or more has degree below seven.
Missing: a degree bound for left multiples $QL$ that holds at every order
(the rank computation in the proof of
Proposition~\ref{shr:h4:prop:higher} treats each range of orders
separately). (Re-scoped 5~October 2026 after the intake's independent
check, which found and proved Proposition~\ref{shr:h4:prop:higher}; this
item first read: ``Still open: whether some recurrence of order three or
more has polynomial coefficients of degree below nine. Missing: a degree
bound for left multiples $QL$ with $Q$ a rational operator, or an
exhaustive search over small orders and degrees.'')
\item \emph{Numerical constants} (Sections~\ref{shr:h4:sec:asymptotics}
and~\ref{shr:h4:sec:inverse}). The manuscript constructs $E_K$, $N_K$,
$C_K$, $H_K$, $W_K$, $X_K$ and the onset \eqref{shr:h4:eq:thresholdonset}
by finite rational arithmetic, but evaluates none of them; its verifier
does not compute them (its README says so). Missing: the evaluation, for
example at $K=0$ and $K=6$, as Part~\ref{shr:h5:part}'s receipt does at
height five \eqref{shr:h5:asy:logexamples}, and a worked certified
threshold.
\item \emph{Even heights six and more} (Section~33.4). ``Higher Pfaffian
sizes contain further trace invariants, so an analogous low-dimensional
reduction would require new work.'' Missing: a reduction of the traces
$\tr(B^j)$, $j\le m/2$, to a few polynomial periods, with their
contiguity; Research question~16 asks the same at odd heights.
\item \emph{Exponentially small content} (Section~31.4, a non-claim). The
homogeneous mode of \eqref{shr:h4:eq:brec} has ratio tending to $-1/27$;
the manuscript selects no canonical exponentially small component and no
Stokes data. Missing: the analysis that Research question~18 asks for at
height five.
\end{enumerate}

\clearpage
\part{A complete exact formula for height five shifted rectangles}
\label{shr:h5:part}

\noindent{\Large A proof of A181199 and Kauers and Koutschan
Conjecture~19\par}
\smallskip
\noindent{\large Research report 233, prepared for Vladimir Reshetnikov,
October 5, 2026\par}
\medskip

\paragraph{Main conclusion.}
The full height-five array count satisfies the nested finite sum printed as
Conjecture 19 of Kauers and Koutschan and the literal third-order,
degree-24 recurrence recorded in OEIS A181199, for every integer $n\ge1$.
All boundary contributions are retained.

\paragraph{Abstract.}
Column interlacing turns the counting problem into an exterior Volterra
transfer. Two odd augmented Pfaffians reduce its exact volume to traces
and three open-chain integrals. A rank-two calculation closes the
length-four open chain and expresses the full count through two polynomial
periods. We derive differential equations, polynomial contiguity, and a
finite moment closure with every endpoint specified. Rational identities
then give a first difference, the actual printed nested sum with its
rising factorials, and an explicit factorization of the prescribed OEIS
operator. The proof is uniform in the width, with separate treatment of
width one. Supplementary exact arithmetic verifies the symbolic
certificates and independent finite counts. Recurrence-based asymptotic
and inverse-index consequences are separated from the exact enumeration;
the earlier fixed-height asymptotic results receive explicit credit.

\paragraph{Status and limits.}
AI-assisted research manuscript, not a peer-reviewed or proof-assistant
formalization. Historical priority is not certified. The recurrence is
proved on $n\ge1$; it fails at $n=0$ under the empty-array convention.

\begin{writenote}[Added 5 October 2026, batch~101: provenance of Part~IV]
Part~IV is a single manuscript, printed in full: Research Report~233 of
the session bundle that arrived in commit \texttt{60f54ea06} and was placed
in commit \texttt{f7c612c72} (batch~101) as an addition to this report.
Its title is \emph{A complete exact formula for height five shifted
rectangles}, subtitled ``A proof of A181199 and Kauers and Koutschan
Conjecture 19'', dated October~5, 2026 (a 30-page PDF; a main file and
eleven section files, 1{,}860 lines with 135 labels). It arrived as the
archive \texttt{Report233.zip} (614{,}266 bytes, 22 files under the doubled
wrapper directory \texttt{Report233/\allowbreak Report233/}). Its title page
reads ``Prepared for Vladimir Reshetnikov'' and its PDF metadata name
``Research report prepared for Vladimir Reshetnikov''; no human author and
no tool are named.

\emph{Pin.} Like Part~\ref{shr:h4:part}, the manuscript names no commit and
pins this report's source by the blob
\texttt{2b40ded9\allowbreak daa6549f\allowbreak 786a52b8\allowbreak 80d5477d\allowbreak e3b5adfc}, the text of
\texttt{6fef5383b} (3~October 2026): it read Part~\ref{shr:part:one} only.
It cites neither the manuscript of Part~\ref{shr:bp:part} nor that of
Part~\ref{shr:h4:part}.

\emph{Files.} Its five programs (its \texttt{exact\_algebra.py} is not the
same file as Part~\ref{shr:h4:part}'s), its receipt and its source record
are shipped with the prefix \texttt{04-height-five-}, in \texttt{code/}, in
\texttt{data/} and at the report root; the note at the end of
Section~42.3 lists them. Its \LaTeX\ source,
delivered README, PDF and \texttt{MANIFEST.sha256} (21 entries, all verified
at placement) are not shipped and survive in the arrival commit's archive.
The manuscript is unrefereed and AI-assisted; nothing in it is formalized,
and its place in this report confers no formal status.

\emph{What it answers.} Theorems~\ref{shr:h5:thm:sum}
and~\ref{shr:h5:thm:operator} prove, from the array definition and for
every $n\ge1$, the nested finite sum printed as Conjecture~19 of
\cite{KK} and the order-three, degree-24 recurrence displayed as
conjectural in OEIS A181199. With Part~\ref{shr:h4:part} this answers
Research question~1 of Part~\ref{shr:part:one} and Research question~8 of
Part~\ref{shr:bp:part}; dated notes there say so. The manuscript also shows
that the operator does not extend to $n=0$ under $a_0=1$ (residual
$5621993879040000$) and that its leading coefficient $c_3$ is positive for
$n\ge1$. The write adds Section~\ref{shr:h5:sec:write}: a remark that the
period and its recurrence are Part~\ref{shr:h4:part}'s, and a proof that
A181199 obeys no first-order recurrence, so the minimal order is two or
three.

\emph{Two independent derivations.} Reports~231 and~233 bear the same date
and neither cites the other. Part~IV re-derives, at height five, the
column transfer, the composition identity, the rank-two form of $B$ and the
scalar integrations by parts of Part~\ref{shr:h4:part}, and it defines the
same period $Z_n$; its recurrence \eqref{shr:h5:eq:Z-shift} is
Theorem~\ref{shr:h4:thm:Z}, proved here by contiguity and a moment closure
instead of a divergence certificate (Remark~\ref{shr:h5:rem:sameZ}). New
in Part~IV: the odd augmented Pfaffian and the enumeration of its 225
products, the open-chain reduction $q_2=16G-12H/N^2+2/N^5$, which has no
analogue at height four, the second period $X_n$, the polynomial
contiguity, the moment closure, the first difference
\eqref{shr:h5:eq:target-diff}, and the factorization
\eqref{shr:h5:eq:operator-factorization}. The two texts share 4.3\,\% of
their word 8-grams (233 against 231), and Part~IV shares 0.2\,\% and
0.5\,\% with Parts~\ref{shr:part:one} and~\ref{shr:bp:part}.

\emph{Overlap with Part~\ref{shr:part:one}, credited.}
Section~\ref{shr:h5:asy:section} recovers, from the recurrence, the case
$m=5$ of Theorem~\ref{shr:thm:main} and
Corollary~\ref{shr:cor:A181199}, as the manuscript says: its
$c_j/c_0$ are $c_{5,j}$ (Table~1), with $c_0=9/32768=J_5/4^{10}$, and
its $d_1,\ldots,d_5$ are those of \eqref{shr:eq:A181199exp}. The sixth
corrections $-6103/16$ and $d_6=-2709616559/3125000$ are new to the
repository, ``without a novelty claim''; at placement
Part~\ref{shr:part:one}'s engine run at order six on a copy gave the same
$b_{5,6}$. The effective constants are new. Section~\ref{shr:h5:inv:section}
uses the same apparatus as Section~\ref{shr:h4:sec:inverse}.

\emph{Source checks.} The manuscript transcribes Conjecture~19 from
arXiv:2303.02793v2 (Section~6.4, printed page~34; ``the rising-factorial
overbars were checked on the printed page'') and the operator from the
A181199 internal record, consulted 5~October 2026. At placement the intake
did not re-inspect the printed page; it checked that the four coefficients
$c_0,\ldots,c_3$ of Appendix~\ref{shr:h5:app:literal} equal those of the
live A181199 entry, checked the $y$-contiguity certificate and the bridge
\eqref{shr:h5:eq:gamma-bridge} in SymPy (residual zero), and checked
Theorem~\ref{shr:h5:thm:sum}, \eqref{shr:h5:eq:full-count} and
\eqref{shr:h5:eq:target-diff} for $n\le20$ against an independent row-state
count; the transcribed operator annihilates those counts for
$n=1,\ldots,15$. Its verifier, rerun at the write on a copy
(Section~42.3), printed its receipt byte for byte.

\emph{Editorial changes.} Section~$k$ of the manuscript is
Section~$k+33$ here and its Appendix~A is Appendix~\ref{shr:h5:app:literal},
printed after Appendix~D; theorem and equation numbers follow their
sections, and its table is Table~\ref{shr:h5:tab:literal-arrays}. Its four Questions are Research
questions~16--19 (one counter runs through the report). Every label
carries the prefix \texttt{shr:h5:}, including its \texttt{asy:} and
\texttt{inv:} labels. Its title, subtitle and title-page paragraphs open
this Part; its table of contents is dropped. Its three citations of its
predecessor now point to Part~\ref{shr:part:one}. The bibliographies are
merged: Kauers--Koutschan, Sun and DLMF~\S\S4.13 and~5.11 are
Part~\ref{shr:part:one}'s entries, and its OEIS entry is
\cite{oeis199int}. Its bibliography cites Sun's article, which its text
does not cite. It said of Part~\ref{shr:part:one}: ``Its all-height
asymptotic results precede the present report; its explicit scope leaves
the particular recurrences unproved. The source is cited for credit, not
redistributed in this package.'' Remarks take this report's theorem style.
The \textbf{[write]} notes and Section~\ref{shr:h5:sec:write} with
Remark~\ref{shr:h5:rem:sameZ}, Proposition~\ref{shr:h5:prop:nothyper} and
Section~\ref{shr:h5:sec:further} are new. No other statement, proof or
number of the manuscript was changed.
\end{writenote}

\begin{writenote}[Added 5 October 2026, batch~101: notation of Part~IV]
Part~IV keeps the manuscript's symbols, including the reuses that it flags
itself ($q_j$, $V$ and $V_j$, $E$, $b$ and $b_n$, $J_0$ and $J_j$). The
count is $a_n=T_5(n)$, with array entries $a_{i,j}$ as in
Part~\ref{shr:part:one}. Within Part~IV each symbol has the meaning below;
the right column gives the tempting false reading from elsewhere in this
report. The kernel functions $W$, $h$, $e$, $w$ and the scalars $c$, $H$,
$U$, $I$, $J_0$ are those of Part~\ref{shr:h4:part} at height five, with
$f=t^n/N$ for Part~III's $\alpha$ and $b=t^{n-1}/(n-1)!$ for its $\beta$.
{\footnotesize
\begin{longtable}{@{}>{\raggedright\arraybackslash}p{0.17\linewidth}>{\raggedright\arraybackslash}p{0.43\linewidth}>{\raggedright\arraybackslash}p{0.34\linewidth}@{}}
\toprule
symbol & meaning in Part~IV & \emph{not}: meaning elsewhere\\
\midrule\endhead
$T_j$; $B$ & $\tr(B^j)$ \eqref{shr:h5:eq:Tq}; the kernel $EA$ & the count $T_m(n)$!; Part~I's $B$, Part~II's $B_i$\\
$P(n)$, $Q(n)$, $D(n)$ & the polynomials \eqref{shr:h5:eq:PQD} & Part~III's field $P$, $Q$ and $D_n$; Part~II's $Q_m(n)$, $D_i$\\
$u(n)$, $v(n)$; $u$, $v$ & \eqref{shr:h5:eq:u}, \eqref{shr:h5:eq:v}; the triangle variables of \eqref{shr:h5:eq:triangle} & ---\\
$\gamma_n$ & \eqref{shr:h5:eq:gamma} & Part~III's certificate constant $\gamma_n$\\
$N$; $M_n$; $F_n$ & $n!$; $M_5(n)$; $a_n/M_n$ (also written $b_n$ in Section~\ref{shr:h5:asy:section}) & $F_m(z)$; Part~III's prefactor $F_n$; $b_{m,j}$\\
$Z_n$; $X_n$ & the period, equal to Part~III's; the second period \eqref{shr:h5:eq:periods} & $Z_m(n)$; Part~II's binomial counts $X_i$\\
$\nu_n$, $\delta_n$ & the period recurrence \eqref{shr:h5:eq:Z-shift} & Part~V's $\nu=\alpha_m$\\
$C_m$ & the ordered simplex & Part~V's leading constant $C_m=K_m$\\
$V$; $V_j$ & the Volterra operator; moments \eqref{shr:h5:eq:moment-defs} & ---\\
$K$, $J$, $E$, $A$; $b$ & kernels \eqref{shr:h5:eq:EAK}; $b=K\one$ & $K_m$; $J_m$; $E_{m,n}(p)$; $A=\binom m3$; $b_{m,j}$\\
$E$ & the sign kernel (Section~\ref{shr:h5:sec:transfer}); the forward shift (Section~\ref{shr:h5:sec:sum-operator}) & ---\\
$q_j$; $q(t)$, $r(t)$ & $\ip b{B^j\one}$ (Sections~\ref{shr:h5:sec:transfer}--\ref{shr:h5:sec:rank}); contiguity polynomials \eqref{shr:h5:eq:qr} with coefficients $q_j$, $r_j$ & $q_{ij}=(y_i+y_j)^2$; $r=m-2$\\
$G$; $L$ & $\ip f{Wh}$ \eqref{shr:h5:eq:HUG}; $\ip b{Ww}$ \eqref{shr:h5:eq:q2-expand} & $L$ a truncation order\\
$y_n$, $x_n$; $k(n)$, $\ell(n)$; $c_n^*$ & polynomials \eqref{shr:h5:eq:periods}; \eqref{shr:h5:eq:kell}; \eqref{shr:h5:eq:cnstar} & Part~I's $Y_i$, $X_i$\\
$\lambda$, $\mu$, $\nu$ & contiguity constants \eqref{shr:h5:eq:lambda-mu-nu}, $\nu=\nu_n$; later $\lambda=\log3125$ & $\lambda=m\log m$\\
$V_j$, $Y_j$, $U_j$, $J_j$; $[p]_V$ & moment families \eqref{shr:h5:eq:moment-defs} and their linear forms & Part~I's $Y_i$; Part~III's $U$\\
$A_5(n)$, $\rho_n$, $\sigma_n$, $\tau_n$, $P_7$, $P_9$ & the $X$ transport \eqref{shr:h5:eq:X-shift}--\eqref{shr:h5:eq:P7P9} & Part~V's cut proportion $\rho$\\
$H_n$, $I_n$ & hypergeometric terms \eqref{shr:h5:eq:HI} & Part~II's $H_m(n)$; Part~III's $I_j$\\
$R(n)$, $S(n)$ & ratios \eqref{shr:h5:eq:RS-ratios} & Part~III's $R(n)$, $S(n)$\\
$c_0,\ldots,c_3$; $c_j$ & the OEIS coefficients (Appendix~\ref{shr:h5:app:literal}); $\beta_{10+j}$ in Section~\ref{shr:h5:asy:section} & $c_{m,j}$; Part~III's $c_j$\\
$r(t)$, $v_*$, $d_*$, $\alpha(t)$, $\eta(t)$, $\mathcal Z$, $\mathcal B$, $z_j$, $\beta_j$ & Section~\ref{shr:h5:asy:section} & $\alpha_m$; $\beta_j$ (Section~6.2)\\
$C_*$, $d_j$, $\ell_j$; $J$ & $K_5$, $b_{5,j}$, logarithmic coefficients; $\lfloor(K+1)/2\rfloor$ in Section~40.3 & $d=\binom m2$; $J_m$\\
$\mathcal N(y)$ & the threshold \eqref{shr:h5:inv:threshold} & Part~III's $N(y)$\\
$B_{21}$, $C_{19}$, $b_j$, $d_j$, $Q_+$, $P_+$ & integer coefficient data of Appendix~\ref{shr:h5:app:literal} & the $d_j$ of Section~\ref{shr:h5:asy:section}\\
\bottomrule
\end{longtable}}
\end{writenote}

\section{The exact target and the scope of the result}
\label{shr:h5:sec:scope}
For a positive integer $n$, let $a_n$ be the number of arrays
$(a_{i,j})_{1\le i\le5,\,1\le j\le n}$ containing $1,\ldots,5n$ exactly
once and satisfying
\begin{equation}
\begin{aligned}
 a_{i,j}&<a_{i,j+1},& a_{i,j}&<a_{i+1,j},\\
 a_{i,j}&<a_{i+1,j+1},& a_{i,j+1}&<a_{i+1,j},
\end{aligned}
\label{shr:h5:eq:arrays}
\end{equation}
whenever the indicated cells exist. These are the literal row, column,
downward-diagonal, and downward-antidiagonal inequalities. The downward
diagonal is redundant; the antidiagonal is essential. Height is fixed at
five throughout. In particular, the width-two count is one, although the
height-two problem at variable width has a different answer.

The sequence is OEIS A181199~\cite{oeis199int}. Kauers and Koutschan state a
specific nested finite sum as Conjecture~19 in Section~6.4 of their
article~\cite{KK}, printed page~34. We first give that formula in a compact
notation which is also used in the proof. Set
\begin{align}
 P(n)={}&25216n^8+9888n^7-14496n^6+11208n^5+23832n^4\\
       &+7383n^3-1522n^2-939n-90,\\
 Q(n)={}&137855872n^{11}+860969696n^{10}+2047036856n^9
          +2032587274n^8\\
       &-24192441n^7-1894061166n^6-1671661480n^5
          -524330624n^4\\
       &+36004789n^3+62751860n^2+13865604n+927360,\\
 D(n)={}&100864n^{11}+241280n^{10}+109376n^9-61744n^8
          +105456n^7\\
       &+247584n^6+103814n^5-40103n^4-33803n^3
          -1331n^2+2445n+306.
 \label{shr:h5:eq:PQD}
\end{align}
For $r\ge0$, the notation $\rise zr=z(z+1)\cdots(z+r-1)$ denotes a
\emph{rising factorial}, with the empty product equal to one. Define
\begin{align}
 u(n)&=\frac{8P(n)}{(2n-1)(4n-1)\rise{3n+1}{3}\rise{4n+1}{4}},
 \label{shr:h5:eq:u}\\
 v(n)&=\frac{(3n+1)(3n+2)(4n+3)Q(n)}
 {(n+1)^2(n+2)^2(2n-1)(2n+1)(2n+3)P(n)P(n+1)},
 \label{shr:h5:eq:v}\\
 \gamma_n&=\frac{(4n+1)D(n)}{16(n+1)^2(2n-1)(2n+1)P(n)}.
 \label{shr:h5:eq:gamma}
\end{align}
The overbars in the source's denominator of $u$ are important: replacing
the rising factorials by powers changes the conjecture. The other
degree-eight factor printed in the denominator of $v$ is exactly
$P(n+1)$, as explicitly verified in Appendix~\ref{shr:h5:app:literal}.

\begin{theorem}[The printed finite sum]\label{shr:h5:thm:sum}
For every integer $n\ge1$,
\begin{equation}
 a_n=1-\frac{27}{4}\sum_{k=1}^{n-1}
 \left\{(-1)^k u(k)\frac{(5k)!}{(3k)!(k!)^2}
 \sum_{i=1}^{k-1}(-1)^i v(i)\frac{(3i)!}{(i!)^3}\right\}.
 \label{shr:h5:eq:main-sum}
\end{equation}
All displayed denominators are nonzero on this range. In particular,
$a_1=a_2=1$ and $a_3=16$.
\end{theorem}

\begin{theorem}[The literal OEIS recurrence]\label{shr:h5:thm:operator}
Let $c_0,c_1,c_2,c_3$ be the polynomials displayed in
Appendix~\ref{shr:h5:app:literal}. They are precisely the four coefficients of the
third-order, degree-24 recurrence in the inspected A181199 entry. Then
\begin{equation}
 c_0(n)a_n+c_1(n)a_{n+1}+c_2(n)a_{n+2}+c_3(n)a_{n+3}=0
 \qquad(n\ge1).
 \label{shr:h5:eq:literal-rec}
\end{equation}
The coefficient $c_3(n)$ is strictly positive there, so these three initial
values determine the sequence uniquely.
\end{theorem}

\subsection{A smaller identity that proves both statements}
Put $N=n!$ and $M_n=(5n)!/N^5$. The proof will establish
\begin{equation}
 a_{n+1}-a_n=2M_nu(n)(Z_n-\gamma_n),
 \label{shr:h5:eq:target-diff}
\end{equation}
where the period is the positive triangle integral
\begin{equation}
 Z_n=n^2\iint_{\substack{0<u,v<1\cr u+v>1}}
      u\{uv(u+v-1)\}^{n-1}\dd u\dd v.
 \label{shr:h5:eq:triangle}
\end{equation}
It obeys the elementary first-order affine recurrence
\begin{align}
 Z_{n+1}&=\nu_nZ_n+\delta_n,\\
 \nu_n&=-\frac{(n+1)^2}{3(3n+1)(3n+2)},&
 \delta_n&=\frac{28n^3+50n^2+29n+5}
 {6(2n+1)(3n+1)(3n+2)}.
 \label{shr:h5:eq:Z-shift}
\end{align}
The initial value is $Z_1=1/3$. Equations~\eqref{shr:h5:eq:target-diff}
and~\eqref{shr:h5:eq:Z-shift}, together with an explicit rational identity for
$\gamma_n$, will prove the actual nested sum, rather than only a compatible
recurrence. A subsequent exact Ore-operator identity proves
Theorem~\ref{shr:h5:thm:operator}.

\subsection{What the present proof adds}
The earlier report \emph{Fixed-height shifted rectangles: asymptotics and
algebraicity} (Part~\ref{shr:part:one}) already established the leading asymptotic for
A181199, all-fixed-order expansions at every fixed height, and further
algebraicity results. Its stated scope explicitly leaves the particular
A181198 and A181199 recurrences unproved. We do not present its
asymptotic conclusions as new. The new mathematical task here is exact
full-count enumeration: the odd augmented Pfaffian, its open-chain
reduction, polynomial period transport, and the specified finite sum and
operator. No asymptotic estimate is used in their proof.

The source comparison is bounded. It covers the cited article, the
inspected OEIS record, and the specified earlier report, pinned in the
bibliography. Exact-identifier searches located the OEIS and original
article descriptions; a conjectural label in a source is not an exhaustive
priority search. No global novelty or external peer-review claim follows.
The standard tools, including order-polytope volumes, determinant
composition, and Pfaffian integration, are not inventions of this report.
Their required special cases are proved below so that the exact result is
self-contained.

\subsection{Indexing and boundary discipline}
The proof concerns positive widths only. At width one, the count is one
and the required period formulas are checked directly. This avoids using a
singular transfer kernel at that index. Moreover, if one adopts the
optional convention $a_0=1$, the left side of~\eqref{shr:h5:eq:literal-rec} at
$n=0$ is
\[
 5621993879040000,
\]
not zero. The recurrence must not be extended to zero without modification.
A boundary term that is exponentially small for large $n$ can still affect
an exact recurrence. The integral proof below retains every such term.

\section{Column transfer and the odd augmented Pfaffian}
\label{shr:h5:sec:transfer}
All integrals in this section are ordinary Lebesgue integrals on compact
cubes or ordered simplices. Equalities among coordinates occupy a set of
measure zero and have no effect on the volumes.

\subsection{The transfer determinant and its normalization}
Write
\[
 C_m=\{0<x_1<\cdots<x_m<1\},\qquad
 V(t,s)=\one_{s<t}.
\]
The same letter $V$ denotes the integral operator on $[0,1]$.
For adjacent ordered columns $x,y$, the inequalities~\eqref{shr:h5:eq:arrays}
are equivalent to
\begin{equation}
 x_1<y_1<x_2<y_2<\cdots<x_5<y_5.
 \label{shr:h5:eq:interlace}
\end{equation}
Indeed the row and antidiagonal inequalities give the displayed chain,
which implies the two within-column orders and the required downward
diagonal comparisons. Thus no extra restriction has been introduced.

\begin{lemma}\label{shr:h5:lem:indicator}
On $C_m\times C_m$, the determinant
$\det[V(y_i,x_j)]_{i,j=1}^m$ is the indicator of
$x_1<y_1<\cdots<x_m<y_m$.
\end{lemma}
\begin{proof}
Each row consists of an initial string of ones followed by zeros, and
these strings are nested. Let their lengths be $r_1\le\cdots\le r_m$.
If any length is zero, or two lengths coincide, the determinant vanishes.
Otherwise the lengths must be $1,\ldots,m$, in which case the matrix is
lower triangular with diagonal one. Those lengths are exactly the asserted
interlacing conditions.
\end{proof}

For kernels $K_1,K_2$ with integrable products, determinant expansion and
Fubini give
\begin{equation}
 \int_{C_m}\det[K_1(y_i,z_j)]\det[K_2(z_i,x_j)]\dd z
   =\det[(K_1K_2)(y_i,x_j)].
 \label{shr:h5:eq:andreief}
\end{equation}
To check the factor, the product of determinants is symmetric in the
$z$ variables. Its cube integral is $m!$ times its ordered-simplex
integral, while expansion of the cube integral gives $m!$ times the
determinant on the right. The two factors cancel.

Every linear extension of the array poset occupies a simplex of volume
$1/(5n)!$ in the unit cube. Applying Lemma~\ref{shr:h5:lem:indicator} between
successive columns and~\eqref{shr:h5:eq:andreief} at the intermediate columns
therefore proves, for $n\ge2$,
\begin{equation}
 \frac{a_n}{(5n)!}
 =\int_{C_5\times C_5}\det[K(y_i,x_j)]\dd x\dd y,
 \qquad K=V^{n-1},
 \label{shr:h5:eq:transfer}
\end{equation}
where
\[
 K(t,s)=\frac{(t-s)_+^{n-2}}{(n-2)!}.
\]
At $n=2$, this exponent-zero truncated-power notation means the indicator
$\one_{s<t}$, consistent with $K=V$; no value of $0^0$ is used.
There is no factorial left over from any intermediate column. Width one
will be handled separately in Section~\ref{shr:h5:sec:periods}.

\subsection{The bordered identity}
Let $J=\one\otimes\one$ be the constant kernel and set
\begin{equation}
 E(t,s)=\sgn(s-t)=J-2V=2V^*-J,\quad
 A=KEK^*,\quad b=K\one=\frac{t^{n-1}}{(n-1)!}.
 \label{shr:h5:eq:EAK}
\end{equation}
The kernel $A$ is skew-symmetric. The odd Pfaffian integral is
\begin{equation}
 \int_{C_5}\det[K(y_i,x_j)]\dd x
 =\Pf\begin{pmatrix}(A(y_i,y_j))_{i,j=1}^5&(b(y_i))_{i=1}^5\\
                   -(b(y_j))_{j=1}^5&0\end{pmatrix}.
 \label{shr:h5:eq:odd-debruijn}
\end{equation}
We spell out its normalization. For an even number $2r$ of rows, put
\[
 M_{ij}=\iint K(y_i,s)E(s,t)K(y_j,t)\dd s\dd t.
\]
The Pfaffian definition is
\[
 \Pf(M)=\frac{1}{2^r r!}\sum_{\sigma\in S_{2r}}
  \sgn(\sigma)\prod_{j=1}^r M_{\sigma(2j-1),\sigma(2j)}.
\]
Expand the pairwise integrals and combine the alternating sum into a
determinant. With $D(z)=\det[K(y_i,z_j)]$, this gives
\[
 \Pf(M)=\frac1{2^rr!}\int_{[0,1]^{2r}}
                 D(z)\prod_{j=1}^rE(z_{2j-1},z_{2j})\dd z.
\]
Relabeling the integration variables in each term of $\Pf(E(z_i,z_j))$
then gives
\[
 \Pf(M)=\frac1{(2r)!}\int_{[0,1]^{2r}}
                         D(z)\Pf(E(z_i,z_j))\dd z.
\]
Both factors change by the same sign under permutation, and the sign
Pfaffian is one on the ordered chamber by its elementary recursion.
Partitioning the cube into its $(2r)!$ ordered chambers proves the even
integration formula with factor one.
For the odd version, adjoin a point $\star$ of mass one, ordered after
$[0,1]$. Extend $K$ by $K(\star,\star)=1$ and zero mixed entries, and use
six rows with the last row at $\star$. In a nonzero determinant integral
the atom appears exactly once. The remaining determinant is the original
five-variable determinant; its pairwise skew integrals form $A$, and the
last column is $K\one=b$. This proves~\eqref{shr:h5:eq:odd-debruijn} with factor
one, including the border sign.

For ordered $y$, the corresponding sign Pfaffian is
\begin{equation}
 \Pf\begin{pmatrix}(E(y_i,y_j))&(1)\\-(1)&0\end{pmatrix}=1.
 \label{shr:h5:eq:sign-pf}
\end{equation}
All entries above its diagonal equal one. The Pfaffian expansion along
the first row gives an alternating sum of five ones, equal to one,
and the same recursion starts from the $2\times2$ case. Multiplying the
right side of~\eqref{shr:h5:eq:odd-debruijn} by~\eqref{shr:h5:eq:sign-pf} allows us to
replace its ordered integral by the cube integral divided by $5!$:
permuting the five ordinary labels changes both Pfaffians by the same sign.

\subsection{The complete enumeration of the 225 products}
Put
\begin{equation}
 B=EA,\qquad T_j=\tr(B^j),\qquad q_j=\ip b{B^j\one}.
 \label{shr:h5:eq:Tq}
\end{equation}
Each six-label perfect matching has one edge incident to the border.
There are $5\cdot3=15$ such matchings, hence $225$ products of terms in the
two bordered Pfaffians. The following argument accounts for every product
and every integrated sign, rather than suppressing the odd open path.

Superpose the two matchings, remove their border vertices, and retain the
border weights $b$ and $1$ at the two exposed ordinary vertices. The
component containing those vertices is an alternating open path. Every
other component is an alternating even cycle. The open path has one,
three, or five ordinary vertices. The four exhaustive classes are
\begin{center}
\begin{tabular}{@{}lrr@{}}
\toprule
Open path and remaining components&Number&Integrated term\\
\midrule
One vertex and two doubled edges&$5\cdot3=15$&$q_0T_1^2$\\
One vertex and one four-cycle&$5\cdot6=30$&$-q_0T_2$\\
Three vertices and one doubled edge&$\binom53\cdot6=60$&$-q_1T_1$\\
Five vertices&$5\cdot4\cdot3!=120$&$q_2$\\
\bottomrule
\end{tabular}
\end{center}
For the three-vertex class, choose its labels and its ordered endpoints;
the middle vertex and remaining doubled edge are forced. For the
five-vertex class, choose the ordered endpoints and the order of the three
internal labels. For the one-vertex class, fix its label; the two matchings
on the other four labels either coincide in one of three ways, or differ
in one of six ordered ways. These counts sum to $225$.

Here are canonical signed representatives of the four classes, with
$E_{ij}=E(y_i,y_j)$ and similarly for $A$:
\begin{align*}
 &E_{12}E_{34}A_{12}A_{34}b_5,\\
 &-E_{12}E_{34}A_{13}A_{24}b_5,\\
 &-E_{12}E_{34}A_{12}A_{35}b_4,\\
 &E_{12}E_{34}A_{13}A_{25}b_4.
\end{align*}
The missing border weight in each case is one. Integrating
$E_{12}A_{12}$ gives $-T_1$, since $A_{21}=-A_{12}$.
The product $E_{12}E_{34}A_{13}A_{24}$ integrates to $T_2$.
Further, $b_4E_{34}A_{35}$ integrates to $-q_1$.
In the last representative, reversing both $E_{34}$ and $A_{13}$ gives
precisely the oriented product defining $\ip b{EAEA\one}=q_2$.
These establish the table's four signs. Relabeling variables changes the
two original Pfaffian signs together, so it preserves their product;
therefore the canonical sign applies to every term in its class.

For complete algorithmic enumeration, pair the smallest unpaired label
$i$ with each possible $j$, multiply the recursive sign by
$(-1)^{p-1}$ when $j$ is in position $p$ among the remaining labels, and
continue. This generates all 15 signed matchings on labels $1,\ldots,6$.
Taking the Cartesian square and following each exposed path generates
exactly the table above. The accompanying verifier implements that finite
procedure; the proof of the counts and signs is the preceding argument.

Division by $120=5!$ gives the crucial full-volume formula
\begin{equation}
 \boxed{\frac{a_n}{(5n)!}
   =q_0\left(\frac{T_1^2}{8}-\frac{T_2}{4}\right)
       -\frac{q_1T_1}{2}+q_2.}
 \label{shr:h5:eq:pf-volume}
\end{equation}
All kernels here are bounded and piecewise polynomial for $n\ge2$.
The notation $\tr(B^j)$ means the finite cyclic kernel integral, so
Fubini is sufficient. No formal infinite-dimensional determinant or
unjustified trace expansion is being used.

\section{The Volterra rank reduction and the open chain}
\label{shr:h5:sec:rank}
The odd Pfaffian formula contains a term $q_2$ that has no analogue in the
height-four trace-only calculation. Its exact reduction is the key extra
step. Throughout this section $n\ge2$; the kernel identities themselves
will also be checked at $n=1$ where needed later.

\subsection{The kernel and its elementary identities}
Write $N=n!$ and introduce the function $f(t)=t^n/N$, so $f'=b$.
Define the symmetric kernel and its boundary functions by
\begin{align}
 W&=V^n(V^*)^n,&
 W(t,s)&=\frac1{((n-1)!)^2}
   \int_0^{\min(t,s)}(t-r)^{n-1}(s-r)^{n-1}\dd r,\\
 h(t)&=W(t,1),& e(t)&=2h(t)-\frac{b(t)}{N},\\
 w&=W\one,&z&=Wf.
 \label{shr:h5:eq:kernel-defs}
\end{align}
Inner products are real integrals on $[0,1]$. Set
\begin{equation}
 H=\ip f h,\quad U=\tr(W^2),\quad
 G=\ip f{Wh}=\ip z h,\quad c=\ip f b=\frac1{2N^2}.
 \label{shr:h5:eq:HUG}
\end{equation}

Preserving the order of noncommuting factors in~\eqref{shr:h5:eq:EAK} gives
\begin{align}
 B=EKEK^*
 &=\bigl(JV^{n-1}-2V^n\bigr)
       \bigl(2(V^*)^n-J(V^*)^{n-1}\bigr)\\
 &=-4W+\one\otimes e+2f\otimes b.
 \label{shr:h5:eq:rank2}
\end{align}
Here $(r\otimes s)g=r\ip s g$. For example,
$V^nJ(V^*)^{n-1}=f\otimes b$, and the first term containing $J$
is $2\one\otimes h$. The remaining double-$J$ term is
$-\one\otimes b/N$. This supplies a direct check of the border
normalization in the rank reduction.

The defining integral gives
\begin{equation}
 W(t,0)=0,\qquad
 (\partial_t+\partial_s)W(t,s)=b(t)b(s),\qquad
 W(rt,rs)=r^{2n-1}W(t,s).
 \label{shr:h5:eq:W-identities}
\end{equation}
For the differential identity, the sum of the external derivatives of
the integrand is its negative derivative in $r$. The endpoint $r=0$
therefore supplies $b(t)b(s)$, while the moving upper endpoint vanishes
for $n\ge2$. At $n=1$, $W(t,s)=\min(t,s)$ and the same identity holds
on both sides of the diagonal directly. The kernel is continuous across
the diagonal; integrations by parts on the two triangles have cancelling
internal boundary fluxes.

Direct integration also gives
\begin{equation}
 \int_0^1h=\frac1{2N^2},\quad
 \int_0^1b=\frac1N,\quad
 \int_0^1e=0,\quad
 \tr W=\frac1{2n(2n-1)((n-1)!)^2}.
 \label{shr:h5:eq:basic-integrals}
\end{equation}
For instance, integrating $t$ first in $\int h$ leaves
$((n-1)!)^{-2}n^{-1}\int_0^1(1-r)^{2n-1}\dd r=1/(2N^2)$.
On the diagonal $W(t,t)=t^{2n-1}/((2n-1)((n-1)!)^2)$, giving the trace.

\subsection{The scalar integration by parts identities}
We need three scalar integrals:
\begin{equation}
 \ip b w=H,\qquad I:=\ip b z=\frac HN-\frac1{8N^4},\qquad
 J_0:=\ip h w=\frac HN-\frac1{8N^4}=I.
 \label{shr:h5:eq:scalar-identities}
\end{equation}
Here $J_0$ is a temporary scalar name, unrelated to the moment $J_j$
introduced later. To prove the first identity, integrate the sum of the
$t$ and $s$ derivatives of $f(t)W(t,s)$ over the square. The interior
integral is
$\ip b w+c/N$. The upper boundaries give
$H+(\int h)/N$; both lower boundaries vanish because $f(0)=0$ and
$W(t,0)=0$. Since $c=\int h$, the first assertion follows.

Next use $f(t)f(s)W(t,s)$. The interior integral is $2I+c^2$,
and the upper boundaries sum to $2H/N$. Thus
$2I+c^2=2H/N$, proving the second assertion. Finally, integrate the
kernel differential identity in $s$:
\begin{equation}
 w'+h=b/N,\qquad w(0)=0,\qquad w(1)=1/(2N^2).
 \label{shr:h5:eq:w-ode}
\end{equation}
Consequently
\[
 J_0=\frac{\ip b w}{N}-\int_0^1w'w
     =\frac HN-\frac{w(1)^2}{2}
     =\frac HN-\frac1{8N^4}.
\]
Every nonzero outer endpoint has been included in these equalities.

\subsection{The two traces}
For rank-one operators, $\tr(r\otimes s)=\ip s r$ and
$\tr((r\otimes s)(u\otimes v))=\ip s u\ip v r$.
Using these rules in~\eqref{shr:h5:eq:rank2} and~\eqref{shr:h5:eq:basic-integrals}
gives
\begin{equation}
 T_1=-4\tr W+2c=-\frac1{(2n-1)N^2}.
 \label{shr:h5:eq:T1}
\end{equation}
The full trace-square expansion, before any simplification, is
\begin{equation}
 T_2=16U-8\ip e w-16I+4c^2+\frac{4\ip e f}{N}.
 \label{shr:h5:eq:T2-expand}
\end{equation}
The square of $\one\otimes e$ has trace zero since $\int e=0$;
the two cross terms between the rank-one summands produce the last term.
Now
\[
 \ip e w=2J_0-H/N,\qquad \ip e f=2H-c/N.
\]
Substitution of~\eqref{shr:h5:eq:scalar-identities} yields
\begin{equation}
 T_2=16U-\frac{16H}{N}+\frac3{N^4}.
 \label{shr:h5:eq:T2}
\end{equation}

\subsection{The open path of length four}
The first two open-chain integrals follow from
$B\one=-4w+2f/N$:
\begin{equation}
 q_0=1/N,\qquad q_1=-4H+1/N^3.
 \label{shr:h5:eq:q01}
\end{equation}
One more application of~\eqref{shr:h5:eq:rank2} gives
\begin{equation}
 q_2=16L-\frac{8I}{N}-\frac{4\ip e w}{N}
            +\frac{2\ip e f}{N^2}+2cq_1,
 \quad L=\ip b{Ww}.
 \label{shr:h5:eq:q2-expand}
\end{equation}
It remains to show $L=G$, without dropping any boundary contribution.
Multiply~\eqref{shr:h5:eq:W-identities} by $f(s)$ and integrate in $s$.
The $s$ derivative has upper flux $h/N$ and lower flux zero, so
\begin{equation}
 Wb=z'-cb+h/N.
 \label{shr:h5:eq:Wb}
\end{equation}
Symmetry of $W$ and a final integration by parts give
\begin{align}
 L&=\ip w{Wb}=\ip w{z'}-cH+J_0/N\\
  &=[wz]_0^1-\ip{w'}z-cH+J_0/N\\
  &=\frac{H}{2N^2}-\frac IN+G-cH+\frac{J_0}{N}=G.
 \label{shr:h5:eq:L-equals-G}
\end{align}
Here $z(1)=H$, $w(1)=1/(2N^2)$, $z(0)=w(0)=0$,
$w'=b/N-h$, $c=1/(2N^2)$, and $I=J_0$.
Thus the new open-chain term is exactly
\begin{equation}
 q_2=16G-\frac{12H}{N^2}+\frac2{N^5}.
 \label{shr:h5:eq:q2}
\end{equation}
This identity is an equality of the full integrals. It is not a leading
term approximation to the path contribution.

\subsection{Eliminating the remaining trace square}
Symmetry of $W$ lets us integrate its square over twice the lower
triangle. Put $s=rt$, or equivalently parametrize the larger coordinate
by $r$ and the smaller by $rt$. Homogeneity gives
\begin{equation}
 U=2\int_0^1r^{4n-1}\dd r\int_0^1h(t)^2\dd t
   =\frac1{2n}\int_0^1h(t)^2\dd t.
 \label{shr:h5:eq:U-h2}
\end{equation}
The factor two includes both triangles; their diagonal has measure zero.
The boundary function satisfies
\begin{equation}
 (1-t)h'+(2n-1)h=\frac{t^{n-1}}{((n-1)!)^2}.
 \label{shr:h5:eq:h-ode}
\end{equation}
A direct derivation, including the width-one case, is given in the next
section. Multiply~\eqref{shr:h5:eq:h-ode} first by $2h$ and then by $t^n$,
and integrate by parts. In both fluxes the upper endpoint is killed by
$1-t$ and the lower endpoint by $h(0)=0$. One obtains
\begin{align}
 (4n-1)\int h^2&=\frac2{((n-1)!)^2}\int t^{n-1}h,\\
 \int t^{n-1}h&=3NH-\frac1{2N^2}.
\end{align}
Combining these with~\eqref{shr:h5:eq:U-h2} proves
\begin{equation}
 U=\frac{3nH}{(4n-1)N}-\frac{n}{2(4n-1)N^4}.
 \label{shr:h5:eq:U-final}
\end{equation}
Equations~\eqref{shr:h5:eq:T1}, \eqref{shr:h5:eq:T2}, \eqref{shr:h5:eq:q01},
\eqref{shr:h5:eq:q2}, and~\eqref{shr:h5:eq:U-final} now express every term in the
Pfaffian volume using $H,G$ and rational functions of $n,N$.

\section{Polynomial periods and their contiguous functions}
\label{shr:h5:sec:periods}
Define the normalized polynomial functions and their periods by
\begin{equation}
 y_n=N^2h,\quad x_n=N^3z,\quad
 Z_n=\int_0^1t^ny_n(t)\dd t=N^3H,\quad
 X_n=\int_0^1x_n(t)y_n(t)\dd t=N^5G.
 \label{shr:h5:eq:periods}
\end{equation}
These definitions make sense for every $n\ge1$, directly through
$W=V^n(V^*)^n$, regardless of the transfer formula's narrower domain.

\begin{proposition}[The exact full-count representation]\label{shr:h5:prop:count}
For every integer $n\ge1$,
\begin{equation}
 F_n:=\frac{a_n}{M_n}=16X_n+k(n)Z_n+\ell(n),
 \label{shr:h5:eq:full-count}
\end{equation}
where
\begin{equation}
 k(n)=-\frac{2(44n^2-26n+3)}{(2n-1)(4n-1)},\quad
 \ell(n)=\frac{224n^3-232n^2+76n-7}{8(2n-1)^2(4n-1)}.
 \label{shr:h5:eq:kell}
\end{equation}
\end{proposition}
\begin{proof}
For $n\ge2$, substitute all the reductions of Section~\ref{shr:h5:sec:rank}
into~\eqref{shr:h5:eq:pf-volume} and multiply by $N^5$. Before eliminating $U$
the result is
\[
 16N^5G-4N^4U-
 \left(8+\frac2{2n-1}\right)N^3H
 +\frac54+\frac1{2(2n-1)}+\frac1{8(2n-1)^2}.
\]
Using~\eqref{shr:h5:eq:U-final} gives precisely~\eqref{shr:h5:eq:full-count}.
At $n=1$, direct integration of $W(t,s)=\min(t,s)$ gives
\begin{equation}
 y_1(t)=t,\qquad x_1(t)=t/2-t^3/6,\qquad
 Z_1=1/3,\qquad X_1=2/15.
 \label{shr:h5:eq:width-one}
\end{equation}
The right side of~\eqref{shr:h5:eq:full-count} is then $1/120=a_1/M_1$,
since $a_1=1$ and $M_1=120$. This proves the stated range.
\end{proof}

The integral for $Z_n$ follows by inserting the kernel definition:
\[
 Z_n=n^2\int_0^1\int_0^t
       t^n(t-r)^{n-1}(1-r)^{n-1}\dd r\dd t.
\]
The change $u=t$, $v=1-r$ gives exactly~\eqref{shr:h5:eq:triangle}, including
its factor $u$ and the triangular domain.

\subsection{Differential identities with all boundary data}
\begin{lemma}\label{shr:h5:lem:odes}
For every $n\ge1$ the functions $x_n,y_n$ are polynomials and satisfy
\begin{align}
 (1-t)y_n'+(2n-1)y_n&=n^2t^{n-1},
 \label{shr:h5:eq:y-ode}\\
 tx_n'-3nx_n&=-y_n,
 \label{shr:h5:eq:x-ode}\\
 x_n(0)=y_n(0)&=0,\qquad x_n(1)=Z_n.
 \label{shr:h5:eq:endpoints}
\end{align}
Also $y_n(1)=n^2/(2n-1)$.
\end{lemma}
\begin{proof}
The kernel formula gives
\[
 y_n(t)=n^2\int_0^t(t-r)^{n-1}(1-r)^{n-1}\dd r.
\]
For $n\ge2$, put
$f_0=(t-r)^{n-1}(1-r)^{n-1}$ and
$g_0=(t-r)^{n-1}(1-r)^n$. Direct differentiation yields
\[
 (1-t)\partial_tf_0+(2n-1)f_0=-\partial_rg_0.
\]
The moving-endpoint contribution $f_0(t,t)$ vanishes, and integrating
the right side gives $g_0(t,0)-g_0(t,t)=t^{n-1}$.
This proves~\eqref{shr:h5:eq:y-ode} for $n\ge2$ and also proves
\eqref{shr:h5:eq:h-ode}. At $n=1$, the moving-endpoint contribution does
\emph{not} vanish: $f_0=1$, $g_0=1-r$. It contributes $1-t$ to
$(1-t)y_1'$, while the integrated $-g_0'$ contributes $t$. Their sum
is one. Equivalently, $y_1=t$ in~\eqref{shr:h5:eq:width-one} verifies the
identity directly. No width-one endpoint is discarded.

Euler homogeneity in~\eqref{shr:h5:eq:W-identities} says
$tW_t+sW_s=(2n-1)W$. Since
$x_n=N^2\int_0^1s^nW(t,s)\dd s$, integration by parts gives
\begin{align*}
 tx_n'&=(2n-1)x_n-N^2\int_0^1s^{n+1}W_s(t,s)\dd s\\
      &=3nx_n-N^2W(t,1)=3nx_n-y_n.
\end{align*}
The upper endpoint is exactly $-y_n$; the lower one is zero. Splitting
at $s=t$ is legitimate even at $n=1$, and the two internal boundary
terms cancel by continuity of $W$. The endpoint identities follow from
$W(t,0)=0$, symmetry, and the definitions; the value of $y_n(1)$ is a
direct integral. Polynomiality also follows from the next explicit form.
\end{proof}

\subsection{The homogeneous coefficient is fixed by the kernel}
For $s\ge t$, beta integration after expanding $(s-r)^{n-1}$ gives
\begin{equation}
 W(t,s)=\sum_{j=0}^{n-1}
 \frac{(-1)^jt^{n+j}s^{n-1-j}}{(n-1-j)!(n+j)!}.
 \label{shr:h5:eq:W-poly}
\end{equation}
In particular,
\begin{equation}
 y_n(t)=N^2\sum_{j=0}^{n-1}
 \frac{(-1)^jt^{n+j}}{(n-1-j)!(n+j)!}.
 \label{shr:h5:eq:y-poly}
\end{equation}
The first-order equation~\eqref{shr:h5:eq:x-ode} alone allows an arbitrary
multiple of $t^{3n}$. The defining kernel fixes it. Indeed
$z=V^n(V^*)^nf$ implies
$z^{(2n)}=(-1)^nf$, by differentiating the two iterated integral operators.
Its leading coefficient is $(-1)^n/(3n)!$. Consequently, writing
\begin{equation}
 c_n^*=\frac{(-1)^n(n!)^3}{(3n)!},
 \label{shr:h5:eq:cnstar}
\end{equation}
we have the exact polynomial
\begin{equation}
 x_n(t)=\sum_{j=n}^{2n-1}\frac{[t^j]y_n(t)}{3n-j}\,t^j
               +c_n^*t^{3n}.
 \label{shr:h5:eq:x-poly}
\end{equation}
Here $[t^j]$ means coefficient extraction. Alternatively,
\eqref{shr:h5:eq:W-poly} on the two halves of the square gives $x_n$ by
ordinary polynomial integration and obtains the same highest coefficient.
Thus~\eqref{shr:h5:eq:x-poly} is not an imposed choice of homogeneous solution.

\subsection{Explicit polynomial contiguity}
Set
\begin{equation}
 \lambda=-\frac{(n+1)^2}{2n(2n+1)},\qquad
 \mu=\frac{(n+1)^2}{2(2n+1)},\qquad
 \nu=-\frac{(n+1)^2}{3(3n+1)(3n+2)}.
 \label{shr:h5:eq:lambda-mu-nu}
\end{equation}
In this subsection coefficients are evaluated at the current $n$. Define
$q(t)=q_0+q_1t+q_2t^2$ and $r(t)=r_0+r_1t+r_2t^2$ by
\begin{align}
 q_0&=-\frac{n+1}{6n(2n+1)},&
 q_1&=\frac{(n+1)(5n+3)}{6n(2n+1)(3n+2)},&
 q_2&=-\frac{n+1}{3(3n+1)(3n+2)},\\
 r_0&=\frac{n+1}{6(2n+1)},&
 r_1&=\frac{7n^2+9n+3}{6(2n+1)(3n+2)},&
 r_2&=-\frac{n^2(n+1)}{3(2n+1)(3n+1)(3n+2)}.
 \label{shr:h5:eq:qr}
\end{align}
These $q_j$ are coefficients of a polynomial, no longer the open-chain
integrals of Section~\ref{shr:h5:sec:transfer}.

\begin{proposition}\label{shr:h5:prop:contiguity}
For every $n\ge1$, identically in $t$,
\begin{align}
 y_{n+1}&=\lambda(1-t)^2y_n+\mu t^n(1+t),
 \label{shr:h5:eq:y-contig}\\
 x_{n+1}&=\nu t^3x_n+(1-t)q(t)y_n+t^nr(t).
 \label{shr:h5:eq:x-contig}
\end{align}
\end{proposition}
\begin{proof}
Substitute the right side of~\eqref{shr:h5:eq:y-contig} into the new
$y$ equation. After dividing its residual by $t^{n-1}$, the coefficient
condition is the elementary polynomial identity
\begin{multline}
 \lambda n^2(1-t)^2+\mu\{n(1-t)(1+t)+t(1-t)\\
 +(2n+1)t(1+t)\}-(n+1)^2t=0.
 \label{shr:h5:eq:y-certificate}
\end{multline}
Its coefficients vanish for~\eqref{shr:h5:eq:lambda-mu-nu}. Both the actual
and candidate $y_{n+1}$ vanish at zero. Their difference satisfies
$(1-t)d'+(2n+1)d=0$, hence is a constant multiple of
$(1-t)^{2n+1}$ on the interior; the value at zero forces that constant
to vanish. Polynomial continuity covers $t=1$.

For the second identity put $Q_0(t)=(1-t)q(t)$. Substitution into the
new $x$ equation reduces, using the two current ODEs, to
\begin{align}
 -\nu t^3+tQ_0'-3(n+1)Q_0-(2n-1)tq
       +\lambda(1-t)^2&=0,\\
 tr'-(2n+3)r+n^2q+\mu(1+t)&=0.
 \label{shr:h5:eq:x-certificate}
\end{align}
These are degree-three and degree-two polynomial identities in $t$,
verified by inserting~\eqref{shr:h5:eq:qr}; they constitute a finite rational
certificate for the substitution. The difference between the candidate
and actual $x_{n+1}$ solves $td'-3(n+1)d=0$, so it is a multiple of
$t^{3n+3}$. Its highest coefficient vanishes because
$c_{n+1}^*/c_n^*=\nu$. All other terms on the proposed right side
have degree at most $2n+2<3n+3$, including when $n=1$.
Thus the homogeneous multiple is zero.
\end{proof}
The next section provides a second check of this normalization at $t=1$:
it derives $Z_{n+1}=\nu Z_n+\delta_n$ from $y$ alone and finds
$r(1)=\delta_n$. Hence the candidate also has the required endpoint
$x_{n+1}(1)=Z_{n+1}$ independently of its leading coefficient.

\section{Finite moment closure and the full first difference}
\label{shr:h5:sec:moments}
The contiguous functions contain only fixed-degree polynomial factors.
Consequently, their product can be integrated through a finite family of
moments. This section derives each reduction with its flux, so that no
endpoint convention is hidden in the symbolic calculation.

\subsection{Four families and four flux identities}
At a fixed positive $n$, abbreviate $x=x_n$, $y=y_n$, $X=X_n$,
$Z=Z_n$, and define
\begin{equation}
 V_j=\int_0^1t^{n+j}y\dd t,\quad
 Y_j=\int_0^1t^jy^2\dd t,\quad
 U_j=\int_0^1t^{n+j}x\dd t,\quad
 J_j=\int_0^1t^jxy\dd t.
 \label{shr:h5:eq:moment-defs}
\end{equation}
We need nonnegative $j$, together with $V_{-1}$; its weight has exponent
$n-1\ge0$, even at $n=1$. The use of $V_j$ for scalar moments is
separate from the Volterra operator $V$.

The differential equations of Lemma~\ref{shr:h5:lem:odes} give the following
polynomial flux identities:
\begin{align}
 [t^{n+j}(1-t)y]'&=(n+j)t^{n+j-1}y-(3n+j)t^{n+j}y
                         +n^2t^{2n+j-1},
 \label{shr:h5:eq:V-flux}\\
 [t^j(1-t)y^2]'&=jt^{j-1}y^2-(4n+j-1)t^jy^2
                         +2n^2t^{n+j-1}y,
 \label{shr:h5:eq:Y-flux}\\
 [t^{n+j+1}x]'&=(4n+j+1)t^{n+j}x-t^{n+j}y,
 \label{shr:h5:eq:U-flux}\\
 [t^{j+1}(1-t)xy]'&=(3n+j+1)t^jxy-(5n+j+1)t^{j+1}xy\\
     &\hspace{1em}-(1-t)t^jy^2+n^2t^{n+j}x.
 \label{shr:h5:eq:J-flux}
\end{align}
In~\eqref{shr:h5:eq:Y-flux} the $j=0$ first term is absent; it is not an
undefined product at zero. The product identity underlying the last line
is worth recording separately:
\begin{equation}
 t(1-t)(xy)'=[3n-(5n-1)t]xy+n^2t^nx-(1-t)y^2.
 \label{shr:h5:eq:xy-ode}
\end{equation}
The final factor $1-t$ cannot be omitted.

Here are every upper and lower flux in the four integrations:
\begin{itemize}
\item For~\eqref{shr:h5:eq:V-flux}, $t^{n+j}(1-t)y$ is zero at one by
$1-t$, and at zero by $y(0)=0$ and $n+j\ge1$.
\item For~\eqref{shr:h5:eq:Y-flux}, $t^j(1-t)y^2$ is zero at one by
$1-t$, and at zero by $y(0)=0$, including $j=0$.
\item For~\eqref{shr:h5:eq:U-flux}, the upper flux is $x(1)=Z$.
The lower flux is zero because $t^{n+j+1}x(t)$ vanishes there.
\item For~\eqref{shr:h5:eq:J-flux}, both fluxes are zero because
$t^{j+1}(1-t)xy$ contains both endpoint factors and $x,y$ are bounded.
\end{itemize}
These are polynomial equalities on the closed interval, also for $n=1$.
Dividing the ODEs by $t$ or $1-t$ on the interior has not suppressed any
nonzero boundary value.

Integration therefore yields the finite closure
\begin{gather}
 V_0=Z,\qquad V_{-1}=3Z-\tfrac12,\qquad
 V_j=\frac{(n+j)V_{j-1}+n^2/(2n+j)}{3n+j}\quad(j\ge1),
 \label{shr:h5:eq:V-rec}\\
 Y_0=\frac{2n^2V_{-1}}{4n-1},\qquad
 Y_j=\frac{jY_{j-1}+2n^2V_{j-1}}{4n+j-1}\quad(j\ge1),
 \label{shr:h5:eq:Y-rec}\\
 U_j=\frac{Z+V_j}{4n+j+1}\quad(j\ge0),\qquad J_0=X,\\
 J_{j+1}=\frac{(3n+j+1)J_j+n^2U_j-Y_j+Y_{j+1}}{5n+j+1}
     \quad(j\ge0).
 \label{shr:h5:eq:J-rec}
\end{gather}
The value $V_{-1}=3Z-1/2$ is simply~\eqref{shr:h5:eq:V-flux} with $j=0$,
integrated and divided by $n$. This division explains again why positive
width is the natural domain.

\subsection{A short certificate for both shifted periods}
For a polynomial $p(t)=\sum p_jt^j$, write
$[p]_V=\sum p_jV_j$ and similarly for $Y,U,J$.
Multiply~\eqref{shr:h5:eq:y-contig} by~\eqref{shr:h5:eq:x-contig} and integrate.
Every term is covered by the four moment families:
\begin{align}
 X_{n+1}={}&[\lambda\nu t^3(1-t)^2]_J
       +[\lambda(1-t)^3q]_Y\\
 &+[\lambda(1-t)^2r+\mu(1+t)(1-t)q]_V\\
 &+[\mu\nu t^3(1+t)]_U
       +\sum_j\frac{[t^j]\{\mu(1+t)r\}}{2n+j+1}.
 \label{shr:h5:eq:X-certificate}
\end{align}
The largest $J$ and $Y$ indices are five; the largest $V$ and $U$
indices are four. Thus finitely many applications of
\eqref{shr:h5:eq:V-rec}--\eqref{shr:h5:eq:J-rec} evaluate the whole expression
as an affine form in the independent symbols $X,Z$.

The simpler calculation from $y$ alone is
\begin{equation}
 Z_{n+1}=\lambda(V_1-2V_2+V_3)
       +\mu\left(\frac1{2n+2}+\frac1{2n+3}\right).
 \label{shr:h5:eq:Z-certificate}
\end{equation}
Its $Z$ coefficient and constant term simplify to exactly $\nu_n$ and
$\delta_n$ in~\eqref{shr:h5:eq:Z-shift}. Also $r(1)=\delta_n$, so
this establishes the independent endpoint check promised after
Proposition~\ref{shr:h5:prop:contiguity}.

For explicit reproducibility, the full $X$ transport is
\begin{equation}
 X_{n+1}=\rho_nX_n+\sigma_nZ_n+\tau_n,
 \label{shr:h5:eq:X-shift}
\end{equation}
where, with $A_5(n)=\prod_{j=1}^4(5n+j)$,
\begin{align}
 \rho_n&=\frac{(n+1)^4}{5A_5(n)}=\frac{M_n}{M_{n+1}},
 \label{shr:h5:eq:rho}\\
 \sigma_n&=-\frac{(n+1)^2P_7(n)}
 {30(2n+1)(3n+1)(3n+2)(4n+1)(4n+3)A_5(n)},\\
 \tau_n&=\frac{P_9(n)}
 {120(2n+1)^2(3n+1)(3n+2)(4n+3)A_5(n)},
 \label{shr:h5:eq:sigma-tau}
\end{align}
with
\begin{align}
 P_7(n)={}&138688n^7+507872n^6+772828n^5+630758n^4\\
         &+296881n^3+80204n^2+11463n+666,\\
 P_9(n)={}&744128n^9+4138176n^8+10020608n^7+13828392n^6\\
         &+11948992n^5+6682909n^4+2410801n^3\\
         &+538853n^2+67431n+3582.
 \label{shr:h5:eq:P7P9}
\end{align}
Equations~\eqref{shr:h5:eq:X-certificate} and~\eqref{shr:h5:eq:Z-certificate},
with the explicit moment recursions, are small rational certificates for
these somewhat larger coefficients. Checking them means expanding
fixed-degree polynomials and testing zero numerators in $\Q(n)$.
No evaluation at finitely many widths, guessed recurrence, or polynomial
interpolation is involved. The supplied standard-library verifier performs
exactly those operations.

\subsection{Eliminating the second period from the first difference}
Substitute~\eqref{shr:h5:eq:X-shift} and~\eqref{shr:h5:eq:Z-shift} into the
full-count representation~\eqref{shr:h5:eq:full-count} at $n+1$.
After subtracting $\rho_nF_n$, the $X_n$ coefficient vanishes because it
is $16\rho_n-16\rho_n$. The other two coefficients are
\begin{align}
 [Z_n](F_{n+1}-\rho_nF_n)
   &=\frac{2(n+1)^2P(n)}
 {15(2n-1)(2n+1)(3n+1)(3n+2)}\\
 &\hspace{1em}\cdot
 \frac1{(4n-1)(4n+1)(4n+3)A_5(n)},
 \label{shr:h5:eq:diff-Z}\\
 [1](F_{n+1}-\rho_nF_n)
   &=-\frac{D(n)}
 {120(2n-1)^2(2n+1)^2(3n+1)(3n+2)}\\
 &\hspace{1em}\cdot\frac1{(4n-1)(4n+3)A_5(n)}.
 \label{shr:h5:eq:diff-constant}
\end{align}
Here $[1]$ denotes the coefficient independent of $X_n,Z_n$, rather than
substitution of their actual values. Direct cancellation with
\eqref{shr:h5:eq:u}, \eqref{shr:h5:eq:gamma}, and~\eqref{shr:h5:eq:rho} identifies these
as $2\rho_nu(n)$ and $-2\rho_nu(n)\gamma_n$. We have proved
\begin{equation}
 F_{n+1}-\rho_nF_n=2\rho_nu(n)(Z_n-\gamma_n).
 \label{shr:h5:eq:F-first-difference}
\end{equation}
Multiplying by $M_{n+1}$ proves~\eqref{shr:h5:eq:target-diff}, including at
$n=1$. This completes the derivation from the literal arrays to the
smaller exact period system announced in Section~\ref{shr:h5:sec:scope}.

As a width-one check independent of this simplification, direct
polynomials give
\[
 y_2=2t^2-2t^3/3,\qquad
 x_2=t^2/2-2t^3/9+t^6/90,
\]
\[
 Z_2=13/45,\qquad X_2=1322/14175,
 \qquad F_2=1/113400.
\]
Since $M_2=113400$, this is $a_2=1$, as the interlacing total order
requires. These values are checks of endpoints and normalization, not
premises of the all-width certificate.

\section{The printed nested sum and the exact OEIS operator}
\label{shr:h5:sec:sum-operator}
The previous sections proved the full-count identities directly. We now
perform the finite algebra that identifies them with the specified source
statements, including their summation constants and the orientation of
the shift operator.

\subsection{Variation of constants gives the actual finite sum}
Recall $c_n^*=(-1)^n(n!)^3/(3n)!$ and $c_{n+1}^*=\nu_nc_n^*$.
The following rational identity links the period's forcing to the printed
summand:
\begin{equation}
 \delta_n+\nu_n\gamma_n-\gamma_{n+1}
   =-\frac{27}{8}\nu_nv(n)
   =\frac{9(n+1)^2v(n)}{8(3n+1)(3n+2)}.
 \label{shr:h5:eq:gamma-bridge}
\end{equation}
Every quantity is explicit in~\eqref{shr:h5:eq:PQD}--\eqref{shr:h5:eq:gamma} and
\eqref{shr:h5:eq:Z-shift}. Multiplication by the displayed denominators reduces
\eqref{shr:h5:eq:gamma-bridge} to a zero polynomial. This is the last rational
certificate required for the finite sum.

Put
\begin{equation}
 H_n=(-1)^nu(n)\frac{(5n)!}{(3n)!(n!)^2}=M_nu(n)c_n^*,\qquad
 I_n=(-1)^nv(n)\frac{(3n)!}{(n!)^3}=\frac{v(n)}{c_n^*}.
 \label{shr:h5:eq:HI}
\end{equation}
From~\eqref{shr:h5:eq:Z-shift} and~\eqref{shr:h5:eq:gamma-bridge},
\begin{equation}
 \frac{Z_{n+1}-\gamma_{n+1}}{c_{n+1}^*}
   -\frac{Z_n-\gamma_n}{c_n^*}=-\frac{27}{8}I_n.
 \label{shr:h5:eq:normalized-Z-diff}
\end{equation}
The constant of summation is determined, rather than chosen: direct
substitution gives $\gamma_1=Z_1=1/3$. Hence
\[
 \frac{Z_k-\gamma_k}{c_k^*}=-\frac{27}{8}\sum_{i=1}^{k-1}I_i.
\]
Inserting this in the proved first difference~\eqref{shr:h5:eq:target-diff}
yields
\begin{equation}
 a_{k+1}-a_k=-\frac{27}{4}H_k\sum_{i=1}^{k-1}I_i.
 \label{shr:h5:eq:sum-increment}
\end{equation}
Summation over $k=1,\ldots,n-1$ with $a_1=1$ gives
\eqref{shr:h5:eq:main-sum}, proving Theorem~\ref{shr:h5:thm:sum}.
This argument matches both nested summation limits in the printed source;
it does not merely construct another solution of a recurrence.

The empty sums give $a_1=a_2=1$. Exact substitution gives
\begin{equation}
 u(1)=4/5,\quad v(1)=5/8208,\quad
 H_1=-16,\quad H_2=608,\quad I_1=-5/1368.
 \label{shr:h5:eq:initial-constants}
\end{equation}
Therefore $a_3=1-(27/4)\,608\,(-5/1368)=16$.

\subsection{The rational factorization and shift order}
Let $E$ denote the forward shift. Coefficients and shifts obey
$Ef(n)=f(n+1)E$. Define
\begin{equation}
 R(n)=\frac{H_{n+1}}{H_n},\qquad
 S(n)=\frac{H_{n+2}I_{n+1}}{H_{n+1}I_n}.
 \label{shr:h5:eq:RS-ratios}
\end{equation}
We use $S$ rather than $U$ for the second ratio to distinguish it from the
kernel trace and the moment family. Explicit cancellation of the
factorials gives
\begin{align}
 R(n)={}&-\frac{5(2n-1)(4n-1)(4n+1)A_5(n)P(n+1)}
 {3(n+2)^2(2n+3)(3n+4)(3n+5)(4n+5)(4n+7)P(n)},
 \label{shr:h5:eq:R}\\
 S(n)={}&\frac{5(2n-1)(2n+1)(3n+4)(3n+5)(4n+5)(4n+7)}
 {(n+3)^4(2n+5)^2(3n+7)(3n+8)(4n+9)(4n+11)}\\
 &\hspace{1em}\cdot
 \frac{\prod_{j=6}^9(5n+j)P(n)Q(n+1)}{P(n+1)Q(n)}.
 \label{shr:h5:eq:S}
\end{align}
The rightmost operator acts first. Equation~\eqref{shr:h5:eq:sum-increment}
therefore implies
\begin{equation}
 (E-R(n))(E-1)a_n=-\frac{27}{4}H_{n+1}I_n.
 \label{shr:h5:eq:inhom-order2}
\end{equation}
The ratio of the forcing at $n+1$ to that at $n$ is $S(n)$. Thus
\begin{equation}
 (E-S(n))(E-R(n))(E-1)a_n=0.
 \label{shr:h5:eq:factored-annihilator}
\end{equation}
For a direct check of Ore multiplication, the coefficients in increasing
powers of $E$ are
\begin{equation}
 -SR,\quad R(n+1)+S(1+R),\quad
 -(1+R(n+1)+S),\quad1.
 \label{shr:h5:eq:monic-coeffs}
\end{equation}
Replacing $R(n+1)$ by $R(n)$ would be an incorrect commutative
multiplication; the shift in this formula is essential.

Define
\begin{multline}
 c_3(n)=3(n+2)^2(2n+3)(2n+5)^2(3n+7)(3n+8)\\
 {}\times(4n+9)(4n+11)Q(n)(n+3)^4.
 \label{shr:h5:eq:c3}
\end{multline}
The exact polynomial identities are
\begin{align}
 c_0&=-c_3SR,\\
 c_1&=c_3\{R(n+1)+S(1+R)\},\\
 c_2&=-c_3\{1+R(n+1)+S\},
 \label{shr:h5:eq:literal-identities}
\end{align}
with the literal coefficient polynomials in Appendix~\ref{shr:h5:app:literal}.
They are verified coefficientwise over $\Q[n]$ after clearing
explicit denominators. That appendix gives the full independent source
transcription, not coefficients defined circularly by these identities.
It follows that
\begin{equation}
 L_{\rm OEIS}=c_3(n)(E-S(n))(E-R(n))(E-1).
 \label{shr:h5:eq:operator-factorization}
\end{equation}
Together with~\eqref{shr:h5:eq:factored-annihilator} this proves
Theorem~\ref{shr:h5:thm:operator}.

There is also a useful explicit forcing version of
\eqref{shr:h5:eq:inhom-order2}:
\begin{equation}
 a_{n+2}-(1+R(n))a_{n+1}+R(n)a_n=M_n\Phi(n),
 \label{shr:h5:eq:order2-forcing}
\end{equation}
where
\begin{align}
 \Phi(n)={}&\frac{15A_5(n)Q(n)}
 {4(n+1)^4(n+2)^4(2n-1)(2n+1)^2(2n+3)^2}\\
 &\hspace{1em}\cdot
 \frac1{(3n+4)(3n+5)(4n+5)(4n+7)P(n)}.
 \label{shr:h5:eq:Phi}
\end{align}
This is obtained directly by cancelling
$-(27/4)H_{n+1}I_n/M_n$. Thus the forcing is fixed by the summands
and its ratio is fixed by~\eqref{shr:h5:eq:RS-ratios}.

\subsection{Positivity and uniqueness on the correct range}
All affine denominators in the rank, differential, and moment formulas
are positive for $n\ge1$. The only further possible zeros occur in
$P(n)$ or $Q(n)$. For $m=n-1\ge0$, the coefficients of $P(m+1)$,
in ascending order, are
\[
 (60480,353502,971955,1683047,1973632,1543976,760768,211616,25216).
\]
Those of $Q(m+1)$ are
\begin{align*}
 (&1077753600,23068181160,128320688700,354332674386,\\
  &588363536007,638861237719,471293180347,238738633847,\\
  &81945774178,18238806776,2377384288,137855872).
\end{align*}
Every coefficient is strictly positive, so both polynomials are positive
for every real $n\ge1$. In particular, $u(n)$, $v(n)$, the ratios, and
all eliminations are defined there, while $c_3(n)>0$. The factorial
$c_n^*$ never vanishes. This proves the denominator assertion in
Theorem~\ref{shr:h5:thm:sum} and the uniqueness assertion in
Theorem~\ref{shr:h5:thm:operator}.

The first terms, with index starting at one, are
\[
 1,\ 1,\ 16,\ 985,\ 141696,\ 36372976,\ 14083834704,
 \ 7372392431849.
\]
These are useful checks of orientation and indexing. The proof is the
uniform integral and rational derivation, not agreement of a finite
prefix. Finally, literal substitution at $n=0$ with the empty-array
values $(a_0,a_1,a_2,a_3)=(1,1,1,16)$ gives the nonzero residual stated
in Section~\ref{shr:h5:sec:scope}; no uniqueness claim from that starting index
is made.

\section{Recurrence-based asymptotics with effective errors}
\label{shr:h5:asy:section}
The leading constant and the existence of all fixed inverse-power orders
were already proved in the earlier all-height report (Part~\ref{shr:part:one}).
This section recovers those results from the exact recurrence and supplies
finite rational residual certificates. It does not claim priority for the
asymptotic expansion.

\subsection{A triangular stable system}
Write $b_n=F_n=a_n/M_n$, where $M_n=(5n)!/(n!)^5$. The proved identities give
\begin{align}
 Z_{n+1}&=\nu_nZ_n+\delta_n,\notag\\
 b_{n+1}&=\rho_nb_n+2\rho_nu(n)(Z_n-\gamma_n),&
 \rho_n&=\frac{(n+1)^4}{5(5n+1)(5n+2)(5n+3)(5n+4)},
 \label{shr:h5:asy:system}
\end{align}
with $Z_1=1/3$ and $b_1=1/120$. Here $b_n$ is a scalar normalized count,
distinct from any function $b(t)$ used in the kernel calculation.
For $n\ge1$,
\begin{equation}
 |\nu_n|\le\frac4{27}<\frac16,\qquad
 0<\rho_n\le\frac{16}{3125}<\frac1{100}.
 \label{shr:h5:asy:contractions}
\end{equation}
These bounds follow by replacing $n+1$ by $2n$ in the numerators and
each positive denominator factor by its leading monomial. The coupling
$2\rho_nu(n)$ is $O(n^{-1})$, which is important for the error orders.

Put $t=1/n$, $f(t)=t/(1+t)$, and define rational functions
\[
 r(t)=\rho_{1/t},\quad v_*(t)=\nu_{1/t},\quad d_*(t)=\delta_{1/t},
 \quad \alpha(t)=2\rho_{1/t}u(1/t),\quad
 \eta(t)=-\alpha(t)\gamma_{1/t}.
\]
All extend analytically to zero; $\alpha(0)=\eta(0)=0$,
$r(0)=1/3125$, and $v_*(0)=-1/27$.
The two formal series $\mathcal Z(t)=\sum_{j\ge0}z_jt^j$ and
$\mathcal B(t)=\sum_{j\ge0}\beta_jt^j$ are specified by
\begin{align}
 \mathcal Z(f(t))&=v_*(t)\mathcal Z(t)+d_*(t),\notag\\
 \mathcal B(f(t))&=r(t)\mathcal B(t)+\alpha(t)\mathcal Z(t)+\eta(t).
 \label{shr:h5:asy:formal}
\end{align}
In the first equation the new coefficient $z_j$ has multiplier $28/27$;
in the second the new coefficient $\beta_j$ has multiplier $3124/3125$.
Moreover, $\alpha(0)=0$ means the second equation at order $j$ needs only
$z_0,\ldots,z_{j-1}$. Thus both series are uniquely determined over $\Q$
by finite coefficient arithmetic, to any prescribed degree.

For completeness, direct expansion gives
\[
 \mathcal Z(t)=\frac14+\frac{t}{16}+\frac{t^2}{32}
 +\frac{t^3}{128}-\frac{t^4}{128}-\frac{9t^5}{512}
 -\frac{21t^6}{1024}-\frac{41t^7}{4096}
 +\frac{29t^8}{1024}+\frac{925t^9}{8192}+O(t^{10}).
\]
Substitution of the displayed rational $\gamma_n$ shows
\[
 \mathcal Z(t)-\gamma_{1/t}
 =\frac{189783}{1613824}t^9+O(t^{10}),\qquad
 \alpha(t)=\frac{197}{84375}t+O(t^2).
\]
Consequently $\beta_0=\cdots=\beta_9=0$ and
$\beta_{10}=9/32768$. These are finite rational identities, not fitted
values. The companion standard-library script checks them coefficientwise.
Set $c_j=\beta_{10+j}$ and $c_0=9/32768$.

\begin{theorem}[Every fixed order with a rational certificate]
\label{shr:h5:asy:fixedorder}
For every integer $K\ge0$, one can compute an integer $N_K\ge2$ and
positive rational constants $E_{Z,K},E_{b,K}$ such that
\begin{align}
 \left|Z_n-\sum_{j=0}^{K+9}z_jn^{-j}\right|
 &\le E_{Z,K}n^{-K-10},\notag\\
 \left|b_n-n^{-10}\sum_{j=0}^Kc_jn^{-j}\right|
 &\le E_{b,K}n^{-K-11}\qquad(n\ge N_K).
 \label{shr:h5:asy:bounds}
\end{align}
In particular,
\begin{align}
 b_n=\frac9{32768}n^{-10}\bigg(&1+\frac{10}{n}
 +\frac{225}{4n^2}+\frac{1819}{8n^3}+\frac{22045}{32n^4}
 +\frac{1353}{n^5}\notag\\
 &-\frac{6103}{16n^6}+O(n^{-7})\bigg).
 \label{shr:h5:asy:bcoeffs}
\end{align}
\end{theorem}
\begin{proof}
Fix $K$, set $L=K+9$, $p_Z=K+10$, $p_b=K+11$, and let
\[
 q_Z(n)=\sum_{j=0}^Lz_jn^{-j},\qquad
 q_b(n)=n^{-10}\sum_{j=0}^Kc_jn^{-j}.
\]
Their exact rational residuals are
\begin{align*}
 D_Z(n)&=q_Z(n+1)-\nu_nq_Z(n)-\delta_n,\\
 D_b(n)&=q_b(n+1)-\rho_nq_b(n)
                   -2\rho_nu(n)(q_Z(n)-\gamma_n).
\end{align*}
Equation~\eqref{shr:h5:asy:formal} gives $D_Z=O(n^{-p_Z})$ and
$D_b=O(n^{-p_b})$. The latter order holds although $q_Z$ stops one degree
before $q_b$, because the coupling vanishes to first order at zero.

Here is an explicit finite procedure to turn these formal cancellations
into bounds. For each of
\[
 t^{-p_Z}D_Z(1/t),\qquad t^{-p_b}D_b(1/t),\qquad \alpha(t)/t,
\]
form an exact ratio $P(t)/H(t)$ of rational polynomials with $H(0)>0$.
Remove any common power of $t$ first. For an integer $N\ge2$, put
\begin{equation}
 h_N=H_0-\sum_{i\ge1}|H_i|N^{-i},\qquad
 \mathcal C_N(P,H)=\frac{\sum_i|P_i|N^{-i}}{h_N}.
 \label{shr:h5:asy:rationalbound}
\end{equation}
Increase $N$, for example by doubling, until all three $h_N$ are positive
and
\[
 \theta_Z=\left(\frac N{N+1}\right)^{p_Z}>\frac16,
 \qquad
 \theta_b=\left(\frac N{N+1}\right)^{p_b}>\frac1{100}.
\]
This search terminates because $H(0)>0$ and both powers tend to one.
Let $C_Z,C_b,A$ be the three respective values of
$\mathcal C_N$. Then for every $n\ge N$,
\[
 |D_Z(n)|\le C_Zn^{-p_Z},\qquad
 |D_b(n)|\le C_bn^{-p_b},\qquad
 |2\rho_nu(n)|\le A/n.
\]
Compute $Z_N,b_N$ exactly from~\eqref{shr:h5:asy:system} and define
\begin{align}
 E_Z&=\max\left\{N^{p_Z}|Z_N-q_Z(N)|,
                       \frac{C_Z}{\theta_Z-1/6}\right\},\notag\\
 E_b&=\max\left\{N^{p_b}|b_N-q_b(N)|,
                       \frac{C_b+AE_Z}{\theta_b-1/100}\right\}.
 \label{shr:h5:asy:constants}
\end{align}
If $E_Z=0$, replace it by any positive rational before defining $E_b$;
if $E_b=0$, replace it by any positive rational as well.

Writing $e_Z=Z-q_Z$ and $e_b=b-q_b$, the exact error equations give
\begin{align*}
 |e_Z(n+1)|&\le |e_Z(n)|/6+C_Zn^{-p_Z},\\
 |e_b(n+1)|&\le |e_b(n)|/100+(A/n)|e_Z(n)|+C_bn^{-p_b}.
\end{align*}
The first estimate propagates $|e_Z(n)|\le E_Zn^{-p_Z}$ because
$E_Z/6+C_Z\le E_Z\theta_Z$. The second propagates
$|e_b(n)|\le E_bn^{-p_b}$ because $p_Z+1=p_b$ and
$E_b/100+AE_Z+C_b\le E_b\theta_b$.
Finally $(n/(n+1))^{p}\ge(N/(N+1))^{p}$ for $n\ge N$.
The initial inequalities in~\eqref{shr:h5:asy:constants} start both inductions.
This proves the claim for the actual counting solution.
\end{proof}

The rational certificate at $K=6$ permits the particularly simple
inequality
\begin{equation}
 |b_n-q_6(n)|\le18n^{-17}\qquad(n\ge64),
 \label{shr:h5:asy:explicit6}
\end{equation}
where $q_6(n)$ is the displayed truncation through $n^{-16}$ in
\eqref{shr:h5:asy:bcoeffs}. This onset and constant are deliberately conservative.

\subsection{Restoring the multinomial factor}
Stirling's expansion and its positive-real remainder bound~\cite{DLMFgamma} give
\begin{align}
 \log M_n\sim{}&n\log3125-2\log n+\log\frac{\sqrt5}{4\pi^2}\notag\\
 &+\sum_{h\ge1}\frac{B_{2h}}{2h(2h-1)}
       \bigl(5^{-(2h-1)}-5\bigr)n^{-(2h-1)}.
 \label{shr:h5:asy:stirling}
\end{align}
Here and below infinite expansions are interpreted at each fixed order.
The first logarithmic correction is $-2/(5n)$. Multiplication with
Theorem~\ref{shr:h5:asy:fixedorder} yields the previously established leading term
and all its rational relative corrections:
\begin{equation}
 a_n=C_*\,3125^n n^{-12}
       \left(1+\sum_{j=1}^Kd_jn^{-j}+O_K(n^{-K-1})\right),
 \qquad C_*:=\frac{9\sqrt5}{2^{17}\pi^2}.
 \label{shr:h5:asy:count}
\end{equation}
For example,
\begin{align*}
 d_1&=\frac{48}{5},& d_2&=\frac{5233}{100},&
 d_3&=\frac{1028391}{5000},\\
 d_4&=\frac{60248377}{100000},&
 d_5&=\frac{273939939}{250000},&
 d_6&=-\frac{2709616559}{3125000}.
\end{align*}
All coefficients follow by exponentiating finite rational series. The first
five displayed normalized and count corrections already appear in the prior
report (Part~\ref{shr:part:one}); the sixth displayed correction is obtained here from
the current recurrence certificate, without a novelty claim.

\subsection{A finite logarithmic remainder}
For inversion we need a bound, not an unspecified $O$ term. Let
$\lambda=\log3125$, and let $\ell_j$ be the rational coefficients in
$\log(1+\sum_{j\ge1}d_jt^j)$. For example,
\begin{equation}
 \ell_1=\frac{48}{5},\qquad \ell_2=\frac{25}{4},\qquad
 \ell_3=-\frac{8889}{5000},\qquad \ell_4=-\frac{335}{8}.
 \label{shr:h5:asy:logcoeffs}
\end{equation}
For each fixed $K$ there are a computable rational $C_K>0$ and integer
$N_K^{\log}$ such that
\begin{equation}
 \left|\log a_n-\left(\lambda n-12\log n+\log C_*
                         +\sum_{j=1}^K\ell_jn^{-j}\right)\right|
 \le C_Kn^{-K-1}\qquad(n\ge N_K^{\log}).
 \label{shr:h5:asy:logerror}
\end{equation}
Here is an explicit construction. Retain $E=E_{b,K}$ from
\eqref{shr:h5:asy:bounds}, put
\[
 u_K(t)=\sum_{j=1}^K\frac{c_j}{c_0}t^j,
 \qquad U=\sum_{j=1}^K\left|\frac{c_j}{c_0}\right|N^{-(j-1)},
\]
and increase $N$ beyond $N_K$ until
$U/N\le1/4$ and $(E/c_0)N^{-K-1}\le1/4$.
Then the mean value theorem for $\log$ on $[1/2,\infty)$ gives
\[
 \left|\log\frac{n^{10}b_n}{c_0}-\log(1+u_K(1/n))\right|
 \le\frac{2E}{c_0}n^{-K-1}.
\]
Define the finite polynomial and coefficient bound
\[
 L_K(t)=\sum_{h=1}^K\frac{(-1)^{h+1}}h u_K(t)^h,
 \qquad V=\sum_{j>K}|[t^j]L_K(t)|N^{-(j-K-1)}.
\]
The scalar logarithm tail is at most
$4U^{K+1}n^{-K-1}/(3(K+1))$, while the terms of $L_K$ of degree above
$K$ contribute at most $Vn^{-K-1}$. Thus the normalized-count contribution
is bounded by $C^{(b)}n^{-K-1}$, where
\[
 C^{(b)}=2E/c_0+V+\frac{4U^{K+1}}{3(K+1)}.
\]
For $K=0$, use $u_K=L_K=0$ and $U=V=0$.

Set $J=\lfloor(K+1)/2\rfloor$ and $q=2J+1$. The Stirling remainder
following $J$ Bernoulli terms in $\log(n!)$ is at most
$|B_{2J+2}|/((2J+2)(2J+1)n^q)$. Applying this separately at $5n$ and at
$n$ gives the explicit additional constant
\[
 C^{(M)}=\frac{|B_{2J+2}|}{(2J+2)(2J+1)}
              (5^{-q}+5)N^{K+1-q}.
\]
Since $q\ge K+1$, this proves~\eqref{shr:h5:asy:logerror} with
$C_K=C^{(b)}+C^{(M)}$ and $N_K^{\log}=N$.
The exact companion certificates, with harmless upward rounding, give
\begin{equation}
 (K,N_K^{\log},C_K)=(0,512,202)
 \quad\hbox{and}\quad (6,64,6100000).
 \label{shr:h5:asy:logexamples}
\end{equation}
Every bound in this construction is rational; the constants $\pi$ and
$\sqrt5$ enter only the displayed leading factor, which is retained exactly.

\begin{remark}[Scope of the expansions]
The contraction proof fixes all algebraic coefficients and controls every
fixed truncation. It does not assert convergence of the infinite formal
series or select a canonical decomposition into exponentially small
sectors. The two homogeneous contraction ratios tend to $1/3125$ and
$-1/27$, but those ratios alone are not Stokes data. None of the finite
error or inverse conclusions requires such a decomposition.
\end{remark}

\section{Finite Lambert models and certified integer thresholds}
\label{shr:h5:inv:section}
The count exists at integer widths. We first invert an explicitly defined
finite real function and then compare it with the exact integer threshold.
This avoids differentiating an unspecified discrete error or rounding an
asymptotic approximation without a margin.

\subsection{A finite monotone model}
Fix $K\ge0$ and use~\eqref{shr:h5:asy:logerror}. For real $x>0$ set
\begin{equation}
 \Psi_K(x)=\lambda x-12\log x+\log C_*
                      +\sum_{j=1}^K\ell_jx^{-j}.
 \label{shr:h5:inv:model}
\end{equation}
Choose $X_K>0$ large enough that
\begin{equation}
 \frac{12}{X_K}+\sum_{j=1}^Kj|\ell_j|X_K^{-j-1}
 \le\frac\lambda2.
 \label{shr:h5:inv:derivative}
\end{equation}
Then $\Psi_K'(x)\ge\lambda/2$ for $x\ge X_K$.
For each $y\ge\exp\Psi_K(X_K)$, there is exactly one
$\xi_K(y)\ge X_K$ satisfying $\Psi_K(\xi_K(y))=\log y$.
This is the inverse of a particular finite function, rather than a choice
of interpolation of the discrete count.

For the leading model the increasing-branch inverse is
\begin{equation}
 w(y):=\xi_0(y)
 =-\frac{12}{\lambda}
 W_{-1}\left(-\frac{\lambda}{12}
                     \left(\frac{C_*}{y}\right)^{1/12}\right).
 \label{shr:h5:inv:Lambert}
\end{equation}
For large $y$, the argument lies in $(-1/e,0)$. Indeed,
$y=C_*e^{\lambda w}w^{-12}$ is equivalent to
$w e^{-\lambda w/12}=(C_*/y)^{1/12}$.
The real branch $W_{-1}$~\cite{DLMFW} gives $w\to\infty$ as $y\to\infty$;
the principal branch describes the small solution of the same leading
model and is not the large-width inverse.

\subsection{Arbitrary fixed inverse corrections}
Let $t=1/w$ and seek the finite approximation
\begin{equation}
 \widehat\xi_K(y)=w+\Delta_K(1/w),\qquad
 \Delta_K(t)=\sum_{j=1}^Ks_jt^j.
 \label{shr:h5:inv:approx}
\end{equation}
Choose $s_1,\ldots,s_K$ to cancel the coefficients of degrees $1$ through
$K$ in
\begin{equation}
 G_K(t)=\lambda\Delta_K(t)-12\log(1+t\Delta_K(t))
       +\sum_{j=1}^K\ell_jt^j(1+t\Delta_K(t))^{-j}.
 \label{shr:h5:inv:residual}
\end{equation}
At degree $j$, the new coefficient $s_j$ occurs only as $\lambda s_j$.
The recursion therefore has a unique solution in
$\Q[\lambda^{-1}]$. The first two coefficients are
\begin{equation}
 s_1=-\frac{48}{5\lambda},\qquad
 s_2=-\frac{25}{4\lambda}-\frac{576}{5\lambda^2}.
 \label{shr:h5:inv:firstcoeffs}
\end{equation}
The logarithmic displacement is already incorporated in $w$, so these
additional corrections start at order $w^{-1}$.

\begin{proposition}[Effective finite inverse error]
\label{shr:h5:inv:finiteerror}
For each fixed $K$, one can compute $H_K\ge0$ and $W_K>0$ such that
\begin{equation}
 |\xi_K(y)-\widehat\xi_K(y)|
 \le\frac{2H_K}{\lambda}w(y)^{-K-1}
 \label{shr:h5:inv:bound}
\end{equation}
whenever $w(y)\ge W_K$ and $y\ge\exp\Psi_K(X_K)$.
\end{proposition}
\begin{proof}
Choose $W_K\ge2X_K$ so that
$\sum_{j=1}^K|s_j|W_K^{-j-1}\le1/2$. Then
$z(t):=1+t\Delta_K(t)\ge1/2$ for $0\le t\le1/W_K$ and
$\widehat\xi_K=wz(1/w)\ge w/2\ge X_K$.
The Taylor coefficients of $G_K$ through degree $K$ vanish.
For $K\ge1$, repeated differentiation of~\eqref{shr:h5:inv:residual} gives
\[
 G_K^{(K+1)}(t)=\frac{P_K(t)}{z(t)^{2K+1}},
\]
where $P_K$ is a finite, computable polynomial with coefficients in
$\Q[\lambda,\lambda^{-1}]$. The largest denominator power comes from
$K+1$ derivatives of $z^{-K}$; derivatives of the logarithm have lower
powers and can be put over the same denominator.
Taylor's theorem therefore gives $|G_K(t)|\le H_Kt^{K+1}$ with
\begin{equation}
 H_K=\frac{2^{2K+1}}{(K+1)!}
                  \sum_i|[t^i]P_K(t)|W_K^{-i}.
 \label{shr:h5:inv:H}
\end{equation}
For $K=0$, take $\Delta_0=0$ and $H_0=0$.
The identity $\Psi_0(w)=\log y$ gives exactly
$\Psi_K(\widehat\xi_K)-\log y=G_K(1/w)$.
Both $\widehat\xi_K$ and $\xi_K$ belong to the interval where
$\Psi_K'\ge\lambda/2$, so the mean value theorem proves the bound.
\end{proof}

This proof uses only finite functions. It makes no convergence assertion
about the infinite inverse series. Certified intervals for $\lambda$
suffice to turn the finite formulas for $W_K,H_K$ into numerical bounds.

\subsection{The exact integer threshold}
For $y>0$, define
\begin{equation}
 \mathcal N(y)=\min\{n\ge1:a_n\ge y\}.
 \label{shr:h5:inv:threshold}
\end{equation}
It exists since~\eqref{shr:h5:asy:count} implies $a_n\to\infty$.
Eventual strict increase follows without differentiating an error term:
from~\eqref{shr:h5:asy:logerror} at $K=0$,
\begin{align*}
 \log a_{n+1}-\log a_n
 &\ge\lambda-12\log(1+1/n)-C_0(n^{-1}+(n+1)^{-1})\\
 &\ge\lambda-\frac{12+2C_0}{n}.
\end{align*}
Thus any integer beyond $N_0^{\log}$ and
$2(12+2C_0)/\lambda$ is an effective monotonicity onset.

Fix $K$, put $p=K+1$, and choose an integer $N_*$ beyond this onset,
$N_K^{\log}$, and $X_K$. Suppose $y$ is large enough that
\begin{equation}
 \begin{gathered}
 y>\max_{1\le n<N_*}a_n,\qquad
 \xi:=\xi_K(y)\ge\max(2N_*+2,4),\\
 \delta:=\frac{2^{p+1}C_K}{\lambda}\xi^{-p}\le1.
 \end{gathered}
 \label{shr:h5:inv:onset}
\end{equation}
Every condition is finite and checkable. The prefix condition is necessary
unless one has separately established monotonicity on the entire prefix.

\begin{theorem}[A qualified ceiling enclosure]
\label{shr:h5:inv:envelope}
Under~\eqref{shr:h5:inv:onset},
\begin{equation}
 \left\lceil\xi_K(y)-\delta\right\rceil
 \le\mathcal N(y)\le
 \left\lceil\xi_K(y)+\delta\right\rceil,
 \qquad \delta=\frac{2^{K+2}C_K}{\lambda}\xi_K(y)^{-K-1}.
 \label{shr:h5:inv:ceilings}
\end{equation}
\end{theorem}
\begin{proof}
Put $m=\lceil\xi+\delta\rceil$. The derivative bound gives
$\Psi_K(m)-\log y\ge\lambda\delta/2$.
Its logarithmic error is at most $C_Km^{-p}\le C_K\xi^{-p}
\le\lambda\delta/2$, so $a_m\ge y$.
Next put $j=\lceil\xi-\delta\rceil-1$. We have
$j<\xi-\delta$ and $j\ge\xi-2\ge\xi/2\ge N_*$.
Hence $\Psi_K(j)-\log y<-\lambda\delta/2$, whereas
$C_Kj^{-p}\le2^pC_K\xi^{-p}=\lambda\delta/2$.
Thus $a_j<y$. Eventual monotonicity and the prefix condition rule out
all earlier crossings. These comparisons prove both inequalities.
\end{proof}

When the two ceilings agree, the integer threshold is certified. For
sufficiently large $y$ one has $2\delta<1$, so at most two adjacent integers
remain. Near an integer boundary, an unconditional rule to round the
asymptotic inverse is not justified. Instead evaluate the remaining exact
recurrence values and compare them with $y$; for an integer or rational
target these are exact terminating rational comparisons, including equality.
For a general real target, the needed exact comparison must be supplied
or certified by the representation of that target. Numerical intervals
alone need not resolve equality with an integer count.

The finite Lambert approximation also yields an explicit enclosure. Write
$E'=2H_K/\lambda$ and $D=2^{p+1}C_K/\lambda$, and enlarge the onset so that
$|\xi_K-\widehat\xi_K|\le E'w^{-p}$ and $\xi_K\ge w/2$.
Then~\eqref{shr:h5:inv:ceilings} implies
\begin{equation}
 \left\lceil\widehat\xi_K-(E'+2^pD)w^{-p}\right\rceil
 \le\mathcal N(y)\le
 \left\lceil\widehat\xi_K+(E'+2^pD)w^{-p}\right\rceil.
 \label{shr:h5:inv:Lambertceilings}
\end{equation}
Here $w$ is~\eqref{shr:h5:inv:Lambert}, $\widehat\xi_K$ is a finite correction
polynomial, and every constant has a specified finite construction.
Interval evaluation must preserve the outward direction of both endpoints.
If an endpoint cannot be separated from an integer, retain both candidates
and use the exact comparison rather than assuming a rounding direction.

\section{Exact verification and questions left open}
\label{shr:h5:sec:verification}
The accompanying code is an auditable supplement to the mathematical
proof. Its principal checks take place in exact polynomial and rational
rings. Finite values are reported in a separate part of the receipt so
that their logical role is not confused with an all-width theorem.

\subsection{What the uniform verifier establishes}
The standard-library program \texttt{code/verify\_certificate.py}
uses dense polynomials over $\Q$ for rational functions of $n$ and sparse
polynomials for the finitely many additional formal variables. Rational
equality is tested by exact zero numerators, not by evaluating at sample
widths. Its certificate covers:
\begin{enumerate}
\item the coefficient identities for both contiguous functions, including
the homogeneous coefficient ratio and the endpoint at $t=1$;
\item the general moment flux identities, all moment reductions used in
the product, and each coefficient of the resulting affine period shift;
\item the trace/open-chain scalar reductions, the full-count expression,
and each independent coefficient in the first-difference identity;
\item the $\gamma$ bridge and initial value, the rising-factorial source
polynomials, and positive coefficients after shifting $P,Q$ by one;
\item all four literal OEIS coefficients, the explicit inhomogeneous
forcing, and a separate direct elimination of the three-state period
system against the literal operator.
\end{enumerate}
The receipt records each reduced moment as a three-component row in the
state $(X,Z,1)$, with polynomial numerators and denominators in ascending
powers of $n$. These rows are one concrete machine-readable form of
Section~\ref{shr:h5:sec:moments}'s certificate. The algebra library has its own
small consistency checks; every essential test uses explicit exceptions
rather than assertions that disappear under Python optimization.

The program also generates all 15 signed six-label matchings and all 225
products. Its path/cycle classification reproduces the four exact counts
and signs in Section~\ref{shr:h5:sec:transfer}. This is a complete enumeration of
a fixed finite Pfaffian calculation. The reduction of arbitrary kernels
to its four classes remains the sign and symmetry argument in the text.

\subsection{Independent finite checks and their limits}
Three separate exact constructions are compared for $1\le n\le20$:
\begin{itemize}
\item A literal-cell-graph dynamic program uses the four local inequalities
in~\eqref{shr:h5:eq:arrays} and counts its linear extensions. It does not use the
proposed recurrence or the Pfaffian formula.
\item Direct polynomial integration of the kernel on both halves of the
square constructs $x_n,y_n$ and evaluates the full-count formula. It does
not obtain the periods by their contiguous recurrence.
\item The printed nested sum is evaluated with exact factorials and
fractions, including the rising-factorial denominator.
\end{itemize}
These agree, and the operator's finite residuals vanish wherever their
terms are available. Separately, three arbitrary integer-matrix examples
compare the discrete odd Pfaffian identities against literal sums of
minors. Such tests are useful for catching transcription, sign,
normalization, and indexing mistakes. They neither replace the integral
argument nor establish a rational identity from finitely many values.
The distinct zero-index failure is tested explicitly.

\subsection{Rebuilding the manuscript and certificates}
The package contains the editable TeX, the frozen PDF, exact-arithmetic
programs and receipts, a source-attribution record, and a frozen SHA-256
inventory. The commands are described in \texttt{README.md}; the verifier
alone requires only Python's standard library. A complete build additionally
requires the stated TeX Live packages and fonts. There is no network access
or downloaded input during reproduction.

The build program verifies every source member, reruns the certificates,
compiles the TeX in a private temporary directory, and checks byte equality
against the frozen receipt and PDF. It creates a deterministic ZIP with
fixed member order, timestamps, and permissions. The output path must be a
new directory outside the source package; existing files, symlink
components, and unsafe paths are rejected. Each Python subprocess has
bytecode writing disabled and preserves the requested optimization level.
The small integer-to-string limit used in the recorded builds is stated
in the README. Source hashes are checked again after reproduction.

Reproducing a fixed PDF byte-for-byte depends on the recorded TeX and font
versions. A different toolchain can produce different bytes from identical
mathematics; it requires a fresh rendering inspection rather than a claim
that the discrepancy is harmless. The inventory guarantees consistency
with the supplied package, not independent authentication or peer review.
No third-party source PDF or private review material is redistributed.

\begin{writenote}[Added 5 October 2026, batch~101: the shipped layout]
``The package'' in this section is the delivered archive. In ProveIt its
programs and receipt carry the prefix \texttt{04-height-five-}:
\texttt{code/}\emph{name}\texttt{.py} (four modules) is shipped as
\texttt{code/\allowbreak 04-\allowbreak height-\allowbreak five-}\emph{name}\texttt{.py},
\texttt{build.py} as \texttt{code/\allowbreak 04-\allowbreak height-\allowbreak five-\allowbreak build.py},
\texttt{code/\allowbreak receipt.json} as \texttt{data/\allowbreak 04-\allowbreak height-\allowbreak five-\allowbreak receipt.json},
and \texttt{SOURCES.md} as \texttt{04-\allowbreak height-\allowbreak five-\allowbreak SOURCES.md}. The editable
source, the frozen PDF, the delivered README and the SHA-256 inventory
named above are not shipped. \texttt{verify\_\allowbreak certificate.py} imports
\texttt{exact\_algebra}, \texttt{printed\_\allowbreak coefficients} and
\texttt{asymptotic\_\allowbreak certificate} under their delivered names, and
\texttt{build.py} needs the inventory and the frozen PDF, so neither runs in
place; the verifier runs on a copy with the delivered names restored, and
the README of the report gives the commands. At the write (5~October 2026,
Python~3.14.4 on Windows, on such a copy) \texttt{python -B
code/\allowbreak verify\_\allowbreak certificate.py} printed
\texttt{data/\allowbreak 04-\allowbreak height-\allowbreak five-\allowbreak receipt.json} byte for byte after removal of
carriage returns, in about 19~seconds; the receipt reports status PASS and
58 named uniform certificates. \texttt{build.py} was not run.
\texttt{code/\allowbreak 04-\allowbreak height-\allowbreak five-\allowbreak printed\_\allowbreak coefficients.py} is a
literal transcription of the A181199 operator (its docstring cites the
OEIS internal entry); as OEIS content that operator is published by The
OEIS Foundation Inc.\ under the Creative Commons Attribution-ShareAlike~4.0
licence. \texttt{code/\allowbreak 04-\allowbreak height-\allowbreak five-\allowbreak asymptotic\_\allowbreak certificate.py}
lists the first eight terms of A181199 as expected values of its own
recurrence run.
\end{writenote}


\subsection{Further questions}
\begin{question}
For heights seven, nine, and higher, can the open paths in the odd
augmented-Pfaffian expansion be reduced to a comparably small family of
polynomial periods?
\end{question}
The present height-five calculation shows why traces alone are
insufficient at odd height: the length-four open path must be retained.
A systematic relation between path length, period dimension, and the
minimal recurrence order would give a useful general theory.

\begin{question}
Can one derive the stated order-three operator directly by a compact
creative-telescoping certificate on the original order polytope, with a
size comparable to the two-period contiguity certificate?
\end{question}
The factorization in~\eqref{shr:h5:eq:operator-factorization} is a factorization
over $\Q(n)$; it is not a proof that order three is minimal among
homogeneous rational-coefficient recurrences for this particular solution.
Minimality is a separate question, not needed for either source statement.

\begin{question}
What is the precise exponentially small information carried by the
homogeneous modes of the triangular period system?
\end{question}
Their limiting multipliers are $1/3125$ and $-1/27$ for the normalized
system. The fixed-order contraction estimates establish Poincar\'e
expansions and effective finite bounds. They do not identify a canonical
exponentially improved expansion, Stokes data, or Borel sum. Those stronger
questions require additional normalization and analysis.

\begin{question}
Can the finite inverse-index bounds be sharpened enough to certify most
threshold ceilings with little or no exact recurrence refinement?
\end{question}
The effective constants here are deliberately conservative. Refining them
is an algorithmic problem distinct from proving the recurrence, and it
must retain the near-integer caveat in any ceiling rule.

\begin{writenote}[Added 5 October 2026, batch~101]
The four questions above are the manuscript's Questions~1--4, numbered
16--19 here. Question~17 continues Research question~1 of
Part~\ref{shr:part:one} and Research question~8 of
Part~\ref{shr:bp:part}, now for a compact certificate rather than a proof;
Question~18 is the height-five form of Research question~4 and of item~(F8);
Question~19 continues Research question~6.
\end{writenote}

\subsection{[write] One period, and no first-order recurrence}
\label{shr:h5:sec:write}

Added 5 October 2026, batch~101. This subsection is the write's, not the
manuscript's.

\begin{remark}[\textbf{[write]} The period of Parts~III and~IV]
\label{shr:h5:rem:sameZ}
The period \eqref{shr:h5:eq:triangle} is the period \eqref{shr:h4:eq:Ztriangle}
of Part~\ref{shr:h4:part}: both are $n^2\int_{\mathcal D}u\,f^{n-1}$ over the
same triangle, and both arise as $N^3H$ with $H=\ip{t^n/N}{h}$
(\eqref{shr:h4:eq:Zdef}, \eqref{shr:h5:eq:periods}). The recurrence
\eqref{shr:h5:eq:Z-shift} is Theorem~\ref{shr:h4:thm:Z} divided by
$3(3n+1)(3n+2)$: indeed $3(3n+1)(3n+2)\nu_n=-(n+1)^2$ and
$3(3n+1)(3n+2)\delta_n=(28n^3+50n^2+29n+5)/(2(2n+1))$. Part~III proves it
by the divergence certificate \eqref{shr:h4:eq:polycert}; Part~IV derives it
from the $y$-contiguity \eqref{shr:h5:eq:y-contig} and the moment closure
\eqref{shr:h5:eq:Z-certificate}, without citing Part~III. These are two
independent proofs of one identity. In the same way \eqref{shr:h5:eq:T1},
\eqref{shr:h5:eq:T2}, \eqref{shr:h5:eq:U-final} and the scalar identities
\eqref{shr:h5:eq:scalar-identities} are \eqref{shr:h4:eq:trB},
\eqref{shr:h4:eq:trB2}, \eqref{shr:h4:eq:UH}, \eqref{shr:h4:eq:I},
\eqref{shr:h4:eq:K0} and~\eqref{shr:h4:eq:J0} with the same proofs, and \eqref{shr:h5:eq:transfer},
\eqref{shr:h5:eq:andreief} and Lemma~\ref{shr:h5:lem:indicator} are
\eqref{shr:h4:eq:transfer}, Lemma~\ref{shr:h4:lem:andreief} and
\eqref{shr:h4:eq:indicator} with $m=5$. The kernels $K$, $W$, $h$ and their
scalars depend on $n$ only, not on the height, so these identities are the
same statements after renaming Part~III's $\alpha$, $\beta$ to Part~IV's
$f$, $b$. Where the printed forms differ they agree by $c=1/(2N^2)$ and
$\int e=0$: thus $c^2/2=1/(8N^4)$ in \eqref{shr:h4:eq:I}, and the term
$\int e$ of \eqref{shr:h4:eq:trB} is absent from \eqref{shr:h5:eq:T1}.
Part~IV's $J_0=I$ in \eqref{shr:h5:eq:scalar-identities} is Part~III's
display~\eqref{shr:h4:eq:J0} combined with \eqref{shr:h4:eq:I}. The proof
of Lemma~\ref{shr:h4:lem:Zirr} uses only the recurrence, which Part~IV
proves for its own $Z_n$; so Proposition~\ref{shr:h5:prop:nothyper} below
does not depend on this identification of the integrals. (Reworded
5~October 2026 after the intake's independent check, recorded at the end
of Section~\ref{shr:sa:sec:further}. The list above first cited only
\eqref{shr:h4:eq:I} and~\eqref{shr:h4:eq:K0} for the scalar identities,
and the last sentence first read: ``The kernels $K$, $W$, $h$ and their
scalars depend on $n$ only, not on the height, so these identities are
literally the same statements.'')
\end{remark}

\begin{proposition}[\textbf{[write]} A181199 is not hypergeometric]
\label{shr:h5:prop:nothyper}
The sequence $a_n=T_5(n)$ has no recurrence of order zero or one, in the
sense of Section~\ref{shr:h4:sec:minimal}. Hence the smallest order of a
recurrence for A181199 is two or three, and the operator of
Theorem~\ref{shr:h5:thm:operator} has minimal order unless A181199 obeys a
recurrence of order two.
\end{proposition}

\begin{proof}
Every $a_n$ with $n\ge1$ is a positive integer, so a recurrence of order
zero is impossible, as in Proposition~\ref{shr:h4:prop:minimal}. Suppose
$q_1(n)a_{n+1}+q_0(n)a_n=0$ for $n\ge n_0$ with $q_1\ne0$; enlarging
$n_0$, $a_{n+1}=\varrho(n)a_n$ for $n\ge n_0$ with $\varrho\in\Q(n)$. By
Section~\ref{shr:h5:sec:sum-operator}, $P(n)>0$ and $Q(n)>0$ for $n\ge1$,
so $u(n)>0$ and $v(n)>0$ there. Put $Y_n=Z_n-\gamma_n$. By
\eqref{shr:h5:eq:Z-shift} and \eqref{shr:h5:eq:gamma-bridge},
\begin{equation}
 Y_{n+1}-\nu_nY_n=\omega(n):=\frac{9(n+1)^2v(n)}{8(3n+1)(3n+2)}>0
 \qquad(n\ge1),
 \label{shr:h5:eq:Yrec}
\end{equation}
and by \eqref{shr:h5:eq:target-diff},
$(\varrho(n)-1)a_n=2M_nu(n)Y_n$ for $n\ge n_0$. If $\varrho(n)=1$ for
infinitely many $n$, then $\varrho=1$ identically, so $Y_n=0$ for all
$n\ge n_0$, contradicting \eqref{shr:h5:eq:Yrec}. Enlarging $n_0$, we
therefore have $\varrho(n)\ne1$, and hence $Y_n\ne0$, for $n\ge n_0$. Then
\[
 \sigma(n):=\frac{Y_{n+1}}{Y_n}
 =\frac{(\varrho(n+1)-1)\,\varrho(n)\,u(n)}{(\varrho(n)-1)\,u(n+1)}
 \cdot\frac{M_n}{M_{n+1}}
\]
is the value of a rational function, since $M_n/M_{n+1}=\rho_n$
\eqref{shr:h5:eq:rho}. By \eqref{shr:h5:eq:Yrec},
$Y_n(\sigma(n)-\nu_n)=\omega(n)\ne0$ for $n\ge n_0$, so $\sigma-\nu_n$ is
a nonzero rational function and
\[
 Z_n=\gamma_n+\frac{\omega(n)}{\sigma(n)-\nu_n}\qquad(n\ge n_0)
\]
coincides with a rational function for all large $n$. By
Remark~\ref{shr:h5:rem:sameZ} this $Z_n$ is the period of
Part~\ref{shr:h4:part}, so Lemma~\ref{shr:h4:lem:Zirr} gives a
contradiction. Theorem~\ref{shr:h5:thm:operator} supplies a recurrence of
order three.
\end{proof}

Whether A181199 obeys a recurrence of order two is item~(F9) below. At
height four, excluding order one settled the question; at height five it
leaves orders two and three.

As at height four, the degree drops at higher order.
Proposition~\ref{shr:h5:prop:higher} was found and proved by the intake's
independent check, 5~October 2026; it is neither the manuscript's nor the
write's. Notation as before Proposition~\ref{shr:h4:prop:higher}, with
$L_{\rm OEIS}=\sum_{j=0}^3c_j(n)E^j$ the operator of
Theorem~\ref{shr:h5:thm:operator} and $Q$ the polynomial of
\eqref{shr:h5:eq:PQD}.

\begin{proposition}[\textbf{[check]} Orders four and five at height five]
\label{shr:h5:prop:higher}
Let $a_n=T_5(n)$ (A181199).
\begin{enumerate}[label=(\alph*),leftmargin=*]
\item There are polynomials $e_0,\ldots,e_4$ of degree $18$ and
$U_0,U_1\in\Q[n]$ of degree five with
\[
 \sum_{j=0}^4e_j(n)E^j=\frac{U_1(n)E+U_0(n)}{Q(n+1)}\,L_{\rm OEIS},
\]
and $\sum_{j=0}^4e_j(n)a_{n+j}=0$ for every $n\ge1$: A181199 obeys a
recurrence of order four and degree $18$.
\item There are polynomials $g_0,\ldots,g_5$ of degree $16$ and
$V_0,V_1,V_2\in\Q[n]$ with
\[
 \sum_{j=0}^5g_j(n)E^j=\Bigl(\frac{V_2(n)E^2}{Q(n+2)}
 +\frac{V_1(n)E}{Q(n+1)Q(n+2)}+\frac{V_0(n)}{Q(n+1)}\Bigr)L_{\rm OEIS},
\]
and $\sum_{j=0}^5g_j(n)a_{n+j}=0$ for every $n\ge1$: A181199 obeys a
recurrence of order five and degree $16$.
\item If A181199 has no recurrence of order two, then every recurrence of
A181199 holds for every $n\ge1$, and the smallest degree of a recurrence is
$24$ at order three, $18$ at order four and $16$ at order five.
\end{enumerate}
The polynomials $e_j$, $g_j$, $U_k$ and $V_k$ are given in
\texttt{data/\allowbreak check-\allowbreak lower-\allowbreak order-\allowbreak recurrences.\allowbreak json}.
\end{proposition}

\begin{proof}
(a) and (b). Expanding the right sides with $Ef(n)=f(n+1)E$, clearing the
denominators and comparing coefficients with the left sides is a finite
identity over $\Q[n]$, checked exactly from the data file and the
coefficients of Appendix~\ref{shr:h5:app:literal}. By
Section~\ref{shr:h5:sec:sum-operator}, $Q(x)>0$ for every real $x\ge1$, so
for $n\ge1$ every coefficient of the rational left factors is finite, and
in (a)
\[
 \sum_{j=0}^4e_j(n)a_{n+j}
 =\frac{U_1(n)\,(L_{\rm OEIS}a)(n+1)+U_0(n)\,(L_{\rm OEIS}a)(n)}{Q(n+1)}=0
\]
by Theorem~\ref{shr:h5:thm:operator}; (b) is the same with three terms.

(c) Let $R$ be a recurrence of order $r$, valid for $n\ge n_0$; by
Proposition~\ref{shr:h5:prop:nothyper} and the hypothesis, $r\ge3$. As in
the proof of Proposition~\ref{shr:h4:prop:higher}, right division by
$L_{\rm OEIS}$ gives
$R=\bigl(\sum_{k=0}^{r-3}W_kE^k\bigr)L_{\rm OEIS}+s_2E^2+s_1E+s_0$, where
the denominators of the $W_k$ and $s_j$ divide
$\prod_{i=0}^{r-3}c_3(n+i)$. The real zeros of $c_3$ \eqref{shr:h5:eq:c3}
are $-3$, $-11/4$, $-8/3$, $-5/2$, $-7/3$, $-9/4$, $-2$, $-3/2$ and those
of $Q$, which lie below one since $Q(x)>0$ for $x\ge1$; so no $c_3(n+i)$
with $i\ge0$ vanishes at an integer $n\ge1$. For $n\ge n_0$ it follows that
$s_2(n)a_{n+2}+s_1(n)a_{n+1}+s_0(n)a_n=0$; after clearing denominators
this is a polynomial relation of order at most two, which the hypothesis
and Proposition~\ref{shr:h5:prop:nothyper} force to be trivial. So
$s_0=s_1=s_2=0$, and $(Ra)(n)=\sum_kW_k(n)(L_{\rm OEIS}a)(n+k)=0$ for every
$n\ge1$. The linear-algebra step of the proof of
Proposition~\ref{shr:h4:prop:higher} then applies, with the exact values
that Theorem~\ref{shr:h5:thm:operator} determines from $a_1=a_2=1$,
$a_3=16$: for $(\bar r,d,N)=(3,23,126)$, $(4,17,120)$ and $(5,15,126)$ the
matrices have full column rank $96$, $90$ and $96$ over $\Q$. So no
recurrence of order at most three has degree below $24$, none of order at
most four has degree below $18$, and none of order at most five has degree
below $16$; Theorem~\ref{shr:h5:thm:operator}, (a) and (b) attain these
degrees.
\end{proof}

Without the hypothesis of (c), the same ranks show only that no recurrence
of those orders and degrees holds for every $n$ in the range of rows.

\subsection{Further questions and research recorded at the write}
\label{shr:h5:sec:further}

Added 5 October 2026, batch~101, under Vladimir Reshetnikov's standing rule.
The manuscript's own open questions are Research questions~16--19 above;
its text after Question~17 states the non-claim that the factorization
over $\Q(n)$ ``is not a proof that order three is minimal''.
\begin{enumerate}[label=(F\arabic*),leftmargin=*,start=9]
\item \emph{Minimality of the A181199 recurrence} (the manuscript's
non-claim after Research question~17; the placement). Proposition~\ref{shr:h5:prop:nothyper}
excludes order one, and Proposition~\ref{shr:h5:prop:higher} gives
recurrences of order four and degree $18$ and of order five and degree
$16$, and the smallest degrees $24$, $18$ and $16$ at orders three, four
and five if order two is excluded. Open: whether some recurrence of order
two exists; the smallest degrees at orders three to five without that
hypothesis, and at orders six and more. Evidence, not proof: no recurrence
of order at most two and degree at most $60$ holds for every
$5\le n\le214$, since the $210\times183$ integer matrix $(n^ia_{n+j})$
($0\le j\le2$, $0\le i\le60$) with those rows has rank $183$ over $\Q$
(exact computation, recorded in \texttt{data/\allowbreak check-\allowbreak lower-\allowbreak order-\allowbreak recurrences.\allowbreak json}). This is not a proof, because a
recurrence need only hold for large $n$, and for a recurrence of order two
nothing like Proposition~\ref{shr:h5:prop:higher}(c) shows that it holds
from $n=5$. (Re-scoped 5~October 2026 after the intake's independent
check, which found and proved Proposition~\ref{shr:h5:prop:higher}; this
item first read: ``Open: whether some recurrence of order two exists, and
the smallest degree at orders three and more.'') Sketch: if
$a_{n+2}=r_1(n)a_{n+1}+r_0(n)a_n$ with rational $r_0,r_1$, then
$\Delta_n=a_{n+1}-a_n=-\tfrac{27}4H_n\sum_{i=1}^{n-1}I_i$
\eqref{shr:h5:eq:sum-increment} satisfies
$\Delta_{n+1}=(r_0+r_1-1)a_n+(r_1-1)\Delta_n$. The coefficient
$r_0+r_1-1$ cannot vanish identically, since $\Delta_n$ would then be
hypergeometric and the proof above applies to $Y_n=\Delta_n/(2M_nu(n))$;
so $a_n$ lies in the $\Q(n)$-span of $\Delta_n$ and $\Delta_{n+1}$.
Conversely such a representation gives order two. Deciding it is a
rational-solution problem for a first-order $2\times2$ difference system
whose coefficients come from $R$ and $S$ \eqref{shr:h5:eq:RS-ratios}.
Missing: that computation, or an argument like
Lemma~\ref{shr:h4:lem:Zirr} for the sum $\sum_{i<n}I_i$.
\item \emph{Numerical constants} (Sections~\ref{shr:h5:asy:section}
and~\ref{shr:h5:inv:section}). The receipt evaluates $E_{Z,K}$, $E_{b,K}$,
$C_K$ and the onsets at $K=0$ and $K=6$
(\eqref{shr:h5:asy:explicit6}, \eqref{shr:h5:asy:logexamples}); the inverse
constants $H_K$, $W_K$, $X_K$ and the onset \eqref{shr:h5:inv:onset} of the
threshold envelope are constructed but not evaluated. Missing: those
values and a worked certified threshold (compare Research question~19).
\end{enumerate}

\clearpage
\part{Two earlier independent routes: all-order asymptotics, rational diagonals and transcendence}
\label{shr:sa:part}

\noindent{\Large Research reports 85 and 91, prepared for Vladimir
Reshetnikov, October 1 and 2, 2026\par}
\medskip

\begin{writenote}[Added 5 October 2026, batch~101: provenance of Part~V]
Part~V prints two manuscripts in full, one section each:
\begin{itemize}
\item Section~\ref{shr:sa:sec} is Research Report~85 of the session
bundle, \emph{Fixed-Height Shifted Strips: All-Orders Asymptotics and
Inversion}, dated 1~October 2026 (a 7-page PDF; one 330-line source with
23 labels and four bibliography entries). It arrived as the archive
\texttt{ProveIt\_\allowbreak Fixed\_\allowbreak Height\_\allowbreak Shifted\_\allowbreak Strips\_\allowbreak Asymptotics.zip}
(357{,}992 bytes, 21 files, no wrapper directory). Its code and records
are shipped with the prefix \texttt{05-\allowbreak strip-\allowbreak asymptotics-}.
\item Section~\ref{shr:sd:sec} is Research Report~91, \emph{Rational
Diagonals and Transcendence for Fixed Height Shifted Strips}, dated
2~October 2026 (a 9-page PDF; one 448-line source with 27 labels and eight
bibliography entries). It arrived as the archive
\texttt{ProveIt\_\allowbreak Shifted\_\allowbreak Strips\_\allowbreak Rational\_\allowbreak Diagonals\_\allowbreak and\_\allowbreak Transcendence.zip}
(371{,}265 bytes, 12 files, no wrapper directory). Its code and records are
shipped with the prefix \texttt{06-\allowbreak strip-\allowbreak diagonals-}.
\end{itemize}
Both arrived in commit \texttt{60f54ea06} and were placed in commit
\texttt{f7c612c72} (batch~101) as additions to this report. Both title
lines and PDF metadata read ``Research report prepared for Vladimir
Reshetnikov with OpenAI'': unlike the sources of Parts~I--IV, they name a
tool. Neither names a ProveIt commit or read the repository; Report~91
cites Report~85 as an ``unpublished project report supplied with this
research task''. Their \LaTeX\ sources, delivered READMEs, PDFs, SHA-256
manifests (20 and 11 entries, verified at placement) and two generic helper
files of Report~85 (a build script and an integrity checker that are byte
copies of files shipped elsewhere in the collection) and the integrity
checker of Report~91 are not shipped; they survive in the arrival commit's
archives. Both manuscripts are unrefereed and AI-assisted; nothing in them
is formalized, and their place in this report confers no formal status.

\emph{Chronology.} The two manuscripts are dated before the manuscripts of
Part~\ref{shr:part:one} (3~October) and Part~\ref{shr:bp:part}
(4~October), neither of which knew them; they reached ProveIt on
5~October 2026. They are printed after Parts~I--IV only because they were
placed later, and no number of those Parts moved. Report~85 proves
Part~\ref{shr:part:one}'s Theorem~\ref{shr:thm:main} (with
\eqref{shr:eq:b12}, the A181199 constant and two corrections) by another
route; Report~91 proves Part~\ref{shr:bp:part}'s
Theorem~\ref{shr:bp:thm:diagonal-main} by the same route, and
Part~\ref{shr:part:one}'s Theorem~\ref{shr:thm:algebraic} with
elementary radial lemmas in place of the Puiseux argument.
Section~\ref{shr:sa:sec:chronology} sets out the correspondence; by their
date lines these are the earliest proofs of those statements in this
report. The dates are the manuscripts' own and certify nothing outside
ProveIt.

\emph{Editorial changes.} Each manuscript's sections are subsections of
its section here, and its subsections are subsubsections. Equation~($k$)
of either manuscript is ($43.k$) or ($44.k$) here. Theorem-like statements
are numbered through the section: Report~85's Theorem~1.1 is
Theorem~\ref{shr:sa:thm:main}; Report~91's Theorems~1.1--1.2, Lemmas~2.1,
3.1 and~3.2, Propositions~4.1 and~6.1, Lemmas~7.1--7.3 and its closing
Remark are 44.1--44.11 in that order.
Labels carry the prefixes \texttt{shr:sa:} and \texttt{shr:sd:}. Their
titles, author lines, dates and abstracts open the two sections. In
Report~85, \verb|\tr| is typeset as $\Tr$, as delivered. Report~91's
citation of Report~85 now points to Section~\ref{shr:sa:sec}. The
bibliographies are merged: the OEIS entries (Report~85 accessed them on
1~October, Report~91 on 2~October 2026), Kauers--Koutschan, Sun and
Bostan--Lairez--Salvy are Part~\ref{shr:part:one}'s and
Part~\ref{shr:bp:part}'s entries, and Lindemann's article \cite{Lindemann}
is added. Remarks take this report's theorem style. The \textbf{[write]}
notes and Section~\ref{shr:sa:sec:chronology} are new. No other statement,
proof or number of either manuscript was changed.
\end{writenote}

\begin{writenote}[Added 5 October 2026, batch~101: notation of Part~V]
Both manuscripts write $g(m,n)$ for $T_m(n)$ and use the same letters for
the same objects in their shared asymptotic argument. Within Part~V each
symbol has the meaning below; the last column gives the tempting false
reading from elsewhere in this report. The worst trap is Report~85's
$c_j(m)$: it is Part~\ref{shr:part:one}'s $b_{m,j}$, the coefficient in the
normalization $K_mm^{mn}n^{-\alpha_m}$, \emph{not} Part~I's $c_{m,j}$.
{\footnotesize
\begin{longtable}{@{}>{\raggedright\arraybackslash}p{0.13\linewidth}>{\raggedright\arraybackslash}p{0.27\linewidth}>{\raggedright\arraybackslash}p{0.27\linewidth}>{\raggedright\arraybackslash}p{0.25\linewidth}@{}}
\toprule
symbol & Report~85 (Section~\ref{shr:sa:sec}) & Report~91 (Section~\ref{shr:sd:sec}) & \emph{not}: meaning elsewhere\\
\midrule\endhead
$g(m,n)$ & $T_m(n)$ & $T_m(n)$, with $g(m,0)=1$ & ---\\
$c_j(m)$ & $b_{m,j}$ of \eqref{shr:eq:mainK} & --- & $c_{m,j}$ of \eqref{shr:eq:mainM}\\
$C_m$ & $K_m$ \eqref{shr:eq:Km} (as $2^{-m(m-1)}=4^{-d}$) & $K_m$ & Part~IV's ordered simplex $C_m$\\
$\nu$; $d$ & $\alpha_m$; $\binom m2$ & the same & Part~IV's $\nu_n$; Parts~III--IV's $d_j$\\
$\alpha_m$; $L_m$ & --- & $2^{-m(m-1)-(m-1)/2}\sqrt m\prod_{j<m}j!$; $\alpha_m\pi^{-(m-2)/2}$ & $\alpha_m=(m^2-1)/2$\\
$G_m(z)$; $A_m(t)$ & --- & $F_m(z)$; $G_m(m^{-m}t)$ & Part~II's Stirling factor $A_m(z)$\\
$R_n$; $R_m$ & $\frac{(mn)!}{(n!)^mf^{(n^m)}}$ & the same $R_n$; a rational function with $G_m=\diag R_m$ & the pair weight $R_n(x)$; $R_m=m^{-m}$; $R_m(n)$\\
$f^\lambda$, $g^\lambda$, $\delta$, $\D$ & ordinary and shifted hook counts; $(m-1,\ldots,0)$; the Vandermonde & the same & $f^{(n^m)}$ as in Corollary~2.4; Part~II's $g^\lambda$\\
$\mu_n$, $\eta_n$, $\rho$, $H(x)$ & the cut law, the hypergeometric law, the cut proportion, binary entropy & the same, with $\rho=1/m$ & $H(\lambda)$; Part~IV's $\rho_n$\\
$y_i$; $\E_m$; $S_r$ & $(\lambda_i-\rho n)/\sqrt n$; the trace-zero Gaussian--Vandermonde expectation; power sums & $y_i$ as in Report~85 & Part~I's $Y_i$, $\E$, $p_r$; Part~II's $S_1$, $S_2$\\
$\beta_k$; $\mathcal L_m$, $\mathcal B$, $L_k$, $H_j$ & $B_{2k}/(2k(2k-1))$; the formal logarithm and its coefficients & --- & recurrence coefficients $\beta_j$\\
$L$ & $m\log m$ (Section~43.6) & a path length; $L_j$ gap lengths & a truncation order; Part~I's $\lambda=m\log m$\\
$d_1$, $d_2$; $w$, $y$ & inverse constants; $Ln$, $\log(x/(C_mL^\nu))$ (Section~43.6) & $d_i=b_i-a_{\sigma(i)}$ & Parts~III--IV's $d_j=b_{m,j}$\\
$N(x)$; $J$; $K$ & the threshold; a truncation order; a gap & --- & Part~III's $N(y)$\\
$B$, $S$; $h_j$ & --- & birth and saturation letters; heights of the Dyck word & Part~II's letters $B$, $D$; Part~I's $h_j$\\
$K_k(a,b)$; $H(t)$, $\del(t)$ & --- & the chamber count; the Heaviside and Kronecker sequences & Part~II's $K_r(u,v)$, $H(a)$\\
$a^j$, $b^j$, $\mathcal D_m$, $\mathcal B_w(n)$, $D_w$, $E_m(n)$ & --- & event vectors, Dyck words, admissible vectors, dimension, the binomial sum & Part~II's $u_s$, $v_s$; $E_{m,n}$\\
$(x)_j$ & the rising factorial & the same & Part~I's falling factorial\\
\bottomrule
\end{longtable}}
\end{writenote}

\section{Report 85: fixed-height shifted strips, all-orders asymptotics and inversion}
\label{shr:sa:sec}

\noindent{\large Research report prepared for Vladimir Reshetnikov with
OpenAI, 1 October 2026\par}

\paragraph{Abstract.}
Let $g(m,n)$ count the $m\times n$ arrays filled by $1,\ldots,mn$ whose rows, columns, diagonals and downward antidiagonals increase. For every fixed height $m$, we prove a complete expansion in powers of $1/n$, with leading constant $2^{-m(m-1)}$ times that for ordinary rectangular standard Young tableaux. At height five this proves the equivalent currently conjectured in OEIS A181199, and gives relative corrections $48/(5n)$ and $5233/(100n^2)$. The proof splits a path at an intermediate shape, applies two shifted hook products, and uses a symmetric lattice expansion. We give a finite exact rule for every coefficient and an inverse expansion, retaining the necessary qualification when rounding to integer indices. No conjectured recurrence is used.

\begin{writenote}[Added 5 October 2026, batch~101]
Every statement of this section was already proved in
Part~\ref{shr:part:one}, two days later by date line and by another route;
Section~\ref{shr:sa:sec:chronology} lists the correspondence. The
abstract's ``currently conjectured'' refers to the A181199 entry as
accessed on 1~October 2026. Three steps of Subsection~43.4 are argued in
outline (item~(F11)); Part~\ref{shr:part:one}'s proof of the same theorem
is complete.
\end{writenote}

\subsection{The result and its scope}
The array counts form the table A181196. Its height-four and height-five rows are A181198 and A181199. The latter entry states the conjecture
\[
 g(5,n)\sim\frac{9\sqrt5}{2^{17}\pi^2}\frac{3125^n}{n^{12}}.
\]
Kauers and Koutschan's recurrence and closed-form conjectures for these two rows remain separate questions \cite{oeis198,oeis199,KK}. Sun proves formulas for $g(3,n)$ and for $g(n,4),g(n,5)$ \cite{Sun}; the last two vary the other dimension.

Put
\[
 d=\binom m2,\qquad \nu=\frac{m^2-1}{2},\qquad
 C_m=2^{-m(m-1)}\sqrt m\,(2\pi)^{-(m-1)/2}
          \prod_{j=0}^{m-1}j!.
\]
\begin{theorem}\label{shr:sa:thm:main}
Fix an integer $m\ge2$. There are effectively computable rational numbers $c_j(m)$ such that, for every fixed $J\ge0$,
\begin{equation}\label{shr:sa:eq:expansion}
 g(m,n)=C_m m^{mn}n^{-\nu}
 \left(1+\sum_{j=1}^J\frac{c_j(m)}{n^j}
                   +O_{m,J}(n^{-J-1})\right).
\end{equation}
The first two coefficients are
\begin{equation}\label{shr:sa:eq:c12}
 c_1(m)=\frac{(m^2-1)^2}{12m},\qquad
 c_2(m)=\frac{(m^2-1)(m^6+3m^2-1)}{288m^2}.
\end{equation}
A single formal coefficient procedure applies to every fixed height.
\end{theorem}
Thus the explicit height-five expansion is
\begin{equation}\label{shr:sa:eq:height5}
 g(5,n)=\frac{9\sqrt5}{2^{17}\pi^2}\frac{3125^n}{n^{12}}
 \left(1+\frac{48}{5n}+\frac{5233}{100n^2}+O(n^{-3})\right).
\end{equation}
For comparison, the height-four row has
\[
 g(4,n)=\frac{3}{1024\sqrt2\,\pi^{3/2}}\frac{256^n}{n^{15/2}}
 \left(1+\frac{75}{16n}+\frac{6905}{512n^2}+O(n^{-3})\right).
\]
The theorem is an all-orders Poincar\'e expansion of the dominant exponential sector. It does not classify every exponentially smaller sector or prove the unused recurrence conjectures. Height one is the immediate identity $g(1,n)=1$.

\subsection{Exact paths and the two hook products}
Fill the array in increasing label order, and let $\lambda_i$ be the number of filled cells in row $i$. Its row-choice word determines the filling uniquely. The state conditions are exactly
\begin{equation}\label{shr:sa:eq:states}
 n\ge\lambda_1\ge\cdots\ge\lambda_m\ge0,
 \quad\text{with repetitions permitted only at $0$ and $n$.}
\end{equation}
Indeed, filling cell $(i+1,j)$ for $j<n$ requires $(i,j+1)$ to be filled already. Filling $(i+1,n)$ requires the entire preceding row. These are the downward-antidiagonal and column conditions; the southeast diagonal conditions follow from row and column increase.

If an endpoint satisfies $\lambda_1<n$, its admissible prefixes are chains of strict partitions, with zero parts omitted. For an interior strict endpoint, the classical shifted hook product is
\begin{equation}\label{shr:sa:eq:shifted}
 g^\lambda=\frac{|\lambda|!}{\prod_i\lambda_i!}
        \prod_{i<j}\frac{\lambda_i-\lambda_j}{\lambda_i+\lambda_j}.
\end{equation}
This formula and the shifted-strip interpretation appear in \cite{Sun}. Reversing the path and its row order gives the complementary partition
\[
 \bar\lambda=(n-\lambda_m,\ldots,n-\lambda_1),
\]
which obeys the same state condition. Consequently the number of admissible complete paths through an interior strict state is exactly
\begin{equation}\label{shr:sa:eq:through}
 g^\lambda g^{\bar\lambda}.
\end{equation}
We do not apply \eqref{shr:sa:eq:shifted} to a saturated boundary state with repeated parts $n$.

Let $f^\lambda$ be the ordinary standard-tableau count, and set $\delta=(m-1,\ldots,0)$. The ordinary hook product and rectangular specialization are
\begin{equation}\label{shr:sa:eq:ordinary}
 f^\lambda=\frac{|\lambda|!\D(\lambda+\delta)}
                  {\prod_i(\lambda_i+m-i)!},\qquad
 f^{(n^m)}=\frac{(mn)!\prod_{j=0}^{m-1}j!}
                    {\prod_{j=0}^{m-1}(n+j)!},
\end{equation}
where $\D(x)=\prod_{i<j}(x_i-x_j)$. One way to obtain the first identity is to shift weakly ordered row-growth paths by $\delta$, apply the reflection principle, and evaluate
$|\lambda|!\det(1/(\lambda_i-i+j)!)$ as a Vandermonde determinant.

Ordinary rectangular branching at any time $k$ gives the probability law
\begin{equation}\label{shr:sa:eq:mu}
 \mu_n(\lambda)=\frac{f^\lambda f^{\bar\lambda}}{f^{(n^m)}},
 \qquad |\lambda|=k.
\end{equation}
Every admissible path is an ordinary path. The ratio of the admissible through-count to the ordinary through-count therefore always belongs to $[0,1]$, including at boundary states.

\subsection{Concentration and the leading constant}
We first cut at $k=n$, so $\rho=k/(mn)=1/m$ is fixed without any parity restriction. Compare \eqref{shr:sa:eq:mu} to the labeled multivariate hypergeometric law
\[
 \eta_n(\ell)=\frac{\prod_i\binom n{\ell_i}}{\binom{mn}n},
 \qquad \sum_i\ell_i=n.
\]
The ratio of unrestricted row words to ordinary tableaux is
\[
 R_n=\frac{(mn)!/(n!)^m}{f^{(n^m)}}
 =\frac{\prod_{j=0}^{m-1}(n+j)!}{(n!)^m\prod_{j=0}^{m-1}j!}
 =O_m(n^d).
\]
Bounding each prefix and suffix by its unrestricted multinomial count yields
\begin{equation}\label{shr:sa:eq:crude}
 \mu_n(\lambda)\le R_n\eta_n(\lambda).
\end{equation}
Stirling bounds and strict concavity of binary entropy imply that a fixed linear deviation from $\lambda_i=\rho n$ has exponentially small probability. The polynomial factor $R_n$ does not affect this conclusion.

A refined bound also excludes bounded row gaps. On a fixed central region
$\epsilon n\le\lambda_i\le(1-\epsilon)n$, put
$y_i=(\lambda_i-\rho n)/\sqrt n$, so $\sum_i y_i=0$. From \eqref{shr:sa:eq:ordinary},
\begin{equation}\label{shr:sa:eq:refined}
 \mu_n(\lambda)=R_n\eta_n(\lambda)
 \frac{\D(\lambda+\delta)^2}
 {\prod_i(\lambda_i+1)_{m-i}(n-\lambda_i+1)_{i-1}},
\end{equation}
where $(x)_j=x(x+1)\cdots(x+j-1)$. Reversing the complementary partition makes its Vandermonde equal to the first one. The denominator is at least a positive constant times $n^{2d}$, and the numerator is at most
$C_m n^d(1+\|y\|)^{2d}$.

Uniform factorial estimates give
\[
 \eta_n(\lambda)\le C n^{-(m-1)/2}e^{-c\|y\|^2}.
\]
To see the uniformity, write each binomial logarithm as
$nH(\lambda_i/n)-\tfrac12\log n+O(1)$. The denominator has the corresponding $n^{-1/2}$ prefactor, the linear entropy terms cancel because the sum is fixed, and $H''$ is bounded above by a negative constant. Thus
\begin{equation}\label{shr:sa:eq:envelope}
 \mu_n(\lambda)\le C n^{-(m-1)/2}
        (1+\|y\|)^{2d}e^{-c\|y\|^2}.
\end{equation}
Gaussian lattice-sum bounds give tightness of $y$. In a fixed box $\|y\|\le M$, a specified gap $\lambda_i-\lambda_j\le K$ occupies at most
$O_M((K+1)n^{(m-2)/2})$ lattice points. Its probability is therefore
$O_M((K+1)/\sqrt n)$. Letting first $n$ and then $M$ tend to infinity proves
\begin{equation}\label{shr:sa:eq:concentration}
 \lambda_i/n\longrightarrow\rho,\qquad
 \min_{i<j}(\lambda_i-\lambda_j)\longrightarrow\infty
 \quad\text{in probability.}
\end{equation}
In particular, interior strict states have probability tending to one.

For such a state, the two hook products give
\[
 \frac{g^\lambda}{f^\lambda}
 =\frac{\prod_i(\lambda_i+m-i)!/\lambda_i!}
        {\prod_{i<j}(\lambda_i+\lambda_j)}
       \frac{\D(\lambda)}{\D(\lambda+\delta)}.
\]
By \eqref{shr:sa:eq:concentration}, its first factor tends to $2^{-d}$ and its second to one. The complementary ratio has the same limit. The exact statewise ratios are bounded by one, so bounded convergence in probability applied to \eqref{shr:sa:eq:mu} proves
\begin{equation}\label{shr:sa:eq:comparison}
 \frac{g(m,n)}{f^{(n^m)}}\longrightarrow2^{-2d}.
\end{equation}
Finally, Stirling's formula in \eqref{shr:sa:eq:ordinary} gives
\[
 f^{(n^m)}\sim\sqrt m\,(2\pi)^{-(m-1)/2}
       \Bigl(\prod_{j=0}^{m-1}j!\Bigr)m^{mn}n^{-\nu}.
\]
This proves the leading assertion and the constant in Theorem~\ref{shr:sa:thm:main}. All the preceding bounds are uniform when $\rho$ ranges over a fixed compact subinterval of $(0,1)$ and $k=\rho mn$ is an integer.

\subsection{A symmetric lattice expansion to every order}
At a general integer cut $k=\rho mn$, the exact interior through-count \eqref{shr:sa:eq:through} becomes
\begin{equation}\label{shr:sa:eq:central}
 \frac{k!(mn-k)!\D(\lambda)^2}
 {\displaystyle\prod_i\lambda_i!(n-\lambda_i)!
  \prod_{i<j}(\lambda_i+\lambda_j)(2n-\lambda_i-\lambda_j)}.
\end{equation}
Restrict all coordinates to a fixed central region as above. Its complement contributes an exponentially small amount relative to $f^{(n^m)}$, hence also to $g(m,n)$ by \eqref{shr:sa:eq:comparison}. The summand \eqref{shr:sa:eq:central} is symmetric. Coincidence terms vanish because of $\D(\lambda)^2$, while its denominators stay positive. Therefore the ordered strict sum is exactly $1/m!$ times the sum over all labeled central integer vectors of sum $k$. There are no chamber-wall correction terms.

Write $\lambda_i=\rho n+\sqrt n\,y_i$. The leading weight on the hyperplane $\sum y_i=0$ is
\begin{equation}\label{shr:sa:eq:gaussian}
 \D(y)^2\exp\left(-\frac{\sum_i y_i^2}{2\rho(1-\rho)}\right).
\end{equation}
The affine lattice has covolume $\sqrt m\,n^{-(m-1)/2}$ in this hyperplane.

Here are the fixed-order error estimates. For a chosen $J$, restrict first to $\|y\|\le n^\varepsilon$, with $\varepsilon>0$ sufficiently small depending on $J$. Uniform Stirling expansion and Taylor expansion of the positive pair-sum factors yield polynomial coefficients at every half-order. Retain the half-orders through $n^{-J-1/2}$ before integration. Taylor's theorem, with a sufficiently long finite expansion, bounds the remaining normalized summand by
\[
 C_J n^{-J-1}(1+\|y\|)^{M_J}e^{-c\|y\|^2}.
\]
Choosing $\varepsilon$ small permits exponentiation of the remainder without losing an integrable Gaussian majorant. Outside this growing window, \eqref{shr:sa:eq:envelope} and \eqref{shr:sa:eq:crude} give a bound of the form a fixed polynomial times $e^{-c n^{2\varepsilon}}$, which is negligible to every algebraic order. These estimates are uniform for $\rho$ in a compact subset of $(0,1)$.

For each polynomial-Gaussian coefficient, Poisson summation replaces the lattice sum by its integral with an error smaller than every power of $1/n$, uniformly in the affine lattice translation. Its Fourier transform is another polynomial times a Gaussian, and every nonzero dual frequency has length at least $c\sqrt n$. The growing-window cutoff can first be removed at a superalgebraically small cost. The coefficient of half-order $j$ has parity $j$ under $y\mapsto-y$; odd coefficients integrate to zero. We obtain
\begin{equation}\label{shr:sa:eq:rhoexp}
 g(m,n)=C_m m^{mn}n^{-\nu}
 \left(\sum_{j=0}^J A_j(m,\rho)n^{-j}+O_{m,J,I}(n^{-J-1})\right),
 \qquad A_0=1.
\end{equation}
The coefficients are analytic in $\rho$ on $(0,1)$.

They are independent of the cut. Choose two rational cut proportions and take $n$ through common multiples for which both cuts are integers. The left side of \eqref{shr:sa:eq:rhoexp} is identical. Successive uniqueness of asymptotic coefficients gives equality at both rational values, at every order. Continuity extends the equality to all cut proportions. Thus coefficients may be calculated at the balanced cut $\rho=1/2$. They hold for all integers $n$ by applying \eqref{shr:sa:eq:rhoexp} at the parity-free cut $k=n$, $\rho=1/m$. This proves the all-orders assertion without extrapolating from an even subsequence.

\subsection{An exact coefficient rule}
Let $\E_m$ denote expectation on $\sum y_i=0$ with normalized density proportional to
$\D(y)^2e^{-2\sum y_i^2}$. Put
$\beta_k=B_{2k}/(2k(2k-1))$, with $B_j$ the Bernoulli numbers. The balanced-cut calculation gives
\begin{equation}\label{shr:sa:eq:recipe}
 c_j(m)=\coef tj\E_m\exp\mathcal L_m(t,y),
\end{equation}
where the formal logarithm is
\begin{align}
 \mathcal L_m(t,y)={}&
 \sum_{k\ge1}\beta_k\bigl(2^{2k}m^{1-2k}-m\bigr)t^{2k-1}
 +\sum_i\mathcal B(t,y_i)\notag\\
 &+\sum_{i<j}\sum_{r\ge1}\frac{(y_i+y_j)^{2r}}r t^r,\label{shr:sa:eq:logrecipe}
\end{align}
with
\begin{align*}
 \mathcal B(t,y)={}&-\sum_{r\ge2}
       \frac{2^{2r}y^{2r}}{2r(2r-1)}t^{r-1}
       +\sum_{r\ge1}\frac{2^{2r-1}y^{2r}}r t^r\\
 &+\sum_{k\ge1}\beta_k t^{2k-1}
 \left\{1-2^{2k-1}\left[(1+2\sqrt t\,y)^{-(2k-1)}
                         +(1-2\sqrt t\,y)^{-(2k-1)}\right]\right\}.
\end{align*}
The braces contain only integer powers of $t$. The first line of \eqref{shr:sa:eq:logrecipe} accounts for the factorial prefactor and central binomials; its last line expands the paired denominators. Every requested coefficient uses finitely many terms.

For an exact rational evaluation, write $S_r=\sum_i y_i^r$, with $S_0=m$ and $S_1=0$. Integration by parts on the trace-zero hyperplane gives, for $k\ge1$ and any polynomial $P$ in the power sums,
\begin{align}\label{shr:sa:eq:ward}
 4\E_m[S_{k+1}P]={}&
 \E_m\left[\left(\sum_{a=0}^{k-1}S_aS_{k-1-a}
                      -\frac{k}{m}S_{k-1}\right)P\right]\notag\\
 &+\E_m\left[\sum_{j\ge1}j\left(S_{k+j-1}
             -\frac{S_kS_{j-1}}m\right)\frac{\partial P}{\partial S_j}\right].
\end{align}
Use the tangent vector field $(y_i^k-S_k/m)P$. The Gaussian derivative contributes $-4S_{k+1}$; the Vandermonde derivative gives the difference-quotient sum, and tangent projection gives the $1/m$ terms. Gaussian decay removes boundary terms, and the squared Vandermonde removes the apparent singularity at a coincidence. Each use of \eqref{shr:sa:eq:ward} decreases total polynomial degree by two. Starting from $\E_m1=1$, it therefore computes every moment by finite rational arithmetic. This also proves the effective and rational coefficient assertion in Theorem~\ref{shr:sa:thm:main}.

For implementing \eqref{shr:sa:eq:logrecipe}, the useful identity
\[
 \sum_{i<j}(y_i+y_j)^{2r}
 =\frac12\left(\sum_{a=0}^{2r}\binom{2r}{a}S_aS_{2r-a}
                              -2^{2r}S_{2r}\right)
\]
expresses every term in power sums. If $\mathcal L=\sum_{k\ge1}L_kt^k$ and
$e^{\mathcal L}=\sum_{j\ge0}H_jt^j$, then
$H_0=1$ and $H_j=j^{-1}\sum_{k=1}^j kL_kH_{j-k}$. This is the finite procedure used by the bundled coefficient generator.

\subsubsection{The first two coefficients}
The first two logarithmic terms are
\[
 L_1=-\frac{m^2-1}{3m}+mS_2-\frac43S_4,\qquad
 L_2=-\frac43S_2+\frac m2S_4-\frac{32}{15}S_6+\frac32S_2^2.
\]
The moments required by $c_1=\E_mL_1$ and
$c_2=\E_mL_2+\tfrac12\E_mL_1^2$ are
\begin{align*}
 \E_mS_2&=\frac{m^2-1}{4},&
 \E_mS_4&=\frac{(m^2-1)(2m^2-3)}{16m},\\
 \E_mS_6&=\frac{5(m^2-1)(m^4-3m^2+3)}{64m^2},&
 \E_mS_2^2&=\frac{(m^2-1)(m^2+1)}{16},\\
 \E_m(S_2S_4)&=\frac{(m^2-1)(m^2+3)(2m^2-3)}{64m},\\
 \E_mS_4^2&=\frac{(m^2-1)(4m^6+20m^4-99m^2+135)}{256m^2}.
\end{align*}
Substitution gives \eqref{shr:sa:eq:c12}; in particular,
\begin{equation}\label{shr:sa:eq:logc2}
 c_2-\frac{c_1^2}{2}=\frac{m^2(m^2-1)}{96}.
\end{equation}
A second calculation uses a traceless Hermitian Gaussian matrix with covariance
\[
 \E[X_{ij}X_{k\ell}]=\frac14
       \left(\delta_{i\ell}\delta_{jk}-\frac{\delta_{ij}\delta_{k\ell}}m\right).
\]
Its eigenvalue density is the one defining $\E_m$. Replacing $S_r$ by $\Tr(X^r)$ and summing Wick pairings gives the same symbolic moment formulas. A separate direct Gaussian integration check verifies the two corrections at heights two, three and four.

There are also exact low-height checks. Since $g(2,n)$ is the $(n-1)$st Catalan number, its corrections are $3/8$ and $25/128$. Sun's proved recurrence simplifies to
\[
 g(3,n+1)+g(3,n)=\frac{(3n)!}{(n!)^3}
                    \frac{7n+1}{2(n+1)^2(4n^2-1)}.
\]
Substituting \eqref{shr:sa:eq:c12} agrees with its expansion through the second correction.

\subsection{Inversion and integer thresholds}
Set $L=m\log m>0$ and define
\[
 y=\log\frac{x}{C_mL^\nu},\qquad d_1=Lc_1(m),\qquad
 d_2=L^2\left(c_2(m)-\frac{c_1(m)^2}{2}\right).
\]
Writing $w=Ln$, the logarithmic forward expansion is
$y=w-\nu\log w+d_1/w+d_2/w^2+\cdots$. Formal substitution gives the smooth inverse
\begin{align}\label{shr:sa:eq:inverse}
 w={}&y+\nu\log y+\frac{\nu^2\log y-d_1}{y}\notag\\
 &+\frac{-\tfrac12\nu^3(\log y)^2
       +(\nu^3+\nu d_1)\log y-\nu d_1-d_2}{y^2}
       +O\left(\frac{(\log y)^3}{y^3}\right).
\end{align}
Every further coefficient is obtained recursively and is a polynomial in $\log y$. The derivative of the logarithmic forward approximation tends to $L$, so the forward remainder transfers to the inverse by the mean-value theorem.

For precision about the discrete sequence, \eqref{shr:sa:eq:expansion} implies
$g(m,n+1)/g(m,n)\to m^m>1$, hence eventual strict increase. Let
$N(x)=\min\{n:g(m,n)\ge x\}$, and let $n_J(x)$ solve the positive $J$-term forward approximation in \eqref{shr:sa:eq:expansion} for sufficiently large real $n$. Comparing the two adjacent threshold integers gives
\begin{equation}\label{shr:sa:eq:rounding}
 N(x)=\left\lceil n_J(x)+\varepsilon_J(x)\right\rceil,
 \qquad |\varepsilon_J(x)|=O_{m,J}((\log x)^{-J-1}).
\end{equation}
Thus taking the ceiling is certified away from an error-sized neighborhood of an integer. An unqualified rounding rule at every $x$ is not asserted. In particular,
\[
 N(x)=\frac{\log x+\nu\log\log x-\log C_m-\nu\log L}{L}+O(1).
\]
The displayed logarithmic formula \eqref{shr:sa:eq:inverse} has its own stated truncation error; the sharper error in \eqref{shr:sa:eq:rounding} refers to solving the corresponding finite forward approximation.

\subsection{Verification and relation to prior results}
The source package contains an endpoint-path enumerator checking sixty small rectangles, thirty published initial terms across A181198 and A181199, and 121 interior hook identities. These are checks of the exact combinatorial correspondence, not a substitute for the uniform proof. Two symbolic moment methods agree on \eqref{shr:sa:eq:c12}; a third, direct Gaussian-integral calculation checks the indicated small heights. A separate formal substitution checks \eqref{shr:sa:eq:inverse}. The generator implements the finite procedure \eqref{shr:sa:eq:recipe}--\eqref{shr:sa:eq:ward} at a requested order.

The hook products, Gaussian expansion and Poisson summation are classical tools. The argument proves the fixed-height comparison and its complete dominant-sector expansion directly for the array counts. The live A181199 entry and the cited primary sources were checked on 1 October 2026; the displayed equivalent remains conjectural there. This is a bounded literature-positioning statement, not an exhaustive priority claim. The order-two and order-three recurrence conjectures in \cite{KK} are not consequences asserted by this report. No proof-assistant certification is claimed.

\subsection{Further questions}
The conjectured exact recurrences and closed forms for heights four and five remain natural targets. The present proof determines their asymptotic behavior without establishing those identities.

A second question is to identify the exponentially smaller sectors omitted by the Poincar\'e expansion, including the contributions of saturated boundary shapes. This would turn the dominant-sector expansion into a fuller transseries.

It is also useful to let the height grow with the width. Since the first relative correction is asymptotic to $m^3/(12n)$, the regime $m^3/n$ of constant order is a concrete candidate for a new uniform approximation. This observation is a proposed scale, not a growing-height theorem. Finally, effective error constants would turn the qualified inverse bounds into numerical certificates for specific large thresholds.

\begin{writenote}[Added 5 October 2026, batch~101: the shipped layout of Report~85]
``The source package'' of Subsection~43.7 is the delivered archive. Its
\texttt{verify.py} and six programs \texttt{code/}\emph{name}\texttt{.py}
are shipped as \texttt{code/\allowbreak 05-\allowbreak strip-\allowbreak asymptotics-\allowbreak verify.py}
and \texttt{code/\allowbreak 05-\allowbreak strip-\allowbreak asymptotics-}\emph{name}\texttt{.py}; the six
records \texttt{code/}\emph{name}\texttt{.json}, \texttt{CHECKS.json} and
\texttt{verification.json} as \texttt{data/\allowbreak 05-\allowbreak strip-\allowbreak asymptotics-}\emph{name}.
\texttt{verify.py} runs each program with working directory
\texttt{code/}, and each program writes its JSON beside itself, so they run
only on a copy with the delivered layout restored, never in place; the
README of the report gives the commands. At placement the whole suite
passed on such a copy (SymPy~1.14.0; six and a half minutes on a loaded
machine), and every refreshed record equalled the shipped one after removal
of carriage returns, except the timing fields of
\texttt{verification.json}. The suite runs its programs with \texttt{-O},
which is harmless: they test with \texttt{raise}, not \texttt{assert}.
\texttt{code/\allowbreak 05-\allowbreak strip-\allowbreak asymptotics-\allowbreak check\_\allowbreak midpoint.py} embeds
fifteen terms each of A181198 and A181199 copied from OEIS; as OEIS content
they are published by The OEIS Foundation Inc.\ under the Creative Commons
Attribution-ShareAlike~4.0 licence.
\end{writenote}

\section{Report 91: rational diagonals and transcendence for fixed-height shifted strips}
\label{shr:sd:sec}

\noindent{\large Research report prepared for Vladimir Reshetnikov with
OpenAI, 2 October 2026\par}

\paragraph{Abstract.}
For every fixed positive height $m$, the number $g(m,n)$ of increasing shifted-strip arrays is a multiple binomial sum in the width $n$. Its ordinary generating function is therefore a rational diagonal and is D-finite. The proof records each row's birth and saturation, leaving a fixed number of strict-chamber walks. Their reflection determinants become everywhere-defined binomial products, and the resulting sum has at most $m^2$ integer coordinates. We also prove that the generating function is algebraic exactly for $m=1,2$. For larger heights, the leading asymptotic gives either logarithmic radial growth after differentiation or an impossible inverse-square-root coefficient. A concise proof of the leading equivalent is included. Exact checks use cell-poset counts, event sums and independent chamber walks. The argument proves existence of polynomial-coefficient recurrences, without asserting either of the particular recurrences conjectured for OEIS A181198 and A181199.

\begin{writenote}[Added 5 October 2026, batch~101]
Theorem~\ref{shr:sd:thm:diagonal} is Part~\ref{shr:bp:part}'s
Theorem~\ref{shr:bp:thm:diagonal-main}, proved by the same route two days
earlier by date line, and Theorem~\ref{shr:sd:thm:trans} is
Part~\ref{shr:part:one}'s Theorem~\ref{shr:thm:algebraic}
(Section~\ref{shr:sa:sec:chronology}). Its bound of $m^2$ coordinates and
Part~\ref{shr:bp:part}'s $m(m-1)$ free indices agree: both proofs count the
same $\sum_jh_{j-1}\le m^2$ slots, and Part~\ref{shr:bp:part} then removes
the $m$ coordinates that the saturation events fix at $n-1$, which
Report~91 keeps and sums against the Kronecker factor $\del(b^j_1-n+1)$.
The $(3,4)$ table of Subsection~44.4 is the table of Section~17.2. Since
5~October 2026 the particular recurrences that the abstract leaves aside
are proved in Parts~\ref{shr:h4:part} and~\ref{shr:h5:part}.
Subsection~44.6 reproduces the argument of Subsection~43.3, as its first
paragraph says, adding the explicit entropy exponent and the finite union
over pairs; it is printed as delivered.
\end{writenote}

\subsection{The statements and their scope}
Let $g(m,n)$ count the $m\times n$ arrays filled bijectively by $1,\ldots,mn$ in which rows increase left to right, columns increase downward, southeast diagonals increase downward, and antidiagonals increase from upper right to lower left. Set $g(m,0)=1$, the empty-array convention, and write
\[
 G_m(z)=\sum_{n\ge0}g(m,n)z^n.
\]
This is the array model in OEIS A181196. Its height-four and height-five sequences are A181198 and A181199 \cite{oeis196,oeis198,oeis199}. It agrees with the truncated shifted-strip tableau model discussed by Sun \cite{Sun}.

\begin{theorem}\label{shr:sd:thm:diagonal}
For every fixed integer $m\ge1$, the sequence $n\mapsto g(m,n)$ is a multiple binomial sum over $\Q$. There is a rational function $R_m$ in finitely many variables, regular at the origin and with rational coefficients, such that
\[
 G_m(z)=\diag R_m(z).
\]
Consequently $G_m$ is D-finite and $g(m,n)$ is P-recursive: there exist $r\ge0$ and polynomials $p_0,\ldots,p_r\in\Q[n]$, not all zero, with
\[
 \sum_{j=0}^r p_j(n)g(m,n+j)=0
\]
for all sufficiently large integers $n$. The binomial-sum construction uses at most $m^2$ summation coordinates for each of finitely many event patterns.
\end{theorem}
Here a diagonal means $[z^n]\diag R=[x_1^n\cdots x_d^n]R$. D-finiteness means that a nonzero linear differential operator with polynomial coefficients annihilates the series.

\begin{theorem}\label{shr:sd:thm:trans}
For positive integer heights, $G_m(z)$ is algebraic over $\C(z)$ exactly when $m=1$ or $m=2$. In those cases,
\[
 G_1(z)=\frac1{1-z},\qquad
 G_2(z)=\frac{3-\sqrt{1-4z}}2.
\]
Thus every fixed height $m\ge3$ gives a transcendental D-finite rational diagonal.
\end{theorem}

Kauers and Koutschan \cite[Section 6.4]{KK} conjecture specific recurrences of order two and degree nine for height four, and order three and degree twenty-four for height five. The OEIS entries still label those formulas conjectural as checked on 2 October 2026. Theorem~\ref{shr:sd:thm:diagonal} proves existence of a recurrence at every height; it does not identify a minimal operator or prove either displayed conjecture. Sun's formulas for $g(n,4)$ and $g(n,5)$ fix the other dimension, while his formula for $g(3,n)$ fixes height three. These distinctions matter because the array is not symmetric in its dimensions.

\subsection{Paths and their boundary events}
Fill an array in increasing label order. Let $\lambda_i$ be the number of cells already filled in row $i$. The row-choice word determines the array uniquely.

\begin{lemma}\label{shr:sd:lem:states}
Legal fillings are in bijection with paths using positive coordinate unit steps from $(0,\ldots,0)$ to $(n,\ldots,n)$ through states satisfying
\begin{equation}\label{shr:sd:eq:states}
 n\ge\lambda_1\ge\cdots\ge\lambda_m\ge0,
 \qquad\text{with repetitions allowed only at $0$ and $n$.}
\end{equation}
\end{lemma}
\begin{proof}
If row $i+1$ contains $j$ filled cells, with $1\le j<n$, its last cell requires the upper-right cell $(i,j+1)$ to be filled, so $\lambda_i\ge j+1$. If $j=n$, column increase requires the upper row to be full. This proves necessity. Conversely, when a cell is inserted along a path satisfying \eqref{shr:sd:eq:states}, its same-row predecessor, upper predecessor and upper-right predecessor, when present, have already been inserted. These are the required local inequalities. Southeast inequalities follow from row and column increase. Thus every path gives a legal array.
\end{proof}

For the moment let $n\ge2$. Mark each row's \emph{birth}, the increment $0\to1$, and its \emph{saturation}, the increment $n-1\to n$. They are distinct events even at $n=2$. Births occur in row order, and saturations occur in row order. At every time the number of births is at least the number of saturations. The event word therefore belongs to the set $\mathcal D_m$ of Dyck words with $m$ letters $B$ and $m$ letters $S$, where $B$ raises height and $S$ lowers it.

After $r$ births and $s$ saturations, the state has the form
\[
 (\underbrace{n,\ldots,n}_{s},a_1,\ldots,a_{r-s},
       \underbrace{0,\ldots,0}_{m-r}),
 \qquad n>a_1>\cdots>a_{r-s}>0.
\]
The active rows are one contiguous strict block. Between events, only its coordinates change.

Fix $w\in\mathcal D_m$, and let $h_j$ be its height after event $j$, with $h_0=h_{2m}=0$. Immediately before event $j$ record a vector $b^j$ of length $h_{j-1}$; immediately afterward record $a^j$. Set $a^0=()$. The event rules are
\begin{align}
 w_j=B:&\quad a^j=(b^j_1,\ldots,b^j_{h_{j-1}},1),
 &&b^j_{h_{j-1}}\ge2\quad\text{if }h_{j-1}>0;\label{shr:sd:eq:birth}\\
 w_j=S:&\quad a^j=(b^j_2,\ldots,b^j_{h_{j-1}}),
 &&b^j_1=n-1.\label{shr:sd:eq:saturate}
\end{align}
All recorded pre-event vectors satisfy
\begin{equation}\label{shr:sd:eq:guards}
 n-1\ge b^j_1>\cdots>b^j_{h_{j-1}}\ge1.
\end{equation}
Empty vectors obey the empty condition. A birth into an empty active block needs no additional guard: any preceding row is already saturated at $n\ge2$. A Dyck word has positive height before a saturation.

\subsection{The chamber transition as a binomial sum}
For strictly decreasing integer vectors $a,b$ of length $k$, let $K_k(a,b)$ count paths from $a$ to $b$, with steps in $\{e_1,\ldots,e_k\}$, staying strictly decreasing. Set $K_0((),())=1$.

\begin{lemma}\label{shr:sd:lem:reflection}
If $L=\sum_i(b_i-a_i)\ge0$, then
\begin{equation}\label{shr:sd:eq:det}
 K_k(a,b)=L!\det_{1\le i,j\le k}\left(\frac1{(b_i-a_j)!}\right),
\end{equation}
where a reciprocal factorial with negative argument is zero. If $L<0$, then $K_k(a,b)=0$.
\end{lemma}
\begin{proof}
Expand the determinant as a signed sum over starting vectors obtained by permuting $a$. For a permutation $\sigma$, the unrestricted walks to $b$ require $b_i-a_{\sigma(i)}$ steps of type $e_i$. Their number is the corresponding multinomial, or zero if one difference is negative; the total length is always $L$.

Cancel walks having a coordinate collision by a sign-reversing involution. Choose the first collision time and the lexicographically first colliding pair at that time. Exchange those two coordinate histories up to and including that collision. The endpoint does not change, while the starting permutation changes by a transposition. Before the chosen time the set of coordinate values is unchanged, and at the collision the whole vector is unchanged, so the chosen first collision is preserved. This is an involution. In an uncancelled walk no two coordinates collide. Since one coordinate changes by one at each step, coordinate ordering cannot change without a collision. The strictly decreasing endpoint forces the initial permutation to be the identity. The survivors are precisely the counted chamber walks. Negative total length is impossible.
\end{proof}

We now remove every problematic factorial from this expression. For all $N,K\in\Z$, use the convention
\begin{equation}\label{shr:sd:eq:binomconv}
 \binom NK=[x^K](1+x)^N,
 \qquad \binom NK=0\quad(K<0).
\end{equation}
For negative $N$, the right side is a formal power series at zero. This is the convention of Bostan, Lairez and Salvy \cite[Section 1.1]{BLS}.

\begin{lemma}\label{shr:sd:lem:binom}
For all strict integer endpoints,
\begin{equation}\label{shr:sd:eq:kernel}
 K_k(a,b)=\sum_{\sigma\in\mathfrak S_k}\sgn(\sigma)
       \prod_{i=1}^k\binom{d_i+\cdots+d_k}{d_i},
 \qquad d_i=b_i-a_{\sigma(i)}.
\end{equation}
Every binomial factor is defined for all integer values of the endpoints.
\end{lemma}
\begin{proof}
If every $d_i$ is nonnegative, the product telescopes to
$(\sum_i d_i)!/\prod_i d_i!$. If some $d_i$ is negative, its own factor vanishes by \eqref{shr:sd:eq:binomconv}, even when an upper argument elsewhere is negative. Thus the signed sum is \eqref{shr:sd:eq:det} whenever $L\ge0$. When $L<0$, every permutation has a negative $d_i$, so every product vanishes. The empty product and the unique empty permutation give $K_0=1$.
\end{proof}

The unrestricted multinomials in this proof already count the ordering of the interior increments. It would be incorrect to multiply the final event sum by another factor for interleaving them.

\subsection{An exact event sum with bounded dimension}
Let $\mathcal B_w(n)$ be the integer pre-event vectors satisfying \eqref{shr:sd:eq:birth}--\eqref{shr:sd:eq:guards}, with $a^j$ defined there. We may include the redundant coordinatewise inequalities $b^j_i\ge a^{j-1}_i$: if they fail, the transition count is already zero.

\begin{proposition}\label{shr:sd:prop:events}
For $n\ge2$,
\begin{equation}\label{shr:sd:eq:eventsum}
 g(m,n)=\sum_{w\in\mathcal D_m}\ \sum_{(b^1,\ldots,b^{2m})\in\mathcal B_w(n)}
       \prod_{j=1}^{2m}K_{h_{j-1}}(a^{j-1},b^j).
\end{equation}
The number of pre-event coordinates for a fixed word is
\begin{equation}\label{shr:sd:eq:dimension}
 D_w=\sum_{j=1}^{2m}h_{j-1}\le m^2.
\end{equation}
\end{proposition}
\begin{proof}
Every complete path determines exactly one event word, exactly one set of pre-event vectors, and exactly one chamber subpath in each gap. Conversely, any chamber subpaths counted on the right join through the forced boundary events to a legal full path. During a gap coordinates increase between endpoints in $[1,n-1]$, so no hidden birth or saturation is possible. The guards are exactly those needed at the junctions.

There is no extra choice of event times. The length of gap $j$ is determined by its endpoints:
\[
 L_j=|b^j|-|a^{j-1}|.
\]
Event $j$ occurs at time $j+\sum_{r\le j}L_r$. Since $a^0=a^{2m}=()$,
\[
 \sum_{j=1}^{2m}L_j
 =\sum_{j=1}^{2m}(|b^j|-|a^j|)
 =-m+m(n-1)=m(n-2).
\]
Adding the $2m$ event increments gives exactly $mn$ steps. No steps are possible before the first birth or after the last saturation.

Finally $h_j\le\min(j,2m-j)$, so
\[
 D_w\le\sum_{j=0}^{2m-1}\min(j,2m-j)=m^2.
\]
Equality holds for $B^mS^m$. The dimension therefore depends on $m$ alone.
\end{proof}

For example, at $(m,n)=(3,4)$ the five event patterns give
\[
\begin{array}{c|rrrrr}
 w&BBBSSS&BBSBSS&BBSSBS&BSBBSS&BSBSBS\\\hline
 \text{count}&4&16&4&4&1
\end{array}
\]
and their sum is $g(3,4)=29$. At width two, only the alternating event word survives, giving $g(m,2)=1$. At width one the events coincide and we do not use the decomposition; directly $g(m,1)=1$.

\subsection{Closure and the rational diagonal theorem}
We use the algebra of multiple binomial sums over $\Q$ as defined in \cite[Definition 1.1]{BLS}. It contains the Kronecker delta, nonzero geometric sequences and \eqref{shr:sd:eq:binomconv}; it is closed under addition, multiplication, integer-affine substitutions and directed indefinite summation. In particular
\[
 H(t)=\mathbf1_{\{t\ge0\}},\qquad \del(t)=\mathbf1_{\{t=0\}}
\]
are binomial sums. Equation (6) of \cite{BLS} gives closure under affine summation bounds; multiplying by $H$ also gives ordinary sums with the empty-sum convention.

For fixed $w$, let every pre-event coordinate range independently from $0$ to $n$. Impose \eqref{shr:sd:eq:guards} by products of
\[
 H(b^j_i-1),\quad H(n-1-b^j_i),\quad
 H(b^j_i-b^j_{i+1}-1).
\]
The birth and saturation guards are respectively
\[
 H(b^j_{h_{j-1}}-2)\quad(h_{j-1}>0),
 \qquad\del(b^j_1-n+1).
\]
The optional monotonicity guards are $H(b^j_i-a^{j-1}_i)$. Because each $a^j$ is obtained from $b^j$ by appending $1$ or deleting its first entry, all these arguments are integer-affine in $n$ and the summation coordinates. By \eqref{shr:sd:eq:kernel}, the same is true of every upper and lower binomial argument. The full guarded summand is therefore a binomial sum, and its $D_w$ nested sums with bounds $0,n$ preserve that property. There are only finitely many $w$.

Write this everywhere-defined finite-sum expression as $E_m(n)$. The zero extension of our counting sequence from $\N$ to $\Z$ is
\begin{equation}\label{shr:sd:eq:extension}
 \del(n)+\del(n-1)+H(n-2)E_m(n).
\end{equation}
Thus it too is a binomial sum. All factors are defined on all integer inputs, so this step does not conceal a negative-factorial singularity. Equivalently, one may use the bounded-fiber polyhedral summation closure in \cite[Proposition 3.12]{BLS}; the direct affine-bound route already suffices.

Theorem 3.5 of \cite{BLS} now supplies a rational power series over $\Q$ whose diagonal is $G_m$. Its Corollary 3.6 supplies a polynomial-coefficient recurrence. The corresponding generating function is D-finite. This proves Theorem~\ref{shr:sd:thm:diagonal}. The closure proof is constructive in principle; it does not assert that the resulting rational representation has few variables, or that a practically small recurrence follows without further simplification.

\subsection{The leading equivalent used for transcendence}
For completeness we give the needed asymptotic argument. It is the leading-order part of the earlier fixed-height report, Report~85 (Section~\ref{shr:sa:sec} of this
Part), reproduced here so that the transcendence conclusion does not rely on an unproved numerical constant. Put
\[
 d=\binom m2,\qquad \nu=\frac{m^2-1}{2},\qquad
 C_m=2^{-m(m-1)}\sqrt m\,(2\pi)^{-(m-1)/2}
       \prod_{j=0}^{m-1}j!.
\]
\begin{proposition}\label{shr:sd:prop:asymptotic}
For every fixed $m\ge2$,
\begin{equation}\label{shr:sd:eq:asymptotic}
 g(m,n)\sim C_m m^{mn}n^{-\nu}.
\end{equation}
\end{proposition}
\begin{proof}
Let $f^\lambda$ be the number of ordinary standard tableaux of partition shape $\lambda$, with at most $m$ rows, and put $\delta=(m-1,\ldots,0)$ and $\D(x)=\prod_{i<j}(x_i-x_j)$. The ordinary hook formula and the shifted hook formula are
\begin{align}
 f^\lambda&=\frac{|\lambda|!\D(\lambda+\delta)}
                  {\prod_i(\lambda_i+m-i)!},&
 f^{(n^m)}&=\frac{(mn)!\prod_{j=0}^{m-1}j!}
                  {\prod_{j=0}^{m-1}(n+j)!},\label{shr:sd:eq:ordinary}\\
 g^\lambda&=\frac{|\lambda|!}{\prod_i\lambda_i!}
                  \prod_{i<j}\frac{\lambda_i-\lambda_j}{\lambda_i+\lambda_j}.
                  \label{shr:sd:eq:shifted}
\end{align}
The last formula applies to positive strict $\lambda$ \cite[equation (1)]{Sun}. It counts our prefixes whenever $\lambda_1<n$, since then no row has saturated. The ordinary formula also follows by shifting weakly ordered paths by $\delta$ and evaluating the reflection determinant.

Cut an ordinary rectangular-tableau path at time $n$. For $|\lambda|=n$, set
\[
 \bar\lambda=(n-\lambda_m,\ldots,n-\lambda_1),\qquad
 \mu_n(\lambda)=\frac{f^\lambda f^{\bar\lambda}}{f^{(n^m)}}.
\]
These are probabilities on the weakly decreasing states. Complementation reverses the path and row order, and preserves \eqref{shr:sd:eq:states}. If $0<\lambda_m<\cdots<\lambda_1<n$, the admissible through-count is exactly $g^\lambda g^{\bar\lambda}$. At any boundary state, the ratio of admissible to ordinary through-count lies in $[0,1]$.

We record the concentration needed for those hook ratios. Let $\rho=1/m$ and compare $\mu_n$ with the labelled multivariate hypergeometric law
\[
 \eta_n(\ell)=\frac{\prod_i\binom n{\ell_i}}{\binom{mn}n},
 \qquad\sum_i\ell_i=n.
\]
Bounding ordinary prefixes and suffixes by unrestricted row words gives
\begin{equation}\label{shr:sd:eq:mucrude}
 \mu_n(\lambda)\le R_n\eta_n(\lambda),\qquad
 R_n=\frac{\prod_{j=0}^{m-1}(n+j)!}
           {(n!)^m\prod_{j=0}^{m-1}j!}=O_m(n^d).
\end{equation}
Stirling bounds and strict concavity of binary entropy show that any fixed linear deviation from $\lambda_i=\rho n$ has exponentially small probability. More explicitly, the exponential part of $\eta_n$ is
$\exp\{n[\sum_i H(\lambda_i/n)-mH(\rho)]\}$, with a polynomial prefactor, where $H(x)=-x\log x-(1-x)\log(1-x)$. On the compact constraint $\sum_i\lambda_i/n=1$ its exponent has a unique maximum at $(\rho,\ldots,\rho)$. Equation~\eqref{shr:sd:eq:mucrude} transfers the concentration to $\mu_n$.

In a central region $\epsilon n\le\lambda_i\le(1-\epsilon)n$, put $y_i=(\lambda_i-\rho n)/\sqrt n$. The exact hook formulas give
\begin{equation}\label{shr:sd:eq:murefine}
 \mu_n(\lambda)=R_n\eta_n(\lambda)
 \frac{\D(\lambda+\delta)^2}
 {\prod_i(\lambda_i+1)_{m-i}(n-\lambda_i+1)_{i-1}},
\end{equation}
where $(x)_j=x(x+1)\cdots(x+j-1)$. The denominator is bounded below by a positive constant times $n^{2d}$, and the numerator by $C n^d(1+\|y\|)^{2d}$. Uniform Stirling bounds on this region and the negative upper bound for $H''$ yield
\[
 \eta_n(\lambda)\le C n^{-(m-1)/2}e^{-c\|y\|^2}.
\]
Hence
\begin{equation}\label{shr:sd:eq:envelope}
 \mu_n(\lambda)\le C n^{-(m-1)/2}(1+\|y\|)^{2d}e^{-c\|y\|^2}.
\end{equation}
Gaussian lattice sums and the exponential outer-region bound imply tightness of $y$. In a fixed box $\|y\|\le M$, requiring one gap $\lambda_i-\lambda_j\le K$ leaves at most
$O_{m,M}((K+1)n^{(m-2)/2})$ lattice points on $\sum\lambda_i=n$. Equation~\eqref{shr:sd:eq:envelope} bounds their probability by $O_{m,M}((K+1)/\sqrt n)$. A finite union over pairs, then tightness, proves
\begin{equation}\label{shr:sd:eq:concentration}
 \lambda_i/n\longrightarrow\rho,\qquad
 \min_{i<j}(\lambda_i-\lambda_j)\longrightarrow\infty
 \quad\text{in probability under }\mu_n.
\end{equation}
In particular the interior strict states have probability tending to one.

At such a state, division of \eqref{shr:sd:eq:shifted} by \eqref{shr:sd:eq:ordinary} gives
\[
 \frac{g^\lambda}{f^\lambda}
 =\frac{\prod_i(\lambda_i+m-i)!/\lambda_i!}
        {\prod_{i<j}(\lambda_i+\lambda_j)}
       \frac{\D(\lambda)}{\D(\lambda+\delta)}.
\]
The first factor tends to $2^{-d}$ by the coordinate concentration, and the second tends to $1$ by the diverging gaps. The complementary ratio has the same limit. All exact through-count ratios lie in $[0,1]$, so bounded convergence in probability gives
\[
 \frac{g(m,n)}{f^{(n^m)}}\longrightarrow2^{-2d}.
\]
Finally Stirling's formula in the rectangular hook product gives
\[
 f^{(n^m)}\sim\sqrt m\,(2\pi)^{-(m-1)/2}
  \Bigl(\prod_{j=0}^{m-1}j!\Bigr)m^{mn}n^{-\nu},
\]
which proves \eqref{shr:sd:eq:asymptotic}.
\end{proof}

In particular the two principal OEIS rows have
\[
 g(4,n)\sim\frac{3}{1024\sqrt2\,\pi^{3/2}}
                      \frac{256^n}{n^{15/2}},\qquad
 g(5,n)\sim\frac{9\sqrt5}{2^{17}\pi^2}
                      \frac{3125^n}{n^{12}}.
\]
The second equivalent is conjectured in A181199. The proof above does not use its conjectural recurrence. No all-orders estimate is needed for what follows.

\subsection{Elementary algebraic obstructions}
The following radial arguments avoid any assumption about the number or nature of other singularities on the circle of convergence.

\begin{lemma}\label{shr:sd:lem:abel}
If $b_n\sim C/n$ with $C>0$, then
\[
 \sum_{n\ge1}b_nt^n\sim C\log\frac1{1-t}\quad(t\uparrow1).
\]
If $b_n\sim Cn^{-1/2}$, then
\[
 \sum_{n\ge1}b_nt^n\sim C\sqrt\pi\,(1-t)^{-1/2}\quad(t\uparrow1).
\]
\end{lemma}
\begin{proof}
For the first claim compare the tail with $(C\pm\varepsilon)/n$ and use $\sum t^n/n=-\log(1-t)$; the finite head is negligible. For the second put $x=-\log t$ and $f(u)=u^{-1/2}e^{-u}$. Since $f$ is positive and decreasing,
\[
 \int_x^\infty f(u)\,du
 \le x\sum_{n\ge1}f(nx)
 \le\int_0^\infty f(u)\,du=\sqrt\pi.
\]
Thus $\sqrt x\sum n^{-1/2}e^{-nx}\to\sqrt\pi$. Now $x\sim1-t$, and tail comparison as before handles $b_n$.
\end{proof}

\begin{lemma}\label{shr:sd:lem:descent}
If $F\in\Q[[t]]$ is algebraic over $\C(t)$, it is algebraic over $\Q(t)$. Algebraic functions in characteristic zero remain algebraic under differentiation and multiplication by $t$.
\end{lemma}
\begin{proof}
A complex polynomial relation is a dependence among finitely many series $t^iF^j$ having rational coefficients. Arrange their coefficients as the columns of a matrix with finitely many columns and possibly infinitely many rows. Its rational row space has finite dimension, so choose finitely many rows spanning it over $\Q$. A nonzero complex kernel implies that this finite rational matrix has rank less than its number of columns. It therefore has a nonzero rational kernel, which annihilates every row and gives a rational polynomial relation. For differentiation, implicit differentiation of a minimal polynomial gives $F'=-P_t(t,F)/P_y(t,F)$; the denominator is nonzero in characteristic zero. The remaining closure is immediate.
\end{proof}

\begin{lemma}\label{shr:sd:lem:obstructions}
An algebraic function over $\C(t)$ cannot satisfy
$F(1-s)\sim C\log(1/s)$ for $C>0$ as $s\downarrow0$.
If a function algebraic over $\Q(t)$ satisfies
$F(t)\sim L(1-t)^{-1/2}$ radially, then $L$ is algebraic.
\end{lemma}
\begin{proof}
For the first assertion, write a nonzero polynomial relation in the form
\[
 0=P(1-s,F(1-s))=\sum_{i\ge i_0}s^i p_i(F(1-s)),
 \qquad p_{i_0}\ne0.
\]
Let $d=\deg p_{i_0}$ and divide by $s^{i_0}F(1-s)^d$. The first term tends to the nonzero leading coefficient of $p_{i_0}$. Every later term tends to zero, since any positive power of $s$ dominates every fixed power of $\log(1/s)$. This is impossible.

For the second, take $P(t,y)\in\Q[t,y]\setminus\{0\}$ with $P(t,F(t))=0$, and let $d=\deg_y P$. The polynomial
\[
 Q_0(u,v)=u^d P(1-u^2,v/u)
\]
is nonzero. Divide it by its largest common power of $u$, obtaining $Q(u,v)\in\Q[u,v]$ with $Q(0,v)\ne0$. For $u>0$ small, put $v(u)=uF(1-u^2)$. Then $Q(u,v(u))=0$ and $v(u)\to L$. Taking limits gives $Q(0,L)=0$.
\end{proof}

\begin{proof}[Proof of Theorem~\ref{shr:sd:thm:trans}]
Normalize the radius by
\[
 A_m(t)=G_m(m^{-m}t),\qquad \tht=t\frac{d}{dt}.
\]
Proposition~\ref{shr:sd:prop:asymptotic} gives
\[
 [t^n]\tht^k A_m(t)\sim C_m n^{k-\nu}.
\]
If $m\ge3$ is odd, then $\nu$ is an integer. Choose $k=\nu-1$. Lemma~\ref{shr:sd:lem:abel} gives
\[
 \tht^{\nu-1}A_m(t)\sim C_m\log\frac1{1-t},
\]
contradicting Lemma~\ref{shr:sd:lem:obstructions} if $G_m$ were algebraic.

If $m\ge4$ is even, put $k=\lfloor\nu\rfloor=(m^2-2)/2$. Then
\[
 \tht^kA_m(t)\sim L_m(1-t)^{-1/2},\qquad
 L_m=C_m\sqrt\pi=\alpha_m\pi^{-(m-2)/2},
\]
where
\[
 \alpha_m=2^{-m(m-1)-(m-1)/2}\sqrt m
             \prod_{j=0}^{m-1}j!\in\overline\Q\setminus\{0\}.
\]
The number $L_m$ is transcendental: otherwise the positive integer power $\pi^{(m-2)/2}$ would be algebraic, contradicting the transcendence of $\pi$ \cite{Lindemann}. If $G_m$ were algebraic over $\C(z)$, its rational coefficients, rational rescaling and Lemma~\ref{shr:sd:lem:descent} would make $\tht^k A_m$ algebraic over $\Q(t)$. Lemma~\ref{shr:sd:lem:obstructions} would then force $L_m$ to be algebraic, a contradiction.

At height one $g(1,n)=1$. At height two the row-choice path is a primitive Dyck path of semilength $n$: equality of its two coordinates is allowed only initially and finally. Removing the first up-step and last down-step leaves a Dyck path of semilength $n-1$. Hence
\[
 g(2,n)=\Cat_{n-1}=\frac1n\binom{2n-2}{n-1}\quad(n\ge1),
\]
and with the constant term $1$ this gives the stated $G_2$.
\end{proof}

\begin{remark}
The same obstruction is reusable whenever a rational-coefficient sequence has an equivalent $C R^{-n}n^{-\nu}$ with rational nonzero $R$. A positive integer $\nu$ produces the logarithmic contradiction. A half-integer $\nu>0$ produces an obstruction when $C\sqrt\pi$ is transcendental. These are sufficient tests, not a classification of algebraic series from arbitrary coefficient asymptotics.
\end{remark}

\subsection{Exact verification and remaining questions}
The accompanying standard-library Python checker uses exact integers and rational numbers. It compares independent implementations rather than a guessed recurrence:
\begin{enumerate}
\item A cell-poset algorithm counts linear extensions from predecessor masks for horizontal, vertical and both downward diagonal edges.
\item A row-state algorithm counts paths satisfying \eqref{shr:sd:eq:states}.
\item The full Dyck-event sum evaluates \eqref{shr:sd:eq:eventsum} with the signed binomial kernel \eqref{shr:sd:eq:kernel}.
\item Direct chamber walks are compared with both the factorial determinant and the binomial kernel, including reversed and zero-length endpoint pairs.
\end{enumerate}
The recorded tests pass for $54$ cell-poset/state instances with $1\le m\le6$ and $0\le n\le8$, including $44$ event-sum instances. There are $2941$ chamber comparisons with up to four coordinates, and $1554$ binomial-product tests including negative lower arguments. All Dyck-word coordinate bounds are checked for $1\le m\le9$. Selected counts are
\[
\begin{array}{c|rrrrrr}
 &n=0&n=1&n=2&n=3&n=4&n=5\\\hline
 m=2&1&1&1&2&5&14\\
 m=3&1&1&1&4&29&290\\
 m=4&1&1&1&8&169&6392\\
 m=5&1&1&1&16&985&141696
\end{array}
\]
The mathematical proof covers every fixed height and all widths; these finite checks audit its indexing and boundary conventions.

The rational diagonal supplied by the closure theorem is not printed as one expanded rational function. Producing a compact representation and reducing it to a manageable differential equation are practical next steps. Proving the specific height-four and height-five conjectures, their minimal orders, and any finer subdominant exponential sectors remain separate problems. No assertion of priority beyond the cited sources and the existing project reports is intended.

\begin{writenote}[Added 5 October 2026, batch~101: the shipped layout of Report~91]
Its \texttt{verify.py}, \texttt{code/\allowbreak check\_\allowbreak exact.py} and
\texttt{build\_local.sh} are shipped as
\texttt{code/\allowbreak 06-\allowbreak strip-\allowbreak diagonals-\allowbreak verify.py}, \texttt{-check\_\allowbreak exact.py} and
\texttt{-build\_\allowbreak local.sh}; \texttt{code/\allowbreak check\_\allowbreak exact.json},
\texttt{CHECKS.json}, \texttt{SOURCES.json} and \texttt{verification.json}
as \texttt{data/\allowbreak 06-\allowbreak strip-\allowbreak diagonals-}\emph{name}. The delivered
\texttt{verify.py} reads and rewrites \texttt{code/\allowbreak check\_\allowbreak exact.json}
and writes \texttt{verification.json} at its root, and \texttt{build\_local.sh} changes
to its own directory, so they run only on a copy with the delivered layout
restored. At the write (5~October 2026, Python~3.14.4, on such a copy)
\texttt{verify.py} passed in under two seconds, reported
\texttt{matches\_\allowbreak distributed\_\allowbreak data: true}, and both refreshed records equal
the shipped ones after removal of carriage returns.
\end{writenote}

\section{[write] Chronology, correspondence and further questions of Part~V}
\label{shr:sa:sec:chronology}

Added 5 October 2026, batch~101. This section is the write's.

\subsection{What Reports~85 and~91 prove, against Parts~I and~II}
Report~85 (1~October 2026) and Report~91 (2~October 2026) were written
without access to this report, and the manuscripts of Part~\ref{shr:part:one}
(3~October) and Part~\ref{shr:bp:part} (4~October) without access to them.
The table pairs their statements. ``Same route'' means the same chain of
ideas, not the same text: the word 8-gram overlaps with
Parts~\ref{shr:part:one} and~\ref{shr:bp:part} are 0.4\,\% and 1.1\,\%
for Report~85, and 0.0\,\% and 1.1\,\% for Report~91 (7.1\,\% of
Report~91's 8-grams recur in Report~85).
{\footnotesize
\begin{longtable}{@{}>{\raggedright\arraybackslash}p{0.27\linewidth}>{\raggedright\arraybackslash}p{0.29\linewidth}>{\raggedright\arraybackslash}p{0.38\linewidth}@{}}
\toprule
Part~V & Parts~I--II & relation\\
\midrule\endhead
Theorem~\ref{shr:sa:thm:main}, \eqref{shr:sa:eq:c12}, \eqref{shr:sa:eq:height5} & Theorem~\ref{shr:thm:main}, \eqref{shr:eq:b12}, Corollary~\ref{shr:cor:A181199} & the same statements ($c_j(m)=b_{m,j}$, $C_m=K_m$); another route: an ordinary-tableau comparison and a cut at $k=n$, Poisson summation and cut independence, against Part~I's random-threshold cut (Theorem~\ref{shr:thm:cut})\\
\eqref{shr:sa:eq:comparison} & Corollary~2.4, leading term & the same statement, proved directly by bounded convergence\\
Subsection~43.5, \eqref{shr:sa:eq:ward} & Section~5 and the integration by parts of Section~6 & a second route to the coefficients: a loop-equation recursion for the trace-zero Gaussian--Vandermonde moments, and a traceless-GUE Wick check (classical tools)\\
\eqref{shr:sa:eq:logc2} & \eqref{shr:eq:b12} & a consequence (Remark~\ref{shr:sa:rem:logc2})\\
Sun's height-three recurrence (Subsection~43.5.1) & --- & not in Parts~I--IV; see item~(F12)\\
\eqref{shr:sa:eq:inverse}, \eqref{shr:sa:eq:rounding} & Section~10.2 & a second route to the $W_{-1}$ inversion with two corrections, with the rounding qualification\\
Theorem~\ref{shr:sd:thm:diagonal} & Theorem~\ref{shr:bp:thm:diagonal-main} & the same theorem by the same route (Dyck event words, reflection determinants, everywhere-defined binomial products, Bostan--Lairez--Salvy); bound $m^2$ against $m(m-1)$ (note at the head of Section~\ref{shr:sd:sec})\\
Theorem~\ref{shr:sd:thm:trans} & Theorem~\ref{shr:thm:algebraic} & the same theorem and the same two obstructions (a logarithm at odd $m$, the amplitude $\alpha_m\pi^{-(m-2)/2}$ at even $m$, the constant of \eqref{shr:eq:amplitude}); elementary radial Lemmas~\ref{shr:sd:lem:abel}--\ref{shr:sd:lem:obstructions} replace the Puiseux argument\\
Proposition~\ref{shr:sd:prop:asymptotic} & Theorem~\ref{shr:thm:main}, leading term & Report~85's Subsection~43.3, reproduced with more detail\\
the $(3,4)$ table of Subsection~44.4 & Section~17.2 & identical entries $4,16,4,4,1$\\
\bottomrule
\end{longtable}}

So, within this report, the first proofs by date line of
Theorems~\ref{shr:thm:main} and~\ref{shr:thm:algebraic} and of
Theorem~\ref{shr:bp:thm:diagonal-main} are those of Reports~85 and~91.
This bears on item~(F4), the priority of
Theorem~\ref{shr:bp:thm:diagonal-main}: that item asked for a search of the
general literature, which is still missing; what is new is that an earlier
manuscript in the same corpus proves the theorem by the same route.
Nothing here is a priority claim outside ProveIt, and the two reports'
own source checks were, in their words, bounded.

\begin{remark}[\textbf{[write]} Equation \eqref{shr:sa:eq:logc2} from Part~I]
\label{shr:sa:rem:logc2}
Report~85's identity $c_2-c_1^2/2=m^2(m^2-1)/96$ follows from
Part~\ref{shr:part:one}'s \eqref{shr:eq:b12}, since $c_j(m)=b_{m,j}$:
\[
 b_{m,2}-\frac{b_{m,1}^2}2
 =\frac{(m^2-1)\bigl((m^6+3m^2-1)-(m^2-1)^3\bigr)}{288m^2}
 =\frac{(m^2-1)\,3m^4}{288m^2}
 =\frac{m^2(m^2-1)}{96}.
\]
\end{remark}

\subsection{Further questions and research recorded at the write}
\label{shr:sa:sec:further}
Under Vladimir Reshetnikov's standing rule, the unproved and outline-level
claims of the two manuscripts are recorded here. None of them is wrong as
far as the intake could check; the intake confirmed the height-five and
height-four displays, \eqref{shr:sa:eq:logc2}, the $(3,4)$ table and
$G_2$.
\begin{enumerate}[label=(F\arabic*),leftmargin=*,start=11]
\item \emph{Outline steps of Report~85} (Subsection~43.4). Three steps are
argued in outline: the growing-window Taylor expansion and its
exponentiation (``Choosing $\varepsilon$ small permits exponentiation of the
remainder without losing an integrable Gaussian majorant''), the uniformity
of Poisson summation over the affine lattice translations, and the passage
from \eqref{shr:sa:eq:rhoexp} to the rounding statement
\eqref{shr:sa:eq:rounding}. The theorem itself is proved in full in
Part~\ref{shr:part:one} (Theorem~\ref{shr:thm:main}, with the explicit
Taylor bound \eqref{shr:eq:taylorbound}), so nothing depends on them.
Missing: explicit window and remainder estimates for this route.
\item \emph{Sun's height-three recurrence} (Subsection~43.5.1). The report
states that ``Sun's proved recurrence simplifies to''
$g(3,n+1)+g(3,n)=\frac{(3n)!}{(n!)^3}\frac{7n+1}{2(n+1)^2(4n^2-1)}$; its
program \texttt{05-\allowbreak strip-\allowbreak asymptotics-\allowbreak ward\_\allowbreak corrections.py} checks the
expansion against this form, and the intake checked the displayed identity
for $n=1,\ldots,24$ against row-state counts. Missing: the derivation of
this form from the recurrence as printed in \cite{Sun}, which neither the
report nor the intake reproduced.
\item \emph{Growing height and smaller sectors} (Subsection~43.8). The scale
$m^3/n$ is ``a proposed scale, not a growing-height theorem'' (Research
question~3); the exponentially smaller sectors and saturated boundary shapes
are Research question~4 and items~(F1)--(F2); effective inverse constants
are Research question~6, now carried out at $m=4$ and $m=5$ in
Parts~\ref{shr:h4:part} and~\ref{shr:h5:part}.
\item \emph{A compact rational diagonal} (Subsection~44.8). ``The rational
diagonal supplied by the closure theorem is not printed as one expanded
rational function''; a compact representation, its differential equation,
the minimal orders and finer subdominant sectors are open. These are
Research questions~9 and~11; at $m=4$ and $m=5$ Parts~\ref{shr:h4:part}
and~\ref{shr:h5:part} supply the operators, and
Propositions~\ref{shr:h4:prop:minimal} and~\ref{shr:h5:prop:nothyper}
bound their orders.
\end{enumerate}

\textbf{Independent check of the write (dated note, 5~October 2026).}
An independent adversarial check of the batch-101 write of Parts~III--V,
made by the intake after the write (\texttt{ca1668212}), reread every
result the write added: Lemma~\ref{shr:h4:lem:Zirr},
Proposition~\ref{shr:h4:prop:minimal} with the remark after it and
item~(F5), Remark~\ref{shr:h5:rem:sameZ},
Proposition~\ref{shr:h5:prop:nothyper} with item~(F9), and
Remark~\ref{shr:sa:rem:logc2}. It found all of them valid, with no
counterexample and no gap in any proof chain. Two changes followed.
Remark~\ref{shr:h5:rem:sameZ} called Part~IV's transfer, rank-two and
scalar identities ``literally the same statements'' as Part~III's; they
are the same after renaming Part~III's $\alpha$, $\beta$ to Part~IV's $f$,
$b$, and Part~IV's $J_0=I$ corresponds to the display
\eqref{shr:h4:eq:J0}, which the remark did not cite. The remark is
reworded where it stands, with a dated note keeping the first wording. And
the question left open after Proposition~\ref{shr:h4:prop:minimal} and in
item~(F5), whether a recurrence of order three or more has degree below
nine, has an answer: the check found and proved
Propositions~\ref{shr:h4:prop:higher} and~\ref{shr:h5:prop:higher}, now
printed with their proofs. Items~(F5) and~(F9) are re-scoped with dated
notes keeping their first wording.

The check used neither the delivered programs nor the write's. Its own
counter, in~C, is a dynamic program over the order ideals of the literal
$m\times n$ poset (row, column and antidiagonal relations), each ideal
given by its row lengths; it runs modulo fourteen primes below $2^{62}$
and recombines by the Chinese remainder theorem, which is exact because
$a_n\le M_n$ and $M_{64}$ is below the product of the primes. It gave
$T_4(n)$ and $T_5(n)$ for $n\le64$. On these counts the operator of
Theorem~\ref{shr:h4:thm:oeis} has residual zero for $n=1,\ldots,62$ and the
identity of Theorem~\ref{shr:h4:thm:main} holds for $n\le63$; the operator
of Theorem~\ref{shr:h5:thm:operator}, with its coefficients typed from
Appendix~\ref{shr:h5:app:literal}, has residual zero for $n=1,\ldots,61$;
and the printed sums of Theorems~\ref{shr:h4:thm:sum}
and~\ref{shr:h5:thm:sum} equal the counts for $2\le n\le64$ and
$1\le n\le64$. Ingredients of the proofs were tested on the same counts:
the affine formula \eqref{shr:h4:eq:affine} for $n\le64$ and the first
difference \eqref{shr:h5:eq:target-diff} for $n\le63$, with $Z_n$ from
\eqref{shr:h4:eq:Zrec} and $Z_1=1/3$ (and $Z_1,Z_2,Z_3=1/3,13/45,461/1680$
by direct triangle integrals); $Y_1=0$ and $Y_n>0$ for $2\le n\le12$,
consistent with the proof of Proposition~\ref{shr:h5:prop:nothyper}, which
needs $Y_n\ne0$ only for large $n$. In SymPy the check confirmed the hinge
identities of both reports. For Report~231: the divergence certificate
\eqref{shr:h4:eq:polycert}, the slanted-edge identity, the edge flux
\eqref{shr:h4:eq:edgeflux} (also against direct integration for $n\le6$),
the cancellation of the $Z_n$ coefficient with the rational return $D_n$,
$q_n=J_{n+1}/J_n$ \eqref{shr:h4:eq:qforcing}, and the three coefficient
identities in the proof of Theorem~\ref{shr:h4:thm:oeis}. For Report~233:
the $\gamma$-bridge \eqref{shr:h5:eq:gamma-bridge} with $\gamma_1=1/3$,
$[Z_n](F_{n+1}-\rho_nF_n)=2\rho_nu(n)$ with the matching constant term,
\eqref{shr:h5:eq:Pplus} and \eqref{shr:h5:eq:Qplus} coefficient by
coefficient, and the factorization
\eqref{shr:h5:eq:operator-factorization}, through the exact right division
of the new operators by $L_{\rm OEIS}$ with zero remainder. For the
minimality statements it computed ranks of the matrices $(n^ia_{n+j})$
modulo $2^{61}-1$ and $2^{62}-57$. At height four: no solution of order one
and degree at most $50$, and at order two a solution space of dimension
exactly $\max(0,d-8)$ for every degree $d\le40$, which confirms the
order-two statement of Proposition~\ref{shr:h4:prop:minimal} in that
range. At height five: no solution of order one or two and degree at most
$60$ on the windows used (for order two, $5\le n\le214$), and at order
three nullity $0$ up to degree $23$ and $1$ at degree $24$. The same
computations, with equations from $n=30$, located the degrees of
Propositions~\ref{shr:h4:prop:higher} and~\ref{shr:h5:prop:higher}, and
the right division by the proved operators turned them into proofs.

Before printing, the session that applied the check re-derived the four
operators of those propositions by exact Ore multiplication and right
division (SymPy~1.14.0) and found zero residual for $n=1,\ldots,234$ on
terms generated by Theorems~\ref{shr:h4:thm:oeis}
and~\ref{shr:h5:thm:operator}; those terms agree with the check's counts
for $n\le64$ and with this report's row-state program
(\texttt{code/tableaux.py}) for $n\le18$ at height four and $n\le14$ at
height five. It redid the rank computations of both proofs exactly over
$\Q$ (python-flint~0.8.0), with equations from $n=1$. The check was a
careful reading with exact computations, not a formal verification or an
external review.

\appendix
\section{A direct proof of the shifted hook product}\label{shr:app:hook}

Let $\lambda_1>\cdots>\lambda_r>0$ be a strict partition of $k$, and put
\[
 H(\lambda)=\frac{k!}{\prod_i\lambda_i!}
 \prod_{i<j}\frac{\lambda_i-\lambda_j}{\lambda_i+\lambda_j}.
\]
We prove that it satisfies the strict-partition path recursion
\begin{equation}
 H(\lambda)=\sum_i H(\lambda-e_i),
 \label{shr:eq:hookrec}
\end{equation}
where only valid strict-partition removals are included and a trailing zero
part is deleted. Together with $H(\varnothing)=1$, this proves
\eqref{shr:eq:shiftedhook} by induction.

For a valid removal,
\begin{equation}
 \frac{H(\lambda-e_i)}{H(\lambda)}
 =\frac{\lambda_i}{k}
 \prod_{j\ne i}
 \frac{\lambda_i-\lambda_j-1}{\lambda_i-\lambda_j}
 \frac{\lambda_i+\lambda_j}{\lambda_i+\lambda_j-1}.
 \label{shr:eq:removalratio}
\end{equation}
The same expression is zero for an invalid removal causing equal positive
parts. If the last part is one, deleting its resulting zero gives the
same ratio, since the extra zero factors cancel. Thus it suffices to prove
\begin{equation}
 \sum_i\lambda_i
 \prod_{j\ne i}
 \frac{\lambda_i-\lambda_j-1}{\lambda_i-\lambda_j}
 \frac{\lambda_i+\lambda_j}{\lambda_i+\lambda_j-1}=k.
 \label{shr:eq:rationalidentity}
\end{equation}

Consider the rational function
\[
 R(z)=\prod_{j=1}^r
 \frac{z-\lambda_j-1}{z-\lambda_j}
 \frac{z+\lambda_j}{z+\lambda_j-1},\qquad
 f(z)=(z-\tfrac12)R(z).
\]
Its possible finite poles are the disjoint sets $\lambda_i$ and
$1-\lambda_i$. Direct substitution gives $R(1-z)=R(z)$ and hence
$f(1-z)=-f(z)$. The residues at the paired poles are equal. At
$z=\lambda_i$, the residue is the negative of the $i$th summand in
\eqref{shr:eq:rationalidentity}.

At infinity, each pair of factors contributes $-2\lambda_j/z^2$ to the
first nonconstant term, so
\[
 R(z)=1-\frac{2k}{z^2}+O(z^{-3}),\qquad
 f(z)=z-\tfrac12-\frac{2k}{z}+O(z^{-2}).
\]
The sum of finite residues is the coefficient of $z^{-1}$, namely $-2k$.
On the other hand it is minus twice the left side of
\eqref{shr:eq:rationalidentity}. The identity follows, proving
\eqref{shr:eq:hookrec} and the hook product. All operations here are rational
identities; the residue language can equivalently be replaced by partial
fractions.

\section{An implementation-independent coefficient specification}

For a fixed height $m$ and requested order $L$, the following finite recipe
fully specifies the coefficient calculation.

First compute $h_0,\ldots,h_L$ from
$\prod_{i<j}(1-u(y_i+y_j)^2)^{-1}$ modulo $u^{L+1}$.
Next compute the sign cumulants $\kappa_r$ from
\[
 \log\cosh t=\sum_{r\ge1}\kappa_r\frac{t^r}{r!}.
\]
The even values begin $\kappa_2=1$, $\kappa_4=-2$, $\kappa_6=16$,
$\kappa_8=-272$; the odd values vanish. The moment polynomials can be
computed recursively as
\begin{equation}
 \mu_0(z)=1,\qquad
 \mu_k(z)=\sum_{\substack{2\le r\le k\\r\text{ even}}}
 \binom{k-1}{r-1}\kappa_r z^{r/2-1}\mu_{k-r}(z)
 \quad(k>0),
 \label{shr:eq:momentrec}
\end{equation}
with $\mu_k=0$ for odd $k$.
Apply \eqref{shr:eq:collapseperms} to each $h_r$, evaluate each monomial by
\eqref{shr:eq:Ez}, retain powers through $z^{L-r}$, and combine them as in
\eqref{shr:eq:coefficientformula}. Finally apply the formal exponential
\eqref{shr:eq:stirling-convert}. This recipe is sufficient to reproduce the
printed coefficients independently of the shipped Python code.

For $m=5$, the two outputs through order five are
\begin{align*}
 (c_{5,j})_{j=0}^5&=
 \left(1,10,\frac{225}{4},\frac{1819}{8},\frac{22045}{32},1353\right),\\
 (b_{5,j})_{j=0}^5&=
 \left(1,\frac{48}{5},\frac{5233}{100},\frac{1028391}{5000},
 \frac{60248377}{100000},\frac{273939939}{250000}\right).
\end{align*}
The first two nonconstant entries in each row also have the independent
symbolic derivation in Section~6. Higher printed entries are exact
evaluations of the proved finite formula, not inferred from observed
numerical errors.

\section{A fully finite recipe for the exact sum}
\label{shr:bp:app:exact}
\begin{writenote}[Added 5 October 2026, batch~98]
Appendices~C and~D belong to Part~\ref{shr:bp:part}: they are the
manuscript's Appendices~A and~B.
\end{writenote}
For clarity, the diagonal theorem can be implemented without making any
choices suggested by numerical data. The following recipe has fixed
size once $m$ is fixed.

First list all Dyck words of semilength $m$. For a chosen word, make one
vector $v_s$ of length $r_s$ for each interval $0\le s<2m$, and replace
the top coordinate before each completion by $n-1$. Sum every remaining
coordinate from $0$ to $n$. Construct $u_0=()$ and the later $u_s$ using
\eqref{shr:bp:eq:skeletonmap}. Insert $H$ factors for
\[
 v_{s,r_s}\ge1,\quad v_{s,1}\le n-1,\quad
 v_{s,i}-v_{s,i+1}\ge1,
\]
where the conditions with coordinates are omitted when $r_s=0$.
For a birth with $r_s>0$, also insert $H(v_{s,r_s}-2)$.
The already substituted completion coordinate enforces its event condition.

For each interval independently choose a permutation
$\sigma_s\in\Sym_{r_s}$, multiply its sign, and set
\[
 a_{s,i}=v_{s,i}-u_{s,\sigma_s(i)}.
\]
Multiply the binomial product \eqref{shr:bp:eq:binomialM} for this displacement
vector. Sum over the finitely many permutations and endpoint coordinates,
and then over the words. This expression equals $T_m(n)$ for every
$n\ge2$ by the bijective proof of Theorem~\ref{shr:bp:thm:skeleton}.

The endpoint-coordinate order in this recipe can be arbitrary; using
nested partial sums and then affine substitution of their upper bounds
realizes the rectangular sums in the binomial-sum algebra. Values outside
the intended polytope are killed by the explicit $H$ factors. In
particular, one never applies a recurrence or divides by a factor that
could vanish at a summation endpoint.

The delivered evaluator uses the same formula more efficiently by
iterating over decreasing endpoint tuples directly and caching suffix
computations. That optimization changes neither the mathematical formula
nor its fixed-dimensional binomial nature. It is not claimed to beat the
prefix dynamic program for every input.

\section{An implementation-independent coefficient algorithm}
\label{shr:bp:app:algorithm}
The symbolic routine can be reproduced using rational polynomial
arithmetic alone. Here are sufficient details to fix every convention.
For a nonnegative integer $a$, form the polynomial in $p$
\begin{equation}
 C_a(n,p)=\sum_{b=0}^a\sum_{k=0}^b
 \binom ab2^b(-n)^{a-b}\Stirling{b}{k}\fall nk\,p^k.
 \label{shr:bp:eq:centeredC}
\end{equation}
It is exactly $\E[(2X-n)^a\mid U=p]$.
For a monomial $y_1^{a_1}\cdots y_r^{a_r}$, set $A=\sum_i a_i$ and
multiply $\prod_i C_{a_i}(n,p)=\sum_{k=0}^A c_k(n)p^k$.
Then
\begin{equation}
 \E\prod_iY_i^{a_i}
 =n^{-A/2}\sum_{k=0}^A c_k(n)
 \frac{\rise nk}{\rise{2n+1}k}.
 \label{shr:bp:eq:moment-rational}
\end{equation}
A convenient common denominator is $\rise{2n+1}A$; its numerator is
\begin{equation}
 \sum_{k=0}^A c_k(n)\rise nk\rise{2n+1+k}{A-k}.
 \label{shr:bp:eq:moment-numerator}
\end{equation}
Reverse the numerator and denominator polynomials and perform ordinary
power-series division at $1/n=0$, retaining the extra factor $n^{-A/2}$.
This is the exact moment functional $\mathcal E_\eps$ used in
\eqref{shr:bp:eq:Qr}. It avoids repeated simplification of many separate
rational functions.

To expand the single-coordinate interaction in \eqref{shr:bp:eq:Psidef}, write
\begin{equation}
 \frac{1-(2+y^2)\eps^2-2y\eps^3}
 {1-(6+y^2)\eps^2/9+2y\eps^3/9}
 =\sum_{k\ge0}c_k(y)\eps^k.
 \label{shr:bp:eq:single-series}
\end{equation}
With $c_k=0$ for $k<0$, its coefficients satisfy
\begin{equation}
 c_k=N_k+\frac{6+y^2}{9}c_{k-2}-\frac{2y}{9}c_{k-3},
 \quad
 N_0=1,\ N_2=-(2+y^2),\ N_3=-2y,
 \label{shr:bp:eq:single-rec}
\end{equation}
and all other $N_k$ are zero. The pair interaction contributes the
geometric series
\[
 \sum_{k\ge0}\frac{(y_i+y_\ell)^{2k}}{4^k}\eps^{2k},\qquad i<\ell.
\]

Multiply these finitely truncated factors, then multiply by
$\Delta(y)^2$. Apply \eqref{shr:bp:eq:moment-rational} to each monomial, divide
by $r!J_r$, and finally multiply by the truncated exponential
\eqref{shr:bp:eq:outer}. No asymptotic fitting or guessed differential equation
enters this calculation. For the height-five coefficients in
\eqref{shr:bp:eq:five}, a truncation at $\eps^6$ suffices.

As intermediate checks, the engine gives
\[
 [z]\mathcal Q_1=-\frac83,\qquad
 [z]\mathcal Q_2=-\frac{119}{18},\qquad
 [z]\mathcal Q_3=-\frac{34}{3},
\]
which also follow from \eqref{shr:bp:eq:q1}. The package stores exact coefficient
strings and reruns their derivation in \texttt{code/verify.py}.

\section{The full literal order three operator}
\label{shr:h5:app:literal}
\begin{writenote}[Added 5 October 2026, batch~101]
Appendix~E belongs to Part~\ref{shr:h5:part}: it is the manuscript's
Appendix~A.
\end{writenote}
This appendix records all four source polynomials without requiring an
external lookup. It also gives the shifted polynomial identities that
connect the compact notation to the printed source. Write
\[
 B_{21}(n)=\sum_{j=0}^{21}b_jn^j,\qquad
 C_{19}(n)=\sum_{j=0}^{19}d_jn^j,
\]
with the integer coefficients in the table below. The distinct notation
$B_{21}$ avoids confusion with the integral operator $B=EA$.
Define
\begin{align}
 Q_+(n)={}&137855872n^{11}+2377384288n^{10}+18238806776n^9\\
 &+81945774178n^8+238738633847n^7+471293180347n^6\\
 &+638861237719n^5+588363536007n^4+354332674386n^3\\
 &+128320688700n^2+23068181160n+1077753600.
 \label{shr:h5:eq:Qplus}
\end{align}
Then the literal coefficient list is
\begin{align}
 c_0(n)={}&25(2n-1)^2(2n+1)(4n-1)(4n+1)Q_+(n)\\
          &\hspace{1em}\cdot\prod_{j\in\{1,2,3,4,6,7,8,9\}}(5n+j),\\
 c_1(n)={}&-15(2n+1)\prod_{j=6}^9(5n+j)\,C_{19}(n),\\
 c_2(n)={}&-(n+2)^2(2n+3)B_{21}(n),\\
 c_3(n)={}&3(n+2)^2(2n+3)(2n+5)^2(3n+7)(3n+8)\\
          &\hspace{1em}\cdot(4n+9)(4n+11)Q(n)(n+3)^4.
 \label{shr:h5:eq:literal-coefficients}
\end{align}
Each has degree exactly 24. These expressions were transcribed from the
A181199 recurrence before comparison with the factorization. Expanding
the products and comparing all coefficients verifies
\eqref{shr:h5:eq:literal-identities} identically, with the four zero residuals
recorded by the exact verifier.

The second degree-eight polynomial printed in the denominator of $v$ is
\begin{align}
 P_+(n)={}&25216n^8+211616n^7+760768n^6+1543976n^5\\
         &+1973632n^4+1683047n^3+971955n^2+353502n+60480.
 \label{shr:h5:eq:Pplus}
\end{align}
The literal expansions satisfy $P_+(n)=P(n+1)$ and
$Q_+(n)=Q(n+1)$. These equalities account for the compact forms used
in~\eqref{shr:h5:eq:v} and~\eqref{shr:h5:eq:S}; neither is an approximation.

\begin{center}
\small
\begin{longtable}{@{}r r r@{}}
\caption{The complete integer coefficient arrays of $B_{21}$ and $C_{19}$}\label{shr:h5:tab:literal-arrays}\\
\toprule
$j$&$b_j$&$d_j$\\\midrule\endfirsthead
\toprule $j$&$b_j$&$d_j$\\\midrule\endhead
\bottomrule\endfoot
0&$-1104708217536000$&$257403484800$\\
1&$-22603977271411200$&$4014999151560$\\
2&$-180089041888657440$&$21777752002716$\\
3&$-685802053101717288$&$50327965592874$\\
4&$-850259283025327524$&$42448206362469$\\
5&$3499276436019940446$&$-54713722756179$\\
6&$19628106098287909191$&$-793918042456325$\\
7&$47097079083848116511$&$-4487751704725767$\\
8&$65692918122110754573$&$-12634296322883951$\\
9&$48833460779191449763$&$-17410620603191989$\\
10&$-1215620047630796049$&$-3351043407516635$\\
11&$-46634308960957698567$&$30740208891967721$\\
12&$-54791308745183340309$&$61042725728660150$\\
13&$-31771892184486994333$&$64302800289940820$\\
14&$-6080171043591548610$&$44061606220301336$\\
15&$5462710847171622988$&$20586994844739808$\\
16&$5556799141672247640$&$6541981632511232$\\
17&$2616700630350746208$&$1357324488921088$\\
18&$760770273998231808$&$166164038596608$\\
19&$139665528127686656$&$9113927409664$\\
20&$14962008250398720$&$0$\\
21&$717088749346816$&$0$\\
\end{longtable}
\end{center}

\begin{thebibliography}{99}
\bibitem{oeis196} The OEIS Foundation Inc.,
\emph{A181196: arrays increasing in four directions}.
\url{https://oeis.org/A181196}. Inspected October 3, 2026.

\bibitem{oeis197} The OEIS Foundation Inc.,
\emph{A181197: the three-row sequence}.
\url{https://oeis.org/A181197}. Inspected October 3, 2026.

\bibitem{oeis198} The OEIS Foundation Inc.,
\emph{A181198: the four-row sequence}.
\url{https://oeis.org/A181198}. Inspected October 3, 2026.

\bibitem{oeis199} The OEIS Foundation Inc.,
\emph{A181199: the five-row sequence}.
\url{https://oeis.org/A181199}. The entry attributes the asymptotic conjecture
to V. Kotesovec and the guessed recurrence to C. Koutschan.
Inspected October 3, 2026.

\bibitem{b198} C. Koutschan, with earlier terms from A. P. Heinz,
\emph{Table for A181198, $n=1,\ldots,100$}.
\url{https://oeis.org/A181198/b181198.txt}. Selected values inspected
October 3, 2026.

\bibitem{b199} C. Koutschan, with earlier terms from A. P. Heinz,
\emph{Table for A181199, $n=1,\ldots,100$}.
\url{https://oeis.org/A181199/b181199.txt}. Selected values inspected
October 3, 2026.

\bibitem{oeis198int} The OEIS Foundation Inc.,
\emph{A181198}, with its internal formula and recurrence record.
\url{https://oeis.org/A181198/internal}. Cited in Part~III: retrieved by
Report~231 on October 5, 2026; ``The inspected internal snapshot retains
the conjectural labels; no change to the entry is part of this report.''

\bibitem{oeis199int} The OEIS Foundation Inc.,
\emph{A181199}, with its internal formula and recurrence record.
\url{https://oeis.org/A181199/internal}. Cited in Part~IV: retrieved by
Report~233 on October 5, 2026; ``This report proves the inspected
recurrence; no edit to the source entry is part of the manuscript or its
reproduction.''

\bibitem{KK} M. Kauers and C. Koutschan,
\emph{Some D-finite and some possibly D-finite sequences in the OEIS},
Journal of Integer Sequences \textbf{26} (2023), article 23.4.5.
\url{https://arxiv.org/abs/2303.02793}.
The arXiv PDF inspected has Section 6.4 on the two shifted-rectangle
sequences; pagination differs from that cited on the OEIS pages.

\bibitem{BLS}
A. Bostan, P. Lairez, and B. Salvy,
\emph{Multiple binomial sums},
Journal of Symbolic Computation \textbf{80}, Part 2 (2017), 351--386.
\url{https://arxiv.org/abs/1510.07487}.
DOI: \href{https://doi.org/10.1016/j.jsc.2016.04.002}{10.1016/j.jsc.2016.04.002}.
See Definition 1.1, Theorem 3.5, and Corollary 3.6.
\bibitem{Sun} P. Sun,
\emph{Enumeration of standard Young tableaux of shifted strips with
constant width}, Electronic Journal of Combinatorics \textbf{24}(2)
(2017), P2.41. Preprint: \url{https://arxiv.org/abs/1506.07256}.
The inspected preprint states the shifted hook product and the
order-statistics model and treats both parameter orientations explicitly.

\bibitem{Chan} B. T. Chan, \emph{Periodic P-partitions},
European Journal of Combinatorics (2023).
Preprint, revised 2020: \url{https://arxiv.org/abs/1803.05594}.
Used here to distinguish fixed row length from fixed number of rows.

\bibitem{FS} P. Flajolet and R. Sedgewick,
\emph{Analytic Combinatorics}, Cambridge University Press, 2009.
Author-maintained companion site: \url{https://ac.cs.princeton.edu/}.
Standard background on algebraic generating functions and Puiseux
singularities; no specific sequence theorem is attributed to this source.

\bibitem{DLMFgamma} NIST Digital Library of Mathematical Functions,
\S5.11, \emph{Asymptotic expansions of the gamma function}.
\url{https://dlmf.nist.gov/5.11}. Inspected October 3, 2026.

\bibitem{DLMFpoly} NIST Digital Library of Mathematical Functions,
\S18.19, \emph{Hahn class: definitions}.
\url{https://dlmf.nist.gov/18.19}. Orthogonal-polynomial background;
the normalization used in this paper is derived explicitly.

\bibitem{DLMFW} NIST Digital Library of Mathematical Functions,
\S4.13, \emph{Lambert $W$-function}.
\url{https://dlmf.nist.gov/4.13}. Inspected October 3, 2026.

\bibitem{Lindemann} F. Lindemann, \emph{Ueber die Zahl $\pi$},
Mathematische Annalen \textbf{20} (1882), 213--225.
\url{https://eudml.org/doc/157031}. Cited in Part~V (Report~91).

\bibitem{ProveIt} V. Reshetnikov, \emph{ProveIt repository}.
\url{https://github.com/VladimirReshetnikov/ProveIt}.
Revision inspected: \texttt{6bf7f30d0352f7596e70928b3d4f304914075907}.
In particular, the README for
\texttt{a189281-path-forest-expansions} under
\texttt{SetTheory/Cardinals/docs/reports/} was read for the repository's
source/status distinctions. It is not a dependency of the mathematical
proofs in this article.
\end{thebibliography}
\end{document}
